% Generated mechanically from frozen input p04/ja/morphisms.tex.
% Do not edit this generated file; edit the bound locale source instead.

% OK, start here.
%

\setcounter{chapter}{28}
\chapter{スキームの射}



\phantomsection
\label{morphisms-section-phantom}


\section{はじめに}
\label{morphisms-section-introduction}

\noindent
本章では、スキームの射のいくつかの型を導入する。
基本的な参考文献は \cite{EGA} である。
































\section{閉埋め込み}
\label{morphisms-section-closed-immersions}

\noindent
本節では、スキームの閉埋め込みについて先に得られた結果を明らかにする。
スキームの射 $i : Z \to X$ が閉埋め込みであるとは、(a) $i$ が $X$ の閉集合への同相写像を誘導し、(b) $i^\sharp : \mathcal{O}_X \to i_*\mathcal{O}_Z$ が全射であり、(c) $i^\sharp$ の核が切断によって局所的に生成される、という条件で定義されることを思い出そう。これはスキームの章の定義 \ref{schemes-definition-immersion} および \ref{schemes-definition-closed-immersion-locally-ringed-spaces} を参照せよ。$Z$ と $X$ がスキームであることを仮定すると、閉埋め込みを特徴づける方法は実に多数ある。

\begin{lemma}
\label{morphisms-lemma-closed-immersion}
$i : Z \to X$ をスキームの射とする。
次は同値である。
\begin{enumerate}
\item 射 $i$ は閉埋め込みである。
\item 任意のアフィン開集合 $\Spec(R) = U \subset X$ に対して、
あるイデアル $I \subset R$ が存在し、
$i^{-1}(U) = \Spec(R/I)$ が $U = \Spec(R)$ 上のスキームとして成り立つ。
\item アフィン開被覆 $X = \bigcup_{j \in J} U_j$ が存在し、
$U_j = \Spec(R_j)$ であり、各 $j \in J$ に対してあるイデアル
$I_j \subset R_j$ が存在して
$i^{-1}(U_j) = \Spec(R_j/I_j)$ が $U_j = \Spec(R_j)$ 上のスキームとして成り立つ。
\item 射 $i$ は $Z$ と $X$ の閉集合との同相写像を誘導し、
$i^\sharp : \mathcal{O}_X \to i_*\mathcal{O}_Z$ は全射である。
\item 射 $i$ は $Z$ と $X$ の閉集合との同相写像を誘導し、
写像 $i^\sharp : \mathcal{O}_X \to i_*\mathcal{O}_Z$ は全射であり、
核 $\Ker(i^\sharp)\subset \mathcal{O}_X$ は準連接イデアル層である。
\item 射 $i$ は $Z$ と $X$ の閉集合との同相写像を誘導し、
写像 $i^\sharp : \mathcal{O}_X \to i_*\mathcal{O}_Z$ は全射であり、
核 $\Ker(i^\sharp)\subset \mathcal{O}_X$ は切断によって局所的に生成される
イデアル層である。
\end{enumerate}
\end{lemma}

\begin{proof}
条件 (6) は閉埋め込みの定義そのものであり、スキームの章の定義
\ref{schemes-definition-closed-immersion-locally-ringed-spaces} および
\ref{schemes-definition-immersion} を参照せよ。
したがって (6) $\Leftrightarrow$ (1) である。スキームの補題
\ref{schemes-lemma-closed-subspace-scheme} により (1) $\Rightarrow$ (2) が成り立つ。
自明に (2) $\Rightarrow$ (3) である。

\medskip\noindent
(3) を仮定する。射のそれぞれ
$\Spec(R_j/I_j) \to \Spec(R_j)$
は閉埋め込みである（スキームの例 \ref{schemes-example-closed-immersion-affines} を参照）。
したがって $i^{-1}(U_j) \to U_j$ はその像への同相写像であり、
$i^\sharp|_{U_j}$ は全射である。ゆえに $i$ はその像への同相写像であり、
$i^\sharp$ は全射である。これは局所的に確認できるからである。
したがって (3) $\Rightarrow$ (4) を得る。

\medskip\noindent
(4) $\Rightarrow$ (1) はスキームの補題
\ref{schemes-lemma-characterize-closed-immersions} である。
(5) $\Rightarrow$ (6) は自明である。そして (6) $\Rightarrow$ (5) は
スキームの補題 \ref{schemes-lemma-closed-subspace-scheme} から従う。
\end{proof}

\begin{lemma}
\label{morphisms-lemma-closed-immersion-ideals}
$X$ をスキームとする。$i : Z \to X$ および $i' : Z' \to X$
を閉埋め込みとし、イデアル層
$\mathcal{I} = \Ker(i^\sharp)$ および $\mathcal{I}' = \Ker((i')^\sharp)$
を $\mathcal{O}_X$ のイデアル層とする。
\begin{enumerate}
\item 射 $i : Z \to X$ が $Z \to Z' \to X$ と因数分解する、すなわち
ある $a : Z \to Z'$ が存在することと、$\mathcal{I}' \subset \mathcal{I}$ であることは同値である。
このとき $a$ は閉埋め込みである。
\item $Z \cong Z'$ が $X$ 上で成り立つことと
$\mathcal{I} = \mathcal{I}'$ は同値である。
\end{enumerate}
\end{lemma}

\begin{proof}
これは閉部分空間についての議論、特にスキームの節
\ref{schemes-section-closed-immersion}、補題
\ref{schemes-lemma-closed-immersion} および
\ref{schemes-lemma-characterize-closed-subspace} から従う。
また、上の補題 \ref{morphisms-lemma-closed-immersion} の特徴づけ (3) からも直接的に従う。
\end{proof}

\begin{lemma}
\label{morphisms-lemma-closed-immersion-bijection-ideals}
$X$ をスキームとする。
$\mathcal{I} \subset \mathcal{O}_X$ をイデアル層とする。
次は同値である。
\begin{enumerate}
\item $\mathcal{I}$ は $\mathcal{O}_X$-加群の層として切断によって局所的に生成される。
\item $\mathcal{I}$ は $\mathcal{O}_X$-加群の層として準連接である。
\item $i : Z \to X$ なるスキームの閉埋め込みが存在し、その対応するイデアル層
$\Ker(i^\sharp)$ が $\mathcal{I}$ に等しい。
\end{enumerate}
\end{lemma}

\begin{proof}
(1) と (2) の同値性はスキームの補題
\ref{schemes-lemma-closed-subspace-scheme} から直ちに分かる。
(1) が成り立つなら、スキームの定義
\ref{schemes-definition-closed-subspace} および例
\ref{schemes-example-closed-subspace} により、閉部分空間 $i : Z \to X$ が存在し、
$\mathcal{I} = \Ker(i^\sharp)$ となる。
スキームの補題 \ref{schemes-lemma-closed-subspace-scheme} により、これは
スキームの閉埋め込みであり、(3) が成り立つ。
逆に、(3) が成り立つなら、スキームの補題
\ref{schemes-lemma-closed-subspace-scheme} により (2) が成り立つ
（スキームの閉埋め込みは、局所環付き空間の閉埋め込みでもあるので、この補題を適用できる）。
\end{proof}

\begin{lemma}
\label{morphisms-lemma-base-change-closed-immersion}
閉埋め込みの基底変換は閉埋め込みである。
\end{lemma}

\begin{proof}
スキームの補題 \ref{schemes-lemma-base-change-immersion} を参照せよ。
\end{proof}

\begin{lemma}
\label{morphisms-lemma-composition-closed-immersion}
閉埋め込みの合成は閉埋め込みである。
\end{lemma}

\begin{proof}
これはスキームの補題 \ref{schemes-lemma-composition-immersion} で既に見たが、
ここでは別の証明を与える。すなわち、補題 \ref{morphisms-lemma-closed-immersion} における
閉埋め込みの特徴づけ (3) から従う。実際、$I \subset R$ がイデアルで、
$\overline{J} \subset R/I$ がイデアルならば、$\overline{J} = J/I$ を満たす
あるイデアル $J \subset R$ が存在し、それは $I$ を含み、$(R/I)/\overline{J} = R/J$ となる。
\end{proof}

\begin{lemma}
\label{morphisms-lemma-closed-immersion-quasi-compact}
閉埋め込みは準コンパクトである。
\end{lemma}

\begin{proof}
この補題はスキームの補題
\ref{schemes-lemma-closed-immersion-quasi-compact} と重複している。
\end{proof}

\begin{lemma}
\label{morphisms-lemma-closed-immersion-separated}
閉埋め込みは分離的である。
\end{lemma}

\begin{proof}
この補題はスキームの補題 \ref{schemes-lemma-immersions-monomorphisms} の特殊例である。
\end{proof}






\section{埋め込み}
\label{morphisms-section-immersions}

\noindent
本節では、埋め込みに関するいくつかの事実をまとめる。

\begin{lemma}
\label{morphisms-lemma-immersion-permanence}
$Z \to Y \to X$ をスキームの射とする。
\begin{enumerate}
\item $Z \to X$ が埋め込みならば、$Z \to Y$ は埋め込みである。
\item $Z \to X$ が準コンパクトな埋め込みであり、$Y \to X$ が
準分離的ならば、$Z \to Y$ は準コンパクトな埋め込みである。
\item $Z \to X$ が閉埋め込みであり、$Y \to X$ が分離的ならば、
$Z \to Y$ は閉埋め込みである。
\end{enumerate}
\end{lemma}

\begin{proof}
いずれの場合も、次の可換図を考察すればよい。
$$
\xymatrix{
Z \ar[r] \ar[rd] & Y \times_X Z \ar[r] \ar[d] & Z \ar[d] \\
& Y \ar[r] & X
}
$$
ここで、上の二つの水平射の合成は恒等射である。
(1) を証明しよう。最初の水平射は
$Y \times_X Z \to Z$ の切断であるから、スキームの補題
\ref{schemes-lemma-section-immersion} により埋め込みである。
射 $Y \times_X Z \to Y$ は $Z \to X$ の基底変換であるから、
埋め込みである（スキームの補題 \ref{schemes-lemma-base-change-immersion}）。
最後に、埋め込みの合成は埋め込みである
（スキームの補題 \ref{schemes-lemma-composition-immersion}）。これで (1) が示された。
残り二つの結果もまったく同じように証明される。
\end{proof}

\begin{lemma}
\label{morphisms-lemma-factor-quasi-compact-immersion}
$h : Z \to X$ を埋め込みとする。
$h$ が準コンパクトならば、$h = i \circ j$ と因数分解できる。ここで
$j : Z \to \overline{Z}$ は開埋め込みであり、
$i : \overline{Z} \to X$ は閉埋め込みである。
\end{lemma}

\begin{proof}
$h$ は準コンパクトかつ準分離的であることに注意せよ
（スキームの補題 \ref{schemes-lemma-immersions-monomorphisms} を参照）。
したがって、スキームの補題 \ref{schemes-lemma-push-forward-quasi-coherent} により、
$h_*\mathcal{O}_Z$ は $\mathcal{O}_X$-加群の準連接層である。
これから
$\mathcal{I} = \Ker(\mathcal{O}_X \to h_*\mathcal{O}_Z)$
は準連接イデアル層であることが従う。スキームの節
\ref{schemes-section-quasi-coherent} を参照せよ。
$\overline{Z} \subset X$ を $\mathcal{I}$ に対応する閉部分スキームとする
（補題 \ref{morphisms-lemma-closed-immersion-bijection-ideals} を参照）。
スキームの補題 \ref{schemes-lemma-characterize-closed-subspace} により、
射 $h$ は $h = i \circ j$ と因数分解する。ここで
$i : \overline{Z} \to X$ は包含射である。
$j$ が開埋め込みであることを見るため、開部分スキーム
$U \subset X$ を、$h$ が $Z$ から $U$ への閉埋め込みを誘導するように選ぶ。
すると、$\mathcal{I}|_U$ が閉埋め込み $Z \to U$ に対応する
イデアル層であることは明らかである。したがって
$Z = \overline{Z} \cap U$ である。
\end{proof}

\begin{lemma}
\label{morphisms-lemma-factor-reduced-immersion}
$h : Z \to X$ を埋め込みとする。
$Z$ が被約ならば、$h = i \circ j$ と因数分解できる。ここで
$j : Z \to \overline{Z}$ は開埋め込みであり、
$i : \overline{Z} \to X$ は閉埋め込みである。
\end{lemma}

\begin{proof}
$\overline{Z} \subset X$ を $h(Z)$ の閉包に被約誘導閉部分スキーム構造を
入れたものとする。スキームの定義
\ref{schemes-definition-reduced-induced-scheme} を参照せよ。
スキームの補題 \ref{schemes-lemma-map-into-reduction} により、射 $h$ は
$h = i \circ j$ と因数分解する。ここで $i : \overline{Z} \to X$ は包含射であり、
$j : Z \to \overline{Z}$ である。埋め込みの定義から、ある開部分スキーム
$U \subset X$ が存在し、$h$ は $U$ への閉埋め込みを経由して因数分解する。
したがって、$\overline{Z} \cap U$ と $h(Z)$ は、ともに $U$ の被約閉部分スキームであり、
同じ基礎閉集合をもつ。ゆえに、スキームの補題
\ref{schemes-lemma-reduced-closed-subscheme} の一意性により
$h(Z) \cong \overline{Z} \cap U$ である。
したがって $j$ は $Z$ と $\overline{Z} \cap U$ との同型を誘導する。
言い換えれば、$j$ は開埋め込みである。
\end{proof}



\begin{example}
\label{morphisms-example-thibaut}
開埋め込みに続く閉埋め込みの合成として表せない埋め込みの例を挙げる。
$k$ を体とする。
$X = \Spec(k[x_1, x_2, x_3, \ldots])$ とする。
$U = \bigcup_{n = 1}^{\infty} D(x_n)$ とする。
すると $U \to X$ は開埋め込みである。
イデアル
$$
I_n =
(x_1^n, x_2^n, \ldots, x_{n - 1}^n, x_n - 1, x_{n + 1}, x_{n + 2}, \ldots)
\subset
k[x_1, x_2, x_3, \ldots][1/x_n].
$$
を考える。$I_n k[x_1, x_2, x_3, \ldots][1/x_nx_m] = (1)$ は
任意の $m \not = n$ に対して成り立つことに注意せよ。
したがって、準連接イデアル $\widetilde I_n$ は $D(x_n)$ 上にあり、
$D(x_nx_m)$ 上で一致する。すなわち、
$\widetilde I_n|_{D(x_nx_m)} = \mathcal{O}_{D(x_n x_m)}$ は
$n \not = m$ のとき成り立つ。
したがって、これらのイデアルは貼り合わさって準連接イデアル層
$\mathcal{I} \subset \mathcal{O}_U$ を与える。
$Z \subset U$ を $\mathcal{I}$ に対応する閉部分スキームとする。
よって $Z \to X$ は埋め込みである。

\medskip\noindent
$Z \to X$ を $Z \to \overline{Z} \to X$ と因数分解することはできないと主張する。
ここで $\overline{Z} \to X$ は閉埋め込みであり、
$Z \to \overline{Z}$ は開埋め込みであるとする。実際、$\overline{Z}$ は、
あるイデアル $I \subset k[x_1, x_2, x_3, \ldots]$ によって定義され、
$I_n = I k[x_1, x_2, x_3, \ldots][1/x_n]$ を満たさなければならない。
しかし、唯一の元 $f \in k[x_1, x_2, x_3, \ldots]$ で、
すべての $I_n$ に入るものは $0$ である！
したがって $I$ は存在しない。
\end{example}

\begin{lemma}
\label{morphisms-lemma-check-immersion}
$f : Y \to X$ をスキームの射とする。すべての $y \in Y$ に対して、
$f(y) \in U \subset X$ なる開部分スキームが存在し、
$f|_{f^{-1}(U)} : f^{-1}(U) \to U$ が埋め込みならば、$f$ は埋め込みである。
\end{lemma}

\begin{proof}
この主張は、スキームの節 \ref{schemes-section-immersions} の末尾にある
閉部分スキームの議論から容易に従うが、詳しい証明も与える。
$Z \subset X$ を $f(Y)$ の閉包とする。閉包をとる操作は開集合への制限と可換なので、
仮定から $f(Y) \subset Z$ は開である。したがって
$Z' = Z \setminus f(Y)$ は閉である。ゆえに
$X' = X \setminus Z'$ は $X$ の開部分スキームであり、
$f$ は $f : Y \to X'$ と包含射との合成に因数分解する。
$y \in Y$ とし、$U \subset X$ を補題の仮定にあるものとする。このとき
$U' = X' \cap U$ は $f'(y)$ の開近傍であり、
$(f')^{-1}(U') \to U'$ は閉な像をもつ埋め込みである
（補題 \ref{morphisms-lemma-immersion-permanence}）。したがって、これは閉埋め込みである。
スキームの補題 \ref{schemes-lemma-immersion-when-closed} を参照せよ。
閉埋め込みであるという性質は目標上局所的なので
（例えば補題 \ref{morphisms-lemma-closed-immersion} による）、
$f'$ は所望の閉埋め込みであると結論する。
\end{proof}








\section{閉埋め込みと準連接層}
\label{morphisms-section-closed-immersions-quasi-coherent}

\noindent
次の補題は、アーベル層に対して加群の補題
\ref{modules-lemma-i-star-exact} が果たす役割を、ついにスキーム上の準連接層に対して果たす。
加群の節 \ref{modules-section-closed-immersion} の議論も参照せよ。

\begin{lemma}
\label{morphisms-lemma-i-star-equivalence}
$i : Z \to X$ をスキームの閉埋め込みとする。
$\mathcal{I} \subset \mathcal{O}_X$ を $Z$ を切り出す準連接イデアル層とする。
関手
$$
i_* :
\QCoh(\mathcal{O}_Z)
\longrightarrow
\QCoh(\mathcal{O}_X)
$$
は完全かつ充満忠実であり、その本質像は準連接
$\mathcal{O}_X$-加群 $\mathcal{G}$ で $\mathcal{I}\mathcal{G} = 0$ を満たすものからなる。
\end{lemma}

\begin{proof}
閉埋め込みは準コンパクトかつ分離的である。補題
\ref{morphisms-lemma-closed-immersion-quasi-compact} および
\ref{morphisms-lemma-closed-immersion-separated} を参照せよ。
したがって、スキームの補題 \ref{schemes-lemma-push-forward-quasi-coherent} を適用でき、
$Z$ 上の準連接層の順像は実際に $X$ 上の準連接層である。

\medskip\noindent
加群の補題 \ref{modules-lemma-i-star-equivalence} により、関手 $i_*$ は充満忠実である。

\medskip\noindent
次に関手 $i_*$ の本質像を記述する。
$\mathcal{I}(i_*\mathcal{F}) = 0$ は任意の準連接 $\mathcal{O}_Z$-加群に対して成り立つ。例えば、
加群の補題 \ref{modules-lemma-i-star-equivalence} による。
次に、$\mathcal{G}$ を任意の準連接 $\mathcal{O}_X$-加群で
$\mathcal{I}\mathcal{G} = 0$ を満たすものとする。標準写像
$$
\mathcal{G} \longrightarrow i_* i^*\mathcal{G}
$$
が同型であることを示せば十分である\footnote{これは加群の補題
\ref{modules-lemma-i-star-equivalence} の証明において、より一般的な状況で示されている。}。
スキームと準連接加群の場合には、$X$ 上でアフィン局所的に考え、補題
\ref{morphisms-lemma-closed-immersion} およびスキームの補題
\ref{schemes-lemma-widetilde-pullback} を用いると、次の代数的主張を証明すればよい。
環 $R$、イデアル $I$、および $R$-加群 $N$ が与えられ、$IN = 0$ を満たすとき、
標準写像
$$
N \longrightarrow N \otimes_R R/I,\quad
n \longmapsto n \otimes 1
$$
は $R$-加群の同型である。この易しい代数的事実の証明は省略する。
\end{proof}

\noindent
$i : Z \to X$ を閉埋め込みとする。上の補題により、記号の濫用ではあるが、
$\mathcal{F}$ で層 $i_*\mathcal{F}$ を $X$ 上の層として表すことがよくある。

\begin{lemma}
\label{morphisms-lemma-largest-quasi-coherent-subsheaf}
$X$ をスキームとする。$\mathcal{F}$ を準連接 $\mathcal{O}_X$-加群とする。
$\mathcal{G} \subset \mathcal{F}$ を $\mathcal{O}_X$-部分加群とする。
次の性質をもつ準連接 $\mathcal{O}_X$-部分加群
$\mathcal{G}' \subset \mathcal{G}$ が一意に存在する。任意の準連接
$\mathcal{O}_X$-加群 $\mathcal{H}$ に対して、写像
$$
\Hom_{\mathcal{O}_X}(\mathcal{H}, \mathcal{G}')
\longrightarrow
\Hom_{\mathcal{O}_X}(\mathcal{H}, \mathcal{G})
$$
は全単射である。特に、$\mathcal{G}'$ は最大の準連接
$\mathcal{O}_X$-部分加群であり、$\mathcal{F}$ の部分加群として $\mathcal{G}$ に含まれる。
\end{lemma}

\begin{proof}
$\mathcal{G}_a$, $a \in A$ を準連接 $\mathcal{O}_X$-部分加群で
$\mathcal{G}$ に含まれるもの全体とする。このとき
$$
\bigoplus\nolimits_{a \in A} \mathcal{G}_a \longrightarrow \mathcal{F}
$$
の像 $\mathcal{G}'$ は準連接である。実際、$X$ 上の準連接層の射の像は準連接であり、
準連接層の直和も準連接である。スキームの節
\ref{schemes-section-quasi-coherent} を参照せよ。
加群 $\mathcal{G}'$ は $\mathcal{G}$ に含まれる。したがって、これは
最大の準連接 $\mathcal{O}_X$-加群であり、$\mathcal{G}$ に含まれる。

\medskip\noindent
公式を証明するため、$\mathcal{H}$ を準連接 $\mathcal{O}_X$-加群とし、
$\alpha : \mathcal{H} \to \mathcal{G}$ を $\mathcal{O}_X$-加群の射とする。
合成 $\mathcal{H} \to \mathcal{G} \to \mathcal{F}$ の像は、準連接層の射の像として
準連接である。したがって、それは $\mathcal{G}'$ に含まれる。
ゆえに、所望のように $\alpha$ は $\mathcal{G}'$ を経由して因数分解する。
\end{proof}

\begin{lemma}
\label{morphisms-lemma-i-upper-shriek}
$i : Z \to X$ をスキームの閉埋め込みとする。
関手\footnote{これは標準的でない記法と思われる。}
$i^! : \QCoh(\mathcal{O}_X) \to \QCoh(\mathcal{O}_Z)$ が存在し、
$i_*$ の右随伴である。（加群の補題 \ref{modules-lemma-i-star-right-adjoint} と比較せよ。）
\end{lemma}

\begin{proof}
準連接 $\mathcal{O}_X$-加群 $\mathcal{G}$ が与えられたとき、
部分層 $\mathcal{H}_Z(\mathcal{G})$、すなわち $\mathcal{G}$ のうち
$\mathcal{I}$ によって零化される局所切断からなる部分層を考える。補題
\ref{morphisms-lemma-largest-quasi-coherent-subsheaf} により、標準的な最大準連接
$\mathcal{O}_X$-部分加群 $\mathcal{H}_Z(\mathcal{G})'$ がある。構成により
$$
\Hom_{\mathcal{O}_X}(i_*\mathcal{F}, \mathcal{H}_Z(\mathcal{G})')
=
\Hom_{\mathcal{O}_X}(i_*\mathcal{F}, \mathcal{G})
$$
が任意の準連接 $\mathcal{O}_Z$-加群 $\mathcal{F}$ に対して成り立つ。
したがって $i^!\mathcal{G} = i^*(\mathcal{H}_Z(\mathcal{G})')$ と定められる。
詳細は省略する。
\end{proof}

\noindent
準連接イデアル層と閉部分スキームとの $1$ 対 $1$ の対応（補題
\ref{morphisms-lemma-closed-immersion-bijection-ideals} を参照）を用いると、
閉部分スキームのスキーム論的交わりと合併を定義できる。

\begin{definition}
\label{morphisms-definition-scheme-theoretic-intersection-union}
$X$ をスキームとする。$Z, Y \subset X$ を準連接イデアル層
$\mathcal{I}, \mathcal{J} \subset \mathcal{O}_X$ に対応する閉部分スキームとする。
$Z$ と $Y$ の {\it スキーム論的交わり}とは、$X$ の閉部分スキームで、
$\mathcal{I} + \mathcal{J}$ によって切り出されるものである。
$Z$ と $Y$ の {\it スキーム論的合併}とは、
$X$ の閉部分スキームで、$\mathcal{I} \cap \mathcal{J}$ によって切り出されるものである。
\end{definition}

\begin{lemma}
\label{morphisms-lemma-scheme-theoretic-intersection}
$X$ をスキームとする。$Z, Y \subset X$ を閉部分スキームとする。
$Z \cap Y$ を $Z$ と $Y$ のスキーム論的交わりとする。
このとき $Z \cap Y \to Z$ および $Z \cap Y \to Y$ は閉埋め込みであり、
$$
\xymatrix{
Z \cap Y \ar[r] \ar[d] & Z \ar[d] \\
Y \ar[r] & X
}
$$
はスキームのカルテジアン図式、すなわち $Z \cap Y = Z \times_X Y$ である。
\end{lemma}

\begin{proof}
射 $Z \cap Y \to Z$ および $Z \cap Y \to Y$ は、補題
\ref{morphisms-lemma-closed-immersion-ideals} により閉埋め込みである。
$U = \Spec(A)$ を $X$ のアフィン開集合とし、$Z \cap U$ および $Y \cap U$ がそれぞれ
イデアル $I \subset A$ および $J \subset A$ に対応するとする。このとき
$Z \cap Y \cap U$ は $I + J \subset A$ に対応する。
$A/I \otimes_A A/J = A/(I + J)$ なので、この図式がカルテジアンであることは、
スキームの節 \ref{schemes-section-fibre-products} における
スキームのファイバー積の記述から分かる。
\end{proof}

\begin{lemma}
\label{morphisms-lemma-scheme-theoretic-union}
$S$ をスキームとする。$X, Y \subset S$ を閉部分スキームとする。
$X \cup Y$ を $X$ と $Y$ のスキーム論的合併とする。
$X \cap Y$ を $X$ と $Y$ のスキーム論的交わりとする。
このとき $X \to X \cup Y$ および $Y \to X \cup Y$ は閉埋め込みであり、
$\mathcal{O}_S$-加群の短完全列
$$
0 \to \mathcal{O}_{X \cup Y} \to \mathcal{O}_X \times \mathcal{O}_Y
\to \mathcal{O}_{X \cap Y} \to 0
$$
が存在し、図式
$$
\xymatrix{
X \cap Y \ar[r] \ar[d] & X \ar[d] \\
Y \ar[r] & X \cup Y
}
$$
はスキームの圏において余カルテジアン、すなわち
$X \cup Y = X \amalg_{X \cap Y} Y$ である。
\end{lemma}

\begin{proof}
射 $X \to X \cup Y$ および $Y \to X \cup Y$ は、補題
\ref{morphisms-lemma-closed-immersion-ideals} により閉埋め込みである。
短完全列では、補題 \ref{morphisms-lemma-i-star-equivalence} の同値性を用いて、
$S$ の閉部分スキーム上の準連接加群を $S$ 上の準連接加群とみなす。
列の最初の射には標準写像 $\mathcal{O}_{X \cup Y} \to \mathcal{O}_X$ および
$\mathcal{O}_{X \cup Y} \to \mathcal{O}_Y$ を用い、二番目の射には標準写像
$\mathcal{O}_X \to \mathcal{O}_{X \cap Y}$ と、標準写像
$\mathcal{O}_Y \to \mathcal{O}_{X \cap Y}$ の負を用いる。
完全性を確認するにはアフィン局所的に考えればよい。
$U = \Spec(A)$ を $S$ のアフィン開集合とし、$X \cap U$ および $Y \cap U$ がそれぞれ
イデアル $I \subset A$ および $J \subset A$ に対応するとする。このとき
$(X \cup Y) \cap U$ は $I \cap J \subset A$ に対応し、
$X \cap Y \cap U$ は $I + J \subset A$ に対応する。
したがって、完全性は
$$
0 \to A/I \cap J \to A/I \times A/J \to A/(I + J) \to 0
$$
の完全性から従う。
図式が余カルテジアンであることを示すため、スキーム $T$ とスキームの射
$f : X \to T$, $g : Y \to T$ が与えられ、それらは射 $X \cap Y \to T$ として
一致するとする。目標は、一意な射 $h : X \cup Y \to T$ が存在し、
$f$ および $g$ と一致することを示すことである。
$h$ を構成するには $X \cup Y$ 上でアフィン局所的に考えればよい。スキームの節
\ref{schemes-section-glueing-schemes} を参照せよ。
$s \in X$, $s \not \in Y$ ならば、$X \to X \cup Y$ は $s$ の近傍で同型であり、
$h$ の構成は明らかである。$s \in Y$, $s \not \in X$ の場合も同様である。
$s \in X \cap Y$ の場合、アフィン開集合 $V = \Spec(B) \subset T$ を選び、
$f(s) = g(s)$ を含むようにできる。次にアフィン開集合
$U = \Spec(A) \subset S$ を、$s$ を含み、$f(X \cap U)$ および $g(Y \cap U)$ が
$V$ に含まれるように選べる。射 $f|_{X \cap U}$ および
$g|_{Y \cap V}$ は $V$ への射であり、環準同型
$$
B \to A/I
\quad\text{and}\quad
B \to A/J
$$
に対応し、これらは $A/(I + J)$ への写像として一致する。
上に表示した短完全列により、これらの環準同型は所望の環準同型
$B \to A/I \cap J$ へ一意に持ち上がる。
\end{proof}









\section{加群の台}
\label{morphisms-section-support}

\noindent
本節では、スキーム上の準連接加群の台に関する初等的な結果をまとめる。
加群層の台は加群の節 \ref{modules-section-support} で定義されていることを思い出そう。
一方、加群の台は代数の節 \ref{algebra-section-support} で定義された。
これらは一致する。

\begin{lemma}
\label{morphisms-lemma-support-affine-open}
$X$ をスキームとする。$\mathcal{F}$ を $X$ 上の準連接層とする。
$\Spec(A) = U \subset X$ をアフィン開集合とし、
$M = \Gamma(U, \mathcal{F})$ とおく。
$x \in U$ とし、$\mathfrak p \subset A$ を対応する素イデアルとする。
次は同値である。
\begin{enumerate}
\item $\mathfrak p$ は $M$ の台に属する。
\item $x$ は $\mathcal{F}$ の台に属する。
\end{enumerate}
\end{lemma}

\begin{proof}
これは等式 $\mathcal{F}_x = M_{\mathfrak p}$ から従う。スキームの補題
\ref{schemes-lemma-spec-sheaves} および定義を参照せよ。
\end{proof}

\begin{lemma}
\label{morphisms-lemma-support-closed-specialization}
$X$ をスキームとする。
$\mathcal{F}$ を $X$ 上の準連接層とする。
$\mathcal{F}$ の台は特殊化について閉じている。
\end{lemma}

\begin{proof}
$x' \leadsto x$ が特殊化であり $\mathcal{F}_x = 0$ ならば、
$\mathcal{F}_{x'}$ は零である。実際、$\mathcal{F}_{x'}$ は加群
$\mathcal{F}_x$ の局所化である。したがって、$\text{Supp}(\mathcal{F})$ の補集合は
一般化について閉じている。
\end{proof}

\noindent
有限型準連接加群については、台は閉であり、ファイバー上で確認でき、
基底変換と可換である。

\begin{lemma}
\label{morphisms-lemma-support-finite-type}
$\mathcal{F}$ をスキーム $X$ 上の有限型準連接加群とする。このとき
\begin{enumerate}
\item $\mathcal{F}$ の台は閉である。
\item $x \in X$ に対して
$$
x \in \text{Supp}(\mathcal{F})
\Leftrightarrow
\mathcal{F}_x \not = 0
\Leftrightarrow
\mathcal{F}_x \otimes_{\mathcal{O}_{X, x}} \kappa(x) \not = 0.
$$
\item 任意のスキームの射 $f : Y \to X$ に対して、引き戻し
$f^*\mathcal{F}$ も有限型であり、
$\text{Supp}(f^*\mathcal{F}) = f^{-1}(\text{Supp}(\mathcal{F}))$ である。
\end{enumerate}
\end{lemma}

\begin{proof}
(1) は加群の補題 \ref{modules-lemma-support-finite-type-closed} の言い換えである。
補題 \ref{morphisms-lemma-support-affine-open}、性質の補題
\ref{properties-lemma-finite-type-module}、および代数の補題
\ref{algebra-lemma-support-closed} を組み合わせても分かる。
(2) の最初の同値は台の定義であり、二番目の同値は中山の補題から従う。
代数の補題 \ref{algebra-lemma-NAK} を参照せよ。
$f : Y \to X$ をスキームの射とする。加群の補題
\ref{modules-lemma-pullback-finite-type} により、$f^*\mathcal{F}$ は有限型である。
最後の主張について、$y \in Y$ とし、その像を $x \in X$ とする。
次を思い出そう。
$$
(f^*\mathcal{F})_y =
\mathcal{F}_x \otimes_{\mathcal{O}_{X, x}} \mathcal{O}_{Y, y},
$$
層の補題 \ref{sheaves-lemma-stalk-pullback-modules} を参照せよ。
したがって、$(f^*\mathcal{F})_y \otimes \kappa(y)$ が零でないことと
$\mathcal{F}_x \otimes \kappa(x)$ が零でないことは同値である。
(2) により、これは $x \in \text{Supp}(\mathcal{F})$ と
$y \in \text{Supp}(f^*\mathcal{F})$ が同値であることを意味し、これが主張 (3) の内容である。
\end{proof}

\begin{lemma}
\label{morphisms-lemma-scheme-theoretic-support}
$\mathcal{F}$ をスキーム $X$ 上の有限型準連接加群とする。
最小の閉部分スキーム $i : Z \to X$ が存在し、ある準連接
$\mathcal{O}_Z$-加群 $\mathcal{G}$ が存在して
$i_*\mathcal{G} \cong \mathcal{F}$ となる。さらに、次が成り立つ。
\begin{enumerate}
\item $\Spec(A) \subset X$ が任意のアフィン開集合で、
$\mathcal{F}|_{\Spec(A)} = \widetilde{M}$ ならば、
$Z \cap \Spec(A) = \Spec(A/I)$ である。ここで $I = \text{Ann}_A(M)$ である。
\item 準連接層 $\mathcal{G}$ は一意な同型を除いて一意である。
\item 準連接層 $\mathcal{G}$ は有限型である。
\item $\mathcal{G}$ および $\mathcal{F}$ の台は $Z$ である。
\end{enumerate}
\end{lemma}

\begin{proof}
$i' : Z' \to X$ を、補題にあるアフィン開集合上の記述を満たす閉部分スキームとする。
すると補題 \ref{morphisms-lemma-i-star-equivalence} により、
$\mathcal{F} \cong i'_*\mathcal{G}'$ が、ある一意な準連接層
$\mathcal{G}'$ について $Z'$ 上で成り立つ。
さらに、$Z'$ がこの性質をもつ最小の閉部分スキームであることは明らかである
（同じ補題による）。最後に、性質の補題
\ref{properties-lemma-finite-type-module} および代数の補題
\ref{algebra-lemma-finite-over-subring} を用いると、$\mathcal{G}'$ が有限型であることが従う。
代数の補題 \ref{algebra-lemma-support-closed} により
$\text{Supp}(\mathcal{G}') = Z$ である。
したがって、補題を証明するには、(1) の特徴づけが実際に閉部分スキームを定めることを
示せば十分である。さらに、そのためには与えられた規則が準連接イデアル層を生み出すことを
示せば十分である。補題 \ref{morphisms-lemma-closed-immersion-bijection-ideals} を参照せよ。
これは次の代数的事実に帰着する。$A$ が環、$f \in A$、
$M$ が有限 $A$-加群ならば、
$\text{Ann}_A(M)_f = \text{Ann}_{A_f}(M_f)$ である。
証明は省略する。
\end{proof}

\begin{definition}
\label{morphisms-definition-scheme-theoretic-support}
$X$ をスキームとする。$\mathcal{F}$ を有限型準連接
$\mathcal{O}_X$-加群とする。$\mathcal{F}$ の {\it スキーム論的な台}とは、
補題 \ref{morphisms-lemma-scheme-theoretic-support} で構成された閉部分スキーム
$Z \subset X$ のことである。
\end{definition}

\noindent
この状況では、$\mathcal{F}$ を $Z$ 上の有限型準連接層とみなすことが多い
（補題 \ref{morphisms-lemma-i-star-equivalence} の圏同値を介して）。












\section{スキーム論的像}
\label{morphisms-section-scheme-theoretic-image}

\noindent
注意：この節の内容の一部はきわめて一般的であり、予想とは異なる挙動を示す。

\begin{lemma}
\label{morphisms-lemma-scheme-theoretic-image}
$f : X \to Y$ をスキームの射とする。閉部分スキーム $Z \subset Y$ であって、
$f$ が $Z$ を経由して因数分解し、かつ任意の他の閉部分スキーム
$Z' \subset Y$ で $f$ が $Z'$ を経由して因数分解するものに対して
$Z \subset Z'$ となるものが存在する。
\end{lemma}

\begin{proof}
$\mathcal{I} = \Ker(\mathcal{O}_Y \to f_*\mathcal{O}_X)$ とおく。
$\mathcal{I}$ が準連接ならば、補題
\ref{morphisms-lemma-closed-immersion-bijection-ideals} を参照して、閉部分スキーム $Z$ を
$\mathcal{I}$ が定めるものとして取ればよい。補題
\ref{morphisms-lemma-closed-immersion-ideals} により、これは機能する。
一般の場合、同じ補題により、最大の準連接イデアル層 $\mathcal{I}'$ で
$\mathcal{I}$ に含まれるものが存在することを示せばよい。これは補題
\ref{morphisms-lemma-largest-quasi-coherent-subsheaf} から従う。
\end{proof}

\begin{definition}
\label{morphisms-definition-scheme-theoretic-image}
$f : X \to Y$ をスキームの射とする。$f$ の {\it スキーム論的像}とは、
最小の閉部分スキーム $Z \subset Y$ であって $f$ がそれを経由して因数分解する
もののことである。
上の補題 \ref{morphisms-lemma-scheme-theoretic-image} を参照せよ。
\end{definition}

\noindent
スキームの射 $f : X \to Y$ とそのスキーム論的像 $Z$ に対し、
$f : X \to Z$ と書いて、$f$ のスキーム論的像を経由する因数分解を表すことが多い。
射 $f$ が準コンパクトでない場合、一般には次が成り立つ。
\begin{enumerate}
\item 集合論的包含 $\overline{f(X)} \subset Z$ は等号でない。すなわち、
$f(X) \subset Z$ は稠密部分集合でない。
\item スキーム論的像の構成は、$Y$ の開部分スキームへの制限と可換でない。
\end{enumerate}
例の章、節 \ref{examples-section-scheme-theoretic-image} には、これら二つの現象の例がある。
これらの現象は埋め込みに対してさえ生じうる。例の章、節
\ref{examples-section-strange-immersion} を参照せよ。
しかし次の補題は、射が準コンパクトならば二つの問題がともに回避されることを示す。

\begin{lemma}
\label{morphisms-lemma-quasi-compact-scheme-theoretic-image}
$f : X \to Y$ をスキームの射とする。$Z \subset Y$ を $f$ のスキーム論的像とする。
$f$ が準コンパクトならば、次が成り立つ。
\begin{enumerate}
\item イデアル層
$\mathcal{I} = \Ker(\mathcal{O}_Y \to f_*\mathcal{O}_X)$
は準連接である。
\item スキーム論的像 $Z$ は $\mathcal{I}$ が定める閉部分スキームである。
\item 任意の開集合 $U \subset Y$ に対し、
$f|_{f^{-1}(U)} : f^{-1}(U) \to U$ のスキーム論的像は $Z \cap U$ に等しい。
\item 像 $f(X) \subset Z$ は $Z$ の稠密部分集合である。言い換えれば、射
$X \to Z$ は支配的である（定義 \ref{morphisms-definition-dominant} を参照）。
\end{enumerate}
\end{lemma}

\begin{proof}
(4) は (3) から従う。$\mathcal{I}$ の形成は $Y$ の開部分スキームへの制限と
可換なので、(3) を示すには (1) を示せば十分である。また、(1) が成り立てば、補題
\ref{morphisms-lemma-scheme-theoretic-image} の証明で (2) を示した。
従って、$\mathcal{I}$ が準連接であることを証明すれば十分である。
準連接性は局所的な性質なので、$Y$ はアフィンであると仮定してよい。
$f$ は準コンパクトだから、有限アフィン開被覆
$X = \bigcup_{i = 1, \ldots, n} U_i$ を取れる。$f'$ を合成
$$
X' = \coprod U_i \longrightarrow X \longrightarrow Y.
$$
とする。このとき $f_*\mathcal{O}_X$ は $f'_*\mathcal{O}_{X'}$ の部分層であり、従って
$\mathcal{I} = \Ker(\mathcal{O}_Y \to f'_*\mathcal{O}_{X'})$ である。
スキームの章、補題 \ref{schemes-lemma-push-forward-quasi-coherent} により、
層 $f'_*\mathcal{O}_{X'}$ は $Y$ 上準連接である。従って主張が従う。
\end{proof}

\begin{example}
\label{morphisms-example-scheme-theoretic-image}
$A \to B$ が核 $I$ をもつ環準同型ならば、
$\Spec(B) \to \Spec(A)$ のスキーム論的像は $\Spec(A/I)$、すなわち
$\Spec(A)$ の閉部分スキームである。これは補題
\ref{morphisms-lemma-quasi-compact-scheme-theoretic-image} から従う。
\end{example}

\noindent
射が準コンパクトならば、そのスキーム論的像は、像の点の特殊化である点のみを付け加える。

\begin{lemma}
\label{morphisms-lemma-reach-points-scheme-theoretic-image}
$f : X \to Y$ を準コンパクト射とする。$Z$ を $f$ のスキーム論的像とする。
$z \in Z$\footnote{補題
\ref{morphisms-lemma-quasi-compact-scheme-theoretic-image} により、集合論的には
$Z$ は $f(X)$ の $Y$ における閉包に一致する。} とする。
付値環 $A$ とその分数体 $K$、および可換図式
$$
\xymatrix{
\Spec(K) \ar[rr] \ar[d] & & X \ar[d] \ar[ld] \\
\Spec(A) \ar[r] & Z \ar[r] & Y
}
$$
であって、$\Spec(A)$ の閉点が $z$ に写るものが存在する。特に、
$Z$ の任意の点は $f(X)$ の点の特殊化である。
\end{lemma}

\begin{proof}
$z \in \Spec(R) = V \subset Y$ を $z$ のアフィン開近傍とする。
補題 \ref{morphisms-lemma-quasi-compact-scheme-theoretic-image} により、
$Z \cap V$ は $f^{-1}(V) \to V$ のスキーム論的像である。従って、$Y$ を $V$ で
置き換えて $Y = \Spec(R)$ はアフィンであると仮定してよい。
この場合、$X$ は $f$ が準コンパクトだから準コンパクトである。
$X = U_1 \cup \ldots \cup U_n$ を有限アフィン開被覆とし、
$U_i = \Spec(A_i)$ と書く。$I = \Ker(R \to A_1 \times \ldots \times A_n)$ とおく。
再び補題 \ref{morphisms-lemma-quasi-compact-scheme-theoretic-image} により、$Z$ は
閉部分スキーム $\Spec(R/I)$（$Y$ のもの）に対応する。
$\mathfrak p \subset R$ を $z$ に対応する素イデアルとすると、
$I_{\mathfrak p} \subset R_{\mathfrak p}$ は等号でない。
従って、局所化は完全なので（代数の章、命題
\ref{algebra-proposition-localization-exact} を参照）、写像
$R_{\mathfrak p} \to
(A_1)_{\mathfrak p} \times \ldots \times (A_n)_{\mathfrak p}$
は零でない。従って環 $(A_i)_{\mathfrak p}$ の少なくとも一つは零でない。
よって、ある $i$ と素イデアル $\mathfrak q_i \subset A_i$ が存在し、これはある
素イデアル $\mathfrak p_i \subset \mathfrak p$ の上にある。
代数の章、補題 \ref{algebra-lemma-dominate} により、付値環
$A \subset K = \kappa(\mathfrak q_i)$ であって、局所環
$R_{\mathfrak p}/\mathfrak p_iR_{\mathfrak p} \subset \kappa(\mathfrak q_i)$
を支配するものを選べる。
これにより所望の図式が得られる。いくつかの詳細は省略する。
\end{proof}

\begin{lemma}
\label{morphisms-lemma-factor-factor}
$$
\xymatrix{
X_1 \ar[d] \ar[r]_{f_1} & Y_1 \ar[d] \\
X_2 \ar[r]^{f_2} & Y_2
}
$$
をスキームの可換図式とする。$Z_i \subset Y_i$, $i = 1, 2$ を $f_i$ の
スキーム論的像とする。このとき射 $Y_1 \to Y_2$ は射 $Z_1 \to Z_2$ および
可換図式
$$
\xymatrix{
X_1 \ar[r] \ar[d] & Z_1 \ar[d] \ar[r] & Y_1 \ar[d] \\
X_2 \ar[r] & Z_2 \ar[r] & Y_2
}
$$
を誘導する。
\end{lemma}

\begin{proof}
$Z_2$ の $Y_1$ におけるスキーム論的逆像は、$Y_1$ の閉部分スキームであり、
$f_1$ はそれを経由して因数分解する。
従って $Z_1$ はこれに含まれる。これで補題が従う。
\end{proof}

\begin{lemma}
\label{morphisms-lemma-scheme-theoretic-image-reduced}
$f : X \to Y$ をスキームの射とする。$X$ が被約ならば、$f$ のスキーム論的像は
$\overline{f(X)}$ 上の被約誘導スキーム構造である。
\end{lemma}

\begin{proof}
$\overline{f(X)}$ 上の被約誘導スキーム構造は、明らかに $Y$ の最小の閉部分スキームであって
$f$ がそれを経由して因数分解するものだからである。スキームの章、補題
\ref{schemes-lemma-map-into-reduction} を参照せよ。
\end{proof}

\begin{lemma}
\label{morphisms-lemma-scheme-theoretic-image-of-partial-section}
$f : X \to Y$ をスキームの分離的射とする。$V \subset Y$ を逆コンパクトな開集合とし、
$s : V \to X$ を $f \circ s = \text{id}_V$ を満たす射とする。
$Y'$ を $s$ のスキーム論的像とする。このとき $Y' \to Y$ は $V$ 上同型である。
\end{lemma}

\begin{proof}
$V$ が $Y$ において逆コンパクトであるという仮定
（位相の章、定義 \ref{topology-definition-quasi-compact}）は、
$V \to Y$ が準コンパクト射であることを意味する。
スキームの章、補題 \ref{schemes-lemma-quasi-compact-permanence} により、射
$s : V \to X$ は準コンパクトである。従って、補題
\ref{morphisms-lemma-quasi-compact-scheme-theoretic-image} により、スキーム論的像
$Y'$（$s$ のもの）の構成は開集合への制限と可換である。特に、
$Y' \cap f^{-1}(V)$ は、分離的射 $f^{-1}(V) \to V$ の切断のスキーム論的像である。
分離的射の切断は閉埋め込みなので
（スキームの章、補題 \ref{schemes-lemma-section-immersion}）、所望の通り
$Y' \cap f^{-1}(V) \to V$ は同型である。
\end{proof}








\section{スキーム論的閉包と稠密性}
\label{morphisms-section-scheme-theoretic-closure}

\noindent
次の定義は \cite[IV, Definition 11.10.2]{EGA} から取る。

\begin{definition}
\label{morphisms-definition-scheme-theoretically-dense}
$X$ をスキームとし、$U \subset X$ を開部分スキームとする。
\begin{enumerate}
\item 射 $U \to X$ のスキーム論的像を {\it $U$ の $X$ におけるスキーム論的閉包}という。
\item $U$ が {\it $X$ においてスキーム論的に稠密}であるとは、任意の開集合
$V \subset X$ に対し、$U \cap V$ の $V$ におけるスキーム論的閉包が $V$ に
等しいことをいう。
\end{enumerate}
\end{definition}

\noindent
この定義では、$U$ が $X$ においてスキーム論的に稠密であることと、
$U$ のスキーム論的閉包が $X$ であることは同値では {\bf ない}。
例 \ref{morphisms-example-scheme-theretically-dense-not-dense} を参照せよ。
これは多少不格好であるが、下の補題
\ref{morphisms-lemma-scheme-theoretically-dense-quasi-compact} および
\ref{morphisms-lemma-reduced-scheme-theoretically-dense} も参照せよ。
一方、この定義では、$U$ が $X$ においてスキーム論的に稠密であることと、
任意の開集合 $V \subset X$ に対して環準同型
$\mathcal{O}_X(V) \to \mathcal{O}_X(U \cap V)$ が単射であることは同値である。
下の補題 \ref{morphisms-lemma-characterize-scheme-theoretically-dense} を参照せよ。
特に、スキーム論的に稠密ならば稠密であることが分かり、これは望ましい。

\begin{example}
\label{morphisms-example-scheme-theretically-dense-not-dense}
スキーム論的閉包が $X$ であっても台位相空間において稠密とは限らない例を挙げる。
$k$ を体とする。
$A = k[x, z_1, z_2, \ldots]/(x^n z_n)$ とおき、
$I = (z_1, z_2, \ldots) \subset A$ とおく。
アフィンスキーム $X = \Spec(A)$ と開部分スキーム $U = X \setminus V(I)$ を考える。
$A \to \prod_n A_{z_n}$ は単射なので、$U$ のスキーム論的閉包は $X$ である。
射 $X \to \Spec(k[x])$ を考える。この射は全射である
（すべての $z_n = 0$ とおけば分かる）。しかし、この射を $U$ に制限したものは、
$x = 0$ という点に写るため全射でない。従って $U$ は $X$ において位相的に稠密ではない。
\end{example}

\begin{lemma}
\label{morphisms-lemma-scheme-theoretically-dense-quasi-compact}
$X$ をスキームとし、$U \subset X$ を開部分スキームとする。
包含射 $U \to X$ が準コンパクトならば、$U$ が $X$ においてスキーム論的に稠密であることと、
$U$ の $X$ におけるスキーム論的閉包が $X$ であることは同値である。
\end{lemma}

\begin{proof}
補題 \ref{morphisms-lemma-quasi-compact-scheme-theoretic-image} の (3) から従う。
\end{proof}

\begin{example}
\label{morphisms-example-scheme-theoretic-closure}
$A$ を環とし、$X = \Spec(A)$ とする。
$f_1, \ldots, f_n \in A$ とし、$U = D(f_1) \cup \ldots \cup D(f_n)$ とおく。
$I = \Ker(A \to \prod A_{f_i})$ とおく。このとき $U$ の $X$ における
スキーム論的閉包は $\Spec(A/I)$、すなわち $X$ の閉部分スキームである。
$U \to X$ は準コンパクトであることに注意せよ。従って補題
\ref{morphisms-lemma-scheme-theoretically-dense-quasi-compact} により、$U$ が $X$ において
スキーム論的に稠密であることと $I = 0$ は同値である。
\end{example}

\begin{lemma}
\label{morphisms-lemma-characterize-scheme-theoretically-dense}
$j : U \to X$ をスキームの開埋め込みとする。このとき、$U$ が $X$ において
スキーム論的に稠密であることと、$\mathcal{O}_X \to j_*\mathcal{O}_U$ が単射であることは同値である。
\end{lemma}

\begin{proof}
$\mathcal{O}_X \to j_*\mathcal{O}_U$ が単射ならば、任意の開集合 $V$（$X$ のもの）への
制限も単射である。従って $U \cap V$ の $V$ におけるスキーム論的閉包は $V$ に等しい。
補題 \ref{morphisms-lemma-scheme-theoretic-image} の証明を参照せよ。
逆に、$U \cap V$ のスキーム論的閉包が $V$ に等しいことを、すべての開集合 $V$ に
対して仮定する。
$\mathcal{O}_X \to j_*\mathcal{O}_U$ が単射でないと仮定する。
すると、あるアフィン開集合 $\Spec(A) = V \subset X$ と非零元 $f \in A$ であって、
$f$ が $\Gamma(V \cap U, \mathcal{O}_X)$ において零に写るものを取れる。
この場合、$V \cap U$ の $V$ におけるスキーム論的閉包は
$\Spec(A/(f))$ に明らかに含まれ、矛盾である。
\end{proof}

\begin{lemma}
\label{morphisms-lemma-intersection-scheme-theoretically-dense}
$X$ をスキームとする。$U$, $V$ が $X$ のスキーム論的に稠密な開部分スキームならば、
$U \cap V$ もそうである。
\end{lemma}

\begin{proof}
$W \subset X$ を任意の開集合とする。写像
$\mathcal{O}_X(W) \to \mathcal{O}_X(W \cap V)
\to \mathcal{O}_X(W \cap V \cap U)$
を考える。補題 \ref{morphisms-lemma-characterize-scheme-theoretically-dense} により、
両方の写像は単射である。従って合成も単射である。
ゆえに再び補題 \ref{morphisms-lemma-characterize-scheme-theoretically-dense} により、
$U \cap V$ は $X$ においてスキーム論的に稠密である。
\end{proof}

\begin{lemma}
\label{morphisms-lemma-quasi-compact-immersion}
$h : Z \to X$ を埋め込みとする。$h$ が準コンパクトであるか、または $Z$ が被約であると仮定する。
$\overline{Z} \subset X$ を $h$ のスキーム論的像とする。このとき射
$Z \to \overline{Z}$ は開埋め込みであり、$Z$ を $\overline{Z}$ のスキーム論的に稠密な
開部分スキームと同一視する。さらに、$Z$ は $\overline{Z}$ において位相的に稠密である。
\end{lemma}

\begin{proof}
補題 \ref{morphisms-lemma-factor-quasi-compact-immersion} または補題
\ref{morphisms-lemma-factor-reduced-immersion} により、$Z \to X$ を
$Z \to \overline{Z}_1 \to X$ と因数分解できる。ここで $Z \to \overline{Z}_1$ は開で、
$\overline{Z}_1 \to X$ は閉である。一方、$Z \to \overline{Z} \subset X$ を
$Z \to X$ のスキーム論的像とする。このとき $\overline{Z} \subset \overline{Z}_1$ である。
$Z$ は $\overline{Z}_1$ の開部分スキームなので、$Z$ は $\overline{Z}$ の開部分スキームでもある。
$Z$ が被約の場合、$Z \subset \overline{Z}_1$ は、補題
\ref{morphisms-lemma-factor-reduced-immersion} の証明における $\overline{Z}_1$ の構成により
位相的に稠密である。
従って $\overline{Z}_1$ と $\overline{Z}$ の台位相空間は同じである。
ゆえに $\overline{Z} \subset \overline{Z}_1$ は、台位相空間上全単射を誘導する
被約スキームへの閉埋め込みであり、従って同型である。
$Z \to X$ が準コンパクトの場合は次のように論じる。
$Z$ が $\overline{Z}$ においてスキーム論的に稠密であるという主張は、補題
\ref{morphisms-lemma-quasi-compact-scheme-theoretic-image} の (3) から従う。
最後の主張は補題 \ref{morphisms-lemma-quasi-compact-scheme-theoretic-image} の (4) から従う。
\end{proof}

\begin{lemma}
\label{morphisms-lemma-reduced-scheme-theoretically-dense}
$X$ を被約スキームとし、$U \subset X$ を開部分スキームとする。次は同値である。
\begin{enumerate}
\item $U$ は $X$ において位相的に稠密である。
\item $U$ の $X$ におけるスキーム論的閉包は $X$ である。
\item $U$ は $X$ においてスキーム論的に稠密である。
\end{enumerate}
\end{lemma}

\begin{proof}
これは補題 \ref{morphisms-lemma-quasi-compact-immersion} と、$Z$ を $X$ の閉部分スキームとするとき、
その台位相空間が $X$ に等しければスキームとして $X$ に等しくなければならないという事実から従う。
\end{proof}

\begin{lemma}
\label{morphisms-lemma-reduced-subscheme-closure}
$X$ をスキームとし、$U \subset X$ を被約な開部分スキームとする。次は同値である。
\begin{enumerate}
\item $U$ の $X$ におけるスキーム論的閉包は $X$ である。
\item $U$ は $X$ においてスキーム論的に稠密である。
\end{enumerate}
これらが成り立つならば、$X$ は被約スキームである。
\end{lemma}

\begin{proof}
これは補題 \ref{morphisms-lemma-quasi-compact-immersion} と、補題
\ref{morphisms-lemma-scheme-theoretic-image-reduced} により $U$ の $X$ における
スキーム論的閉包が被約であるという事実から従う。
\end{proof}


\begin{lemma}
\label{morphisms-lemma-equality-of-morphisms}
$S$ をスキームとし、$X$, $Y$ を $S$ 上のスキームとする。
$f, g : X \to Y$ を $S$ 上のスキームの射とする。
$U \subset X$ を $f|_U = g|_U$ を満たす開部分スキームとする。
$U$ の $X$ におけるスキーム論的閉包が $X$ であり、$Y \to S$ が分離的ならば、
$f = g$ である。
\end{lemma}

\begin{proof}
定義とスキームの章、補題 \ref{schemes-lemma-where-are-they-equal} から従う。
\end{proof}




\section{支配的射}
\label{morphisms-section-dominant}

\noindent
スキームの射が支配的であるという定義は、多様体の支配的射という概念に
慣れている場合に予想するものとは少し異なる。

\begin{definition}
\label{morphisms-definition-dominant}
スキームの射 $f : X \to S$ は、$f$ の像が $S$ の稠密部分集合であるとき
{\it 支配的}であるという。
\end{definition}

\noindent
例えば、$k$ が無限体で、$\lambda_1, \lambda_2, \ldots$ が $k$ の相異なる元からなる
可算集合ならば、射
$$
\coprod\nolimits_{i = 1, 2, \ldots } \Spec(k)
\longrightarrow
\Spec(k[x])
$$
であって、第 $i$ 因子を点 $x = \lambda_i$ に写すものは支配的である。

\begin{lemma}
\label{morphisms-lemma-generic-points-in-image-dominant}
$f : X \to S$ をスキームの射とする。$S$ の各既約成分のすべての生成点が
$f$ の像に属するならば、$f$ は支配的である。
\end{lemma}

\begin{proof}
これは、スキームの台位相空間が sober であるという事実
（スキームの章、補題 \ref{schemes-lemma-scheme-sober}）と、$S$ のすべての点が
$S$ のある既約成分に含まれるという事実
（位相の章、補題 \ref{topology-lemma-irreducible}）から直ちに従う位相的事実である。
\end{proof}

\noindent
射が支配的ならば標的の生成点が像に属するはずだという予想は、射が準コンパクトならば成り立つ。

\begin{lemma}
\label{morphisms-lemma-quasi-compact-dominant}
\begin{slogan}
像が生成点を含む射は支配的である
\end{slogan}
$f : X \to S$ をスキームの準コンパクト射とする。このとき $f$ が支配的であることと、
任意の既約成分 $Z \subset S$ に対して $Z$ の生成点が $f$ の像に属することは同値である。
\end{lemma}

\begin{proof}
$V \subset S$ をアフィン開集合とする。$f$ は準コンパクトなので、有限個のアフィン開集合
$U_i \subset f^{-1}(V)$、$i = 1, \ldots, n$ で $f^{-1}(V)$ を被覆するものを選べる。
アフィンスキームの射
$$
f' :
\coprod\nolimits_{i = 1, \ldots, n} U_i
\longrightarrow
V.
$$
を考える。アフィンスキームの非交和はアフィンである。スキームの章、補題
\ref{schemes-lemma-disjoint-union-affines} を参照せよ。
$V$ の既約成分の生成点は、$S$ の既約成分で $V$ と交わるものの生成点にちょうど一致する。
また、$f$ が支配的であることと $f'$ が支配的であることは、上で
$V, n, U_i$ をどのように選んでも同値である。従って補題はアフィンスキームの射の場合に帰着した。
アフィンの場合は代数の章、補題 \ref{algebra-lemma-image-dense-generic-points} である。
\end{proof}

\begin{lemma}
\label{morphisms-lemma-base-change-dominant}
$f : X \to S$ をスキームの準コンパクトかつ支配的な射とする。
$g : S' \to S$ をスキームの射とし、$f' : X' \to S'$ を $f$ の $g$ による基底変換とする。
$g$ に沿って一般化が持ち上がるならば、$f'$ は支配的である。
\end{lemma}

\noindent
$g$ が開ならば（補題 \ref{morphisms-lemma-open-generizing}）、または平坦ならば
（補題 \ref{morphisms-lemma-generalizations-lift-flat}）、一般化は $g$ に沿って持ち上がることを後で見る。

\begin{proof}
スキームの章、補題 \ref{schemes-lemma-quasi-compact-preserved-base-change} により、
$f'$ は準コンパクトである。$\eta' \in S'$ を $S'$ の既約成分の生成点とする。
$g$ に沿って一般化が持ち上がるならば、$\eta = g(\eta')$ は $S$ の既約成分の生成点である。
補題 \ref{morphisms-lemma-quasi-compact-dominant} により、$\eta$ は $f$ の像に属する。
従ってスキームの章、補題 \ref{schemes-lemma-points-fibre-product} により、
$\eta'$ は $f'$ の像に属する。補題 \ref{morphisms-lemma-quasi-compact-dominant} により、
$f'$ は支配的である。
\end{proof}

\noindent
この主張の一部は、任意の（必ずしも準コンパクトでない）支配的射に対しても成り立つ。

\begin{lemma}
\label{morphisms-lemma-open-base-change-dominant}
$f : X \to S$ をスキームの支配的射とする。
$g : S' \to S$ をスキームの開射とし、$f' : X' \to S'$ を $f$ の $g$ による基底変換とする。
このとき $f'$ は支配的である。
\end{lemma}

\begin{proof}
$f'$ が支配的でないならば、$f'(X')$ はある閉部分集合 $B$ で $S'$ と等しくないものに含まれる。
すると $S'-B$ は $S'$ の空でない開集合である。$g$ は開なので、$g(S'-B)$ は
$S$ の空でない開集合である。$f(X)$ は $S$ において稠密なので、点 $x\in X$ で
$f(x) \in g(S'-B)$ となるものが存在する。
$W = \Spec(\kappa (f(x)))$ とおき、$X_W$, $S_W$, $S'_W$ で、これらのスキームを
$S$ から $W$ へ引き戻したものを表す。このとき $X_W$ と $(S'-B)_W$ は空でない。
従って $X_W \times_W (S'-B)_W$ も空でない。なぜなら $W$ は体のスペクトルだからである。
これは $f'(X \times_S S')$ が $B \subset S'$ に含まれることと矛盾する。
従って実際には $f'$ は支配的である。
\end{proof}

\begin{lemma}
\label{morphisms-lemma-quasi-compact-generic-point-not-in-image}
$f : X \to S$ をスキームの準コンパクト射とする。
$\eta \in S$ を $S$ の既約成分の生成点とする。$\eta \not \in f(X)$ ならば、
開集合 $V \subset S$ であって $\eta$ の近傍であり、$f^{-1}(V) = \emptyset$ となるものが存在する。
\end{lemma}

\begin{proof}
$Z \subset S$ を $f$ のスキーム論的像とする。$\eta \not \in Z$ を示さなければならない。
これは補題 \ref{morphisms-lemma-reach-points-scheme-theoretic-image} から従うが、次のようにも分かる。
補題 \ref{morphisms-lemma-quasi-compact-scheme-theoretic-image} により、射 $X \to Z$ は支配的である。
補題 \ref{morphisms-lemma-quasi-compact-dominant} によれば、これは $Z$ のすべての既約成分の
すべての生成点が $X \to Z$ の像に属することを意味する。
仮定により $\eta \not \in Z$ である。実際、もし $\eta$ が $Z$ に属すれば、
$Z$ のある既約成分の生成点となる。
\end{proof}

\noindent
支配的であることが、すべての既約成分の生成点が像に属することと同値になる別の場合がある。

\begin{lemma}
\label{morphisms-lemma-dominant-finite-number-irreducible-components}
$f : X \to S$ をスキームの射とする。$X$ が有限個の既約成分を持つと仮定する。
このとき、$f$ が支配的であることと、任意の既約成分 $Z \subset S$ に対して
$Z$ の生成点が $f$ の像に属することは同値である。これらが成り立つならば、
$S$ も有限個の既約成分を持つ。
\end{lemma}

\begin{proof}
$f$ が支配的であると仮定する。$X = Z_1 \cup Z_2 \cup \ldots \cup Z_n$ を
$X$ の既約成分への分解とする。$\xi_i \in Z_i$ をその生成点とすれば、
$Z_i = \overline{\{\xi_i\}}$ である。$f(Z_i)$ は $S$ の既約部分集合であることに注意せよ。
従って
$$
S = \overline{f(X)} = \bigcup \overline{f(Z_i)} =
\bigcup \overline{\{f(\xi_i)\}}
$$
は生成点が $f$ の像に属する既約部分集合の有限合併である。補題が従う。
\end{proof}

\begin{lemma}
\label{morphisms-lemma-dominant-between-integral}
$f : X \to Y$ を二つの整スキームの間の射とする。次は同値である。
\begin{enumerate}
\item $f$ は支配的である。
\item $f$ は $X$ の生成点を $Y$ の生成点に写す。
\item ある空でないアフィン開集合 $U \subset X$ および $V \subset Y$ で
$f(U) \subset V$ を満たすものについて、環準同型 $\mathcal{O}_Y(V) \to \mathcal{O}_X(U)$ は単射である。
\item すべての空でないアフィン開集合 $U \subset X$ および $V \subset Y$ で
$f(U) \subset V$ を満たすものについて、環準同型 $\mathcal{O}_Y(V) \to \mathcal{O}_X(U)$ は単射である。
\item ある $x \in X$ とその像 $y = f(x) \in Y$ に対して、局所環準同型
$\mathcal{O}_{Y, y} \to \mathcal{O}_{X, x}$ は単射である。
\item すべての $x \in X$ とその像 $y = f(x) \in Y$ に対して、局所環準同型
$\mathcal{O}_{Y, y} \to \mathcal{O}_{X, x}$ は単射である。
\end{enumerate}
\end{lemma}

\begin{proof}
(1) と (2) の同値性は補題
\ref{morphisms-lemma-dominant-finite-number-irreducible-components} から従う。
$U \subset X$ および $V \subset Y$ を $f(U) \subset V$ を満たす空でないアフィン開集合とする。
環 $A = \mathcal{O}_X(U)$ および $B = \mathcal{O}_Y(V)$ は整域であることを思い出そう。
射 $f|_U : U \to V$ は環準同型 $\varphi : B \to A$ に対応する。
$X$ と $Y$ の生成点は、素イデアル $(0) \subset A$ と $(0) \subset B$ に対応する。
従って (2) は条件 $(0) = \varphi^{-1}((0))$、すなわち $\varphi$ が単射であるという条件と同値である。
このようにして (2)、(3)、(4) が同値であることが分かる。
同様に、(5) にあるような $x$ と $y$ が与えられると、局所環
$\mathcal{O}_{X, x}$ と $\mathcal{O}_{Y, y}$ は整域であり、素イデアル
$(0) \subset \mathcal{O}_{X, x}$ と $(0) \subset \mathcal{O}_{Y, y}$ は
$X$ と $Y$ の生成点に対応する（$x$ における局所環のスペクトルを、
$x$ に特殊化する点の集合と同一視することによる。スキームの章、補題
\ref{schemes-lemma-specialize-points} を参照せよ）。
従って上と全く同じように論じて、(2)、(5)、(6) が同値であることが分かる。
\end{proof}







\section{全射}
\label{morphisms-section-surjective}

\begin{definition}
\label{morphisms-definition-surjective}
スキームの射が台位相空間上全射であるとき、その射を {\it 全射}という。
\end{definition}

\begin{lemma}
\label{morphisms-lemma-composition-surjective}
全射の合成は全射である。
\end{lemma}

\begin{proof}
省略する。
\end{proof}

\begin{lemma}
\label{morphisms-lemma-when-point-maps-to-pair}
$X$ と $Y$ を基礎スキーム $S$ 上のスキームとする。点 $x \in X$ と $y \in Y$ が
与えられたとき、$X \times_S Y$ のある点が射影のもとで $x$ と $y$ に写ることと、
$x$ と $y$ が $S$ の同じ点の上にあることは同値である。
\end{lemma}

\begin{proof}
条件が必要であることは明らかであり、逆はスキームの章、補題
\ref{schemes-lemma-points-fibre-product} の証明から従う。
\end{proof}

\begin{lemma}
\label{morphisms-lemma-base-change-surjective}
全射の基底変換は全射である。
\end{lemma}

\begin{proof}
$f: X \to Y$ を基礎スキーム $S$ 上のスキームの射とする。
$S' \to S$ がスキームの射ならば、$p: X_{S'} \to X$ および
$q: Y_{S'} \to Y$ を標準射影とする。可換正方形
$$
\xymatrix{
X_{S'} \ar[d]_{f_{S'}} \ar[r]_p & X \ar[d]^{f} \\
Y_{S'} \ar[r]^{q} & Y.
}
$$
は、$X_{S'}$ を $X \to Y$ と $Y_{S'} \to Y$ のファイバー積と同一視する。
$Z$ を $X$ の台位相空間の部分集合とする。このとき
$q^{-1}(f(Z)) = f_{S'}(p^{-1}(Z))$ である。実際、
$y' \in q^{-1}(f(Z))$ であることと、$q(y') = f(x)$ となるある $x \in Z$ が存在することは同値であり、
さらに補題 \ref{morphisms-lemma-when-point-maps-to-pair} により、ある $x' \in X_{S'}$ が存在して
$f_{S'}(x') = y'$ および $p(x') = x$ となることと同値である。特に、$Z = X$ と取れば、
$f$ が全射ならば基底変換 $f_{S'}: X_{S'} \to Y_{S'}$ も全射であることが分かる。
\end{proof}

\begin{example}
\label{morphisms-example-injective-not-preserved-base-change}
全単射性は基底変換で安定でなく、従って単射性も安定でない。
例えば全単射
$\Spec(\mathbf{C}) \to \Spec(\mathbf{R})$
を考える。基底変換
$\Spec(\mathbf{C} \otimes_{\mathbf{R}} \mathbf{C}) \to
\Spec(\mathbf{C})$
は単射でない。実際、同型
$\mathbf{C} \otimes_{\mathbf{R}} \mathbf{C} \cong \mathbf{C} \times \mathbf{C}$
が存在し（この分解は冪等元
$\frac{1 \otimes 1 + i \otimes i}{2}$ から生じる）、従って
$\Spec(\mathbf{C} \otimes_{\mathbf{R}} \mathbf{C})$ は二点を持つ。
\end{example}

\begin{lemma}
\label{morphisms-lemma-surjection-from-quasi-compact}
$$
\xymatrix{
X \ar[rr]_f \ar[rd]_p & &
Y \ar[dl]^q \\
& Z
}
$$
をスキームの射の可換図式とする。$f$ が全射で $p$ が準コンパクトならば、
$q$ は準コンパクトである。
\end{lemma}

\begin{proof}
$W \subset Z$ を準コンパクト開集合とする。仮定により $p^{-1}(W)$ は準コンパクトである。
従って位相の章、補題 \ref{topology-lemma-image-quasi-compact} により、逆像
$q^{-1}(W) = f(p^{-1}(W))$ も準コンパクトである。これで補題が従う。
\end{proof}




\section{純非分離射と普遍的単射}
\label{morphisms-section-radicial}

\noindent
本節では、スキームの射が {\it 純非分離}であることと、スキームの射が
{\it 普遍的単射}であることを定義する。次いで、これらの概念が一致することを示す。
両方を導入する理由は、代数空間の場合には対応する概念があり、それらが必ずしも
常に一致するとは限らないからである。

\begin{definition}
\label{morphisms-definition-universally-injective}
$f : X \to S$ を射とする。
\begin{enumerate}
\item $f$ が {\it 普遍的単射}であるとは、スキームの任意の射 $S' \to S$ に対して基底変換
$f' : X_{S'} \to S'$ が（台位相空間上）単射であることをいう。
\item $f$ が {\it 純非分離}であるとは、$f$ が位相空間の写像として単射であり、かつ
すべての $x \in X$ に対して体拡大 $\kappa(x)/\kappa(f(x))$ が純非分離であることをいう。
\end{enumerate}
\end{definition}

\begin{lemma}
\label{morphisms-lemma-universally-injective}
\begin{reference}
\cite[Proposition 3.7.1]{EGA1-second}
\end{reference}
$f : X \to S$ をスキームの射とする。次は同値である。
\begin{enumerate}
\item すべての体 $K$ に対して、誘導される写像
$\Mor(\Spec(K), X) \to \Mor(\Spec(K), S)$
は単射である。
\item 射 $f$ は普遍的単射である。
\item 射 $f$ は純非分離である。
\item 対角射 $\Delta_{X/S} : X \longrightarrow X \times_S X$ は全射である。
\end{enumerate}
\end{lemma}

\begin{proof}
$K$ を体とし、$s : \Spec(K) \to S$ を射とする。
$x : \Spec(K) \to X$ という射で $f \circ x = s$ を満たすものを与えることは、射影
$X_K = \Spec(K) \times_S X \to \Spec(K)$ の切断を与えることと同じであり、さらに
点 $x \in X_K$ で剰余体が $K$ であるものを与えることと同じである。
従って (2) は (1) を含意する。

\medskip\noindent
逆に、(1) が成り立つと仮定する。$x, x' \in X_{S'}$ が同じ点
$s' \in S'$ に写ると仮定する。体の可換図式
$$
\xymatrix{
K & \kappa(x) \ar[l] \\
\kappa(x') \ar[u] & \kappa(s') \ar[l] \ar[u]
}
$$
を選ぶ。スキームの章、補題 \ref{schemes-lemma-characterize-points} により、二つの射
$a, a' : \Spec(K) \to X_{S'}$ を得る。一方は点 $x$ と埋め込み
$\kappa(x) \subset K$ に対応し、他方は点 $x'$ と埋め込み
$\kappa(x') \subset K$ に対応する。また $f' \circ a = f' \circ a'$ である。
条件 (1) により、$a$ と $a'$ を $X_{S'} \to X$ と合成したものは等しい。
$X_{S'}$ は $S'$ と $X$ の $S$ 上のファイバー積なので、$a = a'$ である。
従って $x = x'$ である。ゆえに (1) は (2) を含意する。

\medskip\noindent
$x, x' \in X$ という相異なる二点が同じ点 $s$ に写るならば、(2) に反する。
ある $s = f(x)$、$x \in X$ に対して体拡大 $\kappa(x)/\kappa(s)$ が純非分離でないならば、
体拡大 $K/\kappa(s)$ であって、$\kappa(x)$ が二つの $\kappa(s)$-準同型を
$K$ へ持つものを取れる。スキームの章、補題
\ref{schemes-lemma-characterize-points} により、これは写像
$\Mor(\Spec(K), X) \to \Mor(\Spec(K), S)$
が単射でないことを含意し、従って (1) に反する。
ゆえに同値な条件 (1) と (2) は $f$ が純非分離であること、すなわち (3) を含意する。

\medskip\noindent
(3) を仮定する。スキームの章、補題 \ref{schemes-lemma-characterize-points} により、射
$\Spec(K) \to X$ は組 $(x, \kappa(x) \to K)$ によって与えられる。
性質 (3) は、組 $(x, \kappa(x) \to K)$ に組
$(s, \kappa(s) \to \kappa(x) \to K)$ を対応させることが単射であると正確に述べる。
言い換えれば (1) が成り立つ。この時点で (1)、(2)、(3) はすべて同値である。

\medskip\noindent
最後に、(4) と (1)、(2)、(3) の同値性を証明する。
$X \times_S X$ の点は四つ組 $(x_1, x_2, s, \mathfrak p)$ によって与えられる。ここで
$x_1, x_2 \in X$、$f(x_1) = f(x_2) = s$ であり、
$\mathfrak p \subset \kappa(x_1) \otimes_{\kappa(s)} \kappa(x_2)$ は素イデアルである。
スキームの章、補題 \ref{schemes-lemma-points-fibre-product} を参照せよ。
$f$ が普遍的単射ならば、その定義において $S'=X$ と取ることにより、
$\Delta_{X/S}$ は全射でなければならない。実際、これは単射
$X \times_S X  \longrightarrow X$ の切断である。
逆に、$\Delta_{X/S}$ が全射ならば、常に $x_1 = x_2 = x$ であり、そのような素イデアル
$\mathfrak p$ はちょうど一つしかない。これは $\kappa(s) \subset \kappa(x)$ が
純非分離であることを意味する。従って $f$ は純非分離である。
別の証明として、$\Delta_{X/S}$ が全射ならば、任意の $S' \to S$ に対して基底変換
$\Delta_{X_{S'}/S'}$ は全射であり、これから $f$ は普遍的単射であることが従う。
これで補題の証明が完了する。
\end{proof}

\begin{lemma}
\label{morphisms-lemma-universally-injective-separated}
普遍的単射な射は分離的である。
\end{lemma}

\begin{proof}
補題 \ref{morphisms-lemma-universally-injective} と、$X \to S$ が分離的であることと
$\Delta_{X/S}$ の像が $X \times_S X$ において閉であることが同値であるという注意を合わせればよい。
スキームの章、定義 \ref{schemes-definition-separated} とそれに続く議論を参照せよ。
\end{proof}

\begin{lemma}
\label{morphisms-lemma-base-change-universally-injective}
普遍的単射な射の基底変換は普遍的単射である。
\end{lemma}

\begin{proof}
これは形式的である。
\end{proof}

\begin{lemma}
\label{morphisms-lemma-composition-universally-injective}
純非分離射の合成は純非分離であり、従って同値な条件である普遍的単射性についても同じことが成り立つ。
\end{lemma}

\begin{proof}
省略する。
\end{proof}








\section{アフィン射}
\label{morphisms-section-affine}

\begin{definition}
\label{morphisms-definition-affine}
スキームの射 $f : X \to S$ は、$S$ の任意のアフィン開集合の逆像が
$X$ のアフィン開集合であるとき {\it アフィン}であるという。
\end{definition}

\begin{lemma}
\label{morphisms-lemma-affine-separated}
アフィン射は分離的かつ準コンパクトである。
\end{lemma}

\begin{proof}
$f : X \to S$ をアフィンとする。準コンパクト性はスキームの章、補題
\ref{schemes-lemma-quasi-compact-affine} から直ちに従う。スキームの章、補題
\ref{schemes-lemma-characterize-separated} を用いて $f$ が分離的であることを示す。
$x_1, x_2 \in X$ を $X$ の点で同じ点 $s \in S$ に写るものとする。
任意のアフィン開集合 $W \subset S$ で $s$ を含むものを選ぶ。仮定により $f^{-1}(W)$ は
アフィンである。引用した補題を $U = V = f^{-1}(W)$ として適用すればよい。
\end{proof}

\begin{lemma}
\label{morphisms-lemma-characterize-affine}
\begin{reference}
\cite[II, Corollary 1.3.2]{EGA}
\end{reference}
$f : X \to S$ をスキームの射とする。次は同値である。
\begin{enumerate}
\item 射 $f$ はアフィンである。
\item アフィン開被覆 $S = \bigcup W_j$ で、各 $f^{-1}(W_j)$ がアフィンとなるものが存在する。
\item 準連接 $\mathcal{O}_S$-代数層 $\mathcal{A}$ と同型
$X \cong \underline{\Spec}_S(\mathcal{A})$（$S$ 上のスキームのもの）が存在する。記法については構成の章、節
\ref{constructions-section-spec} を参照せよ。
\end{enumerate}
さらに、この場合 $X = \underline{\Spec}_S(f_*\mathcal{O}_X)$ である。
\end{lemma}

\begin{proof}
(1) が (2) を含意することは明らかである。

\medskip\noindent
$S = \bigcup_{j \in J} W_j$ を、各 $f^{-1}(W_j)$ がアフィンとなるアフィン開被覆と仮定する。
スキームの章、補題 \ref{schemes-lemma-quasi-compact-affine} により $f$ は準コンパクトである。
スキームの章、補題 \ref{schemes-lemma-characterize-quasi-separated} により、射 $f$ は準分離的である。
従ってスキームの章、補題 \ref{schemes-lemma-push-forward-quasi-coherent} により、層
$\mathcal{A} = f_*\mathcal{O}_X$ は準連接 $\mathcal{O}_S$-代数層である。
ゆえにスキーム $g : Y = \underline{\Spec}_S(\mathcal{A}) \to S$（$S$ 上のもの）が得られる。
恒等写像
$\text{id} : \mathcal{A} = f_*\mathcal{O}_X \to f_*\mathcal{O}_X$
は、相対スペクトルの定義を介して射 $can : X \to Y$（$S$ 上のもの）を与える。
構成の章、補題 \ref{constructions-lemma-canonical-morphism} を参照せよ。
仮定と今引用した補題により、制限
$can|_{f^{-1}(W_j)} : f^{-1}(W_j) \to g^{-1}(W_j)$
は同型である。従って $can$ は同型である。これで (2) が (3) を含意することを示した。

\medskip\noindent
(3) を仮定する。構成の章、補題 \ref{constructions-lemma-spec-properties} により、
任意のアフィン開集合の逆像はアフィンであり、従って定義により射はアフィンである。
\end{proof}

\begin{remark}
\label{morphisms-remark-direct-argument}
上の補題 \ref{morphisms-lemma-characterize-affine} において、(2) が (1) を含意することを
次のように直接論じることもできる。
$S = \bigcup W_j$ を、各 $f^{-1}(W_j)$ がアフィンとなるアフィン開被覆と仮定する。
まず上の証明と同様に $\mathcal{A} = f_*\mathcal{O}_X$ が準連接であることを示す。
$\Spec(R) = V \subset S$ をアフィン開集合とする。$f^{-1}(V)$ がアフィンであることを
示さなければならない。
$A = \mathcal{A}(V) = f_*\mathcal{O}_X(V) = \mathcal{O}_X(f^{-1}(V))$ とおく。
スキームの章、補題 \ref{schemes-lemma-morphism-into-affine} により、標準射
$\psi : f^{-1}(V) \to \Spec(A)$（$\Spec(R) = V$ 上のもの）がある。
スキームの章、補題 \ref{schemes-lemma-good-subcover} により、整数 $n \geq 0$、標準開被覆
$V = \bigcup_{i = 1, \ldots, n} D(h_i)$、$h_i \in R$、および写像
$a : \{1, \ldots, n\} \to J$ であって、各 $D(h_i)$ がアフィンスキーム
$W_{a(i)}$ の標準開集合でもあるものが存在する。アフィンスキームの射による標準開集合の
逆像は標準開集合である。代数の章、補題 \ref{algebra-lemma-spec-functorial} を参照せよ。
従って $f^{-1}(D(h_i))$ は $f^{-1}(W_{a(i)})$ の標準開集合であり、特に
$f^{-1}(D(h_i))$ はアフィンである。$\mathcal{A}$ は準連接なので
$A_{h_i} = \mathcal{A}(D(h_i)) = \mathcal{O}_X(f^{-1}(D(h_i)))$ であり、従って
$f^{-1}(D(h_i))$ は $A_{h_i}$ のスペクトルである。
ゆえに射 $\psi$ は開集合 $f^{-1}(D(h_i))$ と $\Spec(A_{h_i})$（$\Spec(A)$ の開集合）との
同型を誘導する。$f^{-1}(V) = \bigcup f^{-1}(D(h_i))$ かつ
$\Spec(A) = \bigcup \Spec(A_{h_i})$ なので主張が従う。
\end{remark}

\begin{lemma}
\label{morphisms-lemma-affine-equivalence-algebras}
$S$ をスキームとする。圏の反同値
$$
\begin{matrix}
\text{Schemes affine} \\
\text{over }S
\end{matrix}
\longleftrightarrow
\begin{matrix}
\text{quasi-coherent sheaves} \\
\text{of }\mathcal{O}_S\text{-algebras}
\end{matrix}
$$
が存在し、$f : X \to S$ に層 $f_*\mathcal{O}_X$ を対応させる。
さらに、この同値は任意の基底変換と可換である。
\end{lemma}

\begin{proof}
右から左への関手は $\underline{\Spec}_S$ によって与えられる。
補題 \ref{morphisms-lemma-characterize-affine} および構成の章、補題
\ref{constructions-lemma-spec-properties} の (3) により、二つの関手は互いに逆である。
最後の主張は構成の章、補題 \ref{constructions-lemma-spec-properties} の (2) である。
\end{proof}

\begin{lemma}
\label{morphisms-lemma-quasi-coherence-scalar-restriction}
\begin{reference}
\cite[I, Proposition 9.6.1]{EGA}
\end{reference}
$S$ をスキームとし、$\mathcal{A}$ を準連接 $\mathcal{O}_S$-代数とする。
$\mathcal{A}$-加群が $\mathcal{O}_S$-加群として準連接であることと、
$\mathcal{A}$-加群として準連接であることは同値である。
\end{lemma}

\begin{proof}
$\mathcal{F}$ を $\mathcal{A}$-加群とする。$\mathcal{F}$ が $\mathcal{A}$-加群として
準連接ならば、すべての $s \in S$ に対して開近傍 $U$（$s$ のもの）と完全列
$$
\bigoplus\nolimits_J\mathcal{A}|_U\to
\bigoplus\nolimits_I\mathcal{A}|_U\to
\mathcal{F}|_U
$$
（$\mathcal{A}|_U$-加群のもの）が存在する。これは $\mathcal{O}_U$-加群の完全列でもある。
従って $\mathcal{F}|_U$ は、スキーム上の準連接 $\mathcal{O}_U$-加群の射の余核として準連接である。
ゆえに $\mathcal{F}$ は $\mathcal{O}_X$-加群として準連接である。

\medskip\noindent
逆に、$\mathcal{F}$ が $\mathcal{O}_X$-加群として準連接であると仮定する。
アフィン開集合 $\Spec(R) = U \subset S$ を選ぶ。$\mathcal{O}_U$-加群の同型
$\mathcal{A}|_U \cong \widetilde{A}$ および $\mathcal{F}|_U \cong \widetilde{M}$ がある。
これは、ある $R$-代数 $A$ とある $R$-加群 $M$ に対してである。
$\mathcal{A}$-加群構造（$\mathcal{F}$ 上のもの）は、$A$-加群構造（$M$ 上のもの）で、
与えられた $R$-加群構造と両立するものに移る（詳細は省略する）。
$A$-加群の完全列
$$
\bigoplus\nolimits_J A \to
\bigoplus\nolimits_I A \to
M \to 0
$$
を選ぶ。関手 $\widetilde{\ }$ は完全なので、これは $\widetilde{A}$-加群の完全列
$$
\bigoplus\nolimits_J \widetilde{A}\to
\bigoplus\nolimits_I \widetilde{A}\to
\widetilde{M} \to 0
$$
を生じる。これは $\mathcal{F}$ が $\mathcal{A}$-加群として準連接であることを意味する。
\end{proof}

\begin{lemma}
\label{morphisms-lemma-affine-equivalence-modules}
$f : X \to S$ をスキームのアフィン射とする。
$\mathcal{A} = f_*\mathcal{O}_X$ とおく。
関手 $\mathcal{F} \mapsto f_*\mathcal{F}$ は圏同値
$$
\left\{
\begin{matrix}
\text{category of quasi-coherent}\\
\mathcal{O}_X\text{-modules}
\end{matrix}
\right\}
\longrightarrow
\left\{
\begin{matrix}
\text{category of quasi-coherent}\\
\mathcal{A}\text{-modules}
\end{matrix}
\right\}
$$
を誘導する。さらに、$\mathcal{A}$-加群が $\mathcal{O}_S$-加群として準連接であることと、
$\mathcal{A}$-加群として準連接であることは同値である。
\end{lemma}

\begin{proof}
最後の主張は補題 \ref{morphisms-lemma-quasi-coherence-scalar-restriction} である。
補題 \ref{morphisms-lemma-affine-separated} とスキームの章、補題
\ref{schemes-lemma-push-forward-quasi-coherent} により、順像
$f_*\mathcal{F}$（準連接 $\mathcal{O}_X$-加群 $\mathcal{F}$ のもの）は
準連接 $\mathcal{O}_S$-加群である。従って補題の主張にある関手が得られる。

\medskip\noindent
この関手の準逆関手 $h$ を構成する。$\mathcal{G}$ を準連接
$\mathcal{A}$-加群とする。このとき
$$
h(\mathcal{G}) = f^*\mathcal{G} \otimes_{f^*\mathcal{A}} \mathcal{O}_X =
f^*\mathcal{G} \otimes_{f^*f_*\mathcal{O}_X} \mathcal{O}_X
$$
とおく。詳しく述べると、引戻し $f^*\mathcal{A} = f^*f_*\mathcal{O}_X$ は
$\mathcal{O}_X$-代数であり、随伴射
$f^*f_*\mathcal{O}_X \to \mathcal{O}_X$ は代数準同型であり、引戻し
$f^*\mathcal{G}$ は $f^*\mathcal{A}$-加群である。準連接性は引戻しと環の変更で
保たれるので、$h(\mathcal{G})$ は準連接である。随伴射
$f^*f_*\mathcal{F} \to \mathcal{F}$ から関手的な写像
$$
h(f_*\mathcal{F}) =
f^*f_*\mathcal{F} \otimes_{f^*f_*\mathcal{O}_X} \mathcal{O}_X \to \mathcal{F}
$$
が得られ、さらに関手的な写像
$$
\mathcal{G} \to f_*h(\mathcal{G}) \quad\text{adjoint to the map}\quad
f^*\mathcal{G} \to f^*\mathcal{G} \otimes_{f^*f_*\mathcal{O}_X} \mathcal{O}_X
$$
が得られる。これは局所切断 $s$（$f^*\mathcal{F}$ のもの）を $s \otimes 1$ に送る。
証明を終えるには、上の $\mathcal{F}$ と $\mathcal{G}$ に対してこれらの写像が同型であることを
示せば十分である。これはアフィン被覆の各要素上、すなわち $X$ と $S$ がアフィンの場合に
検証すればよい。

\medskip\noindent
この構成を成り立たせる鍵となる代数的観察は次の通りである。環準同型
$R \to A$ を取り、$N$ を $A$-加群とする。このとき
$$
(N \otimes_R A) \otimes_{(A \otimes_R A)} A = N
$$
である。実際、この等式の左辺は、$h$ を準連接 $\widetilde{A}$-加群
$\widetilde{N}$（$\Spec(R)$ 上のもの）に適用した結果である。詳細は省略する。
\end{proof}

\begin{lemma}
\label{morphisms-lemma-composition-affine}
アフィン射の合成はアフィンである。
\end{lemma}

\begin{proof}
$f : X \to Y$ と $g : Y \to Z$ をアフィン射とする。
$U \subset Z$ をアフィン開集合とする。$g^{-1}(U)$ は $g$ に関する仮定によりアフィンである。
すると $f^{-1}(g^{-1}(U))$ は $f$ に関する仮定によりアフィンである。
従って $(g \circ f)^{-1}(U)$ はアフィンである。
\end{proof}

\begin{lemma}
\label{morphisms-lemma-base-change-affine}
アフィン射の基底変換はアフィンである。
\end{lemma}

\begin{proof}
$f : X \to S$ をアフィン射とする。$S' \to S$ を任意の射とする。
$f' : X_{S'} = S' \times_S X \to S'$ で $f$ の基底変換を表す。
任意の $s' \in S'$ に対し、あるアフィン開近傍
$s' \in V \subset S'$ で、あるアフィン開集合 $U \subset S$ の中へ写るものが存在する。
仮定により $f^{-1}(U)$ はアフィンである。スキームの章、節
\ref{schemes-section-fibre-products} の内容により、
$f^{-1}(U)_V = V \times_U f^{-1}(U)$ はアフィンであり、
$(f')^{-1}(V)$ に等しい。従って $S'$ は、その各要素の $f'$ による逆像がアフィンとなる
アフィン開被覆を持つ。上の補題 \ref{morphisms-lemma-characterize-affine} により結論を得る。
\end{proof}

\begin{lemma}
\label{morphisms-lemma-closed-immersion-affine}
閉埋め込みはアフィンである。
\end{lemma}

\begin{proof}
これを最初に示唆するのはスキームの章、補題
\ref{schemes-lemma-closed-immersion-affine-case} である。完全な主張についてはスキームの章、補題
\ref{schemes-lemma-closed-subspace-scheme} を参照せよ。
\end{proof}

\begin{lemma}
\label{morphisms-lemma-affine-s-open}
$X$ をスキームとする。
$\mathcal{L}$ を可逆 $\mathcal{O}_X$-加群とする。
$s \in \Gamma(X, \mathcal{L})$ とする。
包含射 $j : X_s \to X$ はアフィンである。
\end{lemma}

\begin{proof}
これは性質の章、補題 \ref{properties-lemma-affine-cap-s-open} と定義から従う。
\end{proof}

\begin{lemma}
\label{morphisms-lemma-affine-permanence}
$g : X \to Y$ を $S$ 上のスキームの射とする。
\begin{enumerate}
\item $X$ が $S$ 上アフィンであり、$\Delta : Y \to Y \times_S Y$ がアフィンならば、
$g$ はアフィンである。
\item $X$ が $S$ 上アフィンであり、$Y$ が $S$ 上分離的ならば、
$g$ はアフィンである。
\item アフィンスキームから、対角射がアフィンであるスキームへの射はアフィンである。
\item アフィンスキームから分離的スキームへの射はアフィンである。
\end{enumerate}
\end{lemma}

\begin{proof}
(1) を証明する。基底変換 $X \times_S Y \to Y$ は補題
\ref{morphisms-lemma-base-change-affine} によりアフィンである。
射 $(1, g) : X \to X \times_S Y$ は、$Y \to Y \times_S Y$ を射
$X \times_S Y \to Y \times_S Y$ により基底変換したものである。従って補題
\ref{morphisms-lemma-base-change-affine} によりアフィンである。
アフィン射の合成はアフィンである（補題 \ref{morphisms-lemma-composition-affine} を参照）ので、(1) が従う。
(2) は (1) から従う。実際、閉埋め込みはアフィンであり
（補題 \ref{morphisms-lemma-closed-immersion-affine} を参照）、$Y/S$ が分離的であるとは
$\Delta$ が閉埋め込みであることを意味する。(3) と (4) は (1) と (2) の特別な場合である。
\end{proof}

\begin{lemma}
\label{morphisms-lemma-morphism-affines-affine}
アフィンスキーム間の射はアフィンである。
\end{lemma}

\begin{proof}
補題 \ref{morphisms-lemma-affine-permanence} で $S = \Spec(\mathbf{Z})$ とすれば直ちに従う。
補題 \ref{morphisms-lemma-characterize-affine} の (1) と (2) の同値からも直接従う。
\end{proof}

\begin{lemma}
\label{morphisms-lemma-Artinian-affine}
$S$ をスキームとする。
$A$ を Artin 環とする。
任意の射 $\Spec(A) \to S$ はアフィンである。
\end{lemma}

\begin{proof}
省略する。
\end{proof}

\begin{lemma}
\label{morphisms-lemma-get-affine}
$j : Y \to X$ をスキームの埋め込みとする。開集合 $U \subset X$ で、その補集合が
$Z = X \setminus U$ であり、次を満たすものが存在すると仮定する。
\begin{enumerate}
\item $U \to X$ はアフィンである。
\item $j^{-1}(U) \to U$ はアフィンである。
\item $j(Y) \cap Z$ は閉である。
\end{enumerate}
このとき $j$ はアフィンである。特に、$X$ がアフィンならば $Y$ もアフィンである。
\end{lemma}

\begin{proof}
スキームの章、定義 \ref{schemes-definition-immersion} により、開部分スキーム
$W \subset X$ で、$j$ が閉埋め込み $i : Y \to W$ と包含射 $W \to X$ の合成として
分解するものが存在する。閉埋め込みはアフィンなので
（補題 \ref{morphisms-lemma-closed-immersion-affine}）、任意の $x \in W$ に対して、
$x$ の $X$ におけるアフィン開近傍で、その $j$ による逆像がアフィンとなるものが存在する。
$x \in U$ ならば、仮定 (2) により同じことが成り立つ。
最後に、$x \in Z$ かつ $x \not \in W$ と仮定する。このとき
$x \not \in j(Y) \cap Z$ である。仮定 (3) により、アフィン開近傍
$V \subset X$（$x$ のもの）で $j(Y) \cap Z$ と交わらないものを選べる。このとき
$j^{-1}(V) = j^{-1}(V \cap U)$ は仮定 (1) と (2) によりアフィンである。
従って補題 \ref{morphisms-lemma-characterize-affine} により $j$ はアフィンである。
\end{proof}

\section{豊富な可逆層の族}
\label{morphisms-section-families-ample-invertible-modules}

\noindent
豊富な可逆層の族という概念についての短い節である。

\begin{definition}
\label{morphisms-definition-family-ample-invertible-modules}
\begin{reference}
\cite[II Definition 2.2.4]{SGA6}
\end{reference}
$X$ をスキームとする。$\{\mathcal{L}_i\}_{i \in I}$ を可逆
$\mathcal{O}_X$-加群の族とする。$\{\mathcal{L}_i\}_{i \in I}$ が
次を満たすとき、これを {\it $X$ 上の豊富な可逆層の族}という。
\begin{enumerate}
\item $X$ は準コンパクトである。
\item 任意の $x \in X$ に対して、ある $i \in I$、ある $n \geq 1$、および
$s \in \Gamma(X, \mathcal{L}_i^{\otimes n})$ で、$x \in X_s$ かつ
$X_s$ がアフィンとなるものが存在する。
\end{enumerate}
\end{definition}

\noindent
$\{\mathcal{L}_i\}_{i \in I}$ がスキーム $X$ 上の豊富な可逆層の族ならば、
有限部分集合 $I' \subset I$ で、$\{\mathcal{L}_i\}_{i \in I'}$ が
$X$ 上の豊富な可逆層の族となるものが存在する（準コンパクト性から直ちに従う）。
豊富な可逆層の族を持つスキームは、次の補題によりアフィン対角射を持ち、
従って特に準分離的である。

\begin{lemma}
\label{morphisms-lemma-affine-diagonal}
$X$ をスキームとし、任意の点 $x \in X$ に対して、可逆
$\mathcal{O}_X$-加群 $\mathcal{L}$ と大域切断
$s \in \Gamma(X, \mathcal{L})$ で、$x \in X_s$ かつ $X_s$ がアフィンとなるものが
存在するとする。このとき $X$ の対角射はアフィン射である。
\end{lemma}

\begin{proof}
可逆 $\mathcal{O}_X$-加群 $\mathcal{L}$、$\mathcal{M}$ と大域切断
$s \in \Gamma(X, \mathcal{L})$、$t \in \Gamma(X, \mathcal{M})$ で、
$X_s$ と $X_t$ がアフィンとなるものが与えられたとき、$X_s \cap X_t$ がアフィンであることを
示せばよい。実際、そのとき補題 \ref{morphisms-lemma-characterize-affine} を
$\Delta : X \to X \times X$ に適用し、
$\Delta^{-1}(X_s \times X_t) = X_s \cap X_t$ であることを用いれば、$\Delta$ が
アフィンであると分かる。$X_s \cap X_t$ がアフィンであることは性質の章、補題
\ref{properties-lemma-affine-cap-s-open} から従う。
\end{proof}

\begin{remark}
\label{morphisms-remark-affine-s-opens-cover-family}
性質の章、補題 \ref{properties-lemma-affine-s-opens-cover-quasi-separated} では、
豊富な可逆層を持つスキームが分離的であることを見た。豊富な可逆層の族を持つスキームについては、
これは成り立たない。実際、$X$ をスキームの章、例
\ref{schemes-example-affine-space-zero-doubled} の $n = 1$ の場合、すなわち原点を二重化した
アフィン直線とする。その例の記法を用いるが、$x$ と書いて $x_1$ を表し、$y$ と書いて $y_1$ を表す。
任意の整数 $n$ に対し、可逆層 $\mathcal{L}_n$（$X$ 上のもの）で、$X_1$ と $X_2$ 上では自明であり、
その遷移関数 $U_{12} \to U_{21}$ が $f(x) \mapsto y^n f(y)$ であるものが存在する。
$\mathcal{L}_n$ の大域切断は組 $(f(x), g(y)) \in k[x] \oplus
k[y]$ で、$y^n f(y) = g(y)$ を満たすものである。切断 $s = (1,
y)$（$\mathcal{L}_1$ のもの）と $t = (x, 1)$（$\mathcal{L}_{-1}$ のもの）は、
$X_s = X_1$ かつ $X_t = X_2$ であるためアフィン開被覆を定める。従って $X$ は
豊富な可逆層の族を持つが、分離的ではない。
\end{remark}

\section{準アフィン射}
\label{morphisms-section-quasi-affine}

\noindent
スキーム $X$ が準コンパクトで、あるアフィンスキームの開部分スキームと同型であるとき、
そのスキームを {\it 準アフィン}と呼ぶことを思い出そう。性質の章、定義
\ref{properties-definition-quasi-affine} を参照せよ。

\begin{definition}
\label{morphisms-definition-quasi-affine}
スキームの射 $f : X \to S$ は、$S$ の任意のアフィン開集合の逆像が準アフィンスキームであるとき
{\it 準アフィン}であるという。
\end{definition}

\begin{lemma}
\label{morphisms-lemma-quasi-affine-separated}
準アフィン射は分離的かつ準コンパクトである。
\end{lemma}

\begin{proof}
$f : X \to S$ を準アフィンとする。準コンパクト性はスキームの章、補題
\ref{schemes-lemma-quasi-compact-affine} から直ちに従う。
$U \subset S$ をアフィン開集合とする。$f^{-1}(U)$ が分離的スキームであることを
示せれば、$f$ は分離的である（スキームの章、補題
\ref{schemes-lemma-characterize-separated} により、分離性は底上局所的である）。
仮定により $f^{-1}(U)$ はアフィンスキームの開部分スキームと同型である。
アフィンスキームは分離的であり、従ってその任意の開部分スキームも分離的である。
\end{proof}

\begin{lemma}
\label{morphisms-lemma-characterize-quasi-affine}
$f : X \to S$ をスキームの射とする。次は同値である。
\begin{enumerate}
\item 射 $f$ は準アフィンである。
\item アフィン開被覆 $S = \bigcup W_j$ で、各 $f^{-1}(W_j)$ が準アフィンとなるものが存在する。
\item 準連接 $\mathcal{O}_S$-代数層 $\mathcal{A}$ と準コンパクトな開埋め込み
$$
\xymatrix{
X \ar[rr] \ar[rd] & & \underline{\Spec}_S(\mathcal{A}) \ar[dl] \\
& S &
}
$$
（$S$ 上のもの）が存在する。
\item (3) と同じであるが、$\mathcal{A} = f_*\mathcal{O}_X$ とし、水平矢印を
構成の章、補題 \ref{constructions-lemma-canonical-morphism} の標準射とする。
\end{enumerate}
\end{lemma}

\begin{proof}
(1) が (2) を含意し、(4) が (3) を含意することは明らかである。

\medskip\noindent
$S = \bigcup_{j \in J} W_j$ を、各 $f^{-1}(W_j)$ が準アフィンとなるアフィン開被覆と仮定する。
スキームの章、補題 \ref{schemes-lemma-quasi-compact-affine} により、$f$ は準コンパクトである。
スキームの章、補題 \ref{schemes-lemma-characterize-quasi-separated} により、射 $f$ は準分離的である。
従ってスキームの章、補題 \ref{schemes-lemma-push-forward-quasi-coherent} により、層
$\mathcal{A} = f_*\mathcal{O}_X$ は準連接 $\mathcal{O}_X$-代数層である。ゆえにスキーム
$g : Y = \underline{\Spec}_S(\mathcal{A}) \to S$（$S$ 上のもの）が得られる。恒等写像
$\text{id} : \mathcal{A} = f_*\mathcal{O}_X \to f_*\mathcal{O}_X$ は、相対スペクトルの定義を介して
射 $can : X \to Y$（$S$ 上のもの）を与える。構成の章、補題
\ref{constructions-lemma-canonical-morphism} を参照せよ。仮定、今引用した補題、および性質の章、補題
\ref{properties-lemma-quasi-affine} により、制限
$can|_{f^{-1}(W_j)} : f^{-1}(W_j) \to g^{-1}(W_j)$ は準コンパクトな開埋め込みである。
従って $can$ は準コンパクトな開埋め込みである。これで (2) が (4) を含意することを示した。

\medskip\noindent
(3) を仮定する。任意のアフィン開集合 $U \subset S$ を選ぶ。構成の章、補題
\ref{constructions-lemma-spec-properties} により、相対スペクトルにおける $U$ の逆像はアフィンである。
従って $f^{-1}(U)$ は準アフィンである（準コンパクト性も (3) に組み込まれていることに注意せよ）。
これで (3) が (1) を含意することを示した。
\end{proof}

\begin{lemma}
\label{morphisms-lemma-composition-quasi-affine}
準アフィン射の合成は準アフィンである。
\end{lemma}

\begin{proof}
$f : X \to Y$ と $g : Y \to Z$ を準アフィン射とする。
$U \subset Z$ をアフィン開集合とする。$g^{-1}(U)$ は $g$ に関する仮定により準アフィンである。
$j : g^{-1}(U) \to V$ を、アフィンスキーム $V$ への準コンパクトな開埋め込みとする。
上の補題 \ref{morphisms-lemma-characterize-quasi-affine} により、$f^{-1}(g^{-1}(U))$ は相対スペクトル
$\underline{\Spec}_{g^{-1}(U)}(\mathcal{A})$ の準コンパクトな開部分スキームである。ここでは準連接
$\mathcal{O}_{g^{-1}(U)}$-代数層 $\mathcal{A}$ を用いた。
スキームの章、補題 \ref{schemes-lemma-push-forward-quasi-coherent} により、層
$\mathcal{A}' = j_*\mathcal{A}$ は準連接 $\mathcal{O}_V$-代数層であり、
$j^*\mathcal{A}' = \mathcal{A}$ を満たす。従って可換図式
$$
\xymatrix{
f^{-1}(g^{-1}(U)) \ar[r] &
\underline{\Spec}_{g^{-1}(U)}(\mathcal{A})
\ar[r] \ar[d] &
\underline{\Spec}_V(\mathcal{A}') \ar[d] \\
& g^{-1}(U) \ar[r]^j & V
}
$$
が得られ、その正方形はファイバー積である。構成の章、補題
\ref{constructions-lemma-spec-properties} を参照せよ。右上隅がアフィンスキームであることに注意せよ。
従って $(g \circ f)^{-1}(U)$ は準アフィンである。
\end{proof}

\begin{lemma}
\label{morphisms-lemma-base-change-quasi-affine}
準アフィン射の基底変換は準アフィンである。
\end{lemma}

\begin{proof}
$f : X \to S$ を準アフィン射とする。上の補題 \ref{morphisms-lemma-characterize-quasi-affine} により、
準連接 $\mathcal{O}_S$-代数層 $\mathcal{A}$ と準コンパクトな開埋め込み
$X \to \underline{\Spec}_S(\mathcal{A})$（$S$ 上のもの）を見つけられる。
$g : S' \to S$ を任意の射とする。$f' : X_{S'} = S' \times_S X \to S'$ で
$f$ の基底変換を表す。準コンパクトな開埋め込みの基底変換は準コンパクトな開埋め込みなので、
$X_{S'} \to \underline{\Spec}_{S'}(g^*\mathcal{A})$ は準コンパクトな開埋め込みである
（スキームの章、補題 \ref{schemes-lemma-quasi-compact-preserved-base-change}、
\ref{schemes-lemma-base-change-immersion}、および構成の章、補題
\ref{constructions-lemma-spec-properties} を用いた）。再び補題
\ref{morphisms-lemma-characterize-quasi-affine} により、$X_{S'} \to S'$ は準アフィンである。
\end{proof}

\begin{lemma}
\label{morphisms-lemma-quasi-compact-immersion-quasi-affine}
準コンパクトな埋め込みは準アフィンである。
\end{lemma}

\begin{proof}
$X \to S$ を準コンパクトな埋め込みとする。任意のアフィン開集合の逆像が準アフィンであることを
示さなければならない。従って $S$ がアフィンスキームであると仮定して、$X$ が準アフィンであることを
示せばよい。補題 \ref{morphisms-lemma-quasi-compact-immersion} により、射 $X \to S$ は
$X \to Z \to S$ と分解し、ここで $Z$ は $S$ の閉部分スキーム、$X \subset Z$ は
準コンパクトな開部分集合である。$S$ はアフィンなので、補題 \ref{morphisms-lemma-closed-immersion} により
$Z$ はアフィンである。これで証明が終わる。
\end{proof}

\begin{lemma}
\label{morphisms-lemma-affine-quasi-affine}
$S$ をスキームとし、$X$ をアフィンスキームとする。射 $f : X \to S$ が準アフィンであることと
準コンパクトであることは同値である。特に、アフィンスキームから準分離的スキームへの任意の射は
準アフィンである。
\end{lemma}

\begin{proof}
$V \subset S$ をアフィン開集合とする。このとき $f^{-1}(V)$ はアフィンスキーム $X$ の
開部分スキームなので、それが準アフィンであることと準コンパクトであることは同値である。
これで最初の主張を示した。任意の $f : X \to S$ に対し、$X$ がアフィンで
$S$ が準分離的であるとき、その準コンパクト性はスキームの章、補題
\ref{schemes-lemma-quasi-compact-permanence} を $X \to S \to \Spec(\mathbf{Z})$ に適用すれば従う。
\end{proof}

\begin{lemma}
\label{morphisms-lemma-quasi-affine-permanence}
$g : X \to Y$ を $S$ 上のスキームの射とする。$X$ が $S$ 上準アフィンで、
$Y$ が $S$ 上準分離的ならば、$g$ は準アフィンである。特に、準アフィンスキームから
準分離的スキームへの任意の射は準アフィンである。
\end{lemma}

\begin{proof}
基底変換 $X \times_S Y \to Y$ は補題 \ref{morphisms-lemma-base-change-quasi-affine} により準アフィンである。
射 $X \to X \times_S Y$ は、$Y \to S$ が準分離的なので準コンパクトな埋め込みである。
スキームの章、補題 \ref{schemes-lemma-section-immersion} を参照せよ。
準コンパクトな埋め込みは補題 \ref{morphisms-lemma-quasi-compact-immersion-quasi-affine} により準アフィンであり、
準アフィン射の合成は準アフィンである（補題 \ref{morphisms-lemma-composition-quasi-affine} を参照）。
これで証明が終わる。
\end{proof}

\section{環準同型の性質によって定義される射の型}
\label{morphisms-section-properties-ring-maps}

\noindent
この節では、スキームの射の局所的性質を定義できるような環準同型の性質を調べる。

\begin{definition}
\label{morphisms-definition-property-local}
$P$ を環準同型の性質とする。
\begin{enumerate}
\item 次の条件が成り立つとき、$P$ は {\it 局所的}であるという。
\begin{enumerate}
\item 任意の環準同型 $R \to A$ と任意の $f \in R$ に対し、
$P(R \to A) \Rightarrow P(R_f \to A_f)$ が成り立つ。
\item 任意の環 $R$, $A$、任意の $f \in R$, $a\in A$、および任意の環準同型
$R_f \to A$ に対し、$P(R_f \to A) \Rightarrow P(R \to A_a)$ が成り立つ。
\item 任意の環準同型 $R \to A$ と $a_i \in A$ で、
$(a_1, \ldots, a_n) = A$ を満たすものに対し、
$\forall i, P(R \to A_{a_i}) \Rightarrow P(R \to A)$ が成り立つ。
\end{enumerate}
\item $P$ は、任意の環準同型 $R \to A$, $R \to R'$ に対して
$P(R \to A) \Rightarrow P(R' \to R' \otimes_R A)$ が成り立つとき、
{\it 基底変換で安定}であるという。
\item $P$ は、任意の環準同型 $A \to B$, $B \to C$ に対して
$P(A \to B) \wedge P(B \to C) \Rightarrow P(A \to C)$ が成り立つとき、
{\it 合成で安定}であるという。
\end{enumerate}
\end{definition}

\begin{definition}
\label{morphisms-definition-locally-P}
$P$ を環準同型の性質とし、$f : X \to S$ をスキームの射とする。
$f$ が {\it 局所的に $P$ 型}であるとは、任意の $x \in X$ に対して、
あるアフィン開近傍 $U$（$x$ の $X$ における近傍）で、あるアフィン開集合
$V \subset S$ に写され、誘導される環準同型
$\mathcal{O}_S(V) \to \mathcal{O}_X(U)$ が性質 $P$ を持つものが存在することをいう。
\end{definition}

\noindent
性質 $P$ が局所的性質でなければ、これは「よい」定義ではない。
$P$ が局所的性質であっても、別の箇所で定義を明示しない限り、
この定義を自動的に用いて射を「局所的に $P$ 型」と呼ぶことはしない。

\begin{lemma}
\label{morphisms-lemma-locally-P}
$f : X \to S$ をスキームの射とし、$P$ を環準同型の性質とする。
$U$ を $X$ のアフィン開集合、$V$ を $S$ のアフィン開集合で
$f(U) \subset V$ を満たすものとする。$f$ が局所的に $P$ 型で、$P$ が局所的ならば、
$P(\mathcal{O}_S(V) \to \mathcal{O}_X(U))$ が成り立つ。
\end{lemma}

\begin{proof}
$f$ は局所的に $P$ 型なので、任意の $u \in U$ に対し、あるアフィン開集合
$U_u \subset X$ と $V_u \subset S$ が存在して、前者は後者へ写され、
$P(\mathcal{O}_S(V_u) \to \mathcal{O}_X(U_u))$ が成り立つ。
$U'_u \subset U \cap U_u$ という $u$ の開近傍で、$U$ と $U_u$ の双方において
標準アフィン開集合となるものを選ぶ。スキームの章、補題
\ref{schemes-lemma-standard-open-two-affines} を参照せよ。定義
\ref{morphisms-definition-property-local} (1)(b) により、
$P(\mathcal{O}_S(V_u) \to \mathcal{O}_X(U'_u))$ が成り立つ。
従って $U_u \subset U$ は標準アフィン開集合であると仮定してよい。
$V'_u \subset V \cap V_u$ という $f(u)$ の開近傍で、$V$ と $V_u$ の双方において
標準アフィン開集合となるものを選ぶ。スキームの章、補題
\ref{schemes-lemma-standard-open-two-affines} を参照せよ。このとき
$U'_u = f^{-1}(V'_u) \cap U_u$ は $U_u$ の（従って $U$ の）標準アフィン開集合であり、
定義 \ref{morphisms-definition-property-local} (1)(a) により
$P(\mathcal{O}_S(V'_u) \to \mathcal{O}_X(U'_u))$ が成り立つ。
従って $U_u \subset U$ と $V_u \subset V$ はともに標準アフィン開集合であると
仮定してよい。定義 \ref{morphisms-definition-property-local} (1)(b) をもう一度適用すると、
$P(\mathcal{O}_S(V) \to \mathcal{O}_X(U_u))$ が成り立つ。
$U$ は準コンパクトなので、有限個の点 $u_1, \ldots, u_n \in U$ を
$$
U = U_{u_1} \cup \ldots \cup U_{u_n}.
$$
となるように選べる。定義 \ref{morphisms-definition-property-local} (1)(c) により、
$P(\mathcal{O}_S(V) \to \mathcal{O}_X(U))$ が成り立つ。
\end{proof}

\begin{lemma}
\label{morphisms-lemma-locally-P-characterize}
$P$ を環準同型の局所的性質とし、$f : X \to S$ をスキームの射とする。
次は同値である。
\begin{enumerate}
\item 射 $f$ は局所的に $P$ 型である。
\item 任意のアフィン開集合 $U \subset X$, $V \subset S$ で $f(U) \subset V$ を満たすものに対し、
$P(\mathcal{O}_S(V) \to \mathcal{O}_X(U))$ が成り立つ。
\item 開被覆 $S = \bigcup_{j \in J} V_j$ と開被覆
$f^{-1}(V_j) = \bigcup_{i \in I_j} U_i$ で、各射
$U_i \to V_j$, $j\in J, i\in I_j$ が局所的に $P$ 型となるものが存在する。
\item アフィン開被覆 $S = \bigcup_{j \in J} V_j$ とアフィン開被覆
$f^{-1}(V_j) = \bigcup_{i \in I_j} U_i$ で、
$P(\mathcal{O}_S(V_j) \to \mathcal{O}_X(U_i))$ がすべての
$j\in J, i\in I_j$ に対して成り立つものが存在する。
\end{enumerate}
さらに、$f$ が局所的に $P$ 型ならば、任意の開部分スキーム
$U \subset X$, $V \subset S$ で $f(U) \subset V$ を満たすものに対し、
制限 $f|_U : U \to V$ は局所的に $P$ 型である。
\end{lemma}

\begin{proof}
これは上の補題 \ref{morphisms-lemma-locally-P} から従う。
\end{proof}

\begin{lemma}
\label{morphisms-lemma-composition-type-P}
$P$ を環準同型の性質とする。$P$ は局所的で、合成で安定であると仮定する。
局所的に $P$ 型である射の合成も局所的に $P$ 型である。
\end{lemma}

\begin{proof}
$f : X \to Y$ と $g : Y \to Z$ を局所的に $P$ 型の射とする。
$x \in X$ とする。アフィン開近傍 $W \subset Z$（$g(f(x))$ のもの）を選ぶ。
アフィン開近傍 $V \subset g^{-1}(W)$（$f(x)$ のもの）を選び、さらにアフィン開近傍
$U \subset f^{-1}(V)$（$x$ のもの）を選ぶ。補題 \ref{morphisms-lemma-locally-P-characterize} により、環準同型
$\mathcal{O}_Z(W) \to \mathcal{O}_Y(V)$ と
$\mathcal{O}_Y(V) \to \mathcal{O}_X(U)$ は $P$ を満たす。
従って $\mathcal{O}_Z(W) \to \mathcal{O}_X(U)$ は $P$ を満たす。これは $P$ が
合成で安定と仮定したためである。
\end{proof}

\begin{lemma}
\label{morphisms-lemma-base-change-type-P}
$P$ を環準同型の性質とする。$P$ は局所的で、基底変換で安定であると仮定する。
局所的に $P$ 型である射の基底変換も局所的に $P$ 型である。
\end{lemma}

\begin{proof}
$f : X \to S$ を局所的に $P$ 型の射とする。$S' \to S$ を任意の射とする。
$f' : X_{S'} = S' \times_S X \to S'$ で $f$ の基底変換を表す。
任意の $s' \in S'$ に対し、アフィン開近傍 $s' \in V' \subset S'$ で、あるアフィン開集合
$V \subset S$ に写されるものが存在する。補題 \ref{morphisms-lemma-locally-P-characterize} により、
開集合 $f^{-1}(V)$ はアフィン開集合 $U_i$ の合併であり、環準同型
$\mathcal{O}_S(V) \to \mathcal{O}_X(U_i)$ は
すべて $P$ を満たす。スキームの章、節 \ref{schemes-section-fibre-products} の内容により、
$f^{-1}(U)_{V'} = V' \times_V f^{-1}(V)$ はアフィン開集合
$V' \times_V U_i$ の合併である。さらに
$\mathcal{O}_{X_{S'}}(V' \times_V U_i) =
\mathcal{O}_{S'}(V') \otimes_{\mathcal{O}_S(V)} \mathcal{O}_X(U_i)$
である。従って環準同型
$\mathcal{O}_{S'}(V') \to \mathcal{O}_{X_{S'}}(V' \times_V U_i)$ は $P$ を満たす。
これは $P$ が基底変換で安定と仮定したためである。
\end{proof}

\begin{lemma}
\label{morphisms-lemma-properties-local}
環準同型 $R \to A$ の次の性質は局所的である。
\begin{enumerate}
\item （局所環上の同型。）
素イデアル $\mathfrak q$（$A$ のもの）が素イデアル $\mathfrak p \subset R$ の上にあるとき、
環準同型 $R \to A$ が誘導する写像
$R_{\mathfrak p} \to A_{\mathfrak q}$ は同型である。
\item （開埋め込み。）
任意の素イデアル $\mathfrak q$（$A$ のもの）に対し、ある $f \in R$ が存在して、
$\varphi(f) \not \in \mathfrak q$ かつ環準同型 $\varphi : R \to A$ が誘導する写像
$R_f \to A_f$ は同型である。
\item （ファイバーが被約。）
任意の素イデアル $\mathfrak p$（$R$ のもの）に対し、ファイバー環
$A \otimes_R \kappa(\mathfrak p)$ は被約である。
\item （ファイバーの次元が高々 $n$。）
任意の素イデアル $\mathfrak p$（$R$ のもの）に対し、ファイバー環
$A \otimes_R \kappa(\mathfrak p)$ の Krull 次元は高々 $n$ である。
\item （標的上局所 Noether。）
環準同型 $R \to A$ は、$A$ が Noether 環であるという性質を持つ。
\item 必要に応じてここに追加する\footnote{ただし、別の箇所で既に個別に扱われていない性質に限る。}。
\end{enumerate}
\end{lemma}

\begin{proof}
省略する。
\end{proof}

\begin{lemma}
\label{morphisms-lemma-properties-base-change}
環準同型の次の性質は基底変換で安定である。
\begin{enumerate}
\item （局所環上の同型。）
素イデアル $\mathfrak q$（$A$ のもの）が素イデアル $\mathfrak p \subset R$ の上にあるとき、
環準同型 $R \to A$ が誘導する写像
$R_{\mathfrak p} \to A_{\mathfrak q}$ は同型である。
\item （開埋め込み。）
任意の素イデアル $\mathfrak q$（$A$ のもの）に対し、ある $f \in R$ が存在して、
$\varphi(f) \not \in \mathfrak q$ かつ環準同型 $\varphi : R \to A$ が誘導する写像
$R_f \to A_f$ は同型である。
\item 必要に応じてここに追加する\footnote{ただし、別の箇所で既に個別に扱われていない性質に限る。}。
\end{enumerate}
\end{lemma}

\begin{proof}
省略する。
\end{proof}

\begin{lemma}
\label{morphisms-lemma-properties-composition}
環準同型の次の性質は合成で安定である。
\begin{enumerate}
\item （局所環上の同型。）
素イデアル $\mathfrak q$（$A$ のもの）が素イデアル $\mathfrak p \subset R$ の上にあるとき、
環準同型 $R \to A$ が誘導する写像
$R_{\mathfrak p} \to A_{\mathfrak q}$ は同型である。
\item （開埋め込み。）
任意の素イデアル $\mathfrak q$（$A$ のもの）に対し、ある $f \in R$ が存在して、
$\varphi(f) \not \in \mathfrak q$ かつ環準同型 $\varphi : R \to A$ が誘導する写像
$R_f \to A_f$ は同型である。
\item （標的上局所 Noether。）
環準同型 $R \to A$ は、$A$ が Noether 環であるという性質を持つ。
\item 必要に応じてここに追加する\footnote{ただし、別の箇所で既に個別に扱われていない性質に限る。}。
\end{enumerate}
\end{lemma}

\begin{proof}
省略する。
\end{proof}

\section{有限型の射}
\label{morphisms-section-finite-type}

\noindent
環準同型 $R \to A$ が有限型であるとは、$A$ が
$R[x_1, \ldots, x_n]$ の商と $R$-代数として同型であることをいう。代数の章、定義
\ref{algebra-definition-finite-type} を参照せよ。

\begin{definition}
\label{morphisms-definition-finite-type}
$f : X \to S$ をスキームの射とする。
\begin{enumerate}
\item $f$ が {\it $x \in X$ において有限型}であるとは、アフィン開近傍
$\Spec(A) = U \subset X$（$x$ のもの）と、アフィン開集合 $\Spec(R) = V \subset S$ で、
$f(U) \subset V$ かつ誘導される環準同型 $R \to A$ が有限型となるものが存在することをいう。
\item $f$ が {\it 局所有限型}であるとは、それが $X$ のすべての点で有限型であることをいう。
\item $f$ が {\it 有限型}であるとは、それが局所有限型かつ準コンパクトであることをいう。
\end{enumerate}
\end{definition}

\begin{lemma}
\label{morphisms-lemma-locally-finite-type-characterize}
$f : X \to S$ をスキームの射とする。次は同値である。
\begin{enumerate}
\item 射 $f$ は局所有限型である。
\item すべてのアフィン開集合 $U \subset X$, $V \subset S$ で $f(U) \subset V$ を満たすものに対し、
環準同型 $\mathcal{O}_S(V) \to \mathcal{O}_X(U)$ は有限型である。
\item 開被覆 $S = \bigcup_{j \in J} V_j$ と開被覆
$f^{-1}(V_j) = \bigcup_{i \in I_j} U_i$ で、各射
$U_i \to V_j$, $j\in J, i\in I_j$ が局所有限型となるものが存在する。
\item アフィン開被覆 $S = \bigcup_{j \in J} V_j$ とアフィン開被覆
$f^{-1}(V_j) = \bigcup_{i \in I_j} U_i$ で、環準同型
$\mathcal{O}_S(V_j) \to \mathcal{O}_X(U_i)$ がすべての
$j\in J, i\in I_j$ に対して有限型となるものが存在する。
\end{enumerate}
さらに、$f$ が局所有限型ならば、任意の開部分スキーム
$U \subset X$, $V \subset S$ で $f(U) \subset V$ を満たすものに対し、
制限 $f|_U : U \to V$ は局所有限型である。
\end{lemma}

\begin{proof}
性質「$R \to A$ が有限型である」が局所的であることを示せば、補題
\ref{morphisms-lemma-locally-P} から従う。定義 \ref{morphisms-definition-property-local} の条件 (a)、(b)、(c) を確認する。
代数の章、補題 \ref{algebra-lemma-base-change-finiteness} により、有限型であることは
基底変換で安定なので (a) が成り立つ。代数の章、補題
\ref{algebra-lemma-compose-finite-type} により有限型であることは合成で安定であり、
任意の環 $R$ に対して環準同型 $R \to R_f$ が有限型であることは明らかである。
従って (b) が成り立つ。最後に (c) は代数の章、補題
\ref{algebra-lemma-cover-upstairs} により成り立つ。
\end{proof}

\begin{lemma}
\label{morphisms-lemma-composition-finite-type}
二つの局所有限型の射の合成は局所有限型である。有限型の射についても同じことが成り立つ。
\end{lemma}

\begin{proof}
補題 \ref{morphisms-lemma-locally-finite-type-characterize} の証明で、有限型であることが
環準同型の局所的性質であることを見た。従って補題の最初の主張は、補題
\ref{morphisms-lemma-composition-type-P} と、有限型であることが合成で安定な環準同型の性質であることを
組み合わせれば従う。代数の章、補題 \ref{algebra-lemma-compose-finite-type} を参照せよ。
以上と、準コンパクトな射の合成が準コンパクトであることから、有限型の射の合成は
有限型である。スキームの章、補題 \ref{schemes-lemma-composition-quasi-compact} を参照せよ。
\end{proof}

\begin{lemma}
\label{morphisms-lemma-base-change-finite-type}
局所有限型の射の基底変換は局所有限型である。有限型の射についても同じことが成り立つ。
\end{lemma}

\begin{proof}
補題 \ref{morphisms-lemma-locally-finite-type-characterize} の証明で、有限型であることが
環準同型の局所的性質であることを見た。従って補題の最初の主張は、補題
\ref{morphisms-lemma-base-change-type-P} と、有限型であることが基底変換で安定な環準同型の性質であることを
組み合わせれば従う。代数の章、補題 \ref{algebra-lemma-base-change-finiteness} を参照せよ。
以上と、準コンパクトな射の基底変換が準コンパクトであることから、有限型の射の基底変換は
有限型の射である。スキームの章、補題
\ref{schemes-lemma-quasi-compact-preserved-base-change} を参照せよ。
\end{proof}

\begin{lemma}
\label{morphisms-lemma-immersion-locally-finite-type}
閉埋め込みは有限型である。埋め込みは局所有限型である。
\end{lemma}

\begin{proof}
開埋め込みは局所同型であり、閉埋め込みは明らかに有限型だからである。
\end{proof}

\begin{lemma}
\label{morphisms-lemma-finite-type-noetherian}
$f : X \to S$ を射とする。$S$ が（局所）Noether で、$f$ が（局所）有限型ならば、
$X$ は（局所）Noether である。
\end{lemma}

\begin{proof}
Noether 環上の有限型環が Noether 環であることから直ちに従う。代数の章、補題
\ref{algebra-lemma-Noetherian-permanence} を参照せよ。
（また、準コンパクトな標的を持つ準コンパクトな射の始域が準コンパクトであることを用いよ。）
\end{proof}

\begin{lemma}
\label{morphisms-lemma-finite-type-Noetherian-quasi-separated}
$f : X \to S$ を局所有限型とし、$S$ を局所 Noether とする。このとき $f$ は準分離的である。
\end{lemma}

\begin{proof}
実際、$X$ が準分離的である。性質の章、補題
\ref{properties-lemma-locally-Noetherian-quasi-separated} と上の補題
\ref{morphisms-lemma-finite-type-noetherian} を参照せよ。そのうえでスキームの章、補題
\ref{schemes-lemma-compose-after-separated} を適用すれば、$f$ が準分離的であると結論できる。
\end{proof}

\begin{lemma}
\label{morphisms-lemma-permanence-finite-type}
$X \to Y$ を基礎スキーム $S$ 上のスキームの射とする。$X$ が $S$ 上局所有限型ならば、
$X \to Y$ は局所有限型である。
\end{lemma}

\begin{proof}
補題 \ref{morphisms-lemma-locally-finite-type-characterize} を介すと、これは次の代数的事実に翻訳される。
環準同型 $A \to B \to C$ が与えられ、$A \to C$ が有限型ならば、$B \to C$ も有限型である。
（代数の章、補題 \ref{algebra-lemma-compose-finite-type} を参照せよ。）
\end{proof}

\section{有限型の点と Jacobson スキーム}
\label{morphisms-section-points-finite-type}

\noindent
$S$ をスキームとする。有限型の点 $s$（$S$ の点）とは、射
$\Spec(\kappa(s)) \to S$ が有限型であるような点をいう。これを調べる理由は、
ある意味で、またある状況では、有限型の点が閉点の代わりを果たせるからである。
例えば、そのような点は常に十分多く存在する。さらに、スキームが Jacobson であることと、
その有限型の点がすべて閉点であることは同値である。

\begin{lemma}
\label{morphisms-lemma-point-finite-type}
$S$ をスキーム、$k$ を体とし、$f : \Spec(k) \to S$ を射とする。
次は同値である。
\begin{enumerate}
\item 射 $f$ は有限型である。
\item 射 $f$ は局所有限型である。
\item アフィン開集合 $U = \Spec(R)$（$S$ のもの）で、$f$ が有限環準同型 $R \to k$ に
対応するものが存在する。
\item アフィン開集合 $U = \Spec(R)$（$S$ のもの）で、$f$ の像が閉点 $u$（$U$ のもの）からなり、
体拡大 $k/\kappa(u)$ が有限となるものが存在する。
\end{enumerate}
\end{lemma}

\begin{proof}
$\Spec(k)$ は一点集合であり、従ってそこからの任意の射は準コンパクトなので、
(1) と (2) の同値性は明らかである。

\medskip\noindent
$f$ が局所有限型であると仮定する。任意のアフィン開集合
$\Spec(R) = U \subset S$ で、$f$ の像が $U$ に含まれ、環準同型
$R \to k$ が有限型となるものを選ぶ。
$\mathfrak p \subset R$ をその核とする。このとき $R/\mathfrak p \subset k$ は有限型である。
代数の章、補題 \ref{algebra-lemma-field-finite-type-over-domain} により、ある
$\overline{f} \in R/\mathfrak p$ が存在して、$(R/\mathfrak p)_{\overline{f}}$ は体であり、
$(R/\mathfrak p)_{\overline{f}} \to k$ は有限次体拡大である。
$f \in R$ が $\overline{f}$ の持ち上げならば、$k$ は有限 $R_f$-加群であることが分かる。
従って (2) $\Rightarrow$ (3) である。

\medskip\noindent
$\Spec(R) = U \subset S$ がアフィン開集合で、$f$ が有限環準同型 $R \to k$ に
対応すると仮定する。このとき補題 \ref{morphisms-lemma-locally-finite-type-characterize} により、
$f$ は局所有限型である。従って (3) $\Rightarrow$ (2) である。

\medskip\noindent
$R \to k$ が有限であると仮定する。代数の章、補題
\ref{algebra-lemma-integral-under-field} により、$R \to k$ の像は体で、その上で $k$ は有限である。
従って $R \to k$ の核は極大イデアルである。ゆえに (3) $\Rightarrow$ (4) である。

\medskip\noindent
(4) $\Rightarrow$ (3) は直ちに従う。
\end{proof}

\begin{lemma}
\label{morphisms-lemma-artinian-finite-type}
$S$ をスキームとする。$A$ を剰余体 $\kappa$ を持つ Artin 局所環とし、
$f : \Spec(A) \to S$ をスキームの射とする。このとき $f$ が有限型であることと、合成
$\Spec(\kappa) \to \Spec(A) \to S$ が有限型であることは同値である。
\end{lemma}

\begin{proof}
射 $\Spec(\kappa) \to \Spec(A)$ は有限型なので、$f$ が有限型ならば合成
$\Spec(\kappa) \to S$ も有限型であることは明らかである（補題
\ref{morphisms-lemma-composition-finite-type} を参照せよ）。逆向きについては、射
$\Spec(A) \to S$ はアフィン開集合 $U = \Spec(B)$（$S$ のもの）に写ることに注意せよ。
これは $\Spec(A)$ が一点しか持たないためである。証明を終えるには、代数の章、補題
\ref{algebra-lemma-essentially-of-finite-type-into-artinian-local} を $B \to A$ に適用する。
\end{proof}

\noindent
点 $s$（スキーム $S$ のもの）が与えられると標準射 $\Spec(\kappa(s)) \to S$ があることを思い出そう。
スキームの章、節 \ref{schemes-section-points} を参照せよ。

\begin{definition}
\label{morphisms-definition-finite-type-point}
$S$ をスキームとする。点 $s$（$S$ のもの）は、標準射 $\Spec(\kappa(s)) \to S$ が有限型であるとき
{\it 有限型の点}であるという。$S_{\text{ft-pts}}$ で有限型の点全体の集合（$S$ のもの）を表す。
\end{definition}

\noindent
有限型の点全体の集合は次のように記述できる。

\begin{lemma}
\label{morphisms-lemma-identify-finite-type-points}
$S$ をスキームとする。このとき
$$
S_{\text{ft-pts}} = \bigcup\nolimits_{U \subset S\text{ open}} U_0
$$
である。ここで $U_0$ は $U$ の閉点全体の集合である。
このとき $U$ を $S$ のすべての開集合にわたらせても、すべてのアフィン開集合に
わたらせてもよい。
\end{lemma}

\begin{proof}
補題 \ref{morphisms-lemma-point-finite-type} から直ちに従う。
\end{proof}

\begin{lemma}
\label{morphisms-lemma-finite-type-points-morphism}
$f : T \to S$ をスキームの射とする。$f$ が局所有限型ならば、
$f(T_{\text{ft-pts}}) \subset S_{\text{ft-pts}}$ である。
\end{lemma}

\begin{proof}
$T$ が体のスペクトルならば、これは補題 \ref{morphisms-lemma-point-finite-type} である。
一般の場合は、局所有限型の射の合成が局所有限型であることから従う
（補題 \ref{morphisms-lemma-composition-finite-type}）。
\end{proof}

\begin{lemma}
\label{morphisms-lemma-finite-type-points-surjective-morphism}
$f : T \to S$ をスキームの射とする。$f$ が局所有限型かつ全射ならば、
$f(T_{\text{ft-pts}}) = S_{\text{ft-pts}}$ である。
\end{lemma}

\begin{proof}
補題 \ref{morphisms-lemma-finite-type-points-morphism} により
$f(T_{\text{ft-pts}}) \subset S_{\text{ft-pts}}$ である。
$s \in S$ を有限型の点とする。$f$ は全射なので、スキーム
$T_s = \Spec(\kappa(s)) \times_S T$ は空でない。従って補題
\ref{morphisms-lemma-identify-finite-type-points} により、それは有限型の点 $t \in T_s$ を持つ。
ここで $T_s \to T$ は $s \to S$ の基底変換として有限型の射である
（補題 \ref{morphisms-lemma-base-change-finite-type}）。従って補題
\ref{morphisms-lemma-finite-type-points-morphism} により、$t$ の $T$ における像は有限型の点であり、
構成により $s$ に写る。
\end{proof}

\begin{lemma}
\label{morphisms-lemma-enough-finite-type-points}
$S$ をスキームとする。任意の局所閉部分集合 $T \subset S$ に対し、
$$
T \not = \emptyset
\Rightarrow
T \cap S_{\text{ft-pts}} \not = \emptyset.
$$
特に、任意の閉部分集合 $T \subset S$ に対し、$T \cap S_{\text{ft-pts}}$ は $T$ で稠密である。
\end{lemma}

\begin{proof}
$T$ は、$T \to S$ が局所閉埋め込みとなるようなスキーム構造を持つことに注意せよ
（スキームの章、補題 \ref{schemes-lemma-reduced-closed-subscheme} を参照）。
任意の局所閉埋め込みは局所有限型である。補題
\ref{morphisms-lemma-immersion-locally-finite-type} を参照せよ。従って補題
\ref{morphisms-lemma-finite-type-points-morphism} により、$T_{\text{ft-pts}} \subset S_{\text{ft-pts}}$ である。
最後に、$T$ の任意の空でないアフィン開集合は少なくとも一つ閉点を持ち、その閉点は補題
\ref{morphisms-lemma-identify-finite-type-points} により $T$ の有限型の点である。
\end{proof}

\noindent
以上から、位相の章、節 \ref{topology-section-space-jacobson} の内容の大部分は、
閉点全体の集合を有限型の点全体の集合に置き換えても成り立つ。
実際、$S$ が Jacobson ならば、有限型の点として閉点が再び得られる。

\begin{lemma}
\label{morphisms-lemma-jacobson-finite-type-points}
$S$ をスキームとする。次は同値である。
\begin{enumerate}
\item スキーム $S$ は Jacobson である。
\item $S_{\text{ft-pts}}$ は $S$ の閉点全体の集合である。
\item 局所有限型のすべての $T \to S$ に対し、閉点は閉点に写る。
\item 局所有限型のすべての $T \to S$ に対し、閉点 $t \in T$ は閉点
$s \in S$ に写り、$\kappa(s) \subset \kappa(t)$ は有限である。
\end{enumerate}
\end{lemma}

\begin{proof}
(4) $\Rightarrow$ (3) $\Rightarrow$ (2) は明らかである。
補題 \ref{morphisms-lemma-enough-finite-type-points} により (2) は (1) を含意する。
従って (1) が (4) を含意することを示せば十分である。
$T \to S$ が局所有限型であると仮定する。閉点 $t \in T$ を選び、その像を $s \in S$ とする。
アフィン開近傍 $\Spec(R) = U \subset S$（$s$ のもの）と
アフィン開近傍 $\Spec(A) = V \subset T$（$t$ のもの）で、$V$ が $U$ に写るものを選ぶ。誘導される環準同型
$R \to A$ は有限型であり（補題 \ref{morphisms-lemma-locally-finite-type-characterize} を参照）、
性質の章、補題 \ref{properties-lemma-locally-jacobson} により $R$ は Jacobson である。
従って結果は代数の章、命題 \ref{algebra-proposition-Jacobson-permanence} から従う。
\end{proof}

\begin{lemma}
\label{morphisms-lemma-Jacobson-universally-Jacobson}
$S$ を Jacobson スキームとする。$S$ 上局所有限型の任意のスキームは Jacobson である。
\end{lemma}

\begin{proof}
これは代数の章、命題 \ref{algebra-proposition-Jacobson-permanence} から明らかである
（性質の章、補題 \ref{properties-lemma-locally-jacobson} と補題
\ref{morphisms-lemma-locally-finite-type-characterize} も参照せよ）。
\end{proof}

\begin{lemma}
\label{morphisms-lemma-ubiquity-Jacobson-schemes}
次の型のスキームは Jacobson である。
\begin{enumerate}
\item 体上局所有限型の任意のスキーム。
\item $\mathbf{Z}$ 上局所有限型の任意のスキーム。
\item 無限個の素イデアルを持つ $1$-次元 Noether 整域上局所有限型の任意のスキーム。
\item $\Spec(R) \setminus \{\mathfrak m\}$ の形のスキームで、
$(R, \mathfrak m)$ が Noether 局所環であるもの、およびその上局所有限型の任意のスキーム。
\end{enumerate}
\end{lemma}

\begin{proof}
補題 \ref{morphisms-lemma-Jacobson-universally-Jacobson} を断りなく用いる。
体のスペクトルは明らかに Jacobson である。$\mathbf{Z}$ のスペクトルは Jacobson である。
代数の章、補題 \ref{algebra-lemma-pid-jacobson} を参照せよ。(3) については代数の章、補題
\ref{algebra-lemma-noetherian-dim-1-Jacobson} を参照せよ。(4) については性質の章、補題
\ref{properties-lemma-complement-closed-point-Jacobson} を参照せよ。
\end{proof}

\section{普遍鎖状スキーム}
\label{morphisms-section-universally-catenary}

\noindent
位相空間 $X$ が {\it 鎖状}であるとは、既約閉部分集合の任意の対
$T \subset T'$ に対して、既約閉部分集合の極大鎖
$$
T = T_0 \subset T_1 \subset \ldots \subset T_e = T'
$$
が存在し、かつそのようなすべての鎖が同じ長さを持つことをいう。位相の章、定義
\ref{topology-definition-catenary} を参照せよ。スキームが鎖状であるとは、その台位相空間が
鎖状であることをいう。性質の章、定義 \ref{properties-definition-catenary} を参照せよ。

\begin{definition}
\label{morphisms-definition-universally-catenary}
$S$ をスキームとし、$S$ は局所 Noether であると仮定する。このとき $S$ は、局所有限型の任意の射
$X \to S$ に対してスキーム $X$ が鎖状であるとき、{\it 普遍鎖状}であるという。
\end{definition}

\noindent
これは鎖状性よりも「よい」概念である。実際、鎖状ではあるが普遍鎖状ではない Noether スキームが
存在する。例の章、節 \ref{examples-section-non-catenary-Noetherian-local} を参照せよ。
多くのスキームが普遍鎖状である。下の補題 \ref{morphisms-lemma-ubiquity-uc} を参照せよ。

\medskip\noindent
環 $A$ が {\it 鎖状}であるとは、素イデアルの任意の対
$\mathfrak p \subset \mathfrak q$ に対して、素イデアルの極大鎖
$$
\mathfrak p =
\mathfrak p_0 \subset \ldots \subset \mathfrak p_e
= \mathfrak q
$$
が存在し、これらがすべて同じ長さを持つことをいう。代数の章、定義
\ref{algebra-definition-catenary} を参照せよ。鎖状スキームと鎖状環の関係は性質の章、節
\ref{properties-section-catenary} で見た。環 $A$ が {\it 普遍鎖状}であるとは、
$A$ が Noether で、任意の有限型環準同型 $A \to B$ に対して環 $B$ が鎖状であることをいう。
代数の章、定義 \ref{algebra-definition-universally-catenary} を参照せよ。
代数幾何に現れる興味深い環の多くはこの性質を満たす。

\begin{lemma}
\label{morphisms-lemma-universally-catenary-local}
$S$ を局所 Noether スキームとする。次は同値である。
\begin{enumerate}
\item $S$ は普遍鎖状である。
\item $S$ の開被覆で、そのすべての要素が普遍鎖状スキームとなるものが存在する。
\item 任意のアフィン開集合 $\Spec(R) = U \subset S$ に対し、環 $R$ は普遍鎖状である。
\item アフィン開被覆 $S = \bigcup U_i$ で、各 $U_i$ が普遍鎖状環のスペクトルとなるものが存在する。
\end{enumerate}
さらに、この場合 $S$ 上局所有限型の任意のスキームも普遍鎖状である。
\end{lemma}

\begin{proof}
補題 \ref{morphisms-lemma-immersion-locally-finite-type} により、開埋め込みは局所有限型である。
局所有限型の射の合成は局所有限型である（補題 \ref{morphisms-lemma-composition-finite-type}）。
従って $S$ が普遍鎖状ならば、その任意の開部分スキーム、および $S$ 上局所有限型の
任意のスキームも普遍鎖状であることは明らかである。これにより補題の最後の主張と、
(1) が (2) を含意することが示された。

\medskip\noindent
$\Spec(R)$ が普遍鎖状スキームならば、任意のスキーム $\Spec(A)$ で、
$A$ が有限型 $R$-代数であるものは鎖状である。従って代数の章、補題 \ref{algebra-lemma-catenary} により、
これらの環 $A$ はすべて鎖状である。ゆえに $R$ は普遍鎖状である。
以上の考察と合わせて、(1) は (3) を含意し、(2) は (4) を含意する。
もちろん (3) は (4) を直ちに含意する。

\medskip\noindent
証明を終えるため、(4) が (1) を含意することを示す。(4) を仮定し、
$X \to S$ を局所有限型の射とする。アフィン開被覆 $X = \bigcup V_j$ で、
各 $V_j \to S$ がいずれかの $U_i$ に写るものを見つけられる。補題
\ref{morphisms-lemma-locally-finite-type-characterize} により、誘導される環準同型
$\mathcal{O}(U_i) \to \mathcal{O}(V_j)$ は有限型である。従って $\mathcal{O}(V_j)$ は鎖状である。
ゆえに性質の章、補題 \ref{properties-lemma-catenary-local} により $X$ は鎖状である。
\end{proof}

\begin{lemma}
\label{morphisms-lemma-universally-catenary-local-rings-universally-catenary}
$S$ を局所 Noether スキームとする。次は同値である。
\begin{enumerate}
\item $S$ は普遍鎖状である。
\item すべての局所環 $\mathcal{O}_{S, s}$（$S$ のもの）は普遍鎖状である。
\end{enumerate}
\end{lemma}

\begin{proof}
$S$ のすべての局所環が普遍鎖状であると仮定する。$f : X \to S$ を局所有限型とする。
$X$ が鎖状であることと、$\mathcal{O}_{X, x}$ がすべての $x \in X$ に対して鎖状であることは
同値である。$f(x) = s$ ならば、$\mathcal{O}_{X, x}$ は $\mathcal{O}_{S, s}$ 上本質的有限型である。
従って $\mathcal{O}_{X, x}$ は鎖状である。これは $\mathcal{O}_{S, s}$ が普遍鎖状であるという
仮定による。

\medskip\noindent
逆に、$S$ が普遍鎖状であると仮定し、$s \in S$ とする。補題
\ref{morphisms-lemma-universally-catenary-local} により、$S$ を $s$ のアフィン開近傍で置き換えてよい。
$S = \Spec(R)$ とし、$s$ は素イデアル $\mathfrak p$ に対応するとする。
任意の有限型 $R_{\mathfrak p}$-代数 $A'$ は $A_{\mathfrak p}$ の形であり、ここで
$R$-代数 $A$ は有限型である。仮定（必要なら補題 \ref{morphisms-lemma-universally-catenary-local}）により
環 $A$ は鎖状であり、従って $A'$（$A$ の局所化）も鎖状である。
ゆえに $R_{\mathfrak p}$ は普遍鎖状である。
\end{proof}

\begin{lemma}
\label{morphisms-lemma-catenary-check-irreducible}
$S$ を局所 Noether スキームとする。このとき $S$ が普遍鎖状であることと、
$S$ の既約成分が普遍鎖状であることは同値である。
\end{lemma}

\begin{proof}
省略する。アフィンの場合については、代数の章、補題
\ref{algebra-lemma-catenary-check-irreducible} を参照せよ。
\end{proof}

\begin{lemma}
\label{morphisms-lemma-ubiquity-uc}
次の型のスキームは普遍鎖状である。
\begin{enumerate}
\item 体上局所有限型の任意のスキーム。
\item Cohen--Macaulay スキーム上局所有限型の任意のスキーム。
\item $\mathbf{Z}$ 上局所有限型の任意のスキーム。
\item $1$-次元 Noether 整域上局所有限型の任意のスキーム。
\item など。
\end{enumerate}
\end{lemma}

\begin{proof}
これらはすべて、Cohen--Macaulay 環が普遍鎖状であることから従う。代数の章、補題
\ref{algebra-lemma-CM-ring-catenary} を参照せよ。また、補題
\ref{morphisms-lemma-universally-catenary-local} の最後の主張を用いる。細部の一部は省略する。
\end{proof}

\section{永田スキーム再論}
\label{morphisms-section-nagata}

\noindent
永田スキームおよび普遍的日本スキームの定義と基本性質については、
性質の章、節 \ref{properties-section-nagata} を参照せよ。

\begin{lemma}
\label{morphisms-lemma-finite-type-nagata}
$f : X \to S$ を射とする。
$S$ が永田スキームで $f$ が局所有限型ならば、$X$ は永田スキームである。
$S$ が普遍的日本スキームで $f$ が局所有限型ならば、
$X$ は普遍的日本スキームである。
\end{lemma}

\begin{proof}
「普遍的日本」については、代数の章、補題
\ref{algebra-lemma-universally-japanese} から従う。
「永田」については、代数の章、命題
\ref{algebra-proposition-nagata-universally-japanese} から従う。
\end{proof}

\begin{lemma}
\label{morphisms-lemma-ubiquity-nagata}
次の型のスキームは永田スキームである。
\begin{enumerate}
\item 体上局所有限型の任意のスキーム。
\item Noether 完備局所環上局所有限型の任意のスキーム。
\item $\mathbf{Z}$ 上局所有限型の任意のスキーム。
\item 標数零の Dedekind 環上局所有限型の任意のスキーム。
\item など。
\end{enumerate}
\end{lemma}

\begin{proof}
補題 \ref{morphisms-lemma-finite-type-nagata} により、上で述べた環が永田環であることを
示せばよい。このことについては、代数の章、命題
\ref{algebra-proposition-ubiquity-nagata} を参照せよ。
\end{proof}

\section{特異点集合再論}
\label{morphisms-section-singular-locus}

\noindent
基底上局所有限型の任意のスキームについて、正則点集合が開となることを保証する
判定条件を求める。定義は次のとおりである。

\begin{definition}
\label{morphisms-definition-J}
$X$ を局所 Noether スキームとする。$X$ が {\it J-2} であるとは、
局所有限型の任意の射 $Y \to X$ に対し、正則点集合
$\text{Reg}(Y)$ が $Y$ の開集合であることをいう。
\end{definition}

\noindent
これは Noether 環に対する対応する概念の類似である。発展代数の章、定義
\ref{more-algebra-definition-J} を参照せよ。

\begin{lemma}
\label{morphisms-lemma-J}
$X$ を局所 Noether スキームとする。次の条件は同値である。
\begin{enumerate}
\item $X$ は J-2 である。
\item $X$ には、各要素が J-2 スキームであるような開被覆が存在する。
\item 任意のアフィン開集合 $\Spec(R) = U \subset X$ に対し、環
$R$ は J-2 である。
\item アフィン開被覆 $S = \bigcup U_i$ であって、各
$\mathcal{O}(U_i)$ がすべての $i$ に対し J-2 であるものが存在する。
\end{enumerate}
さらに、この場合には $X$ 上局所有限型の任意のスキームも J-2 である。
\end{lemma}

\begin{proof}
補題 \ref{morphisms-lemma-immersion-locally-finite-type} により、開埋め込みは局所有限型である。
局所有限型の射の合成は局所有限型である
（補題 \ref{morphisms-lemma-composition-finite-type}）。したがって、$X$ が J-2 ならば、
任意の開部分スキームおよび $X$ 上局所有限型の任意のスキームも J-2 であることは明らかである。
これで補題の最後の主張が証明された。

\medskip\noindent
$\Spec(R)$ が J-2 ならば、任意の有限型 $R$-代数 $A$ に対し、
スキーム $\Spec(A)$ の正則点集合は開である。したがって定義により $R$ は
J-2 である（発展代数の章、定義 \ref{more-algebra-definition-J} を参照）。
上の考察と合わせると、(1) は (3) を、(2) は (4) を含意する。
もちろん (1) $\Rightarrow$ (2) および (3) $\Rightarrow$ (4) は自明である。

\medskip\noindent
証明を完了するため、(4) が (1) を含意することを示す。(4) を仮定し、
$Y \to X$ を局所有限型の射とする。アフィン開被覆 $Y = \bigcup V_j$ であって、
各 $V_j \to X$ の像がいずれか一つの $U_i$ に含まれるものを取ることができる。
補題 \ref{morphisms-lemma-locally-finite-type-characterize} により、誘導される環準同型
$\mathcal{O}(U_i) \to \mathcal{O}(V_j)$ は有限型である。したがって
$V_j = \Spec(\mathcal{O}(V_j))$ の正則点集合は開である。
$\text{Reg}(Y) \cap V_j = \text{Reg}(V_j)$ なので、望んだとおり
$\text{Reg}(Y)$ は開である。
\end{proof}

\begin{lemma}
\label{morphisms-lemma-ubiquity-J-2}
次の型のスキームは J-2 である。
\begin{enumerate}
\item 体上局所有限型の任意のスキーム。
\item Noether 完備局所環上局所有限型の任意のスキーム。
\item $\mathbf{Z}$ 上局所有限型の任意のスキーム。
\item 次元 $1$ の Noether 局所環上局所有限型の任意のスキーム。
\item 次元 $1$ の永田環上局所有限型の任意のスキーム。
\item 標数零の Dedekind 環上局所有限型の任意のスキーム。
\item など。
\end{enumerate}
\end{lemma}

\begin{proof}
補題 \ref{morphisms-lemma-J} により、上で述べた環が J-2 であることを示せばよい。
このことについては、発展代数の章、命題
\ref{more-algebra-proposition-ubiquity-J-2} を参照せよ。
\end{proof}

\section{優秀スキーム}
\label{morphisms-section-excellent}

\noindent
環が G-環かつ J-2 であるとき準優秀であるということを思い出そう。
環が準優秀かつ普遍鎖状であるとき優秀であるという。

\begin{definition}
\label{morphisms-definition-excellent}
$X$ をスキームとする。
\begin{enumerate}
\item $X$ が {\it 準優秀}であるとは、任意の $x \in X$ に対し、
アフィン開近傍 $x \in U \subset X$ であって、その環
$\mathcal{O}_X(U)$ が準優秀となるものが存在することをいう（発展代数の章、定義
\ref{more-algebra-definition-excellent} を参照）。
\item $X$ が {\it 優秀}であるとは、任意の $x \in X$ に対し、
アフィン開近傍 $x \in U \subset X$ であって、その環
$\mathcal{O}_X(U)$ が優秀となるものが存在することをいう（発展代数の章、定義
\ref{more-algebra-definition-excellent} を参照）。
\end{enumerate}
\end{definition}

\noindent
準優秀スキームは局所 Noether である（G-環は Noether 環である）。

\begin{lemma}
\label{morphisms-lemma-characterize-quasi-excellent}
$X$ をスキームとする。次の条件は同値である。
\begin{enumerate}
\item $X$ は準優秀である。
\item $X$ は G-スキームかつ J-2 である。
\end{enumerate}
\end{lemma}

\begin{proof}
(1) $\Rightarrow$ (2) は、準優秀スキームの定義
（定義 \ref{morphisms-definition-excellent}）、準優秀環の定義
（発展代数の章、定義 \ref{more-algebra-definition-excellent}）、
G-スキームの定義（性質の章、定義
\ref{properties-definition-G-scheme}）、および補題 \ref{morphisms-lemma-J}
を組み合わせれば従う。逆に、$X$ が G-スキームかつ J-2 ならば、
任意の $x \in X$ に対し、アフィン開近傍 $x \in U \subset X$ であって、
$\mathcal{O}_X(U)$ が G-環となるものを取ることができる。
補題 \ref{morphisms-lemma-J} により、環 $\mathcal{O}_X(U)$ も J-2 であり、従って準優秀である。
\end{proof}

\begin{lemma}
\label{morphisms-lemma-quasi-excellent-nagata}
準優秀スキームは永田スキームである。
\end{lemma}

\begin{proof}
発展代数の章、補題 \ref{more-algebra-lemma-quasi-excellent-nagata} を参照せよ。
\end{proof}

\begin{lemma}
\label{morphisms-lemma-characterize-excellent}
$X$ をスキームとする。次の条件は同値である。
\begin{enumerate}
\item $X$ は優秀である。
\item $X$ は準優秀かつ普遍鎖状である。
\end{enumerate}
\end{lemma}

\begin{proof}
(1) を仮定する。優秀環は準優秀なので、$X$ が準優秀であることは明らかである。
優秀環は普遍鎖状なので、補題 \ref{morphisms-lemma-universally-catenary-local} により、
$X$ が普遍鎖状であることも従う。(2) を仮定する。
$x \in X$ とする。$X$ は準優秀なので、アフィン開近傍
$x \in U \subset X$ であって、$\mathcal{O}_X(U)$ が準優秀となるものが存在する。
$X$ は普遍的に圏であるので、
補題 \ref{morphisms-lemma-universally-catenary-local} により、環 $\mathcal{O}_X(U)$ も
普遍鎖状であり、従って優秀である。
\end{proof}

\begin{lemma}
\label{morphisms-lemma-locally-excellent}
$X$ をスキームとする。次の条件は同値である。
\begin{enumerate}
\item スキーム $X$ は（準）優秀である。
\item 任意のアフィン開集合 $U \subset X$ に対し、環 $\mathcal{O}_X(U)$ は
（準）優秀である。
\item アフィン開被覆 $X = \bigcup U_i$ であって、各
$\mathcal{O}_X(U_i)$ が（準）優秀であるものが存在する。
\item 開被覆 $X = \bigcup X_j$ であって、各開部分スキーム
$X_j$ が（準）優秀であるものが存在する。
\end{enumerate}
さらに、$X$ が（準）優秀ならば、任意の開部分スキームも（準）優秀である。
\end{lemma}

\begin{proof}
補題 \ref{morphisms-lemma-characterize-quasi-excellent} により、準優秀であることは
G-スキームかつ J-2 であることと同じである。従って準優秀の場合、主張は補題
\ref{morphisms-lemma-J} および性質の章、補題
\ref{properties-lemma-locally-G-scheme} から従う。
優秀の場合には、補題 \ref{morphisms-lemma-characterize-excellent}、準優秀の場合、および補題
\ref{morphisms-lemma-universally-catenary-local} を用いる。
\end{proof}

\begin{lemma}
\label{morphisms-lemma-finite-type-excellent}
$f : X \to S$ を射とする。$S$ が（準）優秀で $f$ が局所有限型ならば、
$X$ は（準）優秀である。
\end{lemma}

\begin{proof}
発展代数の章、補題 \ref{more-algebra-lemma-finite-type-over-excellent} を参照せよ。
\end{proof}

\begin{lemma}
\label{morphisms-lemma-ubiquity-excellent}
次の型のスキームは優秀である。
\begin{enumerate}
\item 体上局所有限型の任意のスキーム。
\item Noether 完備局所環上局所有限型の任意のスキーム。
\item $\mathbf{Z}$ 上局所有限型の任意のスキーム。
\item 標数零の Dedekind 環上局所有限型の任意のスキーム。
\end{enumerate}
\end{lemma}

\begin{proof}
補題 \ref{morphisms-lemma-locally-excellent} および \ref{morphisms-lemma-finite-type-excellent} により、
上で述べた環が優秀であることを示せばよい。このことについては、
発展代数の章、命題 \ref{more-algebra-proposition-ubiquity-excellent} を参照せよ。
\end{proof}

\section{準有限射}
\label{morphisms-section-quasi-finite}

\noindent
準有限射を堅固に扱うことは、後に現れる多くの展開の基礎となる。
これは Zariski の主定理のさまざまな形、ファイバーの次元の振る舞い、
\'etale 射の降下などへつながる。本節を読む前に、代数の章、節
\ref{algebra-section-quasi-finite} の代数的結果に目を通しておくとよいだろう。

\medskip\noindent
有限型環準同型 $R \to A$ が素イデアル $\mathfrak q$ において準有限であるとは、
$\mathfrak q$ がそのファイバーの孤立点を定めることだった。代数の章、定義
\ref{algebra-definition-quasi-finite} を参照せよ。

\begin{definition}
\label{morphisms-definition-quasi-finite}
\begin{reference}
\cite[II Definition 6.2.3]{EGA}
\end{reference}
$f : X \to S$ をスキームの射とする。
\begin{enumerate}
\item $f$ が {\it 点 $x \in X$ において準有限}であるとは、アフィン近傍
$\Spec(A) = U \subset X$（$x$ の近傍）とアフィン開集合
$\Spec(R) = V \subset S$ であって、$f(U) \subset V$、環準同型
$R \to A$ は有限型、かつ $R \to A$ は $A$ の $x$ に対応する素イデアルにおいて
準有限となるものが存在することをいう（上を参照）。
\item $f$ が {\it 局所準有限}であるとは、$f$ がすべての点 $x$（$X$ の点）
において準有限であることをいう。
\item $f$ が {\it 準有限}であるとは、$f$ が有限型であり、すべての点 $x$ が
そのファイバーの孤立点であることをいう。
\end{enumerate}
\end{definition}

\noindent
局所準有限射が局所有限型であることは自明である。局所有限型の射 $f$ が
$x$ において準有限であることと、$x$ がそのファイバーの孤立点であることが
同値であることを後で示す。さらに、射が準有限となる点の集合は開である。
これは Zariski の主定理を扱う節 \ref{morphisms-section-Zariski} で示す。

\begin{lemma}
\label{morphisms-lemma-algebraic-residue-field-extension-closed-point-fibre}
$f : X \to S$ をスキームの射とする。$x \in X$ を点とし、$s = f(x)$ と置く。
$\kappa(x)/\kappa(s)$ が代数的体拡大ならば、
\begin{enumerate}
\item $x$ はそのファイバーの閉点である。
\item さらに $s$ が $S$ の閉点ならば、$x$ は $X$ の閉点である。
\end{enumerate}
\end{lemma}

\begin{proof}
第二の主張は、第一の主張から初等的な位相により従う。
第一の主張を証明するには、スキームの章、補題
\ref{schemes-lemma-fibre-topological} により、$X$ を $X_s$ で、$S$ を
$\Spec(\kappa(s))$ で置き換えてよい。従って $S = \Spec(k)$ が体のスペクトルであると
仮定してよい。この場合、任意のアフィン開集合
$\Spec(A) = U \subset X$ であって $x$ を含むものを取る。点 $x$ は素イデアル
$\mathfrak q \subset A$ に対応し、$\kappa(\mathfrak q)/k$ は代数的体拡大である。
代数の章、補題 \ref{algebra-lemma-finite-residue-extension-closed} により、
$\mathfrak q$ は極大イデアル、すなわち $x \in U$ は閉点である。
アフィン開集合は $X$ の位相の基底をなすので、$\{x\}$ は閉である。
\end{proof}

\noindent
次の補題は Hilbert の零点定理の一つの形である。

\begin{lemma}
\label{morphisms-lemma-closed-point-fibre-locally-finite-type}
$f : X \to S$ をスキームの射とする。$x \in X$ を点とし、$s = f(x)$ と置く。
$f$ が局所有限型であると仮定する。このとき、$x$ がそのファイバーの閉点であることと、
$\kappa(x)/\kappa(s)$ が有限体拡大であることは同値である。
\end{lemma}

\begin{proof}
拡大が有限ならば、上の補題
\ref{morphisms-lemma-algebraic-residue-field-extension-closed-point-fibre} により、$x$ は
そのファイバーの閉点である。逆に、$x$ がそのファイバーの閉点であると仮定する。
アフィン開集合 $\Spec(A) = U \subset X$ および
$\Spec(R) = V \subset S$ であって、$f(U) \subset V$ となるものを取る。
補題 \ref{morphisms-lemma-locally-finite-type-characterize} により、環準同型 $R \to A$ は有限型である。
$\mathfrak q \subset A$、およびそれぞれ $\mathfrak p \subset R$ を、$x$、およびそれぞれ
$s$ に対応する素イデアルとする。ファイバー環
$\overline{A} = A \otimes_R \kappa(\mathfrak p)$ を考える。
$\overline{\mathfrak q}$ を、$\overline{A}$ の素イデアルであって
$\mathfrak q$ に対応するものとする。
$x$ がそのファイバーの閉点であるという仮定から、$\overline{\mathfrak q}$ は
$\overline{A}$ の極大イデアルである。$\overline{A}$ は体
$\kappa(\mathfrak p)$ 上有限型の代数なので、Hilbert の零点定理、すなわち代数の章、定理
\ref{algebra-theorem-nullstellensatz} により、$\kappa(\overline{\mathfrak q})$ は
$\kappa(\mathfrak p)$ の有限拡大である。
$\kappa(s) = \kappa(\mathfrak p)$ かつ
$\kappa(x) = \kappa(\mathfrak q) = \kappa(\overline{\mathfrak q})$ なので、結論を得る。
\end{proof}

\begin{lemma}
\label{morphisms-lemma-base-change-closed-point-fibre-locally-finite-type}
$f : X \to S$ を局所有限型のスキームの射とする。$g : S' \to S$ を任意の射とし、
基底変換を $f' : X' \to S'$ と書く。$x' \in X'$ が点
$x \in X$ であって $X_{f(x)}$ において閉であるものに写るならば、$x'$ は
$X'_{f'(x')}$ において閉である。
\end{lemma}

\begin{proof}
剰余体 $\kappa(x')$ は
$\kappa(f'(x')) \otimes_{\kappa(f(x))} \kappa(x)$ の商である。スキームの章、補題
\ref{schemes-lemma-points-fibre-product} を参照せよ。従って、それは
$\kappa(f'(x'))$ の有限拡大である。実際、補題
\ref{morphisms-lemma-closed-point-fibre-locally-finite-type} により $\kappa(x)$ は
$\kappa(f(x))$ の有限拡大である。
この補題をもう一度適用すれば、$x'$ がそのファイバーにおいて閉であることが分かる。
\end{proof}

\begin{lemma}
\label{morphisms-lemma-residue-field-quasi-finite}
$f : X \to S$ をスキームの射とする。$x \in X$ を点とし、$s = f(x)$ と置く。
$f$ が $x$ において準有限ならば、剰余体拡大
$\kappa(x)/\kappa(s)$ は有限である。
\end{lemma}

\begin{proof}
これは代数の章、定義 \ref{algebra-definition-quasi-finite} から明らかである。
\end{proof}

\begin{lemma}
\label{morphisms-lemma-quasi-finite-at-point-characterize}
$f : X \to S$ をスキームの射とする。$x \in X$ を点とし、$s = f(x)$ と置く。
$X_s$ を $f$ の $s$ におけるファイバーとする。$f$ が局所有限型であると仮定する。
次の条件は同値である。
\begin{enumerate}
\item 射 $f$ は $x$ において準有限である。
\item 点 $x$ は $X_s$ において孤立している。
\item 点 $x$ は $X_s$ において閉であり、点 $x' \in X_s$ であって
$x' \not = x$ かつ $x$ に特殊化するものは存在しない。
\item $\Spec(A) = U \subset X$、$\Spec(R) = V \subset S$ を、
$f(U) \subset V$ かつ $x \in U$ が $\mathfrak q \subset A$ に対応するような
任意の一対のアフィン開集合とするとき、環準同型 $R \to A$ は
$\mathfrak q$ において準有限である。
\end{enumerate}
\end{lemma}

\begin{proof}
$f$ が $x$ において準有限であると仮定する。仮定により、開集合
$U \subset X$、$V \subset S$ であって、$f(U) \subset V$、$x \in U$、かつ
$x$ が $U_s$ の孤立点となるものが存在する。従って $\{x\} \subset U_s$ は開集合である。
$U_s = U \cap X_s \subset X_s$ も開なので、$\{x\} \subset X_s$ も開集合である。
従って $x$ は $X_s$ の孤立点である。

\medskip\noindent
補題 \ref{morphisms-lemma-ubiquity-Jacobson-schemes}（および補題
\ref{morphisms-lemma-base-change-finite-type}）により、$X_s$ は Jacobson スキームであることに注意する。
$x$ が $X_s$ において孤立している、すなわち $\{x\} \subset X_s$ が開ならば、
$\{x\}$ は閉点を含む（Jacobson 性質による）ので、$x$ は $X_s$ において閉である。
点 $x' \in X_s$ であって、$x$ と異なり $x$ に特殊化するものが存在しないことは明らかである。

\medskip\noindent
$x$ が $X_s$ において閉であり、点 $x' \in X_s$ であって、
$x$ と異なり $x$ に特殊化するものが存在しないと仮定する。アフィン開集合の対
$\Spec(A) = U \subset X$、$\Spec(R) = V \subset S$ であって、
$f(U) \subset V$ かつ $x \in U$ となるものを考える。
$\mathfrak q \subset A$ を $x$ に対応するもの、$\mathfrak p \subset R$ を
$s$ に対応するものとする。補題 \ref{morphisms-lemma-locally-finite-type-characterize} により、
環準同型 $R \to A$ は有限型である。ファイバー環
$\overline{A} = A \otimes_R \kappa(\mathfrak p)$ を考える。
$\overline{\mathfrak q}$ を、$\overline{A}$ の素イデアルであって
$\mathfrak q$ に対応するものとする。$\Spec(\overline{A})$ はファイバー
$X_s$ の開部分スキームなので、$\overline{q}$ は $\overline{A}$ の極大イデアルであり、
$\Spec(\overline{A})$ には $\overline{\mathfrak q}$ に特殊化する点が存在しない。
これは $\dim(\overline{A}_{\overline{q}}) = 0$ を含意する。従って代数の章、定義
\ref{algebra-definition-quasi-finite} により、$R \to A$ は $\mathfrak q$ において準有限、
すなわち定義により $X \to S$ は $x$ において準有限である。

\medskip\noindent
ここまでで、条件 (1)--(3) がすべて同値であることを示した。
(4) が (1) を含意することは明らかである。また、$x$ が $X_s$ の孤立点ならば、
それは $U_s$ の孤立点でもある（その点を含む任意の開集合 $U$ に対し）ので、
(2) が (4) を含意することも明らかである。
\end{proof}

\begin{lemma}
\label{morphisms-lemma-finite-fibre}
$f : X \to S$ をスキームの射とする。$s \in S$ とする。次を仮定する。
\begin{enumerate}
\item $f$ は局所有限型である。
\item $f^{-1}(\{s\})$ は有限集合である。
\end{enumerate}
このとき $X_s$ は有限離散位相空間であり、$f$ は $X$ の $s$ 上にある各点において
準有限である。
\end{lemma}

\begin{proof}
$T$ を、(a) 体 $k$ 上局所有限型であり、(b) 有限個の点を持つスキームとする。
補題 \ref{morphisms-lemma-ubiquity-Jacobson-schemes} により、$T$ は Jacobson スキームである。
有限 Jacobson 空間は離散的である。位相の章、補題
\ref{topology-lemma-finite-jacobson} を参照せよ。この考察を、ファイバー
$X_s$（これは $\Spec(\kappa(s))$ 上局所有限型である）に適用すると第一の主張を得る。
最後に、補題 \ref{morphisms-lemma-quasi-finite-at-point-characterize} を適用すると第二の主張を得る。
\end{proof}

\begin{lemma}
\label{morphisms-lemma-locally-quasi-finite-fibres}
\begin{slogan}
離散的なファイバーを持つ有限型射は準有限である。
\end{slogan}
$f : X \to S$ をスキームの射とする。$f$ が局所有限型であると仮定する。
次の条件は同値である。
\begin{enumerate}
\item $f$ は局所準有限である。
\item 任意の $s \in S$ に対し、ファイバー $X_s$ は離散位相空間である。
\item 任意の射 $\Spec(k) \to S$（ここで $k$ は体）に対し、基底変換 $X_k$ の
台位相空間は離散的である。
\end{enumerate}
\end{lemma}

\begin{proof}
(3) が (2) を含意することは直ちに分かる。補題
\ref{morphisms-lemma-quasi-finite-at-point-characterize} により、(2) と (1) は同値である。
(2) を仮定し、$\Spec(k) \to S$ を (3) のような射とする。
$s \in S$ を $\Spec(k) \to S$ の像と書く。このとき $X_k$ は $X_s$ の、射
$\Spec(k) \to \Spec(\kappa(s))$ による基底変換である。
従って、補題 \ref{morphisms-lemma-base-change-closed-point-fibre-locally-finite-type} により、
$X_k$ のすべての点は閉である。$X_k \to \Spec(k)$ は局所有限型なので
（補題 \ref{morphisms-lemma-base-change-finite-type} による）、補題
\ref{morphisms-lemma-quasi-finite-at-point-characterize} を適用すると、$X_k$ のすべての点が
孤立していること、すなわち $X_k$ の台位相空間が離散的であることが分かる。
\end{proof}

\begin{lemma}
\label{morphisms-lemma-quasi-finite-locally-quasi-compact}
$f : X \to S$ をスキームの射とする。$f$ が準有限であることと、$f$ が局所準有限かつ
準コンパクトであることは同値である。
\end{lemma}

\begin{proof}
$f$ が準有限であると仮定する。定義 \ref{morphisms-definition-finite-type} により準コンパクトである。
$x \in X$ とする。補題 \ref{morphisms-lemma-quasi-finite-at-point-characterize} により、$f$ は
$x$ において準有限であることが分かる。従って $f$ は準コンパクトかつ局所準有限である。

\medskip\noindent
$f$ が準コンパクトかつ局所準有限であると仮定する。このとき $f$ は有限型である。
$x \in X$ を点とする。補題 \ref{morphisms-lemma-quasi-finite-at-point-characterize} により、
$x$ がそのファイバーの孤立点であることが分かる。補題が証明された。
\end{proof}

\begin{lemma}
\label{morphisms-lemma-quasi-finite}
$f : X \to S$ をスキームの射とする。次の条件は同値である。
\begin{enumerate}
\item $f$ は準有限である。
\item $f$ は局所有限型かつ準コンパクトで、有限ファイバーを持つ。
\end{enumerate}
\end{lemma}

\begin{proof}
$f$ が準有限であると仮定する。特に $f$ は局所有限型かつ準コンパクトである
（有限型だからである）。$s \in S$ とする。すべての $x \in X_s$ は
$X_s$ において孤立しているので、$X_s = \bigcup_{x \in X_s} \{x\}$ は開被覆である。
$f$ は準コンパクトなので、ファイバー $X_s$ は準コンパクトである。
従って $X_s$ は有限である。

\medskip\noindent
逆に、$f$ が局所有限型かつ準コンパクトで、有限ファイバーを持つと仮定する。
補題 \ref{morphisms-lemma-finite-fibre} により局所準有限である。従って補題
\ref{morphisms-lemma-quasi-finite-locally-quasi-compact} により準有限である。
\end{proof}

\noindent
環準同型 $R \to A$ が準有限であるとは、それが有限型で、$A$ の {\it すべての}
素イデアルにおいて準有限であることだった。代数の章、定義
\ref{algebra-definition-quasi-finite} を参照せよ。

\begin{lemma}
\label{morphisms-lemma-locally-quasi-finite-characterize}
$f : X \to S$ をスキームの射とする。次の条件は同値である。
\begin{enumerate}
\item 射 $f$ は局所準有限である。
\item アフィン開集合の任意の対 $U \subset X$、$V \subset S$ であって
$f(U) \subset V$ となるものに対し、環準同型
$\mathcal{O}_S(V) \to \mathcal{O}_X(U)$ は準有限である。
\item 開被覆 $S = \bigcup_{j \in J} V_j$ および開被覆
$f^{-1}(V_j) = \bigcup_{i \in I_j} U_i$ であって、各射
$U_i \to V_j$（$j\in J, i\in I_j$）が局所準有限となるものが存在する。
\item アフィン開被覆 $S = \bigcup_{j \in J} V_j$ およびアフィン開被覆
$f^{-1}(V_j) = \bigcup_{i \in I_j} U_i$ であって、環準同型
$\mathcal{O}_S(V_j) \to \mathcal{O}_X(U_i)$ が、すべての
$j\in J, i\in I_j$ に対し準有限となるものが存在する。
\end{enumerate}
さらに、$f$ が局所準有限ならば、任意の開部分スキーム
$U \subset X$、$V \subset S$ であって $f(U) \subset V$ となるものに対し、制限
$f|_U : U \to V$ は局所準有限である。
\end{lemma}

\begin{proof}
環準同型 $R \to A$ に対し、$P(R \to A)$ が「$R \to A$ は準有限である」を
意味するものと定義する（補題の前の注意を参照）。$P$ は環準同型の局所的性質であると主張する。
定義 \ref{morphisms-definition-property-local} の条件 (a), (b), (c) を確認する。
補題 \ref{morphisms-lemma-locally-finite-type-characterize} の証明において、「有限型である」という
性質について (a), (b), (c) が成り立つことを見た。有限型環準同型
$R \to A$ に対し、$R \to A$ が $\mathfrak q$ において準有限であるという性質は、局所環
$A_{\mathfrak q}$ の $R_{\mathfrak p}$-代数としての構造のみに依存することに注意する。ここで
$\mathfrak p = R \cap \mathfrak q$ である（通常の記法の濫用をしている）。
これらの注意から、定義 \ref{morphisms-definition-property-local} の (a), (b), (c) は直ちに従う。
例えば、$R \to A$ を環準同型とし、環準同型 $R \to A_{a_i}$ がすべて準有限であって、
$a_1, \ldots, a_n \in A$ が単位イデアルを生成すると仮定する。
$R \to A$ は有限型であると結論できる。また、任意の素イデアル
$\mathfrak q \subset A$ に対し、局所環 $A_{\mathfrak q}$ は $R$-代数として局所環
$(A_{a_i})_{\mathfrak q_i}$ と同型である。ここで $i$ はある添字であり、
$\mathfrak q_i \subset A_{a_i}$ である。従って $R \to A$ は $\mathfrak q$ において準有限である。

\medskip\noindent
以上から、前段落の $P$ に対して補題 \ref{morphisms-lemma-locally-P} を適用できる。
従って、$f$ が局所準有限であることと、$f$ が局所的に型 $P$ であることが同値であると
証明すれば十分である。$P(R \to A)$ は「$R \to A$ は準有限である」、すなわち
$R \to A$ が $A$ のすべての素イデアルにおいて準有限であることを意味するので、これは補題
\ref{morphisms-lemma-quasi-finite-at-point-characterize} から従う。
\end{proof}

\begin{lemma}
\label{morphisms-lemma-composition-quasi-finite}
局所準有限な二つの射の合成は局所準有限である。準有限射についても同じことが成り立つ。
\end{lemma}

\begin{proof}
補題 \ref{morphisms-lemma-locally-quasi-finite-characterize} の証明において、
$P = $「準有限」が環準同型の局所的性質であり、スキームの射が局所準有限であることと、
定義 \ref{morphisms-definition-locally-P} の意味で局所的に型 $P$ であることが同値であると分かった。
従って補題の第一の主張は、補題 \ref{morphisms-lemma-composition-type-P} と、準有限であるという
環準同型の性質が合成の下で安定であるという事実を合わせれば従う。後者については、
代数の章、補題 \ref{algebra-lemma-quasi-finite-composition} を参照せよ。
以上と補題 \ref{morphisms-lemma-quasi-finite-locally-quasi-compact}、および準コンパクト射の合成が
準コンパクトであるという事実（スキームの章、補題
\ref{schemes-lemma-composition-quasi-compact} を参照）から、準有限射の合成が準有限であると分かる。
\end{proof}

\noindent
後で（補題 \ref{morphisms-lemma-quasi-finite-points-open}）、次の補題の集合 $U$ が開であることを示す。

\begin{lemma}
\label{morphisms-lemma-base-change-quasi-finite}
\begin{slogan}
（局所）準有限射は基底変換の下で安定である。
\end{slogan}
$f : X \to S$ をスキームの射とする。$g : S' \to S$ をスキームの射とする。
$f' : X' \to S'$ を $f$ の $g$ による基底変換と書き、射影を
$g' : X' \to X$ と書く。$X$ が $S$ 上局所有限型であると仮定する。
\begin{enumerate}
\item $U \subset X$（それぞれ $U' \subset X'$）を、$f$（それぞれ $f'$）が
準有限となる点の集合とする。このとき $U' = U \times_S S' = (g')^{-1}(U)$ である。
\item 局所準有限射の基底変換は局所準有限である。
\item 準有限射の基底変換は準有限である。
\end{enumerate}
\end{lemma}

\begin{proof}
第一および第二の主張は、対応する代数的結果、すなわち代数の章、補題
\ref{algebra-lemma-quasi-finite-base-change} から従う（補題
\ref{morphisms-lemma-base-change-finite-type} により $f'$ も局所有限型であることと合わせる）。
以上、補題 \ref{morphisms-lemma-quasi-finite-locally-quasi-compact}、および準コンパクト射の基底変換が
準コンパクトであるという事実（スキームの章、補題
\ref{schemes-lemma-quasi-compact-preserved-base-change} を参照）から、準有限射の基底変換が
準有限であることが分かる。
\end{proof}

\begin{lemma}
\label{morphisms-lemma-quasi-finite-at-a-finite-number-of-points}
$f : X \to S$ を有限型のスキームの射とする。$s \in S$ とする。
$X$ の $s$ 上にあって $f$ が準有限となる点は高々有限個である。
\end{lemma}

\begin{proof}
ファイバー $X_s$ は体上有限型のスキームなので Noether である
（補題 \ref{morphisms-lemma-finite-type-noetherian}）。従って $X_s$ の位相は Noether であり
（性質の章、補題 \ref{properties-lemma-Noetherian-topology}）、孤立点を高々有限個しか
持てない（初等的な位相による）。従って補題
\ref{morphisms-lemma-quasi-finite-at-point-characterize} から主張が従う。
\end{proof}

\begin{lemma}
\label{morphisms-lemma-monomorphism-loc-finite-type-loc-quasi-finite}
$f : X \to Y$ をスキームの射とする。$f$ が局所有限型かつモノ射ならば、$f$ は
分離的かつ局所準有限である。
\end{lemma}

\begin{proof}
スキームの章、補題 \ref{schemes-lemma-monomorphism-separated} により、モノ射は分離的である。
モノ射は単射なので、例えば補題 \ref{morphisms-lemma-quasi-finite-at-point-characterize} により、
$f$ はすべての $x \in X$ において準有限である。
\end{proof}

\begin{lemma}
\label{morphisms-lemma-immersion-locally-quasi-finite}
任意の埋め込みは局所準有限である。
\end{lemma}

\begin{proof}
開埋め込みは局所同型であり、閉埋め込みは明らかに準有限なので、これは成り立つ。
\end{proof}

\begin{lemma}
\label{morphisms-lemma-permanence-quasi-finite}
$X \to Y$ を基底スキーム $S$ 上のスキームの射とする。$x \in X$ とする。
$X \to S$ が $x$ において準有限ならば、$X \to Y$ は $x$ において準有限である。
$X$ が $S$ 上局所準有限ならば、$X \to Y$ は局所準有限である。
\end{lemma}

\begin{proof}
補題 \ref{morphisms-lemma-locally-quasi-finite-characterize} により、これは次の代数的事実に翻訳される。
環準同型 $A \to B \to C$ が与えられ、$A \to C$ が準有限ならば、$B \to C$ は準有限である。
これは代数の章、補題 \ref{algebra-lemma-four-rings} において
$R = A$、$S = S' = C$、$R' = B$ とすれば従う。
\end{proof}

\begin{lemma}
\label{morphisms-lemma-quasi-finite-local-source}
$f : X \to Y$ および $g : Y \to S$ をスキームの射とする。$f$ が全射、
$g \circ f$ が局所準有限、かつ $g$ が局所有限型ならば、$g : Y \to S$ は局所準有限である。
\end{lemma}

\begin{proof}
$x \in X$ とし、その像を $y \in Y$ および $s \in S$ とする。
$g \circ f$ は局所準有限なので、補題 \ref{morphisms-lemma-residue-field-quasi-finite} により、拡大
$\kappa(x)/\kappa(s)$ は有限である。従って $\kappa(y)/\kappa(s)$ は有限である。
従って補題 \ref{morphisms-lemma-algebraic-residue-field-extension-closed-point-fibre} により、
$y$ は $Y_s$ の閉点である。$f$ は全射なので、$Y$ のすべての点が $S$ 上の
そのファイバーにおいて閉であることが分かる。従って補題
\ref{morphisms-lemma-quasi-finite-at-point-characterize} により、$g$ はすべての点で準有限である。
\end{proof}

\section{有限表示の射}
\label{morphisms-section-finite-presentation}

\noindent
環準同型 $R \to A$ が有限表示であるとは、$A$ が
$R[x_1, \ldots, x_n]/(f_1, \ldots, f_m)$ と $R$-代数として同型であることだった。
ここで $n, m$ はある整数、$f_j$ はある多項式である。代数の章、定義
\ref{algebra-definition-finite-type} を参照せよ。

\begin{definition}
\label{morphisms-definition-finite-presentation}
$f : X \to S$ をスキームの射とする。
\begin{enumerate}
\item $f$ が {\it $x \in X$ において有限表示} であるとは、アフィン開集合
$\Spec(A) = U \subset X$ であって $x$ の近傍であるものと、アフィン開集合
$\Spec(R) = V \subset S$ であって
$f(U) \subset V$ となり、誘導される環準同型 $R \to A$ が有限表示であるものが
存在することをいう。
\item $f$ が {\it 局所有限表示} であるとは、それが $X$ の各点において
有限表示であることをいう。
\item $f$ が {\it 有限表示} であるとは、それが局所有限表示、準コンパクト、
かつ準分離的であることをいう。
\end{enumerate}
\end{definition}

\noindent
有限表示の射は、局所有限表示である準コンパクト射だけでは
{\bf ない}ことに注意せよ。後に、局所有限表示の射は、任意の有向系
$T_i$（$S$ 上のアフィンスキームからなる）に対して
$$
\colim \Mor_S(T_i, X) = \Mor_S(\lim T_i, X)
$$
を満たす射として特徴付けられる。極限の章、命題
\ref{limits-proposition-characterize-locally-finite-presentation} を参照せよ。
極限の章、節 \ref{limits-section-descending-relative} では、$S = \lim_i S_i$ が
アフィンスキームの極限ならば、任意のスキーム $X$ であって $S$ 上有限表示のものは、
スキーム $X_i$ であって $S_i$ 上のものへ、ある $i$ に対して降下することを示す。

\begin{lemma}
\label{morphisms-lemma-locally-finite-presentation-characterize}
$f : X \to S$ をスキームの射とする。次の条件は同値である。
\begin{enumerate}
\item 射 $f$ は局所有限表示である。
\item 任意のアフィン開集合 $U \subset X$、$V \subset S$ であって
$f(U) \subset V$ となるものに対し、環準同型
$\mathcal{O}_S(V) \to \mathcal{O}_X(U)$ は有限表示である。
\item 開被覆 $S = \bigcup_{j \in J} V_j$ と開被覆
$f^{-1}(V_j) = \bigcup_{i \in I_j} U_i$ であって、各射
$U_i \to V_j$（$j\in J, i\in I_j$）が局所有限表示となるものが存在する。
\item アフィン開被覆 $S = \bigcup_{j \in J} V_j$ とアフィン開被覆
$f^{-1}(V_j) = \bigcup_{i \in I_j} U_i$ であって、環準同型
$\mathcal{O}_S(V_j) \to \mathcal{O}_X(U_i)$ がすべての
$j\in J, i\in I_j$ に対し有限表示となるものが存在する。
\end{enumerate}
さらに、$f$ が局所有限表示ならば、任意の開部分スキーム
$U \subset X$、$V \subset S$ であって $f(U) \subset V$ となるものに対し、制限
$f|_U : U \to V$ は局所有限表示である。
\end{lemma}

\begin{proof}
性質「$R \to A$ が有限表示である」が局所的であることを示せば、補題
\ref{morphisms-lemma-locally-P-characterize} から従う。定義
\ref{morphisms-definition-property-local} の条件 (a), (b), (c) を確認する。
有限表示であることは基底変換の下で安定であるから、代数の章、補題
\ref{algebra-lemma-base-change-finiteness} により (a) が成り立つ。
有限表示であることは合成の下で安定であり、任意の環 $R$ に対して環準同型
$R \to R_f$ は明らかに有限表示であるから、代数の章、補題
\ref{algebra-lemma-compose-finite-type} により (b) が成り立つ。
最後に、代数の章、補題 \ref{algebra-lemma-cover-upstairs} により (c) が成り立つ。
\end{proof}

\begin{lemma}
\label{morphisms-lemma-composition-finite-presentation}
二つの局所有限表示射の合成は局所有限表示である。
有限表示射についても同じことが成り立つ。
\end{lemma}

\begin{proof}
補題 \ref{morphisms-lemma-locally-finite-presentation-characterize} の証明において、
有限表示であることが環準同型の局所的性質であることを見た。
従って補題の第一の主張は、補題 \ref{morphisms-lemma-composition-type-P} と、
有限表示であることが合成の下で安定な環準同型の性質であるという事実との組合せから従う。
代数の章、補題 \ref{algebra-lemma-compose-finite-type} を参照せよ。
以上と、準コンパクトかつ準分離的な射の合成が準コンパクトかつ準分離的であるという事実
（スキームの章、補題 \ref{schemes-lemma-composition-quasi-compact} および
\ref{schemes-lemma-separated-permanence} を参照）により、
有限表示射の合成は有限表示である。
\end{proof}

\begin{lemma}
\label{morphisms-lemma-base-change-finite-presentation}
局所有限表示射の基底変換は局所有限表示である。
有限表示射についても同じことが成り立つ。
\end{lemma}

\begin{proof}
補題 \ref{morphisms-lemma-locally-finite-presentation-characterize} の証明において、
有限表示であることが環準同型の局所的性質であることを見た。
従って補題の第一の主張は、補題 \ref{morphisms-lemma-composition-type-P} と、
有限表示であることが基底変換の下で安定な環準同型の性質であるという事実との組合せから従う。
代数の章、補題 \ref{algebra-lemma-base-change-finiteness} を参照せよ。
以上と、準コンパクトかつ準分離的な射の基底変換が準コンパクトかつ準分離的であるという事実
（スキームの章、補題 \ref{schemes-lemma-quasi-compact-preserved-base-change} および
\ref{schemes-lemma-separated-permanence} を参照）により、
有限表示射の基底変換は有限表示射である。
\end{proof}

\begin{lemma}
\label{morphisms-lemma-open-immersion-locally-finite-presentation}
任意の開埋め込みは局所有限表示である。
\end{lemma}

\begin{proof}
開埋め込みは局所同型だからである。
\end{proof}

\begin{lemma}
\label{morphisms-lemma-quasi-compact-open-immersion-finite-presentation}
任意の開埋め込みが有限表示であることと、それが準コンパクトであることは同値である。
\end{lemma}

\begin{proof}
開埋め込みが局所有限表示であることを、補題
\ref{morphisms-lemma-open-immersion-locally-finite-presentation} で見た。
埋め込みが分離的、従って準分離的であることを、スキームの章、補題
\ref{schemes-lemma-immersions-monomorphisms} で見た。
このことと定義 \ref{morphisms-definition-finite-presentation} から補題が従う。
\end{proof}

\begin{lemma}
\label{morphisms-lemma-closed-immersion-finite-presentation}
\begin{slogan}
有限表示の閉埋め込みは、有限型準連接イデアル層に対応する。
\end{slogan}
閉埋め込み $i : Z \to X$ が有限表示であるための必要十分条件は、
付随する準連接イデアル層
$\mathcal{I} = \Ker(\mathcal{O}_X \to i_*\mathcal{O}_Z)$
が（$\mathcal{O}_X$-加群として）有限型であることである。
\end{lemma}

\begin{proof}
任意のアフィン開集合 $\Spec(R) \subset X$ 上で
$i^{-1}(\Spec(R)) = \Spec(R/I)$ かつ
$\mathcal{I} = \widetilde{I}$ である。さらに、$\mathcal{I}$ が有限型であることと、
$I$ が有限 $R$-加群であることは、そのような各アフィン開集合に対して同値である
（性質の章、補題 \ref{properties-lemma-finite-type-module} を参照）。
また、$R/I$ が $R$ 上有限表示であることと、$I$ が有限 $R$-加群であることは同値である。
従って主張を得る。
\end{proof}

\begin{lemma}
\label{morphisms-lemma-finite-presentation-finite-type}
局所有限表示射は局所有限型である。有限表示射は有限型である。
\end{lemma}

\begin{proof}
省略する。
\end{proof}

\begin{lemma}
\label{morphisms-lemma-noetherian-finite-type-finite-presentation}
\begin{slogan}
局所 Noether 基底上では、有限型は有限表示である。
\end{slogan}
$f : X \to S$ を射とする。
\begin{enumerate}
\item $S$ が局所 Noether で $f$ が局所有限型ならば、$f$ は局所有限表示である。
\item $S$ が局所 Noether で $f$ が有限型ならば、$f$ は有限表示である。
\end{enumerate}
\end{lemma}

\begin{proof}
第一の主張は、Noether 環上有限型の環が有限表示であるという事実から従う。
代数の章、補題 \ref{algebra-lemma-Noetherian-finite-type-is-finite-presentation} を参照せよ。
$f$ が有限型であり、$S$ が局所 Noether であると仮定する。
このとき $f$ は準コンパクトであり、(1) により局所有限表示である。
従って、$f$ が準分離的であることを示せば十分である。
これは補題 \ref{morphisms-lemma-finite-type-Noetherian-quasi-separated}
（および補題 \ref{morphisms-lemma-finite-presentation-finite-type}）から従う。
\end{proof}

\begin{lemma}
\label{morphisms-lemma-finite-presentation-quasi-compact-quasi-separated}
$S$ を準コンパクトかつ準分離的なスキームとする。
$X$ が $S$ 上有限表示ならば、$X$ は準コンパクトかつ準分離的である。
\end{lemma}

\begin{proof}
省略する。
\end{proof}

\begin{lemma}
\label{morphisms-lemma-finite-presentation-permanence}
$f : X \to Y$ を $S$ 上のスキームの射とする。
\begin{enumerate}
\item $X$ が $S$ 上局所有限表示であり、$Y$ が $S$ 上局所有限型ならば、
$f$ は局所有限表示である。
\item $X$ が $S$ 上有限表示であり、$Y$ が準分離的かつ $S$ 上局所有限型ならば、
$f$ は有限表示である。
\end{enumerate}
\end{lemma}

\begin{proof}
(1) の証明。補題 \ref{morphisms-lemma-locally-finite-presentation-characterize} により、
これは次の代数的事実に帰着する。環準同型 $A \to B \to C$ が与えられ、
$A \to C$ が有限表示、$A \to B$ が有限型ならば、$B \to C$ は有限表示である。
代数の章、補題 \ref{algebra-lemma-compose-finite-type} を参照せよ。

\medskip\noindent
(2) は (1) と、スキームの章、補題 \ref{schemes-lemma-compose-after-separated} および
\ref{schemes-lemma-quasi-compact-permanence} から従う。
\end{proof}

\begin{lemma}
\label{morphisms-lemma-diagonal-morphism-finite-type}
$f : X \to Y$ を、対角射 $\Delta : X \to X \times_Y X$ をもつスキームの射とする。
$f$ が局所有限型ならば、$\Delta$ は局所有限表示である。
$f$ が準分離的かつ局所有限型ならば、$\Delta$ は有限表示である。
\end{lemma}

\begin{proof}
$\Delta$ は $X$ 上のスキームの射であることに注意する（その構造射は第二射影
$X \times_Y X \to X$ である）。
$f$ が局所有限型であると仮定する。$X$ は $X$ 上有限表示であり、
$X \times_Y X$ は $X$ 上局所有限型であることに注意する
（補題 \ref{morphisms-lemma-base-change-finite-type} による）。
従って第一の主張は補題 \ref{morphisms-lemma-finite-presentation-permanence} から従う。
第二の主張は、第一の主張、各定義、および対角射がモノ射、従って分離的であるという事実
（スキームの章、補題 \ref{schemes-lemma-monomorphism-separated}）から従う。
\end{proof}

\section{構成可能集合}
\label{morphisms-section-constructible}

\noindent
スキームの構成可能集合および局所構成可能集合については、性質の章、節
\ref{properties-section-constructible} で論じた。本節では、（局所）構成可能集合の
像および逆像に関するいくつかの結果を証明する。主結果は Chevalley の定理であり、
有限表示射による局所構成可能集合の像が局所構成可能であることを述べる。

\begin{lemma}
\label{morphisms-lemma-inverse-image-constructible}
$f : X \to Y$ をスキームの射とする。$E \subset Y$ を部分集合とする。
$E$ が $Y$ において（局所）構成可能ならば、$f^{-1}(E)$ は $X$ において
（局所）構成可能である。
\end{lemma}

\begin{proof}
任意の構成可能部分集合の逆像が構成可能であることを示すには、任意の
逆コンパクトな開集合 $V$ であって $Y$ に含まれるものの逆像が $X$ において
逆コンパクトであることを示せば十分である。
位相の章、補題 \ref{topology-lemma-inverse-images-constructibles} を参照せよ。
$V$ が $Y$ において逆コンパクトであることの意味は、開埋め込み $V \to Y$ が
準コンパクトであることにほかならない。従って基底変換
$f^{-1}(V) = X \times_Y V \to X$ も準コンパクトである。
スキームの章、補題 \ref{schemes-lemma-quasi-compact-preserved-base-change} を参照せよ。
従って $f^{-1}(V)$ は $X$ において逆コンパクトである。
$E$ が $Y$ において局所構成可能であると仮定する。$x \in X$ を選ぶ。
アフィン近傍 $V$（$f(x)$ のもの）と、アフィン近傍 $U \subset X$（$x$ のもの）であって
$f(U) \subset V$ となるものを選ぶ。従って $f|_U : U \to V$ を $V$ への射とみなす。
性質の章、補題 \ref{properties-lemma-locally-constructible} により、
$E \cap V$ は $V$ において構成可能である。構成可能の場合により、
$(f|_U)^{-1}(E \cap V)$ は $U$ において構成可能である。
$(f|_U)^{-1}(E \cap V) = f^{-1}(E) \cap U$ だから、主張を得る。
\end{proof}

\begin{lemma}
\label{morphisms-lemma-chevalley}
$f : X \to Y$ をスキームの射とする。次を仮定する。
\begin{enumerate}
\item $f$ は準コンパクトかつ局所有限表示である。
\item $Y$ は準コンパクトかつ準分離的である。
\end{enumerate}
このとき、$X$ の任意の構成可能部分集合の像は $Y$ において構成可能である。
\end{lemma}

\begin{proof}
性質の章、補題
\ref{properties-lemma-constructible-quasi-compact-quasi-separated} により、
$Y$ がアフィンの場合にこの補題を証明すれば十分である。この場合 $X$ は準コンパクトである。
従って $X = U_1 \cup \ldots \cup U_n$ と書ける。ここで各 $U_i$ は
$X$ のアフィン開集合である。$E \subset X$ が構成可能ならば、各
$E \cap U_i$ も構成可能である。位相の章、補題
\ref{topology-lemma-open-immersion-constructible-inverse-image} を参照せよ。
$f(E) = \bigcup f(E \cap U_i)$ であり、構成可能集合の有限合併は構成可能だから、
$X$ がアフィンの場合に帰着する。この場合の結果は、代数の章、定理
\ref{algebra-theorem-chevalley} である。
\end{proof}

\begin{theorem}[Chevalley の定理]
\label{morphisms-theorem-chevalley}
\begin{reference}
\cite[IV, Theorem 1.8.4]{EGA}
\end{reference}
$f : X \to Y$ をスキームの射とする。$f$ が準コンパクトかつ局所有限表示であると仮定する。
このとき、任意の局所構成可能部分集合の像は局所構成可能である。
\end{theorem}

\begin{proof}
$E \subset X$ を局所構成可能とする。$f(E)$ も局所構成可能であることを示さなければならない。
$f(E) \cap V$ が、任意のアフィン開集合 $V \subset Y$ に対して構成可能であることを示す。
従って $Y$ がアフィンの場合に帰着する。この場合 $X$ は準コンパクトである。
従って $X = U_1 \cup \ldots \cup U_n$ と書ける。ここで各 $U_i$ は
$X$ のアフィン開集合である。$E \subset X$ が局所構成可能ならば、
各 $E \cap U_i$ は構成可能である。性質の章、補題
\ref{properties-lemma-locally-constructible} を参照せよ。
$f(E) = \bigcup f(E \cap U_i)$ であり、構成可能集合の有限合併は構成可能だから、
$X$ がアフィンの場合に帰着する。この場合の結果は、代数の章、定理
\ref{algebra-theorem-chevalley} である。
\end{proof}

\begin{lemma}
\label{morphisms-lemma-constructible-containing-open}
$X$ をスキームとする。$x \in X$ とする。$E \subset X$ を局所構成可能部分集合とする。
$\{x' \mid x' \leadsto x\} \subset E$ ならば、$E$ は $x$ の開近傍を含む。
\end{lemma}

\begin{proof}
$\{x' \mid x' \leadsto x\} \subset E$ と仮定する。$X$ はアフィンであると仮定してよい。
この場合 $E$ は構成可能である。性質の章、補題
\ref{properties-lemma-locally-constructible} を参照せよ。
特に補集合 $E^c$ も構成可能である。代数の章、補題
\ref{algebra-lemma-constructible-is-image} により、アフィンスキームの射
$f : Y \to X$ であって $E^c = f(Y)$ となるものを見いだせる。
$Z \subset X$ を $f$ のスキーム論的像とする。
補題 \ref{morphisms-lemma-reach-points-scheme-theoretic-image} と仮定
$\{x' \mid x' \leadsto x\} \subset E$ により、$x \not \in Z$ である。
従って $X \setminus Z \subset E$ は $x$ の開近傍であり、$E$ に含まれる。
\end{proof}

\section{開射}
\label{morphisms-section-open}

\begin{definition}
\label{morphisms-definition-open}
$f : X \to S$ を射とする。
\begin{enumerate}
\item $f$ が {\it 開} であるとは、台となる位相空間上の写像が開であることをいう。
\item $f$ が {\it 普遍開} であるとは、スキームの任意の射 $S' \to S$ に対し、
基底変換 $f' : X_{S'} \to S'$ が開であることをいう。
\end{enumerate}
\end{definition}

\noindent
位相の章、補題 \ref{topology-lemma-closed-open-map-specialization} によれば、
ある種の位相空間の開写像に沿って一般化が持ち上がる。実際、スキームの任意の
開射に沿って一般化が持ち上がる（補題 \ref{morphisms-lemma-open-generizing} を参照）。
また、スキームの平坦射に沿って一般化が持ち上がることを後で見る
（補題 \ref{morphisms-lemma-generalizations-lift-flat}）。このことから逆に、その射が開であると
結論できる場合がある。

\begin{lemma}
\label{morphisms-lemma-locally-finite-presentation-universally-open}
$f : X \to S$ を射とする。
\begin{enumerate}
\item $f$ が局所有限表示であり、一般化が $f$ に沿って持ち上がるならば、
$f$ は開である。
\item $f$ が局所有限表示であり、一般化が $f$ のすべての基底変換に沿って
持ち上がるならば、$f$ は普遍開である。
\end{enumerate}
\end{lemma}

\begin{proof}
第一の主張を証明すれば十分である。$X$ と $S$ がともにアフィンの場合に帰着する。
この場合、結果は代数の章、補題
\ref{algebra-lemma-going-up-down-specialization} および命題
\ref{algebra-proposition-fppf-open} から従う。
\end{proof}

\noindent
有限表示平坦射の場合については、補題 \ref{morphisms-lemma-fppf-open} も参照せよ。

\begin{lemma}
\label{morphisms-lemma-composition-open}
（普遍）開射の合成は（普遍）開である。
\end{lemma}

\begin{proof}
省略する。
\end{proof}

\begin{lemma}
\label{morphisms-lemma-scheme-over-field-universally-open}
$k$ を体とする。$X$ を $k$ 上のスキームとする。
構造射 $X \to \Spec(k)$ は普遍開である。
\end{lemma}

\begin{proof}
$S \to \Spec(k)$ を射とする。基底変換 $X_S \to S$ が開であることを示さなければならない。
問題は $S$ および $X$ 上局所的だから、$S$ と $X$ はアフィンであると仮定してよい。
この場合、結果は代数の章、補題
\ref{algebra-lemma-map-into-tensor-algebra-open} である。
\end{proof}

\begin{lemma}
\label{morphisms-lemma-open-generizing}
\begin{reference}
含意 (a) $\Rightarrow$ (b)（\cite[IV, Corollary 1.10.4]{EGA}）から従う。
\end{reference}
$\varphi : X \to Y$ をスキームの射とする。$\varphi$ が開ならば、$\varphi$ は
一般化的である（すなわち、一般化が $\varphi$ に沿って持ち上がる）。
$\varphi$ が普遍開ならば、$\varphi$ は普遍一般化的である。
\end{lemma}

\begin{proof}
$\varphi$ が開であると仮定する。$y' \leadsto y$ を $Y$ の点の特殊化とする。
$x \in X$ を $\varphi(x) = y$ となるものとする。
アフィン開集合 $U \subset X$ および $V \subset Y$ であって
$\varphi(U) \subset V$ かつ $x \in U$ となるものを選ぶ。このとき $y' \in V$ でもある。
従って $X$ を $U$ で、$Y$ を $V$ で置き換え、$X$, $Y$ がアフィンであると仮定してよい。
アフィンの場合は、代数の章、補題 \ref{algebra-lemma-open-going-down} と
代数の章、補題 \ref{algebra-lemma-going-up-down-specialization} の組合せである。
\end{proof}

\begin{lemma}
\label{morphisms-lemma-descent-quasi-compact}
$f : X \to Y$ をスキームの射とする。$g : Y' \to Y$ を開かつ全射とし、
基底変換 $f' : X' \to Y'$ が準コンパクトであるとする。このとき $f$ は準コンパクトである。
\end{lemma}

\begin{proof}
$V \subset Y$ を準コンパクト開集合とする。$g$ は開かつ全射だから、
準コンパクト開集合 $W' \subset W$ であって $g(W') = V$ となるものを見いだせる。
仮定により $(f')^{-1}(W')$ は準コンパクトである。$(f')^{-1}(W')$ の $X$ における像は
$f^{-1}(V)$ に等しい。補題 \ref{morphisms-lemma-when-point-maps-to-pair} を参照せよ。
従って $f^{-1}(V)$ は準コンパクト空間の像として準コンパクトである。
位相の章、補題 \ref{topology-lemma-image-quasi-compact} を参照せよ。
ゆえに $f$ は準コンパクトである。
\end{proof}

\section{submersive な射}
\label{morphisms-section-submersive}

\begin{definition}
\label{morphisms-definition-submersive}
$f : X \to Y$ をスキームの射とする。
\begin{enumerate}
\item $f$ が {\it submersive}\footnote{これは微分多様体の沈め込みという概念とは
大きく異なる。} であるとは、台となる位相空間の連続写像が submersive であることをいう。
位相の章、定義 \ref{topology-definition-submersive} を参照せよ。
\item $f$ が {\it 普遍 submersive} であるとは、スキームの任意の射
$Y' \to Y$ に対し、基底変換 $Y' \times_Y X \to Y'$ が submersive であることをいう。
\end{enumerate}
\end{definition}

\noindent
submersive な射は特に全射であることに注意する。

\begin{lemma}
\label{morphisms-lemma-base-change-universally-submersive}
スキームの普遍 submersive な射をスキームの任意の射によって基底変換して得られる射は、
普遍 submersive である。
\end{lemma}

\begin{proof}
定義から直ちに従う。
\end{proof}

\begin{lemma}
\label{morphisms-lemma-composition-universally-submersive}
スキームの二つの（普遍）submersive な射の合成は（普遍）submersive である。
\end{lemma}

\begin{proof}
省略する。
\end{proof}

\section{平坦射}
\label{morphisms-section-flat}

\noindent
平坦性は、代数幾何学における最も重要な技術的道具の一つである。本節でこの概念を導入する。
補題 \ref{morphisms-lemma-fppf-open} を除き、ここでは意図的に議論を直接的な考察に限定する。
非常に重要な一群の結果である平坦性の判定法については、代数の章、節
\ref{algebra-section-criteria-flatness}、
\ref{algebra-section-flatness-artinian}、
\ref{algebra-section-more-flatness-criteria}、および
発展的な射の章、節 \ref{more-morphisms-section-criterion-flat-fibres} で論じる。
スキームの平坦射に関する発展的な内容には独立した章を設けてある。
平坦性詳論の章、節 \ref{flat-section-introduction} を参照せよ。

\medskip\noindent
加群 $M$（環 $R$ 上）が {\it 平坦} であるとは、関手
$-\otimes_R M : \text{Mod}_R \to \text{Mod}_R$ が完全であることだった。
環準同型 $R \to A$ が {\it 平坦} であるとは、$A$ が $R$-加群として平坦であることをいう。
代数の章、定義 \ref{algebra-definition-flat} を参照せよ。

\begin{definition}
\label{morphisms-definition-flat}
$f : X \to S$ をスキームの射とする。
$\mathcal{F}$ を準連接な $\mathcal{O}_X$-加群の層とする。
\begin{enumerate}
\item $f$ が {\it 点 $x \in X$ において平坦} であるとは、
局所環 $\mathcal{O}_{X, x}$ が局所環 $\mathcal{O}_{S, f(x)}$ 上平坦であることをいう。
\item $\mathcal{F}$ が {\it $S$ 上、点 $x \in X$ において平坦} であるとは、
茎 $\mathcal{F}_x$ が平坦な $\mathcal{O}_{S, f(x)}$-加群であることをいう。
\item $f$ が {\it 平坦} であるとは、$f$ が $X$ の各点において平坦であることをいう。
\item $\mathcal{F}$ が {\it $S$ 上平坦} であるとは、
$\mathcal{F}$ が $S$ 上、各点 $x$（$X$ の点）において平坦であることをいう。
\end{enumerate}
\end{definition}

\noindent
従って、$f$ が平坦であることと、構造層 $\mathcal{O}_X$ が $S$ 上平坦であることは同値である。

\begin{lemma}
\label{morphisms-lemma-flat-module-characterize}
$f : X \to S$ をスキームの射とする。
$\mathcal{F}$ を準連接な $\mathcal{O}_X$-加群の層とする。次の条件は同値である。
\begin{enumerate}
\item 層 $\mathcal{F}$ は $S$ 上平坦である。
\item 任意のアフィン開集合 $U \subset X$、$V \subset S$ であって
$f(U) \subset V$ となるものに対し、$\mathcal{O}_S(V)$-加群 $\mathcal{F}(U)$ は平坦である。
\item 開被覆 $S = \bigcup_{j \in J} V_j$ および開被覆
$f^{-1}(V_j) = \bigcup_{i \in I_j} U_i$ であって、各加群 $\mathcal{F}|_{U_i}$ が
$V_j$ 上平坦となるものが存在する。ここで $j\in J, i\in I_j$ は任意である。
\item アフィン開被覆 $S = \bigcup_{j \in J} V_j$ およびアフィン開被覆
$f^{-1}(V_j) = \bigcup_{i \in I_j} U_i$ であって、$\mathcal{F}(U_i)$ が平坦な
$\mathcal{O}_S(V_j)$-加群となるものが存在する。ここで $j\in J, i\in I_j$ は任意である。
\end{enumerate}
さらに、$\mathcal{F}$ が $S$ 上平坦ならば、任意の開部分スキーム
$U \subset X$、$V \subset S$ であって $f(U) \subset V$ となるものに対し、
制限 $\mathcal{F}|_U$ は $V$ 上平坦である。
\end{lemma}

\begin{proof}
$R \to A$ を環準同型とする。$M$ を $A$-加群とする。
$M$ が $R$-平坦ならば、すべての素イデアル $\mathfrak q$ に対し、
加群 $M_{\mathfrak q}$ は $R_{\mathfrak p}$ 上平坦である。ここで
$\mathfrak p$ は $R$ の素イデアルであって $\mathfrak q$ の下にあるものである。
逆に、$M_{\mathfrak q}$ が $R_{\mathfrak p}$ 上平坦であることがすべての素イデアル
$\mathfrak q$（$A$ の素イデアル）に対して成り立つならば、$M$ は $R$ 上平坦である。
代数の章、補題 \ref{algebra-lemma-flat-localization} を参照せよ。
この同値性から補題の各主張は容易に従う。
\end{proof}

\begin{lemma}
\label{morphisms-lemma-flat-characterize}
$f : X \to S$ をスキームの射とする。次の条件は同値である。
\begin{enumerate}
\item 射 $f$ は平坦である。
\item 任意のアフィン開集合 $U \subset X$、$V \subset S$ であって
$f(U) \subset V$ となるものに対し、環準同型
$\mathcal{O}_S(V) \to \mathcal{O}_X(U)$ は平坦である。
\item 開被覆 $S = \bigcup_{j \in J} V_j$ および開被覆
$f^{-1}(V_j) = \bigcup_{i \in I_j} U_i$ であって、各射
$U_i \to V_j$（$j\in J, i\in I_j$）が平坦となるものが存在する。
\item アフィン開被覆 $S = \bigcup_{j \in J} V_j$ およびアフィン開被覆
$f^{-1}(V_j) = \bigcup_{i \in I_j} U_i$ であって、
$\mathcal{O}_S(V_j) \to \mathcal{O}_X(U_i)$ が平坦となるものが存在する。
ここで $j\in J, i\in I_j$ は任意である。
\end{enumerate}
さらに、$f$ が平坦ならば、任意の開部分スキーム
$U \subset X$、$V \subset S$ であって $f(U) \subset V$ となるものに対し、
制限 $f|_U : U \to V$ は平坦である。
\end{lemma}

\begin{proof}
上の補題 \ref{morphisms-lemma-flat-module-characterize} の特別な場合である。
\end{proof}

\begin{lemma}
\label{morphisms-lemma-pushforward-flat-affine}
$f : X \to Y$ を基底スキーム $S$ 上のスキームのアフィン射とする。
$\mathcal{F}$ を準連接な $\mathcal{O}_X$-加群とする。
このとき $\mathcal{F}$ が $S$ 上平坦であることと、
$f_*\mathcal{F}$ が $S$ 上平坦であることは同値である。
\end{lemma}

\begin{proof}
補題 \ref{morphisms-lemma-flat-module-characterize} と、$f$ がアフィン射であることにより、
アフィンの場合に帰着する。$X \to Y \to S$ が環準同型
$C \leftarrow B \leftarrow A$ に対応するとする。
$N$ を $C$-加群であって $\mathcal{F}$ に対応するものとする。
$f_*\mathcal{F}$ は $N$ を $B$-加群とみなしたものに対応することを思い出そう。
スキームの章、補題 \ref{schemes-lemma-widetilde-pullback} を参照せよ。
従って結果は明らかである。
\end{proof}

\begin{lemma}
\label{morphisms-lemma-composition-module-flat}
$X \to Y \to Z$ をスキームの射とする。$\mathcal{F}$ を準連接な
$\mathcal{O}_X$-加群とする。$x \in X$ を、像が $y$（$Y$ の点）である点とする。
$\mathcal{F}$ が $Y$ 上、$x$ において平坦であり、$Y$ が $Z$ 上、
$y$ において平坦ならば、$\mathcal{F}$ は $Z$ 上、$x$ において平坦である。
\end{lemma}

\begin{proof}
代数の章、補題 \ref{algebra-lemma-composition-flat} を参照せよ。
\end{proof}

\begin{lemma}
\label{morphisms-lemma-composition-flat}
平坦射の合成は平坦である。
\end{lemma}

\begin{proof}
補題 \ref{morphisms-lemma-composition-module-flat} の特別な場合である。
\end{proof}

\begin{lemma}
\label{morphisms-lemma-base-change-module-flat}
$f : X \to S$ をスキームの射とする。
$\mathcal{F}$ を準連接な $\mathcal{O}_X$-加群の層とする。
$g : S' \to S$ をスキームの射とする。
$g' : X' = X_{S'} \to X$ を射影と書く。
$x' \in X'$ を、像が $x = g'(x') \in X$ である点とする。
$\mathcal{F}$ が $S$ 上、$x$ において平坦ならば、
$(g')^*\mathcal{F}$ は $S'$ 上、$x'$ において平坦である。
特に、$\mathcal{F}$ が $S$ 上平坦ならば、$(g')^*\mathcal{F}$ は $S'$ 上平坦である。
\end{lemma}

\begin{proof}
代数の章、補題 \ref{algebra-lemma-flat-base-change} を参照せよ。
\end{proof}

\begin{lemma}
\label{morphisms-lemma-base-change-flat}
平坦射の基底変換は平坦である。
\end{lemma}

\begin{proof}
補題 \ref{morphisms-lemma-base-change-module-flat} の特別な場合である。
\end{proof}

\begin{lemma}
\label{morphisms-lemma-generalizations-lift-flat}
$f : X \to S$ をスキームの平坦射とする。このとき $f$ に沿って一般化が持ち上がる。
位相の章、定義 \ref{topology-definition-lift-specializations} を参照せよ。
\end{lemma}

\begin{proof}
代数の章、節 \ref{algebra-section-going-up} を参照せよ。
\end{proof}

\begin{lemma}
\label{morphisms-lemma-fppf-open}
局所有限表示の平坦射は普遍開である。
\end{lemma}

\begin{proof}
上の補題 \ref{morphisms-lemma-generalizations-lift-flat} および補題
\ref{morphisms-lemma-locally-finite-presentation-universally-open} から従う。
次のように直接論じることもできる。

\medskip\noindent
$f : X \to S$ を平坦かつ局所有限表示とする。
補題 \ref{morphisms-lemma-base-change-flat} および
\ref{morphisms-lemma-base-change-finite-presentation} により、$f$ の任意の基底変換は平坦かつ
局所有限表示である。従って $f$ が開であることを示せば十分である。
$f$ が開であることを示すには、$X$ をアフィン開集合で被覆して
$X = \bigcup U_i$ と書き、$U_i \to S$ が開であるようにできることを示せば十分である。
$X$ はアフィン開集合 $U_i \subset X$ により被覆でき、各
$U_i$ がアフィン開集合 $V_i \subset S$ に写り、誘導される環準同型
$\mathcal{O}_S(V_i) \to \mathcal{O}_X(U_i)$ が平坦かつ有限表示となるようにできる
（補題 \ref{morphisms-lemma-flat-characterize} および
\ref{morphisms-lemma-locally-finite-presentation-characterize}）。
このとき $U_i \to V_i$ は、代数の章、命題 \ref{algebra-proposition-fppf-open} により
開である。これで証明が完了する。
\end{proof}

\begin{lemma}
\label{morphisms-lemma-pf-flat-module-open}
$f : X \to Y$ をスキームの射とする。
$\mathcal{F}$ を準連接な $\mathcal{O}_X$-加群とする。
$f$ が局所有限表示、$\mathcal{F}$ が有限型、$X = \text{Supp}(\mathcal{F})$ であり、
$\mathcal{F}$ が $Y$ 上平坦であると仮定する。このとき $f$ は普遍開である。
\end{lemma}

\begin{proof}
補題 \ref{morphisms-lemma-base-change-module-flat}、
\ref{morphisms-lemma-base-change-finite-presentation}、および
\ref{morphisms-lemma-support-finite-type} により、仮定は基底変換の下で保たれる。
補題 \ref{morphisms-lemma-locally-finite-presentation-universally-open} により、
$f$ に沿って一般化が持ち上がることを示せば十分である。
これは代数の章、補題 \ref{algebra-lemma-going-down-flat-module} から従う。
\end{proof}

\begin{lemma}
\label{morphisms-lemma-fpqc-quotient-topology}
\begin{reference}
\cite[Expose VIII, Corollaire 4.3]{SGA1} および
\cite[IV, Corollaire 2.3.12]{EGA}
\end{reference}
$f : X \to Y$ を準コンパクトかつ全射な平坦射とする。
部分集合 $T \subset Y$ が開（それぞれ閉）であることと、
$f^{-1}(T)$ が開（それぞれ閉）であることは同値である。
言い換えれば、$f$ は submersive な射である。
\end{lemma}

\begin{proof}
問題は $Y$ 上局所的だから、$Y$ はアフィンであると仮定してよい。
この場合、$X$ は $f$ が準コンパクトであることにより準コンパクトである。
$X = X_1 \cup \ldots \cup X_n$ と、アフィン開集合の有限合併として書く。
このとき $f' : X' = X_1 \amalg \ldots \amalg X_n \to Y$ は、アフィンスキームの
全射な平坦射である。$T \subset Y$ に対して
$(f')^{-1}(T) = f^{-1}(T) \cap X_1 \amalg \ldots \amalg f^{-1}(T) \cap X_n$
であることに注意する。従って、$f^{-1}(T)$ が開であることと
$(f')^{-1}(T)$ が開であることは同値である。ゆえに $X$ と $Y$ がともにアフィンであると仮定してよい。

\medskip\noindent
$f : \Spec(B) \to \Spec(A)$ を、平坦な環準同型 $A \to B$ に対応する
アフィンスキームの全射とする。$f^{-1}(T)$ は閉であり、
$f^{-1}(T) = V(J)$（$J \subset B$ はイデアル）と書けると仮定する。
このとき $T = f(f^{-1}(T)) = f(V(J))$ は
$\Spec(B/J) \to \Spec(A)$ の像である（ここで $f$ が全射であることを用いた）。
一方、$f$ に沿って一般化が持ち上がる（補題 \ref{morphisms-lemma-generalizations-lift-flat}）。
従って位相の章、補題 \ref{topology-lemma-lift-specializations-images} により、
$Y \setminus T = f(X \setminus f^{-1}(T))$ は一般化の下で安定である。
従って $T$ は特殊化の下で安定である
（位相の章、補題 \ref{topology-lemma-open-closed-specialization}）。
ゆえに $T$ は、代数の章、補題
\ref{algebra-lemma-image-stable-specialization-closed} により閉である。
\end{proof}

\begin{lemma}
\label{morphisms-lemma-flat-permanence}
$h : X \to Y$ を $S$ 上のスキームの射とする。
$\mathcal{G}$ を $Y$ 上の準連接層とする。
$x \in X$ とし、$y = h(x) \in Y$ とする。$h$ が $x$ において平坦ならば、
$$
\mathcal{G}\text{ flat over }S\text{ at }y
\Leftrightarrow
h^*\mathcal{G}\text{ flat over }S\text{ at }x.
$$
特に、$h$ が全射かつ平坦ならば、$\mathcal{G}$ が $S$ 上平坦であることと、
$h^*\mathcal{G}$ が $S$ 上平坦であることは同値である。
$h$ が全射かつ平坦であり、$X$ が $S$ 上平坦ならば、$Y$ は $S$ 上平坦である。
\end{lemma}

\begin{proof}
代数の章、補題 \ref{algebra-lemma-flatness-descends-more-general} を適用して証明できる。
直接の証明を与えよう。$s \in S$ を $y$ の像とする。局所環準同型
$\mathcal{O}_{S, s} \to \mathcal{O}_{Y, y} \to \mathcal{O}_{X, x}$ を考える。
仮定により、環準同型 $\mathcal{O}_{Y, y} \to \mathcal{O}_{X, x}$ は忠実平坦である。
代数の章、補題 \ref{algebra-lemma-local-flat-ff} を参照せよ。
$N = \mathcal{G}_y$ とおく。このとき
$h^*\mathcal{G}_x = N \otimes_{\mathcal{O}_{Y, y}} \mathcal{O}_{X, x}$ である。
層の章、補題 \ref{sheaves-lemma-stalk-pullback-modules} を参照せよ。
$M' \to M$ を $\mathcal{O}_{S, s}$-加群の単射とする。
上に述べた忠実平坦性により、
\begin{align*}
\Ker(
M' \otimes_{\mathcal{O}_{S, s}} N \to M \otimes_{\mathcal{O}_{S, s}} N)
\otimes_{\mathcal{O}_{Y, y}} \mathcal{O}_{X, x} \\
=
\Ker(
M' \otimes_{\mathcal{O}_{S, s}} N
\otimes_{\mathcal{O}_{Y, y}} \mathcal{O}_{X, x}
\to
M \otimes_{\mathcal{O}_{S, s}} N
\otimes_{\mathcal{O}_{Y, y}} \mathcal{O}_{X, x})
\end{align*}
である。従って補題の同値性は、代数の章、補題 \ref{algebra-lemma-flat} の
平坦性の第二の特徴付けから従う。
\end{proof}

\begin{lemma}
\label{morphisms-lemma-flat-pullback-support}
$f : Y \to X$ をスキームの射とする。$\mathcal{F}$ を有限型準連接
$\mathcal{O}_X$-加群とし、そのスキーム論的台を $Z \subset X$ とする。
$f$ が平坦ならば、$f^{-1}(Z)$ は $f^*\mathcal{F}$ のスキーム論的台である。
\end{lemma}

\begin{proof}
補題 \ref{morphisms-lemma-scheme-theoretic-support} によるアフィン上のスキーム論的台の
特徴付けを用いると、代数の章、補題
\ref{algebra-lemma-annihilator-flat-base-change} に帰着する。
\end{proof}

\begin{lemma}
\label{morphisms-lemma-flat-morphism-scheme-theoretically-dense-open}
$f : X \to Y$ をスキームの平坦射とする。
$V \subset Y$ をスキーム論的に稠密な逆コンパクト開集合とする。
このとき $f^{-1}V$ は $X$ においてスキーム論的に稠密である。
\end{lemma}

\begin{proof}
補題 \ref{morphisms-lemma-characterize-scheme-theoretically-dense} の特徴付けを用いる。
任意の開集合 $U \subset X$ に対して、写像
$\mathcal{O}_X(U) \to \mathcal{O}_X(U \cap f^{-1}V)$ が単射であることを示さなければならない。
$U$ が、アフィン開集合 $W \subset Y$ に写るアフィン開集合である場合に証明すれば十分である。
$W = \Spec(A)$、$U = \Spec(B)$ と書く。
このとき $V \cap W = D(f_1) \cup \ldots \cup D(f_n)$ と、ある
$f_i \in A$ を用いて書ける。代数の章、補題 \ref{algebra-lemma-qc-open} を参照せよ。
従って $B \to B_{f_1} \times \ldots \times B_{f_n}$ が単射であることを示せばよい。
$A \to A_{f_1} \times \ldots \times A_{f_n}$ は単射であり、
$A \to B$ は平坦であることが分かっている。
$B_{f_i} = A_{f_i} \otimes_A B$ だから、主張を得る。
\end{proof}

\begin{lemma}
\label{morphisms-lemma-flat-base-change-scheme-theoretic-image}
\begin{slogan}
準コンパクトの場合、スキーム論的像を取る操作は平坦基底変換と可換である
\end{slogan}
$f : X \to Y$ をスキームの平坦射とする。
$g : V \to Y$ をスキームの準コンパクト射とする。
$Z \subset Y$ を $g$ のスキーム論的像とし、$Z' \subset X$ を基底変換
$V \times_Y X \to X$ のスキーム論的像とする。このとき $Z' = f^{-1}Z$ である。
\end{lemma}

\begin{proof}
$Z$ は $\mathcal{I} = \Ker(\mathcal{O}_Y \to g_*\mathcal{O}_V)$ により定まり、
$Z'$ は
$\mathcal{I}' = \Ker(\mathcal{O}_X \to
(V \times_Y X \to X)_*\mathcal{O}_{V \times_Y X})$
により定まることを思い出そう。補題 \ref{morphisms-lemma-quasi-compact-scheme-theoretic-image} を参照せよ。
従って問題は $X$ および $Y$ 上局所的であり、$X$ と $Y$ はアフィンであると仮定してよい。
$V$ を $\coprod V_i$ で置き換えてよいことに注意する。ここで
$V = V_1 \cup \ldots \cup V_n$ は有限アフィン開被覆である。
従って $g$ はアフィンであると仮定してよい。この場合、
$(V \times_Y X \to X)_*\mathcal{O}_{V \times_Y X}$ は
$g_*\mathcal{O}_V$ の $f$ による引き戻しである。
$f$ は平坦だから、$f^*\mathcal{I} = \mathcal{I}'$ を得る。これで補題が成り立つ。
\end{proof}

\section{平坦な閉埋め込み}
\label{morphisms-section-flat-closed-immersions}

\noindent
スキームの連結成分は必ずしも開ではない。しかし常に標準的なスキーム構造をもつ。
本節ではこのことを説明する。

\begin{lemma}
\label{morphisms-lemma-characterize-flat-closed-immersions}
$X$ をスキームとする。$X$ の閉部分スキームにその台となる閉部分集合を対応させる規則は、
次の全単射を定める。
$$
\left\{
\begin{matrix}
\text{closed subschemes }Z \subset X \\
\text{such that }Z \to X\text{ is flat}
\end{matrix}
\right\}
\leftrightarrow
\left\{
\begin{matrix}
\text{closed subsets }Z \subset X \\
\text{closed under generalizations}
\end{matrix}
\right\}
$$
$Z \subset X$ がそのような閉部分スキームならば、スキームの任意の射
$g : Y \to X$ であって集合論的に $g(Y) \subset Z$ を満たすものは、
（スキーム論的に）$Z$ を経由して分解する。
\end{lemma}

\begin{proof}
この全単射のアフィンの場合は、代数の章、補題
\ref{algebra-lemma-pure-open-closed-specializations} である。
一般のスキーム $X$ の場合、$X$ をアフィンで被覆し、貼り合わせることにより全単射を得る。
詳細は省略する。最後の主張については、射影 $Z \times_{X, g} Y \to Y$ が平坦な
（補題 \ref{morphisms-lemma-base-change-flat}）閉埋め込みであり、台となる位相空間上で
全単射であることに注意する。従って、証明の最初の部分で得た全単射により、
この射影は同型でなければならない。
\end{proof}

\begin{lemma}
\label{morphisms-lemma-flat-closed-immersions-finite-presentation}
有限表示の平坦な閉埋め込みは、開かつ閉な部分スキームの開埋め込みである。
\end{lemma}

\begin{proof}
アフィンの場合は、代数の章、補題 \ref{algebra-lemma-finitely-generated-pure-ideal} である。
一般の場合、$X$ をアフィンで被覆することにより補題が従う。詳細は省略する。
\end{proof}

\noindent
連結成分 $T$（スキーム $X$ のもの）は、一般化の下で安定な閉部分集合であることに注意する。
従って次の定義は意味をもつ。

\begin{definition}
\label{morphisms-definition-scheme-structure-connected-component}
$X$ をスキームとする。$T \subset X$ を連結成分とする。
{\it $T$ 上の標準的なスキーム構造}とは、$T$ 上のスキーム構造であって、
閉埋め込み $T \to X$ が平坦となる唯一のものをいう。
補題 \ref{morphisms-lemma-characterize-flat-closed-immersions} を参照せよ。
\end{definition}

\noindent
前の補題を用いると、すべての有限平坦 $\mathcal{O}_X$-加群が有限局所自由となるのは
いつかを判定できる。

\begin{lemma}
\label{morphisms-lemma-finite-flat-is-finite-locally-free}
$X$ をスキームとする。次の条件は同値である。
\begin{enumerate}
\item 任意の有限平坦準連接 $\mathcal{O}_X$-加群は有限局所自由である。
\item 一般化で閉じた任意の閉部分集合 $Z \subset X$ は開である。
\end{enumerate}
\end{lemma}

\begin{proof}
アフィンの場合は、代数の章、補題
\ref{algebra-lemma-finite-flat-module-finitely-presented} である。
スキームの場合はアフィンの場合から直接には従わないので、議論を繰り返す。

\medskip\noindent
(1) を仮定する。閉埋め込み $i : Z \to X$ であって $i$ が平坦であるものを考える。
このとき $i_*\mathcal{O}_Z$ は準連接かつ平坦であり、従って (1) により有限局所自由である。
ゆえに $Z = \text{Supp}(i_*\mathcal{O}_Z)$ は開でもあり、(2) が成り立つ。
従って含意 (1) $\Rightarrow$ (2) は、補題
\ref{morphisms-lemma-characterize-flat-closed-immersions} による平坦な閉埋め込みの特徴付けから従う。

\medskip\noindent
逆を示すため、$X$ が (2) を満たすと仮定する。
$\mathcal{F}$ を有限平坦準連接 $\mathcal{O}_X$-加群とする。
台 $Z = \text{Supp}(\mathcal{F})$（$\mathcal{F}$ の台）は閉である。
加群の章、補題 \ref{modules-lemma-support-finite-type-closed} を参照せよ。
一方、$x \leadsto x'$ が特殊化ならば、代数の章、補題
\ref{algebra-lemma-finite-flat-local} により、加群 $\mathcal{F}_{x'}$ は
$\mathcal{O}_{X, x'}$ 上自由であり、
$$
\mathcal{F}_x =
\mathcal{F}_{x'} \otimes_{\mathcal{O}_{X, x'}} \mathcal{O}_{X, x}.
$$
従って
$x' \in \text{Supp}(\mathcal{F}) \Rightarrow x \in \text{Supp}(\mathcal{F})$、
すなわち台は一般化で閉じている。$X$ は (2) を満たすから、
$\mathcal{F}$ の台は開かつ閉である。加群 $\wedge^i(\mathcal{F})$（$i = 1, 2, 3, \ldots$）
も有限平坦準連接 $\mathcal{O}_X$-加群である。
加群の章、節 \ref{modules-section-symmetric-exterior} を参照せよ。
このとき
$\text{Supp}(\wedge^{i + 1}(\mathcal{F})) \subset
\text{Supp}(\wedge^i(\mathcal{F}))$
であることに注意する。従って、開かつ閉な部分集合による分解
$$
X = U_0 \amalg U_1 \amalg U_2 \amalg \ldots
$$
であって、$\wedge^i(\mathcal{F})$ の台が $U_i \cup U_{i + 1} \cup \ldots$ となるという条件を
すべての $i$ に対して満たすものが存在する。$x$ を $X$ の点とし、$x \in U_r$ とする。
このとき
$\wedge^i(\mathcal{F})_x \otimes \kappa(x) =
\wedge^i(\mathcal{F}_x \otimes \kappa(x))$
であることに注意する。従って、$x \in U_r$ ならば、$\mathcal{F}_x \otimes \kappa(x)$
は次元 $r$ のベクトル空間である。Nakayama の補題、すなわち代数の章、補題
\ref{algebra-lemma-NAK} により、アフィン開近傍 $U \subset U_r \subset X$
（$x$ の近傍）と切断 $s_1, \ldots, s_r \in \mathcal{F}(U)$ を、
誘導される写像
$$
\mathcal{O}_U^{\oplus r} \longrightarrow \mathcal{F}|_U, \quad
(f_1, \ldots, f_r) \longmapsto \sum f_i s_i
$$
が全射となるように選べる。このことは、$\wedge^r(\mathcal{F}|_U)$ が有限平坦準連接
$\mathcal{O}_U$-加群で、その台が $U$ 全体であることを意味する。
上記により、それは単一の元 $s_1 \wedge \ldots \wedge s_r$ により生成される。
従って $\wedge^r(\mathcal{F}|_U) \cong \mathcal{O}_U/\mathcal{I}$ と、
ある準連接イデアル層 $\mathcal{I}$ を用いて書ける。このとき
$\mathcal{O}_U/\mathcal{I}$ は $\mathcal{O}_U$ 上平坦であり、
$V(\mathcal{I}) = U$ である。補題
\ref{morphisms-lemma-characterize-flat-closed-immersions} を適用すると $\mathcal{I} = 0$ が従う。
従って $s_1 \wedge \ldots \wedge s_r$ は $\wedge^r(\mathcal{F}|_U)$ の基底であり、
上に表示した写像は全射であるだけでなく単射でもある。
これで $\mathcal{F}$ が望みどおり有限局所自由であることが示された。
\end{proof}

\section{Generic flatness}
\label{morphisms-section-generic-flatness}

\noindent
整な基底上有限型のスキームは、基底のある稠密開集合上で平坦である。
代数、節 \ref{algebra-section-generic-flatness} では、Noether 的な場合、
有限表示の射の場合、および一般の場合を証明した。ここでは一般の場合のみを述べて証明する。
ただし、実際にはこれをさらに強めた命題
\ref{morphisms-proposition-generic-flatness-reduced} が成り立ち、基底が被約であることのみを
仮定しても同じ結果が成立する。

\begin{proposition}[Generic flatness]
\label{morphisms-proposition-generic-flatness}
$f : X \to S$ をスキームの射とする。
$\mathcal{F}$ を準連接な $\mathcal{O}_X$-加群の層とする。
次を仮定する。
\begin{enumerate}
\item $S$ は整である。
\item $f$ は有限型である。
\item $\mathcal{F}$ は有限型の $\mathcal{O}_X$-加群である。
\end{enumerate}
このとき、ある開稠密部分スキーム $U \subset S$ が存在し、
$X_U \to U$ は平坦かつ有限表示であり、
$\mathcal{F}|_{X_U}$ は $U$ 上平坦かつ
$\mathcal{O}_{X_U}$ 上有限表示である。
\end{proposition}

\begin{proof}
$S$ は整なので既約であり（性質、補題
\ref{properties-lemma-characterize-integral} 参照）、空でない開集合はすべて稠密である。
従って $S$ を $S$ のアフィン開集合で置き換え、
$S = \Spec(A)$ はアフィンであると仮定してよい。
$S$ が整であることから $A$ は整域である。
$f$ は有限型なので準コンパクトであり、従って $X$ は準コンパクトである。
よって有限アフィン開被覆 $X = \bigcup_{i = 1, \ldots, n} X_i$ が取れる。
$X_i = \Spec(B_i)$ と書く。
このとき $B_i$ は有限型 $A$-代数である。補題
\ref{morphisms-lemma-locally-finite-type-characterize} 参照。
さらに有限型の $B_i$-加群 $M_i$ が存在し、$\mathcal{F}|_{X_i}$ は
$B_i$-加群 $M_i$ に付随する準連接層となる。性質、補題
\ref{properties-lemma-finite-type-module} 参照。
次に、添字の各対 $i, j$ に対してイデアル $I_{ij} \subset B_i$ を、
$X_i \setminus X_i \cap X_j = V(I_{ij})$ が
$X_i = \Spec(B_i)$ の中で成り立つように取る。
$M_{ij} = B_i/I_{ij}$ とおき、これを $B_i$-加群とみなす。
すると $V(I_{ij}) = \text{Supp}(M_{ij})$ であり、
$M_{ij}$ は有限 $B_i$-加群である。

\medskip\noindent
ここで代数、補題 \ref{algebra-lemma-generic-flatness} を、
対 $(A \to B_i, M_{ij})$ および対 $(A \to B_i, M_i)$ に適用する。
これにより、零でない $f_{ij}, f_i \in A$ が得られ、次が成り立つ。
(a) $A_{f_{ij}} \to B_{i, f_{ij}}$ は平坦かつ有限表示であり、
$M_{ij, f_{ij}}$ は $A_{f_{ij}}$ 上平坦かつ
$B_{i, f_{ij}}$ 上有限表示である。
(b) $B_{i, f_i}$ は $A_f$ 上平坦かつ有限表示であり、
$M_{i, f_i}$ は $B_{i, f_i}$ 上平坦かつ有限表示である。
$f = (\prod f_i) (\prod f_{ij})$ とおく。
$U = D(f)$ と取ればよいと主張する。

\medskip\noindent
この主張を示すには、$A$ を $A_f$ で置き換えてよい。すなわち、
$U = \Spec(A_f) \to S$ による基底変換を行う。
この基底変換の後では、各 $A \to B_i$ は平坦かつ有限表示であり、
$M_i$, $M_{ij}$ は $A$ 上平坦かつ $B_i$ 上有限表示である。
これにより、$X \to S$ は準コンパクト、局所有限表示かつ平坦であり、
$\mathcal{F}$ は $S$ 上平坦かつ $\mathcal{O}_X$ 上有限表示であることが
既に分かる。補題 \ref{morphisms-lemma-locally-finite-presentation-characterize}
および性質、補題 \ref{properties-lemma-finite-presentation-module} 参照。
$M_{ij}$ は $B_i$ 上有限表示なので、
$X_i \cap X_j = X_i \setminus \text{Supp}(M_{ij})$ は
$X_i$ の準コンパクト開集合である。代数、補題
\ref{algebra-lemma-support-finite-presentation-constructible} 参照。
従って、スキーム、補題 \ref{schemes-lemma-characterize-quasi-separated} により
$X \to S$ は準分離的である。これで命題が証明された。
\end{proof}

\noindent
実際には、任意の被約な基底上での generic flatness も成り立つ。
その定式化は次のとおりである。

\begin{proposition}[Generic flatness、被約の場合]
\label{morphisms-proposition-generic-flatness-reduced}
$f : X \to S$ をスキームの射とする。
$\mathcal{F}$ を準連接な $\mathcal{O}_X$-加群の層とする。
次を仮定する。
\begin{enumerate}
\item $S$ は被約である。
\item $f$ は有限型である。
\item $\mathcal{F}$ は有限型の $\mathcal{O}_X$-加群である。
\end{enumerate}
このとき、ある開稠密部分スキーム $U \subset S$ が存在し、
$X_U \to U$ は平坦かつ有限表示であり、
$\mathcal{F}|_{X_U}$ は $U$ 上平坦かつ
$\mathcal{O}_{X_U}$ 上有限表示である。
\end{proposition}

\begin{proof}
結論を急ぐ読者のために述べると、この証明は命題
\ref{morphisms-proposition-generic-flatness} の証明で、代数、補題
\ref{algebra-lemma-generic-flatness-reduced} を代数、補題
\ref{algebra-lemma-generic-flatness} の代わりに用いたものである。

\medskip\noindent
平坦性と有限表示性は基底について局所的である
（補題 \ref{morphisms-lemma-flat-module-characterize} および
\ref{morphisms-lemma-locally-finite-presentation-characterize} 参照）から、
$S$ 上アフィン局所的に考えてよい。
従って $S = \Spec(A)$ と仮定でき、ここで $A$ は被約環である
（性質、補題 \ref{properties-lemma-characterize-reduced} 参照）。
$f$ は有限型なので準コンパクトであり、従って $X$ は準コンパクトである。
よって有限アフィン開被覆 $X = \bigcup_{i = 1, \ldots, n} X_i$ が取れる。
$X_i = \Spec(B_i)$ と書く。
このとき $B_i$ は有限型 $A$-代数である。補題
\ref{morphisms-lemma-locally-finite-type-characterize} 参照。
さらに有限型の $B_i$-加群 $M_i$ が存在し、$\mathcal{F}|_{X_i}$ は
$B_i$-加群 $M_i$ に付随する準連接層となる。性質、補題
\ref{properties-lemma-finite-type-module} 参照。
次に、添字の各対 $i, j$ に対してイデアル $I_{ij} \subset B_i$ を、
$X_i \setminus X_i \cap X_j = V(I_{ij})$ が
$X_i = \Spec(B_i)$ の中で成り立つように取る。
$M_{ij} = B_i/I_{ij}$ とおき、これを $B_i$-加群とみなす。
すると $V(I_{ij}) = \text{Supp}(M_{ij})$ であり、
$M_{ij}$ は有限 $B_i$-加群である。

\medskip\noindent
ここで代数、補題 \ref{algebra-lemma-generic-flatness-reduced} を、
対 $(A \to B_i, M_{ij})$ および対 $(A \to B_i, M_i)$ に適用する。
これにより、稠密開集合 $U(A \to B_i, M_{ij}) \subset S$ および
稠密開集合 $U(A \to B_i, M_i) \subset S$ が得られる。
記法は代数、式 (\ref{algebra-equation-good-locus}) のとおりである。
有限個の稠密開集合の共通部分は稠密開集合なので、
$$
U =
\bigcap\nolimits_{i, j} U(A \to B_i, M_{ij})
\quad\cap\quad
\bigcap\nolimits_i U(A \to B_i, M_i)
$$
は $S$ の開稠密集合である。この $U$ が求める開集合であると主張する。

\medskip\noindent
$u \in U$ を取る。集合 $U(A \to B_i, M_{ij})$ および
$U(A \to B, M_i)$ の定義により、$f_{ij}, f_i \in A$ が存在して次が成り立つ。
(a) $u \in D(f_i)$ かつ $u \in D(f_{ij})$。
(b) $A_{f_{ij}} \to B_{i, f_{ij}}$ は平坦かつ有限表示であり、
$M_{ij, f_{ij}}$ は $A_{f_{ij}}$ 上平坦かつ
$B_{i, f_{ij}}$ 上有限表示である。
(c) $B_{i, f_i}$ は $A_{f_i}$ 上平坦かつ有限表示であり、
$M_{i, f_i}$ は $B_{i, f_i}$ 上有限表示かつ $A_{f_i}$ 上平坦である。
$f = (\prod f_i) (\prod f_{ij})$ とおく。
すると、$X \to S$ が $D(f)$ 上平坦かつ有限表示であり、
$\mathcal{F}$ を $X_{D(f)}$ に制限したものが $D(f)$ 上平坦かつ
$X_{D(f)}$ の構造層上有限表示であることを示せば十分である。

\medskip\noindent
従って $A$ を $A_f$ で置き換えてよい。すなわち、
$\Spec(A_f) \to S$ による基底変換を行う。
この基底変換の後では、各 $A \to B_i$ は平坦かつ有限表示であり、
$M_i$, $M_{ij}$ は $A$ 上平坦かつ $B_i$ 上有限表示である。
これにより、$X \to S$ は準コンパクト、局所有限表示かつ平坦であり、
$\mathcal{F}$ は $S$ 上平坦かつ $\mathcal{O}_X$ 上有限表示であることが
既に分かる。補題 \ref{morphisms-lemma-locally-finite-presentation-characterize}
および性質、補題 \ref{properties-lemma-finite-presentation-module} 参照。
$M_{ij}$ は $B_i$ 上有限表示なので、
$X_i \cap X_j = X_i \setminus \text{Supp}(M_{ij})$ は
$X_i$ の準コンパクト開集合である。代数、補題
\ref{algebra-lemma-support-finite-presentation-constructible} 参照。
従って、スキーム、補題 \ref{schemes-lemma-characterize-quasi-separated} により
$X \to S$ は準分離的である。これで命題が証明された。
\end{proof}

\begin{remark}
\label{morphisms-remark-flattening}
上の結果は、スキームの射に対するより精密な平坦化の手法への第一歩である。
Raynaud と Gruson の論文 \cite{GruRay} には、
この方向における多くの優れた結果が含まれている。
\end{remark}

\section{射とファイバーの次元}
\label{morphisms-section-dimension-fibres}

\noindent
$X$ を位相空間とし、$x \in X$ とする。
$\dim_x(X)$ は、$x$ の $X$ における開近傍の次元の最小値として
定義したことを思い出そう。位相、定義 \ref{topology-definition-Krull} 参照。

\begin{lemma}
\label{morphisms-lemma-dimension-fibre-at-a-point}
$f : X \to S$ をスキームの射とする。
$x \in X$ とし、$s = f(x)$ とおく。
$f$ は局所有限型であると仮定する。このとき
$$
\dim_x(X_s) =
\dim(\mathcal{O}_{X_s, x}) + \text{trdeg}_{\kappa(s)}(\kappa(x)).
$$
\end{lemma}

\begin{proof}
直ちに $S = s$ かつ $X$ がアフィンの場合に帰着する。
この場合、主張は代数、補題
\ref{algebra-lemma-dimension-at-a-point-finite-type-field} から従う。
\end{proof}

\begin{lemma}
\label{morphisms-lemma-dimension-fibre-at-a-point-additive}
$f : X \to Y$ および $g : Y \to S$ をスキームの射とする。
$x \in X$ とし、$y = f(x)$, $s = g(y)$ とおく。
$f$ および $g$ は局所有限型であると仮定する。このとき
$$
\dim_x(X_s) \leq \dim_x(X_y) + \dim_y(Y_s).
$$
さらに、$\mathcal{O}_{X_s, x}$ が $\mathcal{O}_{Y_s, y}$ 上平坦ならば
等号が成り立つ。この平坦性は、例えば $\mathcal{O}_{X, x}$ が
$\mathcal{O}_{Y, y}$ 上平坦ならば成り立つ。
\end{lemma}

\begin{proof}
$\text{trdeg}_{\kappa(s)}(\kappa(x)) =
\text{trdeg}_{\kappa(y)}(\kappa(x)) + \text{trdeg}_{\kappa(s)}(\kappa(y))$
に注意する。従って補題 \ref{morphisms-lemma-dimension-fibre-at-a-point} により、
主張は
$$
\dim(\mathcal{O}_{X_s, x})
\leq
\dim(\mathcal{O}_{X_y, x}) + \dim(\mathcal{O}_{Y_s, y}).
$$
と同値である。これについては代数、補題
\ref{algebra-lemma-dimension-base-fibre-total} を参照。
平坦な場合については代数、補題
\ref{algebra-lemma-dimension-base-fibre-equals-total} を参照。
\end{proof}

\begin{lemma}
\label{morphisms-lemma-dimension-fibre-after-base-change}
スキームのファイバー積図式
$$
\xymatrix{
X' \ar[r]_{g'} \ar[d]_{f'} & X \ar[d]^f \\
S' \ar[r]^g & S
}
$$
を考える。$f$ は局所有限型であると仮定する。
$x' \in X'$, $x = g'(x')$, $s' = f'(x')$ および
$s = g(s') = f(x)$ とする。このとき
\begin{enumerate}
\item $\dim_x(X_s) = \dim_{x'}(X'_{s'})$。
\item $F$ が射 $X'_{s'} \to X_s$ の $x$ 上のファイバーならば、
$$
\dim(\mathcal{O}_{F, x'}) =
\dim(\mathcal{O}_{X'_{s'}, x'}) - \dim(\mathcal{O}_{X_s, x}) =
\text{trdeg}_{\kappa(s)}(\kappa(x)) -
\text{trdeg}_{\kappa(s')}(\kappa(x'))
$$
特に $\dim(\mathcal{O}_{X'_{s'}, x'}) \geq \dim(\mathcal{O}_{X_s, x})$
かつ $\text{trdeg}_{\kappa(s)}(\kappa(x)) \geq
\text{trdeg}_{\kappa(s')}(\kappa(x'))$。
\item $s', s, x$ が与えられたとき、$x'$ を
$\dim(\mathcal{O}_{X'_{s'}, x'}) = \dim(\mathcal{O}_{X_s, x})$ および
$\text{trdeg}_{\kappa(s)}(\kappa(x)) = \text{trdeg}_{\kappa(s')}(\kappa(x'))$
が成り立つように選べる。
\end{enumerate}
\end{lemma}

\begin{proof}
(1) は代数、補題
\ref{algebra-lemma-dimension-at-a-point-preserved-field-extension}
から直ちに従う。(2) と (3) は代数、補題
\ref{algebra-lemma-inequalities-under-field-extension} から従う。
\end{proof}

\noindent
次の補題は、非自明な代数的結果、すなわち Zariski の主定理の代数的な形から従う。

\begin{lemma}
\label{morphisms-lemma-openness-bounded-dimension-fibres}
\begin{reference}
\cite[IV Theorem 13.1.3]{EGA}
\end{reference}
$f : X \to S$ をスキームの射とする。
$n \geq 0$ とする。$f$ は局所有限型であると仮定する。
集合
$$
U_n = \{x \in X \mid \dim_x X_{f(x)} \leq n\}
$$
は $X$ の開集合である。
\end{lemma}

\begin{proof}
これは代数、補題
\ref{algebra-lemma-dimension-fibres-bounded-open-upstairs} から直ちに従う。
\end{proof}

\begin{lemma}
\label{morphisms-lemma-morphism-finite-type-bounded-dimension}
$f : X \to Y$ を有限型の射とし、$Y$ は準コンパクトとする。
このとき $f$ のファイバーの次元は有界である。
\end{lemma}

\begin{proof}
補題 \ref{morphisms-lemma-openness-bounded-dimension-fibres} により、
集合 $U_n \subset X$、すなわちファイバーの次元が $\leq n$ である点の集合は開集合である。
$f$ は有限型なので、各点はいずれかの $U_n$ に属する
（体上有限型の代数の次元は有限だからである）。
$Y$ は準コンパクトであり、$f$ は有限型なので、$X$ は準コンパクトである。
従って $X = U_n$ がある $n$ に対して成り立つ。
\end{proof}

\begin{lemma}
\label{morphisms-lemma-openness-bounded-dimension-fibres-finite-presentation}
$f : X \to S$ をスキームの射とする。
$n \geq 0$ とする。$f$ は局所有限表示であると仮定する。
補題 \ref{morphisms-lemma-openness-bounded-dimension-fibres} の開集合
$$
U_n = \{x \in X \mid \dim_x X_{f(x)} \leq n\}
$$
は $X$ において逆コンパクトである
（位相、定義 \ref{topology-definition-quasi-compact} 参照）。
\end{lemma}

\begin{proof}
位相空間 $X$ の位相には、アフィン開集合 $U \subset X$ であって
誘導射 $f|_U : U \to S$ があるアフィン開集合 $V \subset S$ を経由するものからなる
基底が存在する。従って $U \cap U_n$ が準コンパクトであることを、
そのような $U$ に対して示せば十分である。
$U_n \cap U$ は開集合 $\{x \in U \mid \dim_x U_{f(x)} \leq n\}$ と同じである。
これにより $X$ および $S$ がアフィンの場合に帰着する。
この場合、補題は代数、補題
\ref{algebra-lemma-dimension-fibres-bounded-quasi-compact-open-upstairs}
（および補題 \ref{morphisms-lemma-locally-finite-presentation-characterize}）から従う。
\end{proof}

\begin{lemma}
\label{morphisms-lemma-dimension-fibre-specialization}
$f : X \to S$ をスキームの射とする。
$x \leadsto x'$ を $X$ の点の非自明な特殊化とし、
これらの点は同じ点 $s \in S$ の上にあるとする。
$f$ は局所有限型であると仮定する。このとき
\begin{enumerate}
\item $\dim_x(X_s) \leq \dim_{x'}(X_s)$。
\item $\dim(\mathcal{O}_{X_s, x}) < \dim(\mathcal{O}_{X_s, x'})$。
\item $\text{trdeg}_{\kappa(s)}(\kappa(x)) >
\text{trdeg}_{\kappa(s)}(\kappa(x'))$。
\end{enumerate}
\end{lemma}

\begin{proof}
(1) は、$X_s$ の開集合で $x'$ を含むものはすべて $x$ も含むことから従う。
(2) について、$\mathcal{O}_{X_s, x}$ は $\mathcal{O}_{X_s, x'}$ を
ある素イデアルで局所化したものである。従って $\mathcal{O}_{X_s, x}$ の
素イデアルの任意の鎖は、$\mathcal{O}_{X_s, x'}$ の真に長い素イデアルの鎖の一部となる。
最後の不等式は代数、補題 \ref{algebra-lemma-tr-deg-specialization} から従う。
\end{proof}

\section{与えられた相対次元をもつ射}
\label{morphisms-section-relative-dimension}

\noindent
与えられた相対次元をもつ射について容易に述べられるよう、次の概念を導入する。

\begin{definition}
\label{morphisms-definition-relative-dimension-d}
$f : X \to S$ をスキームの射とする。
$f$ は局所有限型であると仮定する。
\begin{enumerate}
\item $f$ の{\it 相対次元が $\leq d$ である（点 $x$ において）}とは、
$\dim_x(X_{f(x)}) \leq d$ が成り立つことをいう。
\item $f$ の{\it 相対次元が $\leq d$} であるとは、
$\dim_x(X_{f(x)}) \leq d$ がすべての $x \in X$ に対して成り立つことをいう。
\item $f$ の{\it 相対次元が $d$} であるとは、
空でないファイバー $X_s$ がすべて次元 $d$ の等次元スキームであることをいう。
\end{enumerate}
\end{definition}

\noindent
この概念は特に振る舞いがよいものではないが、多くの状況ではうまく機能する。

\begin{lemma}
\label{morphisms-lemma-base-change-relative-dimension-d}
$f : X \to S$ を局所有限型のスキームの射とする。
$f$ の相対次元が $d$ ならば、$f$ の任意の基底変換についても同じことが成り立つ。
相対次元が $\leq d$ であることについても同様である。
\end{lemma}

\begin{proof}
これは補題 \ref{morphisms-lemma-dimension-fibre-after-base-change} から直ちに従う。
\end{proof}

\begin{lemma}
\label{morphisms-lemma-composition-relative-dimension-d}
$f : X \to Y$, $g : Y \to Z$ は局所有限型とする。
$f$ の相対次元が $\leq d$ であり、$g$ の相対次元が $\leq e$ ならば、
$g \circ f$ の相対次元は $\leq d + e$ である。
また、
\begin{enumerate}
\item $f$ の相対次元は $d$ である。
\item $g$ の相対次元は $e$ である。
\item $f$ は平坦である。
\end{enumerate}
が成り立つならば、$g \circ f$ の相対次元は $d + e$ である。
\end{lemma}

\begin{proof}
これは補題 \ref{morphisms-lemma-dimension-fibre-at-a-point-additive} から直ちに従う。
\end{proof}

\noindent
一般には、射を相対次元が一定となる部分に分解することはできない。
しかし、射のファイバーが Cohen-Macaulay であり、射が平坦かつ有限表示ならば
そのような分解が可能である。

\begin{lemma}
\label{morphisms-lemma-flat-finite-presentation-CM-fibres-relative-dimension}
\begin{slogan}
Cohen-Macaulay 射は、純相対次元をもつ開かつ閉な部分に分解される
\end{slogan}
$f : X \to S$ をスキームの射とする。
次を仮定する。
\begin{enumerate}
\item $f$ は平坦である。
\item $f$ は局所有限表示である。
\item すべての $s \in S$ に対してファイバー $X_s$ は Cohen-Macaulay である
（性質、定義 \ref{properties-definition-Cohen-Macaulay}）。
\end{enumerate}
このとき、開かつ閉な部分スキーム $X_d \subset X$ が存在して、
$X = \coprod_{d \geq 0} X_d$ であり、$f|_{X_d} : X_d \to S$ の
相対次元は $d$ である。
\end{lemma}

\begin{proof}
これは代数、補題 \ref{algebra-lemma-relative-dimension-CM} から直ちに従う。
\end{proof}

\begin{lemma}
\label{morphisms-lemma-locally-quasi-finite-rel-dimension-0}
$f : X \to S$ をスキームの射とする。
$f$ は局所有限型であると仮定する。
$x \in X$ とし、$s = f(x)$ とおく。
このとき、$f$ が $x$ で準有限であることと $\dim_x(X_s) = 0$ は同値である。
特に、$f$ が局所準有限であることと、$f$ の相対次元が $0$ であることは同値である。
\end{lemma}

\begin{proof}
第一の証明。
$f$ が $x$ で準有限ならば、$\kappa(x)$ は $\kappa(s)$ の有限拡大であり
（補題 \ref{morphisms-lemma-residue-field-quasi-finite}）、
$x$ は $X_s$ の孤立点である
（補題 \ref{morphisms-lemma-quasi-finite-at-point-characterize}）。
従って、補題 \ref{morphisms-lemma-dimension-fibre-at-a-point} により $\dim_x(X_s) = 0$ である。
逆に、$\dim_x(X_s) = 0$ ならば、補題
\ref{morphisms-lemma-dimension-fibre-at-a-point} により、
$\kappa(s) \subset \kappa(x)$ は代数拡大であり、
$X_s$ の点で $x$ に特殊化するものはほかに存在しない。
よって、補題 \ref{morphisms-lemma-algebraic-residue-field-extension-closed-point-fibre}
により $x$ はそのファイバーで閉点であり、補題
\ref{morphisms-lemma-quasi-finite-at-point-characterize} (3) により、
$f$ は $x$ で準有限であると結論できる。

\medskip\noindent
第二の証明。ファイバー $X_s$ は体上局所有限型のスキームなので、
局所 Noether 的である（補題 \ref{morphisms-lemma-finite-type-noetherian}）。
従って、主張は補題 \ref{morphisms-lemma-quasi-finite-at-point-characterize}
および性質、補題 \ref{properties-lemma-isolated-dim-0-locally-noetherian} から従う。
\end{proof}

\begin{lemma}
\label{morphisms-lemma-rel-dimension-dimension}
$f : X \to Y$ を局所 Noether スキームの射で、
平坦、局所有限型かつ相対次元 $d$ のものとする。
任意の点 $x$ が $X$ に属し、その像 $y$ が $Y$ に属するとき、
$\dim_x(X) = \dim_y(Y) + d$ が成り立つ。
\end{lemma}

\begin{proof}
$X$ および $Y$ を $x$ および $y$ の開近傍に縮小することで、
$\dim(X) = \dim_x(X)$ かつ $\dim(Y) = \dim_y(Y)$ と仮定してよい。
これは点におけるスキームの次元の定義による
（性質、定義 \ref{properties-definition-dimension}）。
射 $f$ は補題 \ref{morphisms-lemma-noetherian-finite-type-finite-presentation}
および \ref{morphisms-lemma-fppf-open} により開射である。
従って $Y$ を縮小して $f$ が全射であるようにできる。
残るのは $\dim(X) = \dim(Y) + d$ を示すことである。

\medskip\noindent
$a$ を $X$ の点とし、その像 $b$ は $Y$ に属するとする。
代数、補題 \ref{algebra-lemma-dimension-base-fibre-equals-total} により、
$$
\dim(\mathcal{O}_{X,a}) = \dim(\mathcal{O}_{Y,b}) + \dim(\mathcal{O}_{X_b, a}).
$$
点 $a$ を $X$ 全体にわたって動かして上限を取ると、求める
$\dim(X) = \dim(Y) + d$ が従う。性質、補題
\ref{properties-lemma-dimension} 参照。
\end{proof}

\section{Syntomic 射}
\label{morphisms-section-syntomic}

\noindent
代数 $A$ が体 $k$ 上の代数であるとき、
{\it $k$ 上の大域完全交叉}であるとは、
$A \cong k[x_1, \ldots, x_n]/(f_1, \ldots, f_c)$ かつ
$\dim(A) = n - c$ と書けることをいう。
代数 $A$ が体 $k$ 上の代数であるとき、
{\it 局所完全交叉}であるとは、$\Spec(A)$ を標準開集合で被覆でき、
それぞれが $k$ 上の大域完全交叉となることをいう。
代数、節 \ref{algebra-section-lci} 参照。
環準同型 $R \to A$ が{\it syntomic} であるとは、有限表示かつ平坦であり、
ファイバーが局所完全交叉環であることをいう。代数、定義
\ref{algebra-definition-lci} 参照。

\begin{definition}
\label{morphisms-definition-syntomic}
$f : X \to S$ をスキームの射とする。
\begin{enumerate}
\item $f$ が{\it 点 $x \in X$ において syntomic} であるとは、
アフィン開近傍 $\Spec(A) = U \subset X$ が $x$ に対して存在し、
アフィン開集合 $\Spec(R) = V \subset S$ が存在して、
$f(U) \subset V$ であり、誘導される環準同型 $R \to A$ が syntomic となることをいう。
\item $f$ が{\it syntomic} であるとは、$X$ のすべての点で syntomic であることをいう。
\item $S = \Spec(k)$ であり $f$ が syntomic であるとき、
$X$ は{\it $k$ 上の局所完全交叉}であるという。
\item アフィンスキームの射 $f : X \to S$ が{\it 標準 syntomic} であるとは、
大域相対完全交叉 $R \to R[x_1, \ldots, x_n]/(f_1, \ldots, f_c)$
（代数、定義 \ref{algebra-definition-relative-global-complete-intersection} 参照）が存在し、
$X \to S$ が
$$
\Spec(R[x_1, \ldots, x_n]/(f_1, \ldots, f_c)) \to \Spec(R).
$$
と同型であることをいう。
\end{enumerate}
\end{definition}

\noindent
文献によっては、syntomic 射は{\it 平坦な局所完全交叉射}と呼ばれる。
これは扱いやすい射のクラスである。
例えば、これらの射を用いて syntomic 位相を定義できる。この位相は滑らかな位相や
エタール位相より細かいが、それらと多くの形式的性質を共有する。

\medskip\noindent
大域相対完全交叉（標準 syntomic 環準同型の定義に用いたもの）は、特に平坦である。
射の詳論、節 \ref{more-morphisms-section-lci} では、射 $X \to S$ で局所的に
$$
\Spec(R[x_1, \ldots, x_n]/(f_1, \ldots, f_c)) \to \Spec(R).
$$
という形をもち、ある Koszul 正則列 $f_1, \ldots, f_r$ が
$R[x_1, \ldots, x_n]$ の中に存在するものを考える。
そのような射を{\it 局所完全交叉射}と呼ぶ。
この定義を導入すると、射が syntomic であることと、
平坦な局所完全交叉射であることは同値となる。

\medskip\noindent
syntomic 射の定義には、分離性や準コンパクト性の仮定がないことに注意する。
従って、syntomic であるかどうかは始域について局所的な問題である。
正確な結果は次のとおりである。

\begin{lemma}
\label{morphisms-lemma-syntomic-characterize}
$f : X \to S$ をスキームの射とする。次は同値である。
\begin{enumerate}
\item 射 $f$ は syntomic である。
\item 任意のアフィン開集合 $U \subset X$, $V \subset S$ で
$f(U) \subset V$ を満たすものについて、環準同型
$\mathcal{O}_S(V) \to \mathcal{O}_X(U)$ は syntomic である。
\item 開被覆 $S = \bigcup_{j \in J} V_j$ および開被覆
$f^{-1}(V_j) = \bigcup_{i \in I_j} U_i$ が存在し、
各射 $U_i \to V_j$, $j\in J, i\in I_j$ は syntomic である。
\item アフィン開被覆 $S = \bigcup_{j \in J} V_j$ およびアフィン開被覆
$f^{-1}(V_j) = \bigcup_{i \in I_j} U_i$ が存在し、環準同型
$\mathcal{O}_S(V_j) \to \mathcal{O}_X(U_i)$ は
すべての $j\in J, i\in I_j$ に対して syntomic である。
\end{enumerate}
さらに、$f$ が syntomic ならば、任意の開部分スキーム
$U \subset X$, $V \subset S$ で $f(U) \subset V$ を満たすものに対して、
制限 $f|_U : U \to V$ は syntomic である。
\end{lemma}

\begin{proof}
性質「$R \to A$ は syntomic である」が局所的であることを示せば、
補題 \ref{morphisms-lemma-locally-P} から従う。
定義 \ref{morphisms-definition-property-local} の条件 (a), (b), (c) を確認する。
代数、補題 \ref{algebra-lemma-base-change-syntomic} により、
syntomic 性は基底変換で保たれるので、(a) が成り立つ。
代数、補題 \ref{algebra-lemma-composition-syntomic} により、
syntomic 性は合成で保たれる。また、任意の環 $R$ について
環準同型 $R \to R_f$ は明らかに syntomic である。
従って (b) が成り立つ。最後に、性質 (c) は代数、補題
\ref{algebra-lemma-local-syntomic} によって成り立つ。
\end{proof}

\begin{lemma}
\label{morphisms-lemma-composition-syntomic}
二つの syntomic 射の合成は syntomic である。
\end{lemma}

\begin{proof}
補題 \ref{morphisms-lemma-syntomic-characterize} の証明において、
syntomic 性が環準同型の局所的性質であることを見た。
従って、この補題の最初の主張は補題 \ref{morphisms-lemma-composition-type-P} と、
syntomic 性が合成で保たれる環準同型の性質であること
（代数、補題 \ref{algebra-lemma-composition-syntomic} 参照）を合わせると従う。
\end{proof}

\begin{lemma}
\label{morphisms-lemma-base-change-syntomic}
syntomic 射の基底変換は syntomic である。
\end{lemma}

\begin{proof}
補題 \ref{morphisms-lemma-syntomic-characterize} の証明において、
syntomic 性が環準同型の局所的性質であることを見た。
従って、この補題は補題 \ref{morphisms-lemma-composition-type-P} と、
syntomic 性が基底変換で保たれる環準同型の性質であること
（代数、補題 \ref{algebra-lemma-base-change-syntomic} 参照）を合わせると従う。
\end{proof}

\begin{lemma}
\label{morphisms-lemma-open-immersion-syntomic}
任意の開埋め込みは syntomic である。
\end{lemma}

\begin{proof}
開埋め込みは局所同型なので成り立つ。
\end{proof}

\begin{lemma}
\label{morphisms-lemma-syntomic-locally-finite-presentation}
syntomic 射は局所有限表示である。
\end{lemma}

\begin{proof}
syntomic 環準同型は定義により有限表示なので成り立つ。
\end{proof}

\begin{lemma}
\label{morphisms-lemma-syntomic-flat}
syntomic 射は平坦である。
\end{lemma}

\begin{proof}
syntomic 環準同型は定義により平坦なので成り立つ。
\end{proof}

\begin{lemma}
\label{morphisms-lemma-syntomic-open}
syntomic 射は普遍開である。
\end{lemma}

\begin{proof}
補題 \ref{morphisms-lemma-syntomic-locally-finite-presentation},
\ref{morphisms-lemma-syntomic-flat} および \ref{morphisms-lemma-fppf-open} を組み合わせる。
\end{proof}

\noindent
$k$ を体とする。$A$ を局所 $k$-代数であって $k$ 上本質的有限型のものとする。
$A$ が{\it $k$ 上の完全交叉}であるとは、
$A \cong R/(f_1, \ldots, f_c)$ と書けることをいう。ここで
$R$ は $k$ 上本質的有限型の正則局所環であり、
$f_1, \ldots, f_c$ は $R$ の正則列である。
代数、定義 \ref{algebra-definition-lci-local-ring} 参照。

\begin{lemma}
\label{morphisms-lemma-local-complete-intersection}
$k$ を体とする。
$X$ を $k$ 上局所有限型のスキームとする。
次は同値である。
\begin{enumerate}
\item $X$ は $k$ 上の局所完全交叉である。
\item すべての $x \in X$ に対して、アフィン開集合
$U = \Spec(R) \subset X$ が $x$ の近傍として存在し、
$R \cong k[x_1, \ldots, x_n]/(f_1, \ldots, f_c)$ は
$k$ 上の大域完全交叉である。
\item すべての $x \in X$ に対して、局所環 $\mathcal{O}_{X, x}$ は
$k$ 上の完全交叉である。
\end{enumerate}
\end{lemma}

\begin{proof}
対応する代数の結果は、代数、補題 \ref{algebra-lemma-lci-at-prime}
および \ref{algebra-lemma-lci-global} にある。
\end{proof}

\noindent
次の補題によれば、任意の syntomic 射は局所的に標準 syntomic である。
従って、標準 syntomic 射を syntomic 射の{\it 局所モデル}として用いることができる。
さらに、有限表示の平坦射が syntomic であることと、
ファイバーが局所完全交叉スキームであることは同値である。

\begin{lemma}
\label{morphisms-lemma-syntomic-locally-standard-syntomic}
$f : X  \to S$ をスキームの射とする。点 $x \in X$ の像を
$s = f(x)$ とする。$V \subset S$ を $s$ のアフィン開近傍とする。
次は同値である。
\begin{enumerate}
\item 射 $f$ は $x$ で syntomic である。
\item アフィン開集合 $U \subset X$ が存在し、$x \in U$ かつ
$f(U) \subset V$ であり、$f|_U : U \to V$ は標準 syntomic である。
\item 射 $f$ は $x$ において有限表示であり、局所環準同型
$\mathcal{O}_{S, s} \to \mathcal{O}_{X, x}$ は平坦であり、
$\mathcal{O}_{X, x}/\mathfrak m_s \mathcal{O}_{X, x}$ は
$\kappa(s)$ 上の完全交叉である
（代数、定義 \ref{algebra-definition-lci-local-ring} 参照）。
\end{enumerate}
\end{lemma}

\begin{proof}
定義および代数、補題 \ref{algebra-lemma-syntomic} から従う。
\end{proof}

\begin{lemma}
\label{morphisms-lemma-syntomic-flat-fibres}
$f : X \to S$ をスキームの射とする。
$f$ が平坦かつ局所有限表示であり、すべてのファイバー $X_s$ が
局所完全交叉ならば、$f$ は syntomic である。
\end{lemma}

\begin{proof}
補題 \ref{morphisms-lemma-local-complete-intersection} および
\ref{morphisms-lemma-syntomic-locally-standard-syntomic} と、局所環の同型
$
\mathcal{O}_{X, x}/\mathfrak m_s \mathcal{O}_{X, x}
\cong
\mathcal{O}_{X_s, x}
$
から明らかである。
\end{proof}

\begin{lemma}
\label{morphisms-lemma-set-points-where-fibres-lci}
$f : X \to S$ をスキームの射とする。
$f$ は局所有限型であると仮定する。集合
$$
T = \{x \in X \mid \mathcal{O}_{X_{f(x)}, x}
\text{ is a complete intersection over }\kappa(f(x))\}
$$
の形成は任意の基底変換と可換である。
すなわち、任意の射 $g : S' \to S$ に対して、
基底変換 $f' : X' \to S'$ を $f$ から得て、
射影 $g' : X' \to X$ を考える。このとき対応する集合 $T'$ は、
射 $f'$ に対して $T' = (g')^{-1}(T)$ を満たす。
特に、$f$ が平坦かつ局所有限表示であると仮定すると、
$f$ が syntomic である点の開集合についても同じことが成り立つ。
\end{lemma}

\begin{proof}
$s' \in S'$ を点とし、$s = g(s')$ とおく。このとき
$$
X'_{s'} =
\Spec(\kappa(s')) \times_{\Spec(\kappa(s))} X_s
$$
である。言い換えると、基底変換のファイバーはファイバーの基底変換である。
従って最初の主張は、代数、補題 \ref{algebra-lemma-lci-field-change-local} と同値である。
第二の主張は最初の主張から従う。実際、この場合、$T$ は
$f$ が syntomic である点の集合であることが補題
\ref{morphisms-lemma-syntomic-locally-standard-syntomic} により分かる。
\end{proof}

\begin{lemma}
\label{morphisms-lemma-standard-syntomic-relative-dimension}
$R$ を環とする。
$R \to A = R[x_1, \ldots, x_n]/(f_1, \ldots, f_c)$ を相対大域完全交叉とする。
$S = \Spec(R)$ および $X = \Spec(A)$ とおく。
射 $f : X \to S$ で環準同型 $R \to A$ に付随するものを考える。
関数 $x \mapsto \dim_x(X_{f(x)})$ は定数であり、その値は $n - c$ である。
\end{lemma}

\begin{proof}
代数、定義 \ref{algebra-definition-relative-global-complete-intersection} により、
$R \to A$ が相対大域完全交叉であるということは、
零でないすべてのファイバー環の次元が $n - c$ であることを意味する。
従って、$\mathfrak p$ を $R$ の素イデアルとすると、ファイバー環
$\kappa(\mathfrak p)[x_1, \ldots, x_n]/(\overline{f}_1, \ldots, \overline{f}_c)$
は零であるか、次元 $n - c$ の大域完全交叉環である。
代数、定義 \ref{algebra-definition-lci-field} の後の議論により、
このことから、それは次元 $n - c$ の等次元環であることが分かる。
よって補題が従う。
\end{proof}

\begin{lemma}
\label{morphisms-lemma-syntomic-relative-dimension}
$f : X \to S$ を syntomic 射とする。
関数 $x \mapsto \dim_x(X_{f(x)})$ は $X$ 上局所定数である。
\end{lemma}

\begin{proof}
補題 \ref{morphisms-lemma-syntomic-locally-standard-syntomic} により、
射 $f$ は局所的にアフィンスキーム間の標準 syntomic 射の形をしている。
従って主張は補題 \ref{morphisms-lemma-standard-syntomic-relative-dimension} から従う。
\end{proof}

\noindent
補題 \ref{morphisms-lemma-syntomic-relative-dimension} により、次の定義は意味をもつ。

\begin{definition}
\label{morphisms-definition-syntomic-relative-dimension}
$d \geq 0$ を整数とする。スキームの射 $f : X \to S$ が
{\it 相対次元 $d$ の syntomic 射}であるとは、$f$ が syntomic であり、
関数について $\dim_x(X_{f(x)}) = d$ がすべての $x \in X$ に対して成り立つことをいう。
\end{definition}

\noindent
言い換えると、$f$ は syntomic であり、空でないファイバーは次元 $d$ の等次元スキームである。

\begin{lemma}
\label{morphisms-lemma-syntomic-permanence}
スキームの射の可換図式
$$
\xymatrix{
X \ar[rr]_f \ar[rd]_p & &
Y \ar[dl]^q \\
& S
}
$$
を考える。次を仮定する。
\begin{enumerate}
\item $f$ は全射かつ syntomic である。
\item $p$ は syntomic である。
\item $q$ は局所有限表示である\footnote{実際には、$f$ が
全射、平坦かつ局所有限表示であり、$p$ が syntomic ならば、
$q$ と $f$ はともに syntomic である。降下、補題
\ref{descent-lemma-syntomic-permanence} 参照。}。
\end{enumerate}
このとき $q$ は syntomic である。
\end{lemma}

\begin{proof}
補題 \ref{morphisms-lemma-flat-permanence} により $q$ は平坦である。
従って $Y \to S$ のファイバーが局所完全交叉であることを示せば十分である。
補題 \ref{morphisms-lemma-syntomic-flat-fibres} 参照。
$s \in S$ とし、射 $X_s \to Y_s$ を考える。
これは射 $X \to Y$ の基底変換なので、全射かつ syntomic である
（補題 \ref{morphisms-lemma-base-change-syntomic}）。
同じ理由により $X_s$ は $\kappa(s)$ 上 syntomic である。
さらに $Y_s$ は $\kappa(s)$ 上局所有限型である
（補題 \ref{morphisms-lemma-base-change-finite-type}）。
このようにして、$S$ が体のスペクトルである場合に帰着する。

\medskip\noindent
$S = \Spec(k)$ と仮定する。$y \in Y$ とする。
アフィン開集合 $\Spec(A) \subset Y$ を $y$ の近傍として取る。
$\Spec(B) \subset X$ をアフィン開集合であって、
$f(\Spec(B)) \subset \Spec(A)$ であり、点 $x \in X$ で
$f(x) = y$ を満たすものを含むように取る。
全射 $k[x_1, \ldots, x_n] \to A$ を取り、その核を $I$ とする。
全射 $A[y_1, \ldots, y_m] \to B$ を取ると、さらに全射
$k[x_i, y_j] \to B$ が得られ、その核を $J$ とする。
素イデアル $\mathfrak q \subset k[x_i, y_j]$ を
$y \in \Spec(B)$ に対応するものとし、
素イデアル $\mathfrak p \subset k[x_i]$ を $x \in \Spec(A)$ に対応するものとする。
$x$ は $y$ に写るので、$\mathfrak p = \mathfrak q \cap k[x_i]$ である。
次の局所環の可換図式を考える。
$$
\xymatrix{
\mathcal{O}_{X, x} \ar@{=}[r] &
B_{\mathfrak q} &
k[x_1, \ldots, x_n, y_1, \ldots, y_m]_{\mathfrak q} \ar[l] \\
\mathcal{O}_{Y, y} \ar@{=}[r] \ar[u] & A_{\mathfrak p} \ar[u] &
k[x_1, \ldots, x_n]_{\mathfrak p} \ar[l] \ar[u]
}
$$
代数、補題 \ref{algebra-lemma-lci-permanence-initial} の仮定が満たされると主張する。
条件 (1) と (2) は明らかである。条件 (4) は $X \to Y$ が平坦であることから従う。
条件 (3) は、環 $\mathcal{O}_{Y, y}$ および
$\mathcal{O}_{X_y, x} = \mathcal{O}_{X, x}/\mathfrak m_y\mathcal{O}_{X, x}$ が、
$f$ および $p$ は syntomic という仮定により完全交叉環であることから従う。
補題 \ref{morphisms-lemma-syntomic-locally-standard-syntomic} 参照。
代数、補題 \ref{algebra-lemma-lci-permanence-initial} の結論は、まさに
$\mathcal{O}_{Y, y}$ が完全交叉環であるということである！
従って再び補題 \ref{morphisms-lemma-syntomic-locally-standard-syntomic} により、
$Y$ は $k$ 上、点 $y$ で syntomic であり、求める結果が得られる。
\end{proof}

\section{埋め込みの余法層}
\label{morphisms-section-conormal-sheaf}

\noindent
$i : Z \to X$ を閉埋め込みとする。
$\mathcal{I} \subset \mathcal{O}_X$ を対応する準連接イデアル層とする。
短完全列
$$
0 \to \mathcal{I}^2 \to \mathcal{I} \to \mathcal{I}/\mathcal{I}^2 \to 0
$$
を $X$ 上の準連接層の列として考える。
層 $\mathcal{I}/\mathcal{I}^2$ は $\mathcal{I}$ により零化されるので、
補題 \ref{morphisms-lemma-i-star-equivalence} により $Z$ 上の層に対応する。
この準連接 $\mathcal{O}_Z$-加群を
{\it $Z$ の $X$ における余法層}と呼ぶ。
しばしば、節 \ref{morphisms-section-closed-immersions-quasi-coherent} で述べた
記法の濫用によって、単に $\mathcal{I}/\mathcal{I}^2$ と書く。

\medskip\noindent
$i : Z \to X$ が（局所閉）埋め込みである場合には、
$i$ の余法層を、閉埋め込み $i : Z \to X \setminus \partial Z$ の余法層として定義する。
ここで $\partial Z = \overline{Z} \setminus Z$ である。
これをしばしば $\mathcal{I}/\mathcal{I}^2$ と書く。ただし $\mathcal{I}$ は
閉埋め込み $i : Z \to X \setminus \partial Z$ のイデアル層である。

\begin{definition}
\label{morphisms-definition-conormal-sheaf}
$i : Z \to X$ を埋め込みとする。
{\it 余法層 $\mathcal{C}_{Z/X}$（$Z$ の $X$ におけるもの）}、
あるいは{\it $i$ の余法層}とは、上で述べた準連接
$\mathcal{O}_Z$-加群 $\mathcal{I}/\mathcal{I}^2$ である。
\end{definition}

\noindent
\cite[IV Definition 16.1.2]{EGA} では、この層を $\mathcal{N}_{Z/X}$ と書く。
ここではこの慣習には従わない。記号 $\mathcal{N}_{Z/X}$ は
{\it 埋め込みの法層}のために取っておきたいからである。
これは
$$
\mathcal{N}_{Z/X} =
\SheafHom_{\mathcal{O}_Z}(\mathcal{C}_{Z/X}, \mathcal{O}_Z) =
\SheafHom_{\mathcal{O}_Z}(\mathcal{I}/\mathcal{I}^2, \mathcal{O}_Z)
$$
と定義される。ただし余法層は有限表示であるとする
（そうでなければ、法層は準連接ですらないかもしれない）。
法層については後で再び扱う（ここに後の参照を挿入）。

\begin{lemma}
\label{morphisms-lemma-affine-conormal}
$i : Z \to X$ を埋め込みとする。
$i$ の余法層は次の性質をもつ。
\begin{enumerate}
\item $U \subset X$ を任意の開部分スキームであって、
$i$ が $Z \xrightarrow{i'} U \to X$ と分解し、
$i'$ が閉埋め込みとなるものとする。
$\mathcal{I} = \Ker((i')^\sharp) \subset \mathcal{O}_U$ とおく。このとき
$$
\mathcal{C}_{Z/X} = (i')^*\mathcal{I}\quad\text{and}\quad
i'_*\mathcal{C}_{Z/X} = \mathcal{I}/\mathcal{I}^2
$$
\item 任意のアフィン開集合 $\Spec(R) = U \subset X$ で
$Z \cap U = \Spec(R/I)$ を満たすものに対して、標準同型
$\Gamma(Z \cap U, \mathcal{C}_{Z/X}) = I/I^2$ が存在する。
\end{enumerate}
\end{lemma}

\begin{proof}
ほとんど定義から明らかである。環 $R$ とイデアル $I$ が $R$ に与えられたとき、
$I/I^2 = I \otimes_R R/I$ であることに注意する。詳細は省略する。
\end{proof}

\begin{lemma}
\label{morphisms-lemma-conormal-functorial}
スキームの圏における可換図式
$$
\xymatrix{
Z \ar[r]_i \ar[d]_f & X \ar[d]^g \\
Z' \ar[r]^{i'} & X'
}
$$
を考える。$i$, $i'$ は埋め込みであると仮定する。
$\mathcal{O}_Z$-加群の標準写像
$$
f^*\mathcal{C}_{Z'/X'}
\longrightarrow
\mathcal{C}_{Z/X}
$$
が存在し、次の性質で特徴づけられる。任意のアフィン開集合の対
$(\Spec(R) = U \subset X, \Spec(R') = U' \subset X')$ で
$g(U) \subset U'$ かつ
$Z \cap U = \Spec(R/I)$ および $Z' \cap U' = \Spec(R'/I')$
を満たすものについて、誘導写像
$$
\Gamma(Z' \cap U', \mathcal{C}_{Z'/X'}) = I'/I'^2
\longrightarrow
I/I^2 = \Gamma(Z \cap U, \mathcal{C}_{Z/X})
$$
は環準同型 $g^\sharp : R' \to R$ から誘導される写像である。
この環準同型は $g^\sharp(I') \subset I$ を満たす。
\end{lemma}

\begin{proof}
$\partial Z' = \overline{Z'} \setminus Z'$ および
$\partial Z = \overline{Z} \setminus Z$ とおく。
これらはそれぞれ $X'$ および $X$ の閉部分集合である。
$X'$ を $X' \setminus \partial Z'$ で、$X$ を
$X \setminus \big(g^{-1}(\partial Z') \cup \partial Z\big)$ で置き換えることで、
$i$ および $i'$ は閉埋め込みであると仮定してよい。

\medskip\noindent
$g \circ i$ が $i'$ を経由することから、
$g^*\mathcal{I}'$ は $\mathcal{I}$ に、標準写像
$g^*\mathcal{I}' \to \mathcal{O}_X$ の下で写る。
スキーム、補題 \ref{schemes-lemma-characterize-closed-subspace} および
\ref{schemes-lemma-restrict-map-to-closed} 参照。
従って準連接層の誘導写像
$g^*(\mathcal{I}'/(\mathcal{I}')^2) \to \mathcal{I}/\mathcal{I}^2$ が得られる。
$i$ によって引き戻すと、
$i^*g^*(\mathcal{I}'/(\mathcal{I}')^2) \to i^*(\mathcal{I}/\mathcal{I}^2)$ を得る。
$i^*(\mathcal{I}/\mathcal{I}^2) = \mathcal{C}_{Z/X}$ に注意する。一方、
$i^*g^*(\mathcal{I}'/(\mathcal{I}')^2) =
f^*(i')^*(\mathcal{I}'/(\mathcal{I}')^2) = f^*\mathcal{C}_{Z'/X'}$ である。
これにより、求める写像が得られる。

\medskip\noindent
この写像が局所的に上記の写像 $I'/(I')^2 \to I/I^2$ で記述されることの確認は、
定義を展開するだけなので省略する。
また、任意の $x \in i(Z)$ に対し、アフィン開近傍 $U$, $U'$ が存在し、
$f(U) \subset U'$ であり、$Z \cap U$ および $U' \cap Z'$ は閉であり、
$x \in U$ となる。証明は省略する。
従って、補題の条件は確かに写像を特徴づける
（そしてこの条件を用いて写像を定義することもできた）。
\end{proof}

\begin{lemma}
\label{morphisms-lemma-conormal-functorial-flat}
スキームの圏におけるファイバー積図式
$$
\xymatrix{
Z \ar[r]_i \ar[d]_f & X \ar[d]^g \\
Z' \ar[r]^{i'} & X'
}
$$
を考え、$i$, $i'$ は埋め込みとする。
このとき、標準写像 $f^*\mathcal{C}_{Z'/X'} \to \mathcal{C}_{Z/X}$、
すなわち補題 \ref{morphisms-lemma-conormal-functorial} の写像は全射である。
$g$ が平坦ならば、これは同型である。
\end{lemma}

\begin{proof}
$R' \to R$ を環準同型とし、$I' \subset R'$ をイデアルとする。
$I = I'R$ とおく。このとき $I'/(I')^2 \otimes_{R'} R \to I/I^2$ は全射である。
$R' \to R$ が平坦ならば、$I = I' \otimes_{R'} R$ および
$I^2 = (I')^2 \otimes_{R'} R$ であり、この写像は同型であることが分かる。
\end{proof}

\begin{lemma}
\label{morphisms-lemma-transitivity-conormal}
$Z \to Y \to X$ をスキームの埋め込みとする。このとき標準完全列
$$
i^*\mathcal{C}_{Y/X} \to
\mathcal{C}_{Z/X} \to
\mathcal{C}_{Z/Y} \to 0
$$
が存在する。ここで写像は補題 \ref{morphisms-lemma-conormal-functorial} から得られ、
$i : Z \to Y$ は最初の射である。
\end{lemma}

\begin{proof}
補題 \ref{morphisms-lemma-conormal-functorial} により、これは次の代数的事実に言い換えられる。
$C \to B \to A$ を全射環準同型とする。
$I = \Ker(B \to A)$, $J = \Ker(C \to A)$ および $K = \Ker(C \to B)$ とおく。
このとき完全列
$$
K/K^2 \otimes_B A \to J/J^2 \to I/I^2 \to 0.
$$
が存在する。これは $I = J/K$ という観察から直ちに従う。
\end{proof}

\section{射の微分層}
\label{morphisms-section-sheaf-differentials}

\noindent
可換代数の章の対応する節（代数、節 \ref{algebra-section-differentials}）
および加群の層の章の対応する節（加群、節
\ref{modules-section-differentials}）に目を通すことを勧める。

\begin{definition}
\label{morphisms-definition-sheaf-differentials}
$f : X \to S$ をスキームの射とする。
{\it 微分層 $\Omega_{X/S}$（$X$ の $S$ 上のもの）}とは、
$f$ を環付き空間の射とみなしたときの微分層
（加群、定義 \ref{modules-definition-differentials}）であり、
その{\it 普遍 $S$-導分}
$$
\text{d}_{X/S} : \mathcal{O}_X \longrightarrow \Omega_{X/S}.
$$
を備えたものである。
\end{definition}

\noindent
$\Omega_{X/S}$ は準連接 $\mathcal{O}_X$-加群である。
これは例えば、対角射 $\Delta : X \to X \times_S X$ の余法層と同型であること
（補題 \ref{morphisms-lemma-differentials-diagonal}）から分かる。
ここでは $X$ の $S$ 上の微分加群を、普遍性、すなわち普遍導分の値域であることを
用いて定義した。相対微分層の別の構成がこの普遍性を満たすならば、
米田の補題により、それは上で定義したものと標準的に同型となる。
便宜のため、ここで普遍性を再掲する。

\begin{lemma}
\label{morphisms-lemma-universal-derivation-universal}
$f : X \to S$ をスキームの射とする。写像
$$
\Hom_{\mathcal{O}_X}(\Omega_{X/S}, \mathcal{F})
\longrightarrow
\text{Der}_S(\mathcal{O}_X, \mathcal{F}),\quad
\alpha \longmapsto \alpha \circ \text{d}_{X/S}
$$
は関手 $\textit{Mod}(\mathcal{O}_X) \to \textit{Sets}$ の同型である。
\end{lemma}

\begin{proof}
これは定義の言い換えにすぎない。
\end{proof}

\begin{lemma}
\label{morphisms-lemma-differentials-restrict-open}
$f : X \to S$ をスキームの射とする。
$U \subset X$, $V \subset S$ を開部分スキームで
$f(U) \subset V$ を満たすものとする。このとき、同型
$\Omega_{X/S}|_U = \Omega_{U/V}$ が $\mathcal{O}_U$-加群の同型として一意に存在し、
$\text{d}_{X/S}|_U = \text{d}_{U/V}$ を満たす。
\end{lemma}

\begin{proof}
標準的な同一視
$f^{-1}\mathcal{O}_S|_U = (f|_U)^{-1}\mathcal{O}_V$ を用いれば、
これは加群、補題 \ref{modules-lemma-localize-differentials} の特別な場合である。
\end{proof}

\noindent
今後はこれらの標準的な同一視を用い、$\Omega_{U/S}$ あるいは $\Omega_{U/V}$ と
単に書くことで、$\Omega_{X/S}$ の $U$ への制限を表す。

\begin{lemma}
\label{morphisms-lemma-affine-case-derivation}
$R \to A$ を環準同型とする。$\mathcal{F}$ を
$\mathcal{O}_X$-加群の層で $X = \Spec(A)$ 上のものとする。
$S = \Spec(R)$ とおく。
$S$-導分で $\mathcal{F}$ に値をもつものに、その大域切断への作用を対応させる規則は、
$S$-導分で $\mathcal{F}$ に値をもつものの集合と、
$R$-導分で $M = \Gamma(X, \mathcal{F})$ に値をもつものの集合との間の全単射を定める。
\end{lemma}

\begin{proof}
$D : A \to M$ を $R$-導分とする。
一意な $S$-導分で $\mathcal{F}$ に値をもち、
大域切断上で $D$ を与えるものが存在することを示す。
$U = D(f) \subset X$ を標準アフィン開集合とする。
$\Gamma(U, \mathcal{O}_X)$ の任意の元は $a/f^n$ という形であり、
ここで $a \in A$ および $n \geq 0$ である。Leibniz 則により
$$
D(a)|_U = a/f^n D(f^n)|_U + f^n D(a/f^n)
$$
が $\Gamma(U, \mathcal{F})$ において成り立つ。
$f$ は $\Gamma(U, \mathcal{F})$ に可逆に作用するので、
これは $D(a/f^n) \in \Gamma(U, \mathcal{F})$ の値を完全に決定する。
これで一意性が示された。存在については、単に
$$
D(a/f^n) := (1/f^n) D(a)|_U - a/f^{2n} D(f^n)|_U
$$
と定義し、これが必要なすべての性質を満たすことを
（$X$ の標準開集合からなる基底上で）証明すればよい。詳細は省略する。
\end{proof}

\begin{lemma}
\label{morphisms-lemma-differentials-affine}
$f : X \to S$ をスキームの射とする。任意のアフィン開集合の対
$\Spec(A) = U \subset X$, $\Spec(R) = V \subset S$ で $f(U) \subset V$
を満たすものに対して、一意な同型
$$
\Gamma(U, \Omega_{X/S}) = \Omega_{A/R}.
$$
が存在し、$\text{d}_{X/S}$ および $\text{d} : A \to \Omega_{A/R}$ と両立する。
\end{lemma}

\begin{proof}
補題 \ref{morphisms-lemma-differentials-restrict-open} により、
$X$ および $S$ を $U$ および $V$ で置き換えてよい。
従って $X = \Spec(A)$ および $S = \Spec(R)$ と仮定でき、
$U = X$ かつ $V = S$ として補題を示せばよい。
$A$-加群 $M = \Gamma(X, \Omega_{X/S})$ を、
$R$-導分 $\text{d}_{X/S} : A \to M$ とともに考える。
$N$ を別の $A$-加群とし、$\widetilde{N}$ を準連接
$\mathcal{O}_X$-加群で $N$ に付随するものとする。
スキーム、節 \ref{schemes-section-quasi-coherent-affine} 参照。
$\text{d}_{X/S} : A \to M$ との前合成によって、矢
$$
\alpha : \Hom_A(M, N) \longrightarrow \text{Der}_R(A, N)
$$
を得る。補題 \ref{morphisms-lemma-universal-derivation-universal} および
\ref{morphisms-lemma-affine-case-derivation} を用いると、同一視
$$
\Hom_{\mathcal{O}_X}(\Omega_{X/S}, \widetilde{N}) =
\text{Der}_S(\mathcal{O}_X, \widetilde{N}) =
\text{Der}_R(A, N)
$$
を得る。大域切断を取ることにより、矢
$\Hom_{\mathcal{O}_X}(\Omega_{X/S}, \widetilde{N}) \to \Hom_R(M, N)$ が定まる。
この矢と上の同一視を組み合わせると、矢
$$
\beta : \text{Der}_R(A, N) \longrightarrow \Hom_R(M, N)
$$
を得る。大域切断上での作用を調べると、$\alpha$ と $\beta$ は互いに逆である。
従って $\text{d}_{X/S} : A \to M$ は
$\text{d} : A \to \Omega_{A/R}$ と同じ普遍性を満たす。
代数、補題 \ref{algebra-lemma-universal-omega} 参照。
よって米田の補題（圏、補題 \ref{categories-lemma-yoneda}）から、
$A$-加群の一意な同型 $M \cong \Omega_{A/R}$ で導分と両立するものが存在する。
\end{proof}

\begin{remark}
\label{morphisms-remark-differentials-glue}
上の補題は、微分加群を構成する第二の方法を与える。
すなわち、$f : X \to S$ をスキームの射とする。
すべてのアフィン開集合 $U \subset X$ で、$S$ のあるアフィン開集合に写るものを考える。
これらは $X$ の位相の基底をなす。
従って $\Gamma(U, \Omega_{X/S})$ をそのような $U$ に対して定めれば十分である。
単に $\Gamma(U, \Omega_{X/S}) = \Omega_{A/R}$ とおく。
ここで $A$, $R$ は上の補題 \ref{morphisms-lemma-differentials-affine} のとおりである。
この方法でも構成できるが、制限写像を構成し層条件を確認するには、
やや多くの代数的準備が必要となる。
\end{remark}

\noindent
次の補題は、微分層を定義するさらに別の方法を与える。
特に、$\Omega_{X/S}$ が準連接であることを、$X$ と $S$ がスキームの場合に示す。

\begin{lemma}
\label{morphisms-lemma-differentials-diagonal}
$f : X \to S$ をスキームの射とする。
$\Omega_{X/S}$ と対角射
$\Delta_{X/S} : X \longrightarrow X \times_S X$ の余法層との間に標準同型が存在する。
\end{lemma}

\begin{proof}
まず幾つかの「大域的な」層と層の大域的写像の存在を示し、
その後で、幾つかのアフィン開集合上でこの構成を記述する。

\medskip\noindent
$\Delta = \Delta_{X/S} : X \to X \times_S X$ は埋め込みであることを思い出そう。
スキーム、補題 \ref{schemes-lemma-diagonal-immersion} 参照。
$\mathcal{J}$ をこの埋め込みのイデアル層とする。
これはある開部分スキーム $W$（$X \times_S X$ の中のもの）上にあり、
$\Delta(X) \subset W$ は閉である。
スキーム、補題 \ref{schemes-lemma-diagonal-immersion} の証明で
得られたものを取ることにする。
環の層 $\mathcal{O}_W/\mathcal{J}^2$ の台は $\Delta(X)$ に含まれることに注意する。
さらに、この層は層の短完全列
$$
0 \to \mathcal{J}/\mathcal{J}^2
\to \mathcal{O}_W/\mathcal{J}^2
\to \Delta_*\mathcal{O}_X
\to 0.
$$
の中に現れる。$\Delta^{-1}$ を用いると、これを
$f^{-1}\mathcal{O}_S$-代数の層の全射で、核が
$\Delta$ の余法層であるものとみなせる
（定義 \ref{morphisms-definition-conormal-sheaf} および補題 \ref{morphisms-lemma-affine-conormal} 参照）。
$$
0 \to \mathcal{C}_{X/X \times_S X}
\to \Delta^{-1}(\mathcal{O}_W/\mathcal{J}^2)
\to \mathcal{O}_X
\to 0
$$
これは加群、補題 \ref{modules-lemma-double-structure-gives-derivation} の状況である。
射影 $p_i : X \times_S X \to X$, $i = 1, 2$ は環の層の写像
$(p_i)^\sharp : (p_i)^{-1}\mathcal{O}_X \to \mathcal{O}_{X \times_S X}$ を誘導する。
$W$ に制限し、$\mathcal{J}^2$ による商を取ることで、
$(p_i)^{-1}\mathcal{O}_X \to \mathcal{O}_W/\mathcal{J}^2$ を得る。
$\Delta^{-1}p_i^{-1}\mathcal{O}_X = \mathcal{O}_X$ なので、写像
$$
s_i : \mathcal{O}_X \to \Delta^{-1}(\mathcal{O}_W/\mathcal{J}^2).
$$
を得る。$s_1$ と $s_2$ はともに写像
$\Delta^{-1}(\mathcal{O}_W/\mathcal{J}^2) \to \mathcal{O}_X$ の切断であり、
加群、補題 \ref{modules-lemma-double-structure-gives-derivation} の状況になっている。
従って $S$-導分
$\text{d} = s_2 - s_1 : \mathcal{O}_X \to \mathcal{C}_{X/X \times_S X}$ を得る。
微分加群の普遍性により、一意な $\mathcal{O}_X$-線形写像
$$
\Omega_{X/S} \longrightarrow \mathcal{C}_{X/X \times_S X},\quad
f\text{d}g \longmapsto fs_2(g) - fs_1(g)
$$
を得る。この写像が同型であることを見るため、適当なアフィン開集合上で計算しよう。
$X$ をアフィン開集合 $\Spec(A) = U \subset X$ で、
像がアフィン開集合 $\Spec(R) = V \subset S$ に含まれるものにより被覆できる。
スキーム、補題 \ref{schemes-lemma-diagonal-immersion} の証明により、
$U \times_V U \subset X \times_S X$ は、上の開集合 $W$ に含まれるアフィン開集合である。
また $U \times_V U = \Spec(A \otimes_R A)$ である。
層 $\mathcal{J}$ はイデアル $J = \Ker(A \otimes_R A \to A)$ に対応する。
短完全列は、$A \otimes_R A$-加群の短完全列
$$
0 \to J/J^2 \to (A \otimes_R A)/J^2 \to A \to 0
$$
に対応する。切断 $s_i$ は環準同型
$$
A \longrightarrow (A \otimes_R A)/J^2,\quad
s_1 : a \mapsto a \otimes 1,\quad
s_2 : a \mapsto 1 \otimes a.
$$
に対応する。補題 \ref{morphisms-lemma-affine-conormal} により
$\Gamma(U, \mathcal{C}_{X/X \times_S X}) = J/J^2$ であり、
補題 \ref{morphisms-lemma-differentials-affine} により
$\Gamma(U, \Omega_{X/S}) = \Omega_{A/R}$ である。
上の写像は $a \text{d}b \mapsto a \otimes b - ab \otimes 1$ という写像であり、
これが同型であることは代数、補題 \ref{algebra-lemma-differentials-diagonal} で示されている。
\end{proof}

\begin{lemma}
\label{morphisms-lemma-functoriality-differentials}
スキームの可換図式
$$
\xymatrix{
X' \ar[d] \ar[r]_f & X \ar[d] \\
S' \ar[r] & S
}
$$
を考える。標準写像 $\mathcal{O}_X \to f_*\mathcal{O}_{X'}$ と写像
$f_*\text{d}_{X'/S'} : f_*\mathcal{O}_{X'} \to f_*\Omega_{X'/S'}$ の合成は
$S$-導分である。従って標準的な $\mathcal{O}_X$-加群の写像
$\Omega_{X/S} \to f_*\Omega_{X'/S'}$ を得る。
$f_*$ と $f^*$ の随伴性により、標準的な $\mathcal{O}_{X'}$-加群準同型
$$
c_f : f^*\Omega_{X/S} \longrightarrow \Omega_{X'/S'}.
$$
を得る。これは、$f^*\text{d}_{X/S}(h)$ が $\text{d}_{X'/S'}(f^* h)$ に写るという性質が、
任意の局所切断 $h$（$\mathcal{O}_X$ のもの）に対して成り立つことにより一意に特徴づけられる。
\end{lemma}

\begin{proof}
これは加群、補題 \ref{modules-lemma-functoriality-differentials-ringed-spaces}
の特別な場合である。スキームの場合には、余法層の関手性
（補題 \ref{morphisms-lemma-conormal-functorial} 参照）と
補題 \ref{morphisms-lemma-differentials-diagonal} を用いて $c_f$ を定義することもできる。
あるいは、補題の最後の行の特徴づけを用いて、アフィン開集合上で定義された写像を
貼り合わせることもできる（代数、式 (\ref{algebra-equation-functorial-omega}) 参照）。
\end{proof}

\begin{lemma}
\label{morphisms-lemma-triangle-differentials}
$f : X \to Y$, $g : Y \to S$ をスキームの射とする。
このとき標準完全列
$$
f^*\Omega_{Y/S} \to \Omega_{X/S} \to \Omega_{X/Y} \to 0
$$
が存在する。ここで写像は補題 \ref{morphisms-lemma-functoriality-differentials} の適用から得られる。
\end{lemma}

\begin{proof}
これは代数、補題 \ref{algebra-lemma-exact-sequence-differentials} を層の言葉にしたものである。
また、環付き空間の射に対する一般形もある。
加群、補題 \ref{modules-lemma-triangle-differentials} 参照。
\end{proof}

\begin{lemma}
\label{morphisms-lemma-base-change-differentials}
$X \to S$ をスキームの射とする。
$g : S' \to S$ をスキームの射とする。
$X' = X_{S'}$ を $X$ の基底変換とする。
$g' : X' \to X$ を射影とする。このとき写像
$$
(g')^*\Omega_{X/S} \to \Omega_{X'/S'}
$$
は、補題 \ref{morphisms-lemma-functoriality-differentials} によるものとして同型である。
\end{lemma}

\begin{proof}
これは代数、補題 \ref{algebra-lemma-differentials-base-change} を層の言葉にしたものである。
\end{proof}

\begin{lemma}
\label{morphisms-lemma-differential-product}
$f : X \to S$ および $g : Y \to S$ を同じ終域をもつスキームの射とする。
$p : X \times_S Y \to X$ および $q : X \times_S Y \to Y$ を射影とする。
補題 \ref{morphisms-lemma-functoriality-differentials} から得られる写像
$$
p^*\Omega_{X/S} \oplus q^*\Omega_{Y/S}
\longrightarrow
\Omega_{X \times_S Y/S}
$$
は同型を与える。
\end{lemma}

\begin{proof}
補題 \ref{morphisms-lemma-base-change-differentials} により、合成
$p^*\Omega_{X/S} \to \Omega_{X \times_S Y/S} \to \Omega_{X \times_S Y/Y}$
は同型であり、$q$ についても同様である。
さらに $p^*\Omega_{X/S} \to \Omega_{X \times_S Y/S}$ の余核は
$\Omega_{X \times_S Y/X}$ であることが、
補題 \ref{morphisms-lemma-triangle-differentials} により分かる。これで主張が従う。
\end{proof}

\begin{lemma}
\label{morphisms-lemma-finite-type-differentials}
$f : X \to S$ をスキームの射とする。
$f$ が局所有限型ならば、$\Omega_{X/S}$ は有限型の $\mathcal{O}_X$-加群である。
\end{lemma}

\begin{proof}
代数、補題 \ref{algebra-lemma-differentials-finitely-generated}、
補題 \ref{morphisms-lemma-differentials-affine}、
補題 \ref{morphisms-lemma-locally-finite-type-characterize} および
性質、補題 \ref{properties-lemma-finite-type-module} から直ちに従う。
\end{proof}

\begin{lemma}
\label{morphisms-lemma-finite-presentation-differentials}
$f : X \to S$ をスキームの射とする。
$f$ が局所有限表示ならば、$\Omega_{X/S}$ は
有限表示の $\mathcal{O}_X$-加群である。
\end{lemma}

\begin{proof}
代数、補題 \ref{algebra-lemma-differentials-finitely-presented}、
補題 \ref{morphisms-lemma-differentials-affine}、
補題 \ref{morphisms-lemma-locally-finite-presentation-characterize} および
性質、補題 \ref{properties-lemma-finite-presentation-module} から直ちに従う。
\end{proof}

\begin{lemma}
\label{morphisms-lemma-immersion-differentials}
$X \to S$ が埋め込み、より一般にモノ射であれば、
$\Omega_{X/S}$ は零である。
\end{lemma}

\begin{proof}
この場合、$\Delta_{X/S}$ は同型であり、従ってその余法層は自明である。
よって $\Omega_{X/S} = 0$ が補題 \ref{morphisms-lemma-differentials-diagonal} により従う。
代数的な形は、代数、補題 \ref{algebra-lemma-trivial-differential-surjective} である。
\end{proof}

\begin{lemma}
\label{morphisms-lemma-differentials-relative-immersion}
$i : Z \to X$ を $S$ 上のスキームの埋め込みとする。
標準完全列
$$
\mathcal{C}_{Z/X} \to i^*\Omega_{X/S} \to \Omega_{Z/S} \to 0
$$
が存在する。ここで最初の矢は $\text{d}_{X/S}$ から誘導され、
第二の矢は補題 \ref{morphisms-lemma-functoriality-differentials} から得られる。
\end{lemma}

\begin{proof}
これは代数、補題 \ref{algebra-lemma-differential-seq} を層の言葉にしたものである。
ただし、最初の矢を大域的に定義できることを確認する必要がある。
そこで「$\text{d}_{X/S}$ から誘導される」という意味をここで説明する。
すなわち、$i$ は閉埋め込みであると、$X$ を縮小することで仮定してよい。
$\mathcal{I} \subset \mathcal{O}_X$ を $Z \subset X$ に対応するイデアル層とする。
このとき $\text{d}_{X/S} : \mathcal{I} \to \Omega_{X/S}$ は
部分層 $\mathcal{I}^2 \subset \mathcal{I}$ を $\mathcal{I}\Omega_{X/S}$ に写す。
従って、写像
$\mathcal{I}/\mathcal{I}^2 \to \Omega_{X/S}/\mathcal{I}\Omega_{X/S}$ が誘導され、
これは $\mathcal{O}_X/\mathcal{I}$-線形である。
補題 \ref{morphisms-lemma-i-star-equivalence} により、これは求める写像
$\mathcal{C}_{Z/X} \to i^*\Omega_{X/S}$ に対応する。
\end{proof}

\begin{lemma}
\label{morphisms-lemma-differentials-relative-immersion-section}
$i : Z \to X$ を $S$ 上のスキームの埋め込みとし、
$i$ は（局所的に）左逆をもつと仮定する。
このとき標準列
$$
0 \to \mathcal{C}_{Z/X} \to i^*\Omega_{X/S} \to \Omega_{Z/S} \to 0
$$
は、補題 \ref{morphisms-lemma-differentials-relative-immersion} の列として、
（局所的に）分裂完全である。
特に、$s : S \to X$ が構造射 $X \to S$ の切断ならば、
写像 $\mathcal{C}_{S/X} \to s^*\Omega_{X/S}$ で $\text{d}_{X/S}$ から誘導されるものは同型である。
\end{lemma}

\begin{proof}
代数、補題 \ref{algebra-lemma-differential-seq-split} により、
$X, Z, S$ がアフィンの場合には主張が成り立つ。
一般の場合には、列は局所的に分裂完全であると結論できる。
$g : X \to Z$ が $i$ の左逆ならば、
$i^*c_g$ は写像 $i^*\Omega_{X/S} \to \Omega_{Z/S}$ の右逆である。
これは加群、補題 \ref{modules-lemma-check-functoriality-differentials} による。
従って、ホモロジー、補題 \ref{homology-lemma-ses-split} により列は分裂する。
最後に、$s$ が切断ならば、それは埋め込み $s : Z = S \to X$ であり、
$S$ 上の射である（スキーム、補題 \ref{schemes-lemma-section-immersion} 参照）。
この場合、$\Omega_{Z/S} = 0$ である。
\end{proof}

\begin{remark}
\label{morphisms-remark-differentials-diagonal}
$X \to S$ をスキームの射とする。
補題 \ref{morphisms-lemma-differential-product} により、
$$
\Omega_{X \times_S X/S} =
\text{pr}_1^*\Omega_{X/S} \oplus \text{pr}_2^*\Omega_{X/S}
$$
である。一方、対角射 $\Delta : X \to X \times_S X$ は埋め込みであり、
局所的に左逆をもつ。従って、補題
\ref{morphisms-lemma-differentials-relative-immersion-section} により、標準短完全列
$$
0 \to \mathcal{C}_{X/X \times_S X} \to \Omega_{X/S} \oplus \Omega_{X/S}
\to \Omega_{X/S} \to 0
$$
を得る。右の矢は $(1, 1)$ であり、確かに分裂全射であることに注意する。
一方、補題 \ref{morphisms-lemma-differentials-diagonal} により、
同一視 $\Omega_{X/S} = \mathcal{C}_{X/X \times_S X}$ がある。
この同一視で $\text{d}_{X/S}(f) = s_2(f) - s_1(f)$ と選んだので、
左の矢は写像 $(-1, 1)$ となる\footnote{すなわち、
局所切断 $\text{d}_{X/S}(f) = 1 \otimes f - f \otimes 1$ は
$\Delta$ のイデアル層の切断であり、$\text{d}_{X \times_S X/X}$ によって局所切断
$1 \otimes 1 \otimes 1 \otimes f - 1 \otimes f \otimes 1 \otimes 1
-1 \otimes 1 \otimes f \otimes 1 + f \otimes 1 \otimes 1 \otimes 1 =
\text{pr}_2^*\text{d}_{X/S}(f) - \text{pr}_1^*\text{d}_{X/S}(f)$ に写る。}。
\end{remark}

\begin{lemma}
\label{morphisms-lemma-two-immersions}
スキームの可換図式
$$
\xymatrix{
Z \ar[r]_i \ar[rd]_j & X \ar[d] \\
& Y
}
$$
を考え、$i$ と $j$ は埋め込みとする。
このとき標準完全列
$$
\mathcal{C}_{Z/Y} \to
\mathcal{C}_{Z/X} \to
i^*\Omega_{X/Y} \to 0
$$
が存在する。最初の矢は補題 \ref{morphisms-lemma-conormal-functorial} から、
第二の矢は補題 \ref{morphisms-lemma-differentials-relative-immersion} から得られる。
\end{lemma}

\begin{proof}
この代数的な形は、代数、補題 \ref{algebra-lemma-application-NL} である。
\end{proof}

\section{有限階微分作用素}
\label{morphisms-section-differential-operators}

\noindent
可換代数の章の対応する節（代数、節
\ref{algebra-section-differential-operators}）および加群の層の章の
対応する節（加群、節 \ref{modules-section-differential-operators}）に
目を通すことを勧める。

\begin{lemma}
\label{morphisms-lemma-affine-case-differential-operators}
$R \to A$ を環準同型とする。$f : X \to S$ を対応するアフィンスキームの射とする。
$\mathcal{F}$ および $\mathcal{G}$ を $\mathcal{O}_X$-加群とする。
$\mathcal{F}$ が準連接ならば、写像
$$
\text{Diff}^k_{X/S}(\mathcal{F}, \mathcal{G}) \to
\text{Diff}^k_{A/R}(\Gamma(X, \mathcal{F}), \Gamma(X, \mathcal{G}))
$$
で微分作用素にその大域切断への作用を対応させるものは全単射である。
\end{lemma}

\begin{proof}
$\mathcal{F} = \widetilde{M}$ と、ある $A$-加群 $M$ を用いて書く。
$N = \Gamma(X, \mathcal{G})$ とおく。
$D : M \to N$ を $k$ 階微分作用素とする。
一意な微分作用素 $\mathcal{F} \to \mathcal{G}$ で、$k$ 階であり
大域切断上で $D$ を与えるものが存在することを示す。
$U = D(f) \subset X$ を標準アフィン開集合とする。
このとき $\mathcal{F}(U) = M_f$ は局所化である。
代数、補題 \ref{algebra-lemma-invert-system-differential-operators} により、
微分作用素 $D$ は一意な微分作用素
$$
D_f : \mathcal{F}(U) = \widetilde{M}(U) = M_f \to N_f = \widetilde{N}(U)
$$
に拡張される。一意性から、これらの写像 $D_f$ は貼り合わさって層の写像
$\widetilde{M} \to \widetilde{N}$ を、$X$ のすべての標準開集合からなる基底上で与える。
従って、一意な層の写像 $\widetilde{D} : \widetilde{M} \to \widetilde{N}$ で、
これらの写像と一致するものが得られる。層、節
\ref{sheaves-section-bases} の議論による。
$\widetilde{D}$ は標準開集合上で $k$ 階微分作用素として与えられるので、
$\widetilde{D}$ は $k$ 階微分作用素である（細部は省略する）。
最後に、標準的な $\mathcal{O}_X$-加群の写像
$c : \widetilde{N} \to \mathcal{G}$
（スキーム、補題 \ref{schemes-lemma-compare-constructions}）を後から合成して、
$c \circ \widetilde{D} : \mathcal{F} \to \mathcal{G}$ を得る。
これは $k$ 階微分作用素であることが、加群、補題
\ref{modules-lemma-composition-differential-operators} により分かる。
これで存在が示された。一意性の証明は省略する。
\end{proof}

\begin{lemma}
\label{morphisms-lemma-base-change-differential-operators}
$a : X \to S$ および $b : Y \to S$ をスキームの射とする。
$\mathcal{F}$ および $\mathcal{F}'$ を準連接 $\mathcal{O}_X$-加群とする。
$D : \mathcal{F} \to \mathcal{F}'$ を $k$ 階の $X/S$ 上の微分作用素とする。
$\mathcal{G}$ を準連接 $\mathcal{O}_Y$-加群とする。
このとき一意な微分作用素
$$
D' :
\text{pr}_1^*\mathcal{F} \otimes_{\mathcal{O}_{X \times_S Y}}
\text{pr}_2^*\mathcal{G}
\longrightarrow
\text{pr}_1^*\mathcal{F}' \otimes_{\mathcal{O}_{X \times_S Y}}
\text{pr}_2^*\mathcal{G}
$$
で、$k$ 階であり $X \times_S Y / Y$ 上で定義され、
$
D'(s \otimes t) = D(s) \otimes t
$
が局所切断 $s$（$\mathcal{F}$ のもの）および
$t$（$\mathcal{G}$ のもの）に対して成り立つものが存在する。
\end{lemma}

\begin{proof}
$X$, $Y$ および $S$ がアフィンの場合、補題
\ref{morphisms-lemma-affine-case-differential-operators} を通じて、
対応する代数の結果から従う。代数、補題
\ref{algebra-lemma-base-change-differential-operators} 参照。
一般の場合には、アフィン被覆（例えば、スキーム、補題
\ref{schemes-lemma-affine-covering-fibre-product} のようなもの）を用いて、
$D'$ を大域的に構成する。詳細は省略する。
\end{proof}

\begin{remark}
\label{morphisms-remark-base-change-differential-operators}
$a : X \to S$ および $b : Y \to S$ をスキームの射とする。
$p : X \times_S Y \to X$ および $q : X \times_S Y \to Y$ を射影とする。
この注意では、$\mathcal{O}_X$-加群 $\mathcal{F}$ および
$\mathcal{O}_Y$-加群 $\mathcal{G}$ が与えられたとき、
$$
\mathcal{F} \boxtimes \mathcal{G} =
p^*\mathcal{F} \otimes_{\mathcal{O}_{X \times_S Y}} q^*\mathcal{G}
$$
とおく。$\mathcal{A}_{X/S}$ を、対象が準連接 $\mathcal{O}_X$-加群であり、
射が $X/S$ 上の有限階微分作用素である加法圏とする。
$\mathcal{A}_{Y/S}$ および $\mathcal{A}_{X \times_S Y/S}$ についても同様とする。
補題 \ref{morphisms-lemma-base-change-differential-operators} の構成は、関手
$$
\boxtimes : 
\mathcal{A}_{X/S} \times \mathcal{A}_{Y/S} \longrightarrow
\mathcal{A}_{X \times_S Y/S},
\quad
(\mathcal{F}, \mathcal{G}) \longmapsto \mathcal{F} \boxtimes \mathcal{G}
$$
を定め、これは射に関して双線形である。
$X = \Spec(A)$, $Y = \Spec(B)$ および $S = \Spec(R)$ ならば、
準連接層と加群との同一視の下で、この関手は対象に対して
$(M, N) \mapsto M \otimes_R N$ で与えられ、
射 $(D, D') : (M, N) \to (M', N')$ を
$D \otimes D' : M \otimes_R N \to M' \otimes_R N'$ に写す。
\end{remark}

\section{滑らかな射}
\label{morphisms-section-smooth}

\noindent
$f : X \to Y$ を位相空間の連続写像とする。
次の条件を考える。すべての $x \in X$ に対して開近傍
$x \in U \subset X$ および $f(x) \in V \subset Y$ と整数 $d$ が存在し、
$f(U) \subset V$ であり、可換図式
$$
\xymatrix{
X \ar[d] &
U \ar[l] \ar[d] \ar[r]_-\pi &
V \times \mathbf{R}^d \ar[ld] \\
Y & V \ar[l]
}
$$
が得られ、ここで $\pi$ はある開部分集合への同相写像である。
スキームの滑らかな射は、スキームの圏におけるこれらの写像の類似物である。
補題 \ref{morphisms-lemma-smooth-locally-standard-smooth} および
補題 \ref{morphisms-lemma-smooth-etale-over-affine-space} 参照。

\medskip\noindent
（おそらく）予想に反して、滑らかな環準同型の概念は微分加群だけを用いて
定義されるわけではない。すなわち、$R \to A$ が{\it 滑らかな環準同型}であるとは、
$A$ が $R$ 上有限表示であり、$A$ の $R$ 上のナイーブ余接複体が、
次数 $0$ に置かれた射影加群と擬同型であることをいう。
代数、定義 \ref{algebra-definition-smooth} 参照。

\begin{definition}
\label{morphisms-definition-smooth}
$f : X \to S$ をスキームの射とする。
\begin{enumerate}
\item $f$ が{\it 点 $x \in X$ において滑らか}であるとは、
アフィン開近傍 $\Spec(A) = U \subset X$ が $x$ に対して存在し、
アフィン開集合 $\Spec(R) = V \subset S$ が存在して、
$f(U) \subset V$ であり、誘導される環準同型 $R \to A$ が滑らかとなることをいう。
\item $f$ が{\it 滑らか}であるとは、$X$ のすべての点で滑らかであることをいう。
\item アフィンスキームの射 $f : X \to S$ が{\it 標準滑らか}であるとは、
標準滑らかな環準同型 $R \to R[x_1, \ldots, x_n]/(f_1, \ldots, f_c)$
（代数、定義 \ref{algebra-definition-standard-smooth} 参照）が存在し、
$X \to S$ が
$$
\Spec(R[x_1, \ldots, x_n]/(f_1, \ldots, f_c)) \to \Spec(R).
$$
と同型であることをいう。
\end{enumerate}
\end{definition}

\noindent
この定義の好ましい点は、射が滑らかである点の集合が自動的に開集合となることである。

\medskip\noindent
定義には分離性や準コンパクト性の仮定がないことに注意する。
従って、滑らかであるかどうかは始域について局所的な問題である。
正確な結果は次のとおりである。

\begin{lemma}
\label{morphisms-lemma-smooth-characterize}
$f : X \to S$ をスキームの射とする。次は同値である。
\begin{enumerate}
\item 射 $f$ は滑らかである。
\item 任意のアフィン開集合 $U \subset X$, $V \subset S$ で
$f(U) \subset V$ を満たすものについて、環準同型
$\mathcal{O}_S(V) \to \mathcal{O}_X(U)$ は滑らかである。
\item 開被覆 $S = \bigcup_{j \in J} V_j$ および開被覆
$f^{-1}(V_j) = \bigcup_{i \in I_j} U_i$ が存在し、
各射 $U_i \to V_j$, $j\in J, i\in I_j$ は滑らかである。
\item アフィン開被覆 $S = \bigcup_{j \in J} V_j$ およびアフィン開被覆
$f^{-1}(V_j) = \bigcup_{i \in I_j} U_i$ が存在し、
環準同型 $\mathcal{O}_S(V_j) \to \mathcal{O}_X(U_i)$ は、
すべての $j\in J, i\in I_j$ に対して滑らかである。
\end{enumerate}
さらに、$f$ が滑らかならば、任意の開部分スキーム
$U \subset X$, $V \subset S$ で $f(U) \subset V$ を満たすものに対して、
制限 $f|_U : U \to V$ は滑らかである。
\end{lemma}

\begin{proof}
性質「$R \to A$ は滑らかである」が局所的であることを示せば、
補題 \ref{morphisms-lemma-locally-P} から従う。
定義 \ref{morphisms-definition-property-local} の条件 (a), (b), (c) を確認する。
代数、補題 \ref{algebra-lemma-base-change-smooth} により、
滑らかさは基底変換で保たれるので、(a) が成り立つ。
代数、補題 \ref{algebra-lemma-compose-smooth} により、
滑らかさは合成で保たれる。また、任意の環 $R$ に対して
環準同型 $R \to R_f$ は（標準）滑らかである。
従って (b) が成り立つ。最後に、性質 (c) は代数、補題
\ref{algebra-lemma-locally-smooth} によって成り立つ。
\end{proof}

\noindent
次の補題は、滑らかな射を、滑らかなファイバーをもつ平坦な有限表示射として特徴づける。
体上滑らかなスキームについては、多様体、節 \ref{varieties-section-smooth}
で詳しく論じることに注意する。

\begin{lemma}
\label{morphisms-lemma-smooth-flat-smooth-fibres}
$f : X \to S$ をスキームの射とする。
$f$ が平坦かつ局所有限表示であり、すべてのファイバー $X_s$ が滑らかならば、
$f$ は滑らかである。
\end{lemma}

\begin{proof}
代数、補題 \ref{algebra-lemma-flat-fibre-smooth} から従う。
\end{proof}

\begin{lemma}
\label{morphisms-lemma-composition-smooth}
二つの滑らかな射の合成は滑らかである。
\end{lemma}

\begin{proof}
補題 \ref{morphisms-lemma-smooth-characterize} の証明において、
滑らかさが環準同型の局所的性質であることを見た。
従って、この補題の最初の主張は補題 \ref{morphisms-lemma-composition-type-P} と、
滑らかさが合成で保たれる環準同型の性質であること
（代数、補題 \ref{algebra-lemma-compose-smooth} 参照）を合わせると従う。
\end{proof}

\begin{lemma}
\label{morphisms-lemma-base-change-smooth}
滑らかな射の基底変換は滑らかである。
\end{lemma}

\begin{proof}
補題 \ref{morphisms-lemma-smooth-characterize} の証明において、
滑らかさが環準同型の局所的性質であることを見た。
従って、この補題は補題 \ref{morphisms-lemma-composition-type-P} と、
滑らかさが基底変換で保たれる環準同型の性質であること
（代数、補題 \ref{algebra-lemma-base-change-smooth} 参照）を合わせると従う。
\end{proof}

\begin{lemma}
\label{morphisms-lemma-open-immersion-smooth}
任意の開埋め込みは滑らかである。
\end{lemma}

\begin{proof}
開埋め込みは局所同型なので成り立つ。
\end{proof}

\begin{lemma}
\label{morphisms-lemma-smooth-syntomic}
滑らかな射は syntomic である。
\end{lemma}

\begin{proof}
代数、補題 \ref{algebra-lemma-smooth-syntomic} を参照。
\end{proof}

\begin{lemma}
\label{morphisms-lemma-smooth-locally-finite-presentation}
滑らかな射は局所有限表示である。
\end{lemma}

\begin{proof}
滑らかな環準同型は定義により有限表示なので成り立つ。
\end{proof}

\begin{lemma}
\label{morphisms-lemma-smooth-flat}
滑らかな射は平坦である。
\end{lemma}

\begin{proof}
補題 \ref{morphisms-lemma-syntomic-flat} と \ref{morphisms-lemma-smooth-syntomic} を組み合わせる。
\end{proof}

\begin{lemma}
\label{morphisms-lemma-smooth-open}
滑らかな射は普遍開である。
\end{lemma}

\begin{proof}
補題 \ref{morphisms-lemma-smooth-flat},
\ref{morphisms-lemma-smooth-locally-finite-presentation} および
\ref{morphisms-lemma-fppf-open} を組み合わせる。
あるいは、補題 \ref{morphisms-lemma-smooth-syntomic},
\ref{morphisms-lemma-syntomic-open} を組み合わせてもよい。
\end{proof}

\noindent
次の補題によれば、任意の滑らかな射は局所的に標準滑らかである。
従って、標準滑らかな射を滑らかな射の{\it 局所モデル}として用いることができる。

\begin{lemma}
\label{morphisms-lemma-smooth-locally-standard-smooth}
\begin{slogan}
滑らかな射は局所的に標準滑らかである。
\end{slogan}
$f : X  \to S$ をスキームの射とする。
$x \in X$ を点とする。
$V \subset S$ を $f(x)$ のアフィン開近傍とする。
次は同値である。
\begin{enumerate}
\item 射 $f$ は $x$ において滑らかである。
\item アフィン開集合 $U \subset X$ が存在し、
$x \in U$ かつ $f(U) \subset V$ であり、
誘導射 $f|_U : U \to V$ は標準滑らかである。
\end{enumerate}
\end{lemma}

\begin{proof}
定義と、代数、補題 \ref{algebra-lemma-standard-smooth} および
\ref{algebra-lemma-smooth-syntomic} から従う。
\end{proof}

\begin{lemma}
\label{morphisms-lemma-smooth-omega-finite-locally-free}
$f : X \to S$ をスキームの射とする。
$f$ は滑らかであると仮定する。
このとき微分加群 $\Omega_{X/S}$（$X$ の $S$ 上のもの）は有限局所自由であり、
$$
\text{rank}_x(\Omega_{X/S}) = \dim_x(X_{f(x)})
$$
がすべての $x \in X$ に対して成り立つ。
\end{lemma}

\begin{proof}
主張は $X$ および $S$ について局所的である。
上の補題 \ref{morphisms-lemma-smooth-locally-standard-smooth} により、
$f$ はアフィンスキーム間の標準滑らかな射であると仮定してよい。
この場合、主張は代数、補題 \ref{algebra-lemma-standard-smooth}
（および相対大域完全交叉の定義、代数、定義
\ref{algebra-definition-relative-global-complete-intersection} 参照）から従う。
\end{proof}

\noindent
補題 \ref{morphisms-lemma-smooth-omega-finite-locally-free} により、次の定義は意味をもつ。

\begin{definition}
\label{morphisms-definition-smooth-relative-dimension}
$d \geq 0$ を整数とする。スキームの射 $f : X \to S$ が
{\it 相対次元 $d$ の滑らかな射}であるとは、
$f$ が滑らかであり、$\Omega_{X/S}$ が一定の階数 $d$ の有限局所自由加群であることをいう。
\end{definition}

\noindent
言い換えると、$f$ は滑らかであり、空でないファイバーは次元 $d$ の等次元スキームである。
以下の補題 \ref{morphisms-lemma-smooth-at-point} により、これは次の条件を課すこととも同値である。
(a) $f$ は局所有限表示である。
(b) $f$ は平坦である。
(c) 空でないファイバーはすべて次元 $d$ の等次元スキームである。
(d) $\Omega_{X/S}$ は階数 $d$ の有限局所自由加群である。
単に $f$ が平坦かつ有限表示であり、
$\Omega_{X/S}$ が階数 $d$ の有限局所自由加群であると仮定するだけでは不十分である。
反例は $\Spec(\mathbf{F}_p[t]) \to \Spec(\mathbf{F}_p[t^p])$ で与えられる。

\medskip\noindent
点における滑らかさの微分による判定法を述べる。
この結果には多くの変種があり、どれも何らかの場面で役立ちうる。
ここには必要に応じてそれらを加えていく。

\begin{lemma}
\label{morphisms-lemma-smooth-at-point}
$f : X \to S$ をスキームの射とする。
$x \in X$ とする。
$s = f(x)$ とおく。
$f$ は局所有限表示であると仮定する。
次は同値である。
\begin{enumerate}
\item 射 $f$ は $x$ において滑らかである。
\item 局所環準同型 $\mathcal{O}_{S, s} \to \mathcal{O}_{X, x}$ は平坦であり、
$X_s \to \Spec(\kappa(s))$ は $x$ において滑らかである。
\item 局所環準同型 $\mathcal{O}_{S, s} \to \mathcal{O}_{X, x}$ は平坦であり、
$\mathcal{O}_{X, x}$-加群 $\Omega_{X/S, x}$ は
高々 $\dim_x(X_{f(x)})$ 個の元で生成できる。
\item 局所環準同型 $\mathcal{O}_{S, s} \to \mathcal{O}_{X, x}$ は平坦であり、
$\kappa(x)$-ベクトル空間
$$
\Omega_{X_s/s, x} \otimes_{\mathcal{O}_{X_s, x}} \kappa(x) =
\Omega_{X/S, x} \otimes_{\mathcal{O}_{X, x}} \kappa(x)
$$
は高々 $\dim_x(X_{f(x)})$ 個の元で生成できる。
\item アフィン開集合 $U \subset X$ および $V \subset S$ が存在し、
$x \in U$, $f(U) \subset V$ であり、
誘導射 $f|_U : U \to V$ は標準滑らかである。
\item アフィン開集合 $\Spec(A) = U \subset X$ および
$\Spec(R) = V \subset S$ が存在し、$x \in U$ は
$\mathfrak q \subset A$ に対応し、$f(U) \subset V$ であり、表示
$$
A = R[x_1, \ldots, x_n]/(f_1, \ldots, f_c)
$$
で、
$$
g =
\det
\left(
\begin{matrix}
\partial f_1/\partial x_1 &
\partial f_2/\partial x_1 &
\ldots &
\partial f_c/\partial x_1 \\
\partial f_1/\partial x_2 &
\partial f_2/\partial x_2 &
\ldots &
\partial f_c/\partial x_2 \\
\ldots & \ldots & \ldots & \ldots \\
\partial f_1/\partial x_c &
\partial f_2/\partial x_c &
\ldots &
\partial f_c/\partial x_c
\end{matrix}
\right)
$$
が $A$ の元で $\mathfrak q$ に属さないものに写るようなものが存在する。
\end{enumerate}
\end{lemma}

\begin{proof}
$f$ が $x$ において滑らかならば、補題
\ref{morphisms-lemma-smooth-locally-standard-smooth} により (5) が成り立ち、
(6) は (5) を少し弱めたものである。
さらに、$f$ が滑らかならば環準同型
$\mathcal{O}_{S, s} \to \mathcal{O}_{X, x}$ は平坦であり
（補題 \ref{morphisms-lemma-smooth-flat} 参照）、
$\Omega_{X/S}$ は階数 $\dim_x(X_s)$ の有限局所自由加群である
（補題 \ref{morphisms-lemma-smooth-omega-finite-locally-free} 参照）。
従って (1) は (3) と (4) を導く。
また、補題 \ref{morphisms-lemma-base-change-smooth} により (1) は (2) を導く。

\medskip\noindent
補題 \ref{morphisms-lemma-base-change-differentials} により、
微分加群 $\Omega_{X_s/s}$（ファイバー $X_s$ の $\kappa(s)$ 上のもの）は、
微分加群 $\Omega_{X/S}$（$X$ の $S$ 上のもの）の引き戻しである。
これにより、補題の (4) に表示した等式が得られる。
補題 \ref{morphisms-lemma-finite-type-differentials} により、これらの加群は有限型である。
従って、加群 $\Omega_{X/S, x}$ と $\Omega_{X_s/s, x}$ の生成元の最小個数は同じであり、
この $\kappa(x)$-ベクトル空間の次元に等しい。これは中山の補題
（代数、補題 \ref{algebra-lemma-NAK}）による。
特に (3) と (4) は同値である。

\medskip\noindent
代数、補題 \ref{algebra-lemma-flat-fibre-smooth} により (2) は (1) を導く。
代数、補題 \ref{algebra-lemma-characterize-smooth-over-field} により
(3) と (4) は (2) を導く。最後に、(6) は (5) を導く。
例えば代数、例 \ref{algebra-example-make-standard-smooth} を参照。
また、(5) は (1) を代数、補題 \ref{algebra-lemma-standard-smooth} によって導く。
\end{proof}

\begin{lemma}
\label{morphisms-lemma-set-points-where-fibres-smooth}
スキームのカルテジアン図式
$$
\xymatrix{
X' \ar[r]_{g'} \ar[d]_{f'} & X \ar[d]^f \\
S' \ar[r]^g & S
}
$$
を考える。$W \subset X$ および $W' \subset X'$ を、それぞれ
$f$ および $f'$ が滑らかである点からなる開部分スキームとする。
次のいずれかが成り立つならば、$W' = (g')^{-1}(W)$ である。
\begin{enumerate}
\item $f$ は平坦かつ局所有限表示である。
\item $f$ は局所有限表示であり、$g$ は平坦である。
\end{enumerate}
\end{lemma}

\begin{proof}
まず $f$ は局所有限型であると仮定する。集合
$$
T = \{x \in X \mid X_{f(x)}\text{ is smooth over }\kappa(f(x))\text{ at }x\}
$$
および対応する集合 $T' \subset X'$（$f'$ に対するもの）を考える。
$T' = (g')^{-1}(T)$ と主張する。実際、$s' \in S'$ を点とし、
$s = g(s')$ とおく。このとき
$$
X'_{s'} =
\Spec(\kappa(s')) \times_{\Spec(\kappa(s))} X_s
$$
である。言い換えると、基底変換のファイバーはファイバーの基底変換である。
従って主張は代数、補題 \ref{algebra-lemma-smooth-field-change-local} と同値である。

\medskip\noindent
これにより (1) の場合が従う。実際、(1) の場合には $T$ は
$f$ が滑らかである点の（開）集合であることが補題
\ref{morphisms-lemma-smooth-at-point} により分かる。

\medskip\noindent
(2) の場合、$x' \in W'$ とする。
このとき $g'$ は $x'$ で平坦であり（補題 \ref{morphisms-lemma-base-change-module-flat}）、
$g \circ f'$ は $x'$ で平坦である（補題 \ref{morphisms-lemma-composition-module-flat}）。
従って $f$ は $x = g'(x')$ で平坦であることが補題 \ref{morphisms-lemma-flat-permanence} により従う。
一方、$x' \in T'$（補題 \ref{morphisms-lemma-base-change-smooth}）なので、$x \in T$ である。
よって $f$ は $x$ で滑らかであることが補題 \ref{morphisms-lemma-smooth-at-point} により分かる。
\end{proof}

\noindent
次の補題では、滑らかな環準同型のナイーブ余接複体の $H^{-1}$ の消滅を実際に用いる。

\begin{lemma}
\label{morphisms-lemma-triangle-differentials-smooth}
$f : X \to Y$, $g : Y \to S$ をスキームの射とする。
$f$ は滑らかであると仮定する。このとき
$$
0 \to f^*\Omega_{Y/S} \to \Omega_{X/S} \to \Omega_{X/Y} \to 0
$$
（補題 \ref{morphisms-lemma-triangle-differentials} 参照）は短完全列である。
\end{lemma}

\begin{proof}
この補題の代数的な形は次のとおりである。
環準同型 $A \to B \to C$ が与えられ、$B \to C$ が滑らかならば、列
$$
0 \to C \otimes_B \Omega_{B/A} \to \Omega_{C/A} \to \Omega_{C/B} \to 0
$$
（代数、補題 \ref{algebra-lemma-exact-sequence-differentials} の列）は完全である。
これは代数、補題 \ref{algebra-lemma-triangle-differentials-smooth} である。
\end{proof}

\begin{lemma}
\label{morphisms-lemma-differentials-relative-immersion-smooth}
$i : Z \to X$ を $S$ 上のスキームの埋め込みとする。
$Z$ は $S$ 上滑らかであると仮定する。このとき標準完全列
$$
0 \to \mathcal{C}_{Z/X} \to i^*\Omega_{X/S} \to \Omega_{Z/S} \to 0
$$
（補題 \ref{morphisms-lemma-differentials-relative-immersion} の列）は短完全列である。
\end{lemma}

\begin{proof}
この補題の代数的な形は次のとおりである。
環準同型 $A \to B \to C$ が与えられ、$A \to C$ は滑らかであり、
$B \to C$ は全射でその核が $J$ ならば、列
$$
0 \to J/J^2 \to C \otimes_B \Omega_{B/A} \to \Omega_{C/A} \to 0
$$
（代数、補題 \ref{algebra-lemma-differential-seq} の列）は完全である。
これは代数、補題 \ref{algebra-lemma-differential-seq-smooth} である。
\end{proof}

\begin{lemma}
\label{morphisms-lemma-two-immersions-smooth}
スキームの可換図式
$$
\xymatrix{
Z \ar[r]_i \ar[rd]_j & X \ar[d] \\
& Y
}
$$
を考え、$i$ および $j$ は埋め込みであり、$X \to Y$ は滑らかとする。
このとき標準完全列
$$
0 \to  \mathcal{C}_{Z/Y} \to \mathcal{C}_{Z/X} \to i^*\Omega_{X/Y} \to 0
$$
（補題 \ref{morphisms-lemma-two-immersions} の列）は完全である。
\end{lemma}

\begin{proof}
この補題の代数的な形は次のとおりである。
環準同型 $A \to B \to C$ が与えられ、$A \to C$ は全射であり、
$A \to B$ は滑らかならば、列
$$
0 \to I/I^2 \to J/J^2 \to C \otimes_B \Omega_{B/A} \to 0
$$
（代数、補題 \ref{algebra-lemma-application-NL} の列）は完全である。
これは代数、補題 \ref{algebra-lemma-application-NL-smooth} である。
\end{proof}

\begin{lemma}
\label{morphisms-lemma-smooth-permanence}
スキームの射の可換図式
$$
\xymatrix{
X \ar[rr]_f \ar[rd]_p & &
Y \ar[dl]^q \\
& S
}
$$
を考える。次を仮定する。
\begin{enumerate}
\item $f$ は全射かつ滑らかである。
\item $p$ は滑らかである。
\item $q$ は局所有限表示である\footnote{実際には、これは (1) と (2) から従う。
降下、補題 \ref{descent-lemma-flat-finitely-presented-permanence} 参照。
さらに、$f$ は全射、平坦かつ局所有限表示であると仮定すれば十分である。
降下、補題 \ref{descent-lemma-smooth-permanence} 参照。}。
\end{enumerate}
このとき $q$ は滑らかである。
\end{lemma}

\begin{proof}
補題 \ref{morphisms-lemma-flat-permanence} により $q$ は平坦である。
点 $y \in Y$ を取る。点 $x \in X$ を $y$ に写るように取る。
$f$ の相対次元は $a$（点 $x$ において）であり、
$p$ の相対次元は $b$（点 $x$ において）であるとする。
補題 \ref{morphisms-lemma-smooth-omega-finite-locally-free} により、
これは $\Omega_{X/S, x}$ が階数 $b$ の自由加群であり、
$\Omega_{X/Y, x}$ が階数 $a$ の自由加群であることを意味する。
補題 \ref{morphisms-lemma-triangle-differentials-smooth} の短完全列により、
これは $(f^*\Omega_{Y/S})_x$ が階数 $b - a$ の自由加群であることを意味する。
中山の補題から、$\Omega_{Y/S, y}$ は $b - a$ 個の元で生成できる。
また、補題 \ref{morphisms-lemma-dimension-fibre-at-a-point-additive} により
$\dim_y(Y_s) = b - a$ である。
従って、$Y \to S$ は $y$ において滑らかであることが、
補題 \ref{morphisms-lemma-smooth-at-point} の (2) により結論される。
\end{proof}

\noindent
次の補題の状況では、$\sigma$ の像は、$X$ 上局所的に正則列によって定義される。
因子、補題 \ref{divisors-lemma-section-smooth-regular-immersion} 参照。

\begin{lemma}
\label{morphisms-lemma-section-smooth-morphism}
$f : X \to S$ をスキームの射とする。
$\sigma : S \to X$ を $f$ の切断とする。
$s \in S$ を点であって $f$ が $x = \sigma(s)$ において滑らかとなるものとする。
このとき、アフィン開近傍 $\Spec(A) = U \subset S$ が $s$ に対して、
および $\Spec(B) = V \subset X$ が $x$ に対して存在し、次が成り立つ。
\begin{enumerate}
\item $f(V) \subset U$ かつ $\sigma(U) \subset V$。
\item $I = \Ker(\sigma^\# : B \to A)$ とおくと、加群 $I/I^2$ は自由 $A$-加群である。
\item $B^\wedge \cong A[[x_1, \ldots, x_d]]$ が $A$-代数の同型として成り立つ。
ここで $B^\wedge$ は $B$ の $I$ に関する完備化を表す。
\end{enumerate}
\end{lemma}

\begin{proof}
アフィン開集合 $U \subset S$ を $s$ を含むように取る。
アフィン開集合 $V \subset f^{-1}(U)$ を $x$ を含むように取る。
アフィン開集合 $U' \subset \sigma^{-1}(V)$ を $s$ を含むように取る。
$V' = f^{-1}(U') \cap V$ は、ファイバー積 $V' = U' \times_U V$ に等しいので
アフィンである。すると $U'$ および $V'$ は (1) を満たす。
$U' = \Spec(A')$ および $V' = \Spec(B')$ と書く。
代数、補題 \ref{algebra-lemma-section-smooth} により、
加群 $I'/(I')^2$ は $A'$-加群として有限局所自由である。
従って $U'$ をより小さいアフィン開集合 $U'' \subset U'$ で、
また $V'$ を $V'' = V' \cap f^{-1}(U'')$ で置き換えることで、
$I''/(I'')^2$ が自由となる状況を得る。すなわち (2) が成り立つ。
この場合、(3) も代数、補題 \ref{algebra-lemma-section-smooth} によって成り立つ。
\end{proof}

\noindent
スキーム $X$ の点 $x$ における次元（性質、定義
\ref{properties-definition-dimension}）は、単に $X$ を位相空間とみなしたときの
$x$ における次元である。位相、定義 \ref{topology-definition-Krull} 参照。
これは一般には局所環 $\mathcal{O}_{X,x}$ の次元ではない。

\begin{lemma}
\label{morphisms-lemma-smoothness-dimension}
$f : X \to Y$ を局所 Noether スキームの滑らかな射とする。
任意の点 $x$ が $X$ に属し、その像 $y$ が $Y$ に属するとき、
$$
\dim_x(X) = \dim_y(Y) + \dim_x(X_y),
$$
が成り立つ。ここで $X_y$ は $y$ 上のファイバーを表す。
\end{lemma}

\begin{proof}
$X$ を $x$ の開近傍で置き換えると、自然数 $d$ が存在して、
$X \to Y$ のすべてのファイバーは各点において次元 $d$ をもつ。
補題 \ref{morphisms-lemma-smooth-omega-finite-locally-free} 参照。
このとき $f$ は平坦であり（補題 \ref{morphisms-lemma-smooth-flat}）、
局所有限型であり（補題 \ref{morphisms-lemma-smooth-locally-finite-presentation}）、
相対次元は $d$ である。
従って主張は補題 \ref{morphisms-lemma-rel-dimension-dimension} から従う。
\end{proof}

\section{不分岐射}
\label{morphisms-section-unramified}

\noindent
（おそらく）より興味深いエタール射のクラスに進む前に、
不分岐射について簡単に論じる。環準同型 $R \to A$ が
有限型であり、かつ $\Omega_{A/R} = 0$ であるとき、これを
{\it 不分岐}であるという（これは \cite{Henselian} の定義である）。
環準同型 $R \to A$ が有限表示であり、かつ $\Omega_{A/R} = 0$
であるとき、これを {\it G-不分岐}であるという（これは
\cite{EGA} の定義である）。代数、定義
\ref{algebra-definition-unramified} 参照。

\begin{definition}
\label{morphisms-definition-unramified}
$f : X \to S$ をスキームの射とする。
\begin{enumerate}
\item $f$ が {\it $x \in X$ において不分岐}であるとは、
$\Spec(A) = U \subset X$ で表される $x$ のアフィン開近傍とアフィン開集合
$\Spec(R) = V \subset S$ が存在して $f(U) \subset V$ を満たし、
誘導される環準同型 $R \to A$ が不分岐であることをいう。
\item $f$ が {\it $x \in X$ において G-不分岐}であるとは、
$\Spec(A) = U \subset X$ で表される $x$ のアフィン開近傍とアフィン開集合
$\Spec(R) = V \subset S$ が存在して $f(U) \subset V$ を満たし、
誘導される環準同型 $R \to A$ が G-不分岐であることをいう。
\item $f$ が $X$ のすべての点において不分岐であるとき、
これを {\it 不分岐}であるという。
\item $f$ が $X$ のすべての点において G-不分岐であるとき、
これを {\it G-不分岐}であるという。
\end{enumerate}
\end{definition}

\noindent
G-不分岐射は不分岐であることに注意せよ。従って、不分岐射に
関する任意の結果から、G-不分岐射に関する対応する結果が従う。
さらに $S$ が局所 Noether ならば、G-不分岐射と不分岐射の間に
相違はない。補題 \ref{morphisms-lemma-noetherian-unramified} 参照。
この定義の好ましい特徴は、射が不分岐（resp.\ G-不分岐）である
点の集合が自動的に開になることである。

\begin{lemma}
\label{morphisms-lemma-unramified-omega-zero}
$f : X \to S$ をスキームの射とする。このとき
\begin{enumerate}
\item $f$ が不分岐であることと、$f$ が局所有限型であり、
$\Omega_{X/S} = 0$ であることは同値であり、
\item $f$ が G-不分岐であることと、$f$ が局所有限表示であり、
$\Omega_{X/S} = 0$ であることは同値である。
\end{enumerate}
\end{lemma}

\begin{proof}
定義により、環準同型 $R \to A$ が不分岐（resp.\ G-不分岐）である
ことと、それが有限型（resp.\ 有限表示）であり、かつ
$\Omega_{A/R} = 0$ であることは同値である。従って補題は定義と
補題 \ref{morphisms-lemma-differentials-affine} から直ちに従う。
\end{proof}

\noindent
定義には分離性または準コンパクト性に関する仮定がないことに
注意せよ。従って、不分岐であるという問題は始域上局所的な性質を
もつ。正確な結果は次のとおりである。

\begin{lemma}
\label{morphisms-lemma-unramified-characterize}
$f : X \to S$ をスキームの射とする。次は同値である。
\begin{enumerate}
\item 射 $f$ は不分岐（resp.\ G-不分岐）である。
\item すべてのアフィン開集合 $U \subset X$、$V \subset S$ で
$f(U) \subset V$ を満たすものに対し、環準同型
$\mathcal{O}_S(V) \to \mathcal{O}_X(U)$ は不分岐
（resp.\ G-不分岐）である。
\item 開被覆 $S = \bigcup_{j \in J} V_j$ と開被覆
$f^{-1}(V_j) = \bigcup_{i \in I_j} U_i$ が存在し、各射
$U_i \to V_j$、$j\in J, i\in I_j$ は不分岐
（resp.\ G-不分岐）である。
\item アフィン開被覆 $S = \bigcup_{j \in J} V_j$ とアフィン開被覆
$f^{-1}(V_j) = \bigcup_{i \in I_j} U_i$ が存在し、環準同型
$\mathcal{O}_S(V_j) \to \mathcal{O}_X(U_i)$ は、すべての
$j\in J, i\in I_j$ に対して不分岐
（resp.\ G-不分岐）である。
\end{enumerate}
さらに $f$ が不分岐（resp.\ G-不分岐）ならば、
任意の開部分スキーム $U \subset X$、$V \subset S$ で
$f(U) \subset V$ を満たすものに対して、その制限
$f|_U : U \to V$ も不分岐
（resp.\ G-不分岐）である。
\end{lemma}

\begin{proof}
「$R \to A$ は不分岐である」という性質が局所的であることを
示せば、補題 \ref{morphisms-lemma-locally-P} から従う。
定義 \ref{morphisms-definition-property-local} の条件 (a)、(b)、(c) を
確認する。これらの性質は代数、補題
\ref{algebra-lemma-unramified} で証明されている。
\end{proof}

\begin{lemma}
\label{morphisms-lemma-composition-unramified}
二つの不分岐射の合成は不分岐である。
G-不分岐射についても同じことが成り立つ。
\end{lemma}

\begin{proof}
補題 \ref{morphisms-lemma-unramified-characterize} の証明は、不分岐
（resp.\ G-不分岐）であることが環準同型の局所的性質であることを
示している。従って補題の第一の主張は、補題
\ref{morphisms-lemma-composition-type-P} と、不分岐（resp.\ G-不分岐）である
ことが合成に関して安定な環準同型の性質であるという事実を
組み合わせると従う。代数、補題
\ref{algebra-lemma-unramified} 参照。
\end{proof}

\begin{lemma}
\label{morphisms-lemma-base-change-unramified}
不分岐射の基底変換は不分岐である。
G-不分岐射についても同じことが成り立つ。
\end{lemma}

\begin{proof}
補題 \ref{morphisms-lemma-unramified-characterize} の証明は、不分岐
（resp.\ G-不分岐）であることが環準同型の局所的性質であることを
示している。従って補題は、補題 \ref{morphisms-lemma-base-change-type-P} と、
不分岐（resp.\ G-不分岐）であることが基底変換に関して安定な
環準同型の性質であるという事実を組み合わせると従う。
代数、補題 \ref{algebra-lemma-unramified} 参照。
\end{proof}

\begin{lemma}
\label{morphisms-lemma-noetherian-unramified}
$f : X \to S$ をスキームの射とし、$S$ は局所 Noether であると
仮定する。このとき $f$ が不分岐であることと、$f$ が G-不分岐で
あることは同値である。
\end{lemma}

\begin{proof}
定義と補題 \ref{morphisms-lemma-noetherian-finite-type-finite-presentation} から従う。
\end{proof}


\begin{lemma}
\label{morphisms-lemma-open-immersion-unramified}
任意の開埋め込みは G-不分岐である。
\end{lemma}

\begin{proof}
開埋め込みは局所同型なので、これは成り立つ。
\end{proof}

\begin{lemma}
\label{morphisms-lemma-closed-immersion-unramified}
閉埋め込み $i : Z \to X$ は不分岐である。
それが G-不分岐であることと、付随する準連接イデアル層
$\mathcal{I} = \Ker(\mathcal{O}_X \to i_*\mathcal{O}_Z)$ が
（$\mathcal{O}_X$-加群として）有限型であることは同値である。
\end{lemma}

\begin{proof}
補題 \ref{morphisms-lemma-closed-immersion-finite-presentation} と
代数、補題 \ref{algebra-lemma-unramified} から従う。
\end{proof}

\begin{lemma}
\label{morphisms-lemma-unramified-locally-finite-type}
不分岐射は局所有限型である。
G-不分岐射は局所有限表示である。
\end{lemma}

\begin{proof}
不分岐な環準同型は定義により有限型である。
G-不分岐な環準同型は定義により有限表示である。
\end{proof}

\begin{lemma}
\label{morphisms-lemma-unramified-quasi-finite}
$f : X \to S$ をスキームの射とする。
$f$ が $x$ において不分岐ならば、$f$ は $x$ において準有限である。
特に、不分岐射は局所準有限である。
\end{lemma}

\begin{proof}
代数、補題 \ref{algebra-lemma-unramified-quasi-finite} 参照。
\end{proof}

\begin{lemma}
\label{morphisms-lemma-unramified-over-field}
不分岐射のファイバーについて、次が成り立つ。
\begin{enumerate}
\item $X$ を体 $k$ 上のスキームとする。構造射
$X \to \Spec(k)$ が不分岐であることと、$X$ が $k$ の有限分離拡大の
スペクトルの非交和であることは同値である。
\item $f : X \to S$ が不分岐射ならば、任意の $s \in S$ に対して
ファイバー $X_s$ は $\kappa(s)$ の有限分離拡大のスペクトルの
非交和である。
\end{enumerate}
\end{lemma}

\begin{proof}
(2) は (1) と補題 \ref{morphisms-lemma-base-change-unramified} から従う。
(1) を証明しよう。まず代数、補題
\ref{algebra-lemma-characterize-unramified} を用いる。この補題により、
$X$ が $k$ の有限分離拡大のスペクトルの非交和ならば
$X \to \Spec(k)$ は不分岐である。逆に、$X \to \Spec(k)$ が
不分岐であると仮定する。代数、補題
\ref{algebra-lemma-unramified-at-prime} により、任意の $x \in X$ に
対して剰余体拡大 $\kappa(x)/k$ は有限分離拡大である。
$X \to \Spec(k)$ は局所準有限なので（補題
\ref{morphisms-lemma-unramified-quasi-finite}）、補題
\ref{morphisms-lemma-quasi-finite-at-point-characterize} により、$X$ のすべての
点は孤立した閉点である。従って $X$ は離散空間であり、特にその
局所環のスペクトルの非交和である。代数、補題
\ref{algebra-lemma-unramified-at-prime} を再び用いると、これらの
局所環は体であり、主張を得る。
\end{proof}

\noindent
次の補題は、不分岐射を、不分岐ファイバーをもつ局所有限型射として
特徴づける。

\begin{lemma}
\label{morphisms-lemma-unramified-etale-fibres}
$f : X \to S$ をスキームの射とする。
\begin{enumerate}
\item $f$ が不分岐ならば、任意の $x \in X$ に対して体拡大
$\kappa(x)/\kappa(f(x))$ は有限分離拡大である。
\item $f$ が局所有限型であり、任意の $s \in S$ に対してファイバー
$X_s$ が $\kappa(s)$ の有限分離拡大のスペクトルの非交和ならば、
$f$ は不分岐である。
\item $f$ が局所有限表示であり、任意の $s \in S$ に対してファイバー
$X_s$ が $\kappa(s)$ の有限分離拡大のスペクトルの非交和ならば、
$f$ は G-不分岐である。
\end{enumerate}
\end{lemma}

\begin{proof}
代数、補題 \ref{algebra-lemma-unramified-at-prime} および
\ref{algebra-lemma-characterize-unramified} から従う。
\end{proof}

\noindent
対角射を用いた不分岐射の特徴づけは次のとおりである。

\begin{lemma}
\label{morphisms-lemma-diagonal-unramified-morphism}
$f : X \to S$ を射とする。
\begin{enumerate}
\item $f$ が不分岐ならば、対角射
$\Delta : X \to X \times_S X$ は開埋め込みである。
\item $f$ が局所有限型であり、$\Delta$ が開埋め込みならば、
$f$ は不分岐である。
\item $f$ が局所有限表示であり、$\Delta$ が開埋め込みならば、
$f$ は G-不分岐である。
\end{enumerate}
\end{lemma}

\begin{proof}
第一の主張は代数、補題
\ref{algebra-lemma-diagonal-unramified} から従う。
第二の主張は、$\Omega_{X/S}$ が対角射の余法層であり
（補題 \ref{morphisms-lemma-differentials-diagonal}）、従って $\Delta$ が
開埋め込みならば明らかに零であることから従う。
\end{proof}

\begin{lemma}
\label{morphisms-lemma-unramified-at-point}
$f : X \to S$ をスキームの射とし、$x \in X$ とする。
$s = f(x)$ とおく。$f$ は局所有限型（resp.\ 局所有限表示）であると
仮定する。次は同値である。
\begin{enumerate}
\item 射 $f$ は $x$ において不分岐（resp.\ G-不分岐）である。
\item ファイバー $X_s$ は $\kappa(s)$ 上で $x$ において不分岐である。
\item $\mathcal{O}_{X, x}$-加群 $\Omega_{X/S, x}$ は零である。
\item $\mathcal{O}_{X_s, x}$-加群 $\Omega_{X_s/s, x}$ は零である。
\item $\kappa(x)$-ベクトル空間
$$
\Omega_{X_s/s, x} \otimes_{\mathcal{O}_{X_s, x}} \kappa(x) =
\Omega_{X/S, x} \otimes_{\mathcal{O}_{X, x}} \kappa(x)
$$
は零である。
\item $\mathfrak m_s\mathcal{O}_{X, x} = \mathfrak m_x$ が成り立ち、
体拡大 $\kappa(x)/\kappa(s)$ は有限分離拡大である。
\end{enumerate}
\end{lemma}

\begin{proof}
$f$ が $x$ において不分岐ならば、定義と節
\ref{morphisms-section-sheaf-differentials} における微分加群の結果により、
$\Omega_{X/S} = 0$ であることが $x$ のある近傍でわかる。従って (1) から
(3) と、(5) の右辺のベクトル空間の消滅が従う。また補題
\ref{morphisms-lemma-base-change-differentials} により、微分加群
$\Omega_{X_s/s}$（ファイバー $X_s$ の $\kappa(s)$ 上のもの）は、
微分加群 $\Omega_{X/S}$（$X$ の $S$ 上のもの）の引き戻しなので、
(2) も従う。
微分加群に関するこの事実から、項 (4) におけるベクトル空間の
表示された等式も従う。補題 \ref{morphisms-lemma-finite-type-differentials} により、
加群 $\Omega_{X/S, x}$ と $\Omega_{X_s/s, x}$ は有限型である。
従って、Nakayama の補題（代数、補題 \ref{algebra-lemma-NAK}）により、
加群 $\Omega_{X/S, x}$ と $\Omega_{X_s/s, x}$ が零であることと、
(4) の対応する $\kappa(x)$-ベクトル空間が零であることは同値である。
特に、これは (3)、(4)、(5) が同値であることを示す。
$\Omega_{X/S}$ の台は $X$ で閉である。加群、補題
\ref{modules-lemma-support-finite-type-closed} 参照。仮定 (3) により、
$x$ はその台に属さない。従って $\Omega_{X/S}$ は $x$ の近傍で
零であり、(1) が従う。これまでの (1) と (3) の同値性を
$X_s \to s$ に適用すると、(2) と (4) の同値性が従う。
以上により (1) -- (5) が同値であることを示した。

\medskip\noindent
別の方法として、代数、補題 \ref{algebra-lemma-unramified} を用いれば、
(1) -- (5) の同値性をより直接に示せる。

\medskip\noindent
(1) と (6) の同値性は補題
\ref{morphisms-lemma-unramified-etale-fibres} から従う。また、代数、補題
\ref{algebra-lemma-unramified-at-prime} および
\ref{algebra-lemma-characterize-unramified} からも、より直接に従う。
\end{proof}

\begin{lemma}
\label{morphisms-lemma-set-points-where-fibres-unramified}
$f : X \to S$ をスキームの射とし、$f$ は局所有限型であると仮定する。
次の開集合の構成は任意の基底変換と可換である。
\begin{align*}
T
& =
\{x \in X \mid X_{f(x)}\text{ is unramified over }\kappa(f(x))\text{ at }x\} \\
& =
\{x \in X \mid X\text{ is unramified over }S\text{ at }x\}
\end{align*}
すなわち、任意の射 $g : S' \to S$ に対し、射
$f' : X' \to S'$（$f$ の基底変換）と射影 $g' : X' \to X$ を考えると、
対応する集合 $T'$（射 $f'$ に対するもの）について等式
$T' = (g')^{-1}(T)$ が成り立つ。
$f$ が局所有限表示であると仮定すれば、$f$ が G-不分岐である点の
開集合についても同じことが成り立つ。
\end{lemma}

\begin{proof}
$s' \in S'$ を点とし、$s = g(s')$ とおく。このとき
$$
X'_{s'} =
\Spec(\kappa(s')) \times_{\Spec(\kappa(s))} X_s
$$
が成り立つ。言い換えれば、基底変換のファイバーはファイバーの
基底変換である。特に
$$
\Omega_{X_s/s, x} \otimes_{\mathcal{O}_{X_s, x}} \kappa(x')
=
\Omega_{X'_{s'}/s', x'} \otimes_{\mathcal{O}_{X'_{s'}, x'}} \kappa(x')
$$
が成り立つ。補題 \ref{morphisms-lemma-base-change-differentials} 参照。
従って補題 \ref{morphisms-lemma-unramified-at-point} により、$x' \in T'$ である
ことと $x \in T$ であることは同値である。第二の主張は第一の主張
から従う。実際、この場合、補題 \ref{morphisms-lemma-unramified-at-point} により、
$T$ は $f$ が G-不分岐である点の（開）集合だからである。
\end{proof}

\begin{lemma}
\label{morphisms-lemma-unramified-permanence}
$f : X \to Y$ を $S$ 上のスキームの射とする。
\begin{enumerate}
\item $X$ が $S$ 上不分岐ならば、$f$ は不分岐である。
\item $X$ が $S$ 上 G-不分岐であり、$Y$ が $S$ 上局所有限型ならば、
$f$ は G-不分岐である。
\end{enumerate}
\end{lemma}

\begin{proof}
$X$ が $S$ 上不分岐であると仮定する。補題
\ref{morphisms-lemma-permanence-finite-type} により、$f$ は局所有限型である。
仮定により $\Omega_{X/S} = 0$ である。従って補題
\ref{morphisms-lemma-triangle-differentials} により $\Omega_{X/Y} = 0$ である。
ゆえに $f$ は不分岐である。$X$ が $S$ 上 G-不分岐であり、$Y$ が
$S$ 上局所有限型ならば、補題
\ref{morphisms-lemma-finite-presentation-permanence} により、$f$ は局所有限表示で
ある。従って $f$ は G-不分岐である。
\end{proof}

\begin{lemma}
\label{morphisms-lemma-value-at-one-point}
$S$ をスキームとし、$X$、$Y$ を $S$ 上のスキームとする。
$f, g : X \to Y$ を $S$ 上の射とし、$x \in X$ とする。
次を仮定する。
\begin{enumerate}
\item 構造射 $Y \to S$ は不分岐である。
\item $f(x) = g(x)$ が $Y$ において成り立つ。$y = f(x) = g(x)$ とおく。
\item 誘導される写像 $f^\sharp, g^\sharp : \kappa(y) \to \kappa(x)$ は
等しい。
\end{enumerate}
このとき、$x$ の $X$ におけるある開近傍上で $f$ と $g$ は等しい。
\end{lemma}

\begin{proof}
射 $(f, g) : X \to Y \times_S Y$ を考える。仮定 (1) と補題
\ref{morphisms-lemma-diagonal-unramified-morphism} により、
$\Delta_{Y/S}(Y)$ の逆像は $X$ で開である。また仮定 (2) と (3) により、
$x$ はこの開部分集合に属する。
\end{proof}

\section{エタール射}
\label{morphisms-section-etale}

\noindent
スキームの Zariski 位相は非常に粗い位相である。
これは特に $\mathbf{C}$ 上の多様体を見ると明らかである。
エタール射を位相における局所同型の類似物とみなすと、はるかに
細かい位相が得られることがわかる。$\mathbf{C}$ 上の多様体では、
この位相を用いて有限係数のコホモロジーを計算すると「正しい」
Betti 数が得られる。別の観察として、$f : X \to Y$ が
$\mathbf{C}$ 上の多様体のエタール射であり、$x$ が $X$ の閉点ならば、
$f$ は完備局所環の同型
$\mathcal{O}^{\wedge}_{Y, f(x)} \to \mathcal{O}^{\wedge}_{X, x}$
を誘導する。

\medskip\noindent
本節では、これらの事柄の研究を始める。実際、より興味深い結果は
別の場所で論じるので、ここでは意図的に議論を最小限にとどめる。
環準同型 $R \to A$ が滑らかであり、かつ $\Omega_{A/R} = 0$ である
とき、それを {\it エタール}であるという。代数、定義
\ref{algebra-definition-etale} 参照。

\begin{definition}
\label{morphisms-definition-etale}
$f : X \to S$ をスキームの射とする。
\begin{enumerate}
\item $f$ が {\it $x \in X$ においてエタール}であるとは、
$\Spec(A) = U \subset X$ で表される $x$ のアフィン開近傍とアフィン開集合
$\Spec(R) = V \subset S$ が存在して $f(U) \subset V$ を満たし、
誘導される環準同型 $R \to A$ がエタールであることをいう。
\item $f$ が $X$ のすべての点においてエタールであるとき、
これを {\it エタール}であるという。
\item アフィンスキームの射 $f : X \to S$ が {\it 標準エタール}で
あるとは、$X \to S$ が
$$
\Spec(R[x]_h/(g)) \to \Spec(R)
$$
と同型であることをいう。ここで $R \to R[x]_h/(g)$ は標準エタールな
環準同型である。代数、定義 \ref{algebra-definition-standard-etale} 参照。
すなわち、$g$ はモニックであり、$g'$ は $R[x]_h/(g)$ で可逆である。
\end{enumerate}
\end{definition}

\noindent
射がエタールであることと、それが相対次元 $0$ の滑らかな射である
ことは同値である（定義 \ref{morphisms-definition-smooth-relative-dimension} 参照）。
この定義の好ましい特徴は、射がエタールである点の集合が自動的に
開になることである。

\medskip\noindent
定義には分離性または準コンパクト性に関する仮定がないことに
注意せよ。従って、エタールであるという問題は始域上局所的な性質を
もつ。正確な結果は次のとおりである。

\begin{lemma}
\label{morphisms-lemma-etale-characterize}
$f : X \to S$ をスキームの射とする。次は同値である。
\begin{enumerate}
\item 射 $f$ はエタールである。
\item すべてのアフィン開集合 $U \subset X$、$V \subset S$ で
$f(U) \subset V$ を満たすものに対し、環準同型
$\mathcal{O}_S(V) \to \mathcal{O}_X(U)$ はエタールである。
\item 開被覆 $S = \bigcup_{j \in J} V_j$ と開被覆
$f^{-1}(V_j) = \bigcup_{i \in I_j} U_i$ が存在し、各射
$U_i \to V_j$、$j\in J, i\in I_j$ はエタールである。
\item アフィン開被覆 $S = \bigcup_{j \in J} V_j$ とアフィン開被覆
$f^{-1}(V_j) = \bigcup_{i \in I_j} U_i$ が存在し、環準同型
$\mathcal{O}_S(V_j) \to \mathcal{O}_X(U_i)$ は、すべての
$j\in J, i\in I_j$ に対してエタールである。
\end{enumerate}
さらに $f$ がエタールならば、任意の開部分スキーム
$U \subset X$、$V \subset S$ で $f(U) \subset V$ を満たすものに対し、
その制限 $f|_U : U \to V$ はエタールである。
\end{lemma}

\begin{proof}
「$R \to A$ はエタールである」という性質が局所的であることを
示せば、補題 \ref{morphisms-lemma-locally-P} から従う。
定義 \ref{morphisms-definition-property-local} の条件 (a)、(b)、(c) を確認する。
これらはすべて代数、補題 \ref{algebra-lemma-etale} から従う。
\end{proof}

\begin{lemma}
\label{morphisms-lemma-composition-etale}
二つのエタール射の合成はエタールである。
\end{lemma}

\begin{proof}
補題 \ref{morphisms-lemma-etale-characterize} の証明で、エタールであることが
環準同型の局所的性質であることを示した。従って補題の第一の主張は、
補題 \ref{morphisms-lemma-composition-type-P} と、エタールであることが合成に
関して安定な環準同型の性質であるという事実を組み合わせると従う。
代数、補題 \ref{algebra-lemma-etale} 参照。
\end{proof}

\begin{lemma}
\label{morphisms-lemma-base-change-etale}
エタール射の基底変換はエタールである。
\end{lemma}

\begin{proof}
補題 \ref{morphisms-lemma-etale-characterize} の証明で、エタールであることが
環準同型の局所的性質であることを示した。従って補題は、補題
\ref{morphisms-lemma-composition-type-P} と、エタールであることが基底変換に
関して安定な環準同型の性質であるという事実を組み合わせると従う。
代数、補題 \ref{algebra-lemma-etale} 参照。
\end{proof}

\begin{lemma}
\label{morphisms-lemma-etale-smooth-unramified}
$f : X \to S$ をスキームの射とし、$x \in X$ とする。このとき
$f$ が $x$ においてエタールであることと、$f$ が $x$ において
滑らかかつ不分岐であることは同値である。
\end{lemma}

\begin{proof}
定義から直ちに従う。
\end{proof}

\begin{lemma}
\label{morphisms-lemma-etale-locally-quasi-finite}
エタール射は局所準有限である。
\end{lemma}

\begin{proof}
補題 \ref{morphisms-lemma-etale-smooth-unramified} により、エタール射は
不分岐である。補題 \ref{morphisms-lemma-unramified-quasi-finite} により、
不分岐射は局所準有限である。
\end{proof}

\begin{lemma}
\label{morphisms-lemma-etale-over-field}
\begin{slogan}
体上のエタールスキームとエタール射のファイバーの記述。
\end{slogan}
エタール射のファイバーについて、次が成り立つ。
\begin{enumerate}
\item $X$ を体 $k$ 上のスキームとする。構造射
$X \to \Spec(k)$ がエタールであることと、$X$ が $k$ の有限分離拡大の
スペクトルの非交和であることは同値である。
\item $f : X \to S$ がエタール射ならば、任意の $s \in S$ に対して
ファイバー $X_s$ は $\kappa(s)$ の有限分離拡大のスペクトルの
非交和である。
\end{enumerate}
\end{lemma}

\begin{proof}
補題 \ref{morphisms-lemma-unramified-over-field} からこれを導ける。その際、上の補題
\ref{morphisms-lemma-etale-smooth-unramified} を用いる。直接の証明も与える。

\medskip\noindent
代数、補題 \ref{algebra-lemma-etale-over-field} を用いる。
従って、$X$ が $k$ の有限分離拡大のスペクトルの非交和ならば、
$X \to \Spec(k)$ がエタールであることは明らかである。逆に、
$X \to \Spec(k)$ がエタールであると仮定する。このとき任意の
アフィン開集合 $U \subset X$ に対して、$U$ は $k$ の有限分離拡大の
スペクトルの有限非交和である。従って $X$ のすべての点は閉点である
（例えば補題
\ref{morphisms-lemma-algebraic-residue-field-extension-closed-point-fibre} 参照）。
ゆえに $X$ は離散空間であり、主張を得る。
\end{proof}

\noindent
次の補題は、エタール射を「エタールなファイバー」をもつ平坦かつ
有限表示な射として特徴づける。

\begin{lemma}
\label{morphisms-lemma-etale-flat-etale-fibres}
$f : X \to S$ をスキームの射とする。$f$ が平坦かつ局所有限表示であり、
任意の $s \in S$ に対してファイバー $X_s$ が $\kappa(s)$ の
有限分離拡大のスペクトルの非交和ならば、$f$ はエタールである。
\end{lemma}

\begin{proof}
代数、補題 \ref{algebra-lemma-characterize-etale} からこれを導ける。
別の証明を与える。

\medskip\noindent
補題 \ref{morphisms-lemma-etale-over-field} により、ファイバー $X_s$ は
$s$ 上エタールであり、従って滑らかである。補題
\ref{morphisms-lemma-smooth-flat-smooth-fibres} により、$X \to S$ は滑らかである。
補題 \ref{morphisms-lemma-unramified-etale-fibres} により、$f$ は不分岐である。
補題 \ref{morphisms-lemma-etale-smooth-unramified} により結論を得る。
\end{proof}

\begin{lemma}
\label{morphisms-lemma-open-immersion-etale}
任意の開埋め込みはエタールである。
\end{lemma}

\begin{proof}
開埋め込みは局所同型なので、これは成り立つ。
\end{proof}

\begin{lemma}
\label{morphisms-lemma-etale-syntomic}
エタール射は syntomic である。
\end{lemma}

\begin{proof}
代数、補題 \ref{algebra-lemma-smooth-syntomic} を参照し、エタール射が
相対次元 $0$ の滑らかな射と同じものであることを用いる。
\end{proof}

\begin{lemma}
\label{morphisms-lemma-etale-locally-finite-presentation}
エタール射は局所有限表示である。
\end{lemma}

\begin{proof}
エタールな環準同型は定義により有限表示なので成り立つ。
\end{proof}

\begin{lemma}
\label{morphisms-lemma-etale-flat}
エタール射は平坦である。
\end{lemma}

\begin{proof}
補題 \ref{morphisms-lemma-syntomic-flat} と \ref{morphisms-lemma-etale-syntomic} を
組み合わせる。
\end{proof}

\begin{lemma}
\label{morphisms-lemma-etale-open}
エタール射は開である。
\end{lemma}

\begin{proof}
補題 \ref{morphisms-lemma-etale-flat}、
\ref{morphisms-lemma-etale-locally-finite-presentation}、および
\ref{morphisms-lemma-fppf-open} を組み合わせる。
\end{proof}

\noindent
次の補題は、局所的には任意のエタール射が標準エタールであることを
述べる。完全な一般性のもとでこれを証明するのは、実のところ少し
技巧を要する。難しい部分は可換代数の章に隠されている。従って、
標準エタール射は一般のエタール射の {\it 局所モデル} である。

\begin{lemma}
\label{morphisms-lemma-etale-locally-standard-etale}
$f : X  \to S$ をスキームの射とし、$x \in X$ を点とする。
$V \subset S$ を $f(x)$ のアフィン開近傍とする。次は同値である。
\begin{enumerate}
\item 射 $f$ は $x$ においてエタールである。
\item アフィン開集合 $U \subset X$ が存在して $x \in U$ と
$f(U) \subset V$ を満たし、誘導される射 $f|_U : U \to V$ は
標準エタールである（定義 \ref{morphisms-definition-etale} 参照）。
\end{enumerate}
\end{lemma}

\begin{proof}
定義と代数、命題
\ref{algebra-proposition-etale-locally-standard} から従う。
\end{proof}

\noindent
点におけるエタール性の微分による判定法を述べる。この結果には、
いずれもどこかで有用となりうる多くの変種がある。必要になった
ものをここに追加していく。

\begin{lemma}
\label{morphisms-lemma-etale-at-point}
$f : X \to S$ をスキームの射とし、$x \in X$ とする。
$s = f(x)$ とおく。$f$ は局所有限表示であると仮定する。
次は同値である。
\begin{enumerate}
\item 射 $f$ は $x$ においてエタールである。
\item 局所環準同型 $\mathcal{O}_{S, s} \to \mathcal{O}_{X, x}$ は
平坦であり、$X_s \to \Spec(\kappa(s))$ は $x$ において
エタールである。
\item 局所環準同型 $\mathcal{O}_{S, s} \to \mathcal{O}_{X, x}$ は
平坦であり、$X_s \to \Spec(\kappa(s))$ は $x$ において
不分岐である。
\item 局所環準同型 $\mathcal{O}_{S, s} \to \mathcal{O}_{X, x}$ は
平坦であり、$\mathcal{O}_{X, x}$-加群 $\Omega_{X/S, x}$ は零である。
\item 局所環準同型 $\mathcal{O}_{S, s} \to \mathcal{O}_{X, x}$ は
平坦であり、$\kappa(x)$-ベクトル空間
$$
\Omega_{X_s/s, x} \otimes_{\mathcal{O}_{X_s, x}} \kappa(x) =
\Omega_{X/S, x} \otimes_{\mathcal{O}_{X, x}} \kappa(x)
$$
は零である。
\item 局所環準同型 $\mathcal{O}_{S, s} \to \mathcal{O}_{X, x}$ は
平坦であり、$\mathfrak m_s\mathcal{O}_{X, x} = \mathfrak m_x$ が成り立ち、
体拡大 $\kappa(x)/\kappa(s)$ は有限分離拡大である。
\item アフィン開集合 $U \subset X$、$V \subset S$ が存在して
$x \in U$、$f(U) \subset V$ を満たし、誘導される射
$f|_U : U \to V$ は相対次元 $0$ の標準滑らかな射である。
\item アフィン開集合 $\Spec(A) = U \subset X$ と
$\Spec(R) = V \subset S$ が存在し、$x \in U$ は
$\mathfrak q \subset A$ に対応し、$f(U) \subset V$ が成り立ち、かつ表示
$$
A = R[x_1, \ldots, x_n]/(f_1, \ldots, f_n)
$$
が存在して
$$
g =
\det
\left(
\begin{matrix}
\partial f_1/\partial x_1 &
\partial f_2/\partial x_1 &
\ldots &
\partial f_n/\partial x_1 \\
\partial f_1/\partial x_2 &
\partial f_2/\partial x_2 &
\ldots &
\partial f_n/\partial x_2 \\
\ldots & \ldots & \ldots & \ldots \\
\partial f_1/\partial x_n &
\partial f_2/\partial x_n &
\ldots &
\partial f_n/\partial x_n
\end{matrix}
\right)
$$
の $A$ における像が $\mathfrak q$ に属さない。
\item アフィン開集合 $U \subset X$、$V \subset S$ が存在して
$x \in U$、$f(U) \subset V$ を満たし、誘導される射
$f|_U : U \to V$ は標準エタールである。
\item アフィン開集合 $\Spec(A) = U \subset X$ と
$\Spec(R) = V \subset S$ が存在し、$x \in U$ は
$\mathfrak q \subset A$ に対応し、$f(U) \subset V$ が成り立ち、かつ表示
$$
A = R[x]_Q/(P) = R[x, 1/Q]/(P)
$$
が存在する。ここで $P, Q \in R[x]$ であり、$P$ はモニックで、
$P' = \text{d}P/\text{d}x$ の $A$ における像は $\mathfrak q$ に属さない。
\end{enumerate}
\end{lemma}

\begin{proof}
補題 \ref{morphisms-lemma-etale-locally-standard-etale} と定義を用いると、
(1) から他のすべての条件が従うことがわかる。各条件 (2) -- (10) について、
補題 \ref{morphisms-lemma-smooth-at-point} と \ref{morphisms-lemma-unramified-at-point} を
組み合わせると、$f$ が $x$ において滑らかかつ不分岐であることを
示せる。従って補題 \ref{morphisms-lemma-etale-smooth-unramified} を適用すれば
(1) が成り立つ。詳細の一部は省略する。
\end{proof}

\begin{lemma}
\label{morphisms-lemma-flat-unramified-etale}
射が点においてエタールであることと、その点において平坦かつ
G-不分岐であることは同値である。射がエタールであることと、
それが平坦かつ G-不分岐であることは同値である。
\end{lemma}

\begin{proof}
補題 \ref{morphisms-lemma-etale-at-point} と \ref{morphisms-lemma-unramified-at-point} から
明らかである。
\end{proof}

\begin{lemma}
\label{morphisms-lemma-set-points-where-fibres-etale}
スキームの Cartesian 図式
$$
\xymatrix{
X' \ar[r]_{g'} \ar[d]_{f'} & X \ar[d]^f \\
S' \ar[r]^g & S
}
$$
を考える。$W \subset X$、resp.\ $W' \subset X'$ を、$f$、resp.\ $f'$ が
エタールである点からなる開部分スキームとする。このとき、次の
いずれかが成り立てば $W' = (g')^{-1}(W)$ である。
\begin{enumerate}
\item $f$ は平坦かつ局所有限表示である。
\item $f$ は局所有限表示であり、$g$ は平坦である。
\end{enumerate}
\end{lemma}

\begin{proof}
まず $f$ が局所有限型であると仮定する。集合
$$
T = \{x \in X \mid f\text{ is unramified at }x\}
$$
と、対応する集合 $T' \subset X'$（$f'$ に対するもの）を考える。このとき、
補題 \ref{morphisms-lemma-set-points-where-fibres-unramified} により
$T' = (g')^{-1}(T)$ である。

\medskip\noindent
従って (1) の場合、補題 \ref{morphisms-lemma-flat-unramified-etale} により
$T$ は $f$ がエタールである点の（開）集合なので、(1) が従う。

\medskip\noindent
(2) の場合、$x' \in W'$ とする。このとき $g'$ は $x'$ において平坦で
あり（補題 \ref{morphisms-lemma-base-change-module-flat}）、$g \circ f'$ も $x'$ に
おいて平坦である（補題 \ref{morphisms-lemma-composition-module-flat}）。従って
補題 \ref{morphisms-lemma-flat-permanence} により、$f$ は $x = g'(x')$ において
平坦である。一方、$x' \in T'$ なので（補題
\ref{morphisms-lemma-base-change-smooth}）、$x \in T$ である。従って補題
\ref{morphisms-lemma-etale-at-point} により、$f$ は $x$ においてエタールである。
\end{proof}

\begin{lemma}
\label{morphisms-lemma-etale-permanence-better}
\begin{reference}
例えば \cite[Remark 18.30 (2).]{GWII} 参照。
\end{reference}
$f : X \to Y$ を $S$ 上のスキームの射とする。$X$ が $S$ 上エタールで
あり、$Y$ が $S$ 上不分岐ならば、$f$ はエタールである。
\end{lemma}

\begin{proof}
分解 $X \to X \times_S Y \to Y$ を考える。第一の矢印は
$\text{id}_X$ と $f$ によって与えられ、第二の矢印は射影である。
両方の矢印がエタールであることを示せば、補題
\ref{morphisms-lemma-composition-etale} により $f$ がエタールであると従う。
実際、射影は $X \to S$ の基底変換なのでエタールである。補題
\ref{morphisms-lemma-base-change-etale} 参照。第一の矢印は対角射
$Y \to Y \times_S Y$ の基底変換である。なぜなら、図式
$$
\xymatrix{
X \ar[d] \ar[r] & X \times_S Y \ar[d] \\
Y \ar[r] & Y \times_S Y
}
$$
は Cartesian だからである。補題
\ref{morphisms-lemma-diagonal-unramified-morphism} により、対角射
$Y \to Y \times_S Y$ は開埋め込みである。開埋め込みの基底変換は
開埋め込みであり（スキーム、補題
\ref{schemes-lemma-base-change-immersion}）、開埋め込みはエタールである
（補題 \ref{morphisms-lemma-open-immersion-etale}）。
\end{proof}

\begin{lemma}
\label{morphisms-lemma-etale-permanence}
\begin{slogan}
エタール射の消去法則。
\end{slogan}
$f : X \to Y$ を $S$ 上のスキームの射とする。$X$ と $Y$ が
$S$ 上エタールならば、$f$ はエタールである。
\end{lemma}

\begin{proof}
補題 \ref{morphisms-lemma-etale-smooth-unramified} により $Y \to S$ は不分岐なので、
補題 \ref{morphisms-lemma-etale-permanence-better} を適用できる。代数、補題
\ref{algebra-lemma-map-between-etale} も参照。
\end{proof}

\begin{lemma}
\label{morphisms-lemma-etale-permanence-two}
スキームの射の可換図式
$$
\xymatrix{
X \ar[rr]_f \ar[rd]_p & &
Y \ar[dl]^q \\
& S
}
$$
があり、次を仮定する。
\begin{enumerate}
\item $f$ は全射かつエタールである。
\item $p$ はエタールである。
\item $q$ は局所有限表示である\footnote{実際、これは (1) と (2) から
従う。降下、補題
\ref{descent-lemma-flat-finitely-presented-permanence} 参照。
さらに、$f$ が全射、平坦かつ局所有限表示であると仮定するだけで十分で
ある。降下、補題 \ref{descent-lemma-smooth-permanence} 参照。}。
\end{enumerate}
このとき $q$ はエタールである。
\end{lemma}

\begin{proof}
補題 \ref{morphisms-lemma-smooth-permanence} により、$q$ は滑らかである。
従って $q$ の相対次元が $0$ であることだけを示せばよい。
これは補題 \ref{morphisms-lemma-dimension-fibre-at-a-point-additive} と、$f$ と $p$ の
相対次元が $0$ であるという事実から従う。
\end{proof}

\noindent
滑らかな射の最後の特徴づけは、滑らかな射 $f : X \to S$ が局所的に、
エタール射と射影 $\mathbf{A}_S^d \to S$ の合成であるというものである。

\begin{lemma}
\label{morphisms-lemma-smooth-etale-over-affine-space}
\begin{slogan}
滑らかなスキームは、エタール局所的にはアフィン空間のように見える。
\end{slogan}
$\varphi : X \to Y$ をスキームの射とし、$x \in X$ とする。
$V \subset Y$ を $\varphi(x)$ のアフィン開近傍とする。
$\varphi$ が $x$ において滑らかならば、整数 $d \geq 0$ とアフィン開集合
$U \subset X$ が存在して $x \in U$ と $\varphi(U) \subset V$ を満たし、
さらに可換図式
$$
\xymatrix{
X \ar[d] & U \ar[l] \ar[d] \ar[r]_-\pi & \mathbf{A}^d_V \ar[ld] \\
Y & V \ar[l]
}
$$
が存在する。ここで $\pi$ はエタールである。
\end{lemma}

\begin{proof}
補題 \ref{morphisms-lemma-smooth-locally-standard-smooth} により、補題にあるような
アフィン開集合 $U$ で、$\varphi|_U : U \to V$ が標準滑らかとなるものを
とれる。$U = \Spec(A)$、$V = \Spec(R)$ と書けば、
$$
A = R[x_1, \ldots, x_n]/(f_1, \ldots, f_c)
$$
と書け、さらに
$$
g =
\det
\left(
\begin{matrix}
\partial f_1/\partial x_1 &
\partial f_2/\partial x_1 &
\ldots &
\partial f_c/\partial x_1 \\
\partial f_1/\partial x_2 &
\partial f_2/\partial x_2 &
\ldots &
\partial f_c/\partial x_2 \\
\ldots & \ldots & \ldots & \ldots \\
\partial f_1/\partial x_c &
\partial f_2/\partial x_c &
\ldots &
\partial f_c/\partial x_c
\end{matrix}
\right)
$$
の $A$ における像は可逆である。このとき
$R[x_{c + 1}, \ldots, x_n] \to A$ が相対次元 $0$ の標準滑らかな
環準同型であることは明らかである。従って、それは相対次元 $0$ の
滑らかな環準同型である。言い換えれば、環準同型
$R[x_{c + 1}, \ldots, x_n] \to A$ はエタールである。
$\mathbf{A}^{n - c}_V = \Spec(R[x_{c + 1}, \ldots, x_n])$ なので、
$d = n - c$ として補題が成り立つ。
\end{proof}

\section{相対的に豊富な層}
\label{morphisms-section-relatively-ample}

\noindent
$X$ をスキームとし、$\mathcal{L}$ を $X$ 上の可逆層とする。
このとき $\mathcal{L}$ が $X$ 上豊富であるとは、$X$ が準コンパクトで
あり、$X$ のすべての点が $X_s$ の形のアフィン開集合に含まれることを
いう。ここで
$s \in \Gamma(X, \mathcal{L}^{\otimes n})$ かつ $n \geq 1$ である。
性質、定義 \ref{properties-definition-ample} 参照。
これを次のように相対的な概念にする。

\begin{definition}
\label{morphisms-definition-relatively-ample}
\begin{reference}
\cite[II Definition 4.6.1]{EGA}
\end{reference}
$f : X \to S$ をスキームの射とし、$\mathcal{L}$ を可逆
$\mathcal{O}_X$-加群とする。$\mathcal{L}$ が {\it 相対的に豊富}、
または {\it $f$-相対的に豊富}、または {\it $X/S$ 上豊富}、または
{\it $f$-豊富}であるとは、$f : X \to S$ が準コンパクトであり、
任意のアフィン開集合 $V \subset S$ に対して、$\mathcal{L}$ の
開部分スキーム $f^{-1}(V)$（$X$ のもの）への制限が豊富であることをいう。
\end{definition}

\noindent
$X$ 上に相対的に豊富な層が存在しても、射 $X \to S$ が有限型である
とは限らないことに注意せよ。

\begin{lemma}
\label{morphisms-lemma-ample-power-ample}
$X \to S$ をスキームの射とし、$\mathcal{L}$ を可逆
$\mathcal{O}_X$-加群とする。$n \geq 1$ とする。このとき
$\mathcal{L}$ が $f$-豊富であることと、$\mathcal{L}^{\otimes n}$ が
$f$-豊富であることは同値である。
\end{lemma}

\begin{proof}
性質、補題 \ref{properties-lemma-ample-power-ample} から従う。
\end{proof}

\begin{lemma}
\label{morphisms-lemma-relatively-ample-separated}
$f : X \to S$ をスキームの射とする。$f$-豊富な可逆層が存在すれば、
$f$ は分離的である。
\end{lemma}

\begin{proof}
分離的であることは基底上局所的である（例えばスキーム、補題
\ref{schemes-lemma-characterize-separated} 参照。定義からも容易に従う）。
従って $S$ がアフィンで、$X$ が豊富な可逆層をもつと仮定してよい。
この場合、結果は性質、補題
\ref{properties-lemma-ample-separated} から従う。
\end{proof}

\noindent
相対的に豊富な可逆層を特徴づける方法は多くあり、性質、命題
\ref{properties-proposition-characterize-ample} の同値条件に類似している。
必要に応じてここに追加していく。

\begin{lemma}
\label{morphisms-lemma-characterize-relatively-ample}
\begin{reference}
\cite[II, Proposition 4.6.3]{EGA}
\end{reference}
$f : X \to S$ をスキームの準コンパクト射とし、$\mathcal{L}$ を
$X$ 上の可逆層とする。次は同値である。
\begin{enumerate}
\item 可逆層 $\mathcal{L}$ は $f$-豊富である。
\item 開被覆 $S = \bigcup V_i$ が存在し、各
$\mathcal{L}|_{f^{-1}(V_i)}$ は $f^{-1}(V_i) \to V_i$ に関して
相対的に豊富である。
\item アフィン開被覆 $S = \bigcup V_i$ が存在し、各
$\mathcal{L}|_{f^{-1}(V_i)}$ は豊富である。
\item 準連接次数付き $\mathcal{O}_S$-代数 $\mathcal{A}$ と、次数付き
$\mathcal{O}_X$-代数の写像
$\psi : f^*\mathcal{A} \to \bigoplus_{d \geq 0} \mathcal{L}^{\otimes d}$
が存在して、$U(\psi) = X$ であり、
$$
r_{\mathcal{L}, \psi} :
X
\longrightarrow
\underline{\text{Proj}}_S(\mathcal{A})
$$
は開埋め込みである（記法については構成、補題
\ref{constructions-lemma-invertible-map-into-relative-proj} 参照）。
\item 射 $f$ は準分離的であり、(4) が
$\mathcal{A} = f_*(\bigoplus_{d \geq 0} \mathcal{L}^{\otimes d})$ と
随伴写像 $\psi$ に対して成り立つ。
\item (4) と同じであるが、$r_{\mathcal{L}, \psi}$ が埋め込みであること
だけを要求する。
\end{enumerate}
\end{lemma}

\begin{proof}
定義から、(1) ならば (2)、(2) ならば (3) であることは直ちにわかる。
(5) ならば (4) であることも明らかである。

\medskip\noindent
(3) がアフィン開被覆 $S = \bigcup V_i$ に対して成り立つと仮定し、
(5) を示す。各 $f^{-1}(V_i)$ は豊富な可逆層をもつので、
$f^{-1}(V_i)$ は分離的である（性質、補題
\ref{properties-lemma-ample-separated}）。従って $f$ は分離的である。
スキーム、補題 \ref{schemes-lemma-push-forward-quasi-coherent} により、
$\mathcal{A} = f_*(\bigoplus_{d \geq 0} \mathcal{L}^{\otimes d})$ は
準連接次数付き $\mathcal{O}_S$-代数である。随伴写像を
$\psi : f^*\mathcal{A} \to \bigoplus_{d \geq 0} \mathcal{L}^{\otimes d}$
と記す。構成、節
\ref{constructions-section-invertible-relative-proj} における開集合
$U(\psi)$ の記述と、$\mathcal{L}|_{f^{-1}(V_i)}$ の豊富性の定義から、
$U(\psi) = X$ が従う。さらに構成、補題
\ref{constructions-lemma-invertible-map-into-relative-proj} の (3) により、
$r_{\mathcal{L}, \psi}$ の $f^{-1}(V_i)$ への制限は、性質、補題
\ref{properties-lemma-map-into-proj} の射と同じである。この射は性質、補題
\ref{properties-lemma-ample-immersion-into-proj} により開埋め込みである。
従って (5) が成り立つ。

\medskip\noindent
(4) ならば (1) であることを示そう。(4) を仮定する。構造射を
$\pi : \underline{\text{Proj}}_S(\mathcal{A}) \to S$ と記し、
アフィン開集合 $V \subset S$ を選ぶ。構成、定義
\ref{constructions-definition-relative-proj} により、
$\pi^{-1}(V) \subset \underline{\text{Proj}}_S(\mathcal{A})$ は
$\text{Proj}(A)$ に等しく、ここで次数付き環として
$A = \mathcal{A}(V)$ である。従って $r_{\mathcal{L}, \psi}$ は
$f^{-1}(V)$ を $\text{Proj}(A)$ の準コンパクトな開集合へ同型に写す。
さらに、$\mathcal{L}^{\otimes d}$ は $\mathcal{O}_{\text{Proj}(A)}(d)$ の
引き戻しと同型である。ただし、ある $d \geq 1$ に対してである。
（構成、補題 \ref{constructions-lemma-invertible-map-into-relative-proj} の
(3) と、構成、補題
\ref{constructions-lemma-invertible-map-into-proj} の最後の主張を参照。）
従って性質、補題 \ref{properties-lemma-open-in-proj-ample} と
\ref{properties-lemma-ample-power-ample} により、
$\mathcal{L}|_{f^{-1}(V)}$ は豊富である。

\medskip\noindent
(6) を仮定する。上で示した (1) - (5) の同値性により、$X/S$ 上
相対的に豊富であるという性質は $S$ 上局所的である。従って $S$ は
アフィンであると仮定してよく、$\mathcal{L}$ が $X$ 上豊富であることを
示せばよい。この場合、射 $r_{\mathcal{L}, \psi}$ は、同じ記号で表す射
$r_{\mathcal{L}, \psi} : X \to \text{Proj}(A)$、すなわち写像
$\psi : A = \mathcal{A}(V) \to \Gamma_*(X, \mathcal{L})$ に付随する射と
同一視される。
（上の参照先を参照。）また上と同様に、$\mathcal{L}^{\otimes d}$ は層
$\mathcal{O}_{\text{Proj}(A)}(d)$ の引き戻しである。ただし、ある
$d \geq 1$ に対してである。さらに $X$ は準コンパクトなので、$X$ は
準コンパクトな開部分スキーム $Y \subset \text{Proj}(A)$ の
閉部分スキームと同一視される。
構成、補題 \ref{constructions-lemma-ample-on-proj}（性質、補題
\ref{properties-lemma-open-in-proj-ample} も参照）により、
$\mathcal{O}_Y(d')$ は $Y$ 上の豊富な可逆層である。ただし、ある
$d' \geq 1$ に対してである。
豊富な層の閉部分スキームへの制限は豊富なので（性質、補題
\ref{properties-lemma-ample-on-closed} 参照）、
$\mathcal{O}_Y(d')$ の引き戻しは豊富である。これらの結果を性質、補題
\ref{properties-lemma-ample-power-ample} と組み合わせると、望むとおり
$\mathcal{L}$ は豊富である。
\end{proof}

\begin{lemma}
\label{morphisms-lemma-ample-over-affine}
\begin{reference}
\cite[II Corollary 4.6.6]{EGA}
\end{reference}
$f : X \to S$ をスキームの射とし、$\mathcal{L}$ を可逆
$\mathcal{O}_X$-加群とする。$S$ はアフィンであると仮定する。
このとき $\mathcal{L}$ が $f$-相対的に豊富であることと、
$\mathcal{L}$ が $X$ 上豊富であることは同値である。
\end{lemma}

\begin{proof}
補題 \ref{morphisms-lemma-characterize-relatively-ample} と定義から直ちに従う。
\end{proof}

\begin{lemma}
\label{morphisms-lemma-quasi-affine-O-ample}
\begin{reference}
\cite[II Proposition 5.1.6]{EGA}
\end{reference}
$f : X \to S$ をスキームの射とする。このとき $f$ が準アフィンである
ことと、$\mathcal{O}_X$ が $f$-相対的に豊富であることは同値である。
\end{lemma}

\begin{proof}
性質、補題 \ref{properties-lemma-quasi-affine-O-ample} と定義から従う。
\end{proof}

\begin{lemma}
\label{morphisms-lemma-pullback-ample-tensor-relatively-ample}
$f : X \to Y$ をスキームの射、$\mathcal{M}$ を可逆
$\mathcal{O}_Y$-加群、$\mathcal{L}$ を可逆 $\mathcal{O}_X$-加群とする。
\begin{enumerate}
\item $\mathcal{L}$ が $f$-豊富であり、$\mathcal{M}$ が豊富ならば、
$\mathcal{L} \otimes f^*\mathcal{M}^{\otimes a}$ は $a \gg 0$ に対して
豊富である。
\item $\mathcal{M}$ が豊富で、$f$ が準アフィンならば、
$f^*\mathcal{M}$ は豊富である。
\end{enumerate}
\end{lemma}

\begin{proof}
$\mathcal{L}$ が $f$-豊富であり、$\mathcal{M}$ が豊富であると仮定する。
仮定により $Y$ と $f$ は準コンパクトである（定義
\ref{morphisms-definition-relatively-ample} および性質、定義
\ref{properties-definition-ample} 参照）。従って $X$ は準コンパクトである。
性質、補題 \ref{properties-lemma-ample-separated} によりスキーム $Y$ は
分離的であり、補題 \ref{morphisms-lemma-relatively-ample-separated} により射 $f$ は
分離的である。従ってスキーム、補題
\ref{schemes-lemma-separated-permanence} により $X$ は分離的である。
$x \in X$ を選ぶ。$m \geq 1$ と
$t \in \Gamma(Y, \mathcal{M}^{\otimes m})$ を、$Y_t$ がアフィンで
$f(x) \in Y_t$ となるように選べる。$\mathcal{L}$ の
$f^{-1}(Y_t) = X_{f^*t}$ への制限は豊富な可逆層なので、$n \geq 1$ と
$s \in \Gamma(X_{f^*t}, \mathcal{L}^{\otimes n})$ を、
$x \in (X_{f^*t})_s$ かつ $(X_{f^*t})_s$ がアフィンとなるように選べる。
上の仮定を満たす性質、補題
\ref{properties-lemma-invert-s-sections} の (2) により、整数 $e \geq 1$ と
切断
$s' \in \Gamma(X, \mathcal{L}^{\otimes n} \otimes f^*\mathcal{M}^{\otimes em})$
が存在し、$s(f^*t)^e$ へ $X_{f^*t}$ 上で制限される。任意の $b > 0$ に
対して、切断 $s'' = s'(f^*t)^b$（
$\mathcal{L}^{\otimes n} \otimes f^*\mathcal{M}^{\otimes (e + b)m}$ のもの）
を考える。このとき
$X_{s''} = (X_{f^*t})_s$ は $X$ の $x$ を含むアフィン開集合である。
$b$ を、$n$ が $e + b$ を割り切るように選ぶと、
$\mathcal{L}^{\otimes n} \otimes f^*\mathcal{M}^{\otimes (e + b)m}$ は、
$n$ 乗であり、その底は $\mathcal{L} \otimes f^*\mathcal{M}^{\otimes a}$
（ある $a$ に対するもの）である。また、十分大きい任意の $a$ であって
$m$ で割り切れるものを得られる。
$X$ は準コンパクトなので、これらのアフィン開集合の有限個が $X$ を
被覆する。従って、十分に割り切れ、かつ十分大きいある $a$ に対して、
可逆層 $\mathcal{L} \otimes f^*\mathcal{M}^{\otimes a}$ は $X$ 上豊富である。
一方、性質、命題 \ref{properties-proposition-characterize-ample} により、
$\mathcal{M}^{\otimes c}$（従ってその $X$ への引き戻し）は、すべての
$c \gg 0$ に対して大域生成される。従って
$\mathcal{L} \otimes f^*\mathcal{M}^{\otimes a + c}$ は $c \gg 0$ に対して
豊富である
（性質、補題 \ref{properties-lemma-ample-tensor-globally-generated}）。
これで (1) が証明された。

\medskip\noindent
(2) は、補題 \ref{morphisms-lemma-quasi-affine-O-ample}、性質、補題
\ref{properties-lemma-ample-power-ample}、および (1) から従う。
\end{proof}

\begin{lemma}
\label{morphisms-lemma-ample-composition}
$g : Y \to S$ と $f : X \to Y$ をスキームの射とする。
$\mathcal{M}$ を可逆 $\mathcal{O}_Y$-加群、$\mathcal{L}$ を可逆
$\mathcal{O}_X$-加群とする。$S$ が準コンパクトで、$\mathcal{M}$ が
$g$-豊富、$\mathcal{L}$ が $f$-豊富ならば、
$\mathcal{L} \otimes f^*\mathcal{M}^{\otimes a}$ は $g \circ f$-豊富であり、
これは $a \gg 0$ に対して成り立つ。
\end{lemma}

\begin{proof}
$S = \bigcup_{i = 1, \ldots, n} V_i$ を有限アフィン開被覆とする。
補題 \ref{morphisms-lemma-characterize-relatively-ample} により、
$\mathcal{L} \otimes f^*\mathcal{M}^{\otimes a}$ が
$(g \circ f)^{-1}(V_i)$ 上豊富であることを、すべての
$i = 1, \ldots, n$ に対して
示せば十分である。従って補題
\ref{morphisms-lemma-pullback-ample-tensor-relatively-ample} から従う。
\end{proof}

\begin{lemma}
\label{morphisms-lemma-ample-base-change}
$f : X \to S$ をスキームの射とし、$\mathcal{L}$ を可逆
$\mathcal{O}_X$-加群とする。$S' \to S$ をスキームの射とする。
$f' : X' \to S'$ を $f$ の基底変換とし、$\mathcal{L}'$ を
$\mathcal{L}$ の $X'$ への引き戻しと記す。$\mathcal{L}$ が $f$-豊富ならば、
$\mathcal{L}'$ は $f'$-豊富である。
\end{lemma}

\begin{proof}
補題 \ref{morphisms-lemma-characterize-relatively-ample} により、アフィン開被覆
$S' = \bigcup U'_i$ で、$\mathcal{L}'$ の $(f')^{-1}(U_i')$ への制限が
豊富な可逆層となるものを、すべての $i$ に対して見つければ十分である。
$U'_i$ がアフィン開集合 $U_i \subset S$ に写るように選んでよい。
この場合、射 $(f')^{-1}(U'_i) \to f^{-1}(U_i)$ は、アフィン射
$U'_i \to U_i$ の基底変換なのでアフィンである（補題
\ref{morphisms-lemma-base-change-affine}）。従って補題
\ref{morphisms-lemma-pullback-ample-tensor-relatively-ample} により、
$\mathcal{L}'|_{(f')^{-1}(U'_i)}$ は豊富である。
\end{proof}

\begin{lemma}
\label{morphisms-lemma-ample-permanence}
$g : Y \to S$ と $f : X \to Y$ をスキームの射とし、$\mathcal{L}$ を
可逆 $\mathcal{O}_X$-加群とする。$\mathcal{L}$ が $g \circ f$-豊富で、
$f$ が準コンパクトならば\footnote{$g$ が準分離的ならば、これは
スキーム、補題 \ref{schemes-lemma-quasi-compact-permanence} から従う。}、
$\mathcal{L}$ は $f$-豊富である。
\end{lemma}

\begin{proof}
$f$ は準コンパクトで、$\mathcal{L}$ は $g \circ f$-豊富であると仮定する。
$U \subset S$ をアフィン開集合、$V \subset Y$ を
$g(V) \subset U$ を満たすアフィン開集合とする。仮定により、
$\mathcal{L}|_{(g \circ f)^{-1}(U)}$ は $(g \circ f)^{-1}(U)$ 上豊富である。
$f^{-1}(V) \subset (g \circ f)^{-1}(U)$ なので、性質、補題
\ref{properties-lemma-ample-on-locally-closed} により、
$\mathcal{L}|_{f^{-1}(V)}$ は $f^{-1}(V)$ 上豊富である。実際、
$f^{-1}(V) \to (g \circ f)^{-1}(U)$ は、スキーム、補題
\ref{schemes-lemma-quasi-compact-permanence} により準コンパクトな
開埋め込みである。実際、$(g \circ f)^{-1}(U)$ は分離的であり
（性質、補題 \ref{properties-lemma-ample-separated}）、$f^{-1}(V)$ は
準コンパクトである（$f$ が準コンパクトだから）。
従って補題 \ref{morphisms-lemma-characterize-relatively-ample} により、
$\mathcal{L}$ は $f$-豊富である。
\end{proof}

\section{非常に豊富な層}
\label{morphisms-section-very-ample}

\noindent
準連接層 $\mathcal{E}$ がスキーム $S$ 上に与えられたとき、
$\mathcal{E}$ に付随する {\it 射影束}とは射
$\mathbf{P}(\mathcal{E}) \to S$ のことである。ここで
$\mathbf{P}(\mathcal{E}) = \underline{\text{Proj}}_S(\text{Sym}(\mathcal{E}))$
である。構成、定義
\ref{constructions-definition-projective-bundle} 参照。

\begin{definition}
\label{morphisms-definition-very-ample}
$f : X \to S$ をスキームの射とし、$\mathcal{L}$ を可逆
$\mathcal{O}_X$-加群とする。$\mathcal{L}$ が
{\it 相対的に非常に豊富}、より正確には
{\it $f$-相対的に非常に豊富}、または {\it $X/S$ 上非常に豊富}、
または {\it $f$-非常に豊富}であるとは、準連接
$\mathcal{O}_S$-加群 $\mathcal{E}$ と浸入
$i : X \to \mathbf{P}(\mathcal{E})$（$S$ 上のもの）が存在して、
$\mathcal{L} \cong i^*\mathcal{O}_{\mathbf{P}(\mathcal{E})}(1)$
となることをいう。
\end{definition}



\noindent
この定義には準コンパクト性の仮定がないので、一般には、相対的に
非常に豊富な可逆層が相対的に豊富な可逆層であるとは限らない。

\begin{lemma}
\label{morphisms-lemma-ample-very-ample}
\begin{reference}
\cite[II, Proposition 4.6.2]{EGA}
\end{reference}
$f : X \to S$ をスキームの射とし、$\mathcal{L}$ を可逆
$\mathcal{O}_X$-加群とする。$f$ が準コンパクトで、$\mathcal{L}$ が
相対的に非常に豊富な可逆層ならば、$\mathcal{L}$ は相対的に豊富な
可逆層である。
\end{lemma}

\begin{proof}
定義により、準連接 $\mathcal{O}_S$-加群 $\mathcal{E}$ と
浸入 $i : X \to \mathbf{P}(\mathcal{E})$ が $S$ 上に存在して、
$\mathcal{L} \cong i^*\mathcal{O}_{\mathbf{P}(\mathcal{E})}(1)$
となる。$\mathcal{A} = \text{Sym}(\mathcal{E})$ とおくと、定義により
$\mathbf{P}(\mathcal{E}) = \underline{\text{Proj}}_S(\mathcal{A})$
である。次数付き $\mathcal{O}_S$-代数 $\mathcal{A}$ には写像
$$
\psi :
\mathcal{A} \to
\bigoplus\nolimits_{n \geq 0}
\pi_*\mathcal{O}_{\mathbf{P}(\mathcal{E})}(n) \to
\bigoplus\nolimits_{n \geq 0}
f_*\mathcal{L}^{\otimes n}
$$
が備わっている。第二の矢印では同一視
$\mathcal{L} \cong i^*\mathcal{O}_{\mathbf{P}(\mathcal{E})}(1)$
を用いる。$f_*$ と $f^*$ の随伴性により射
$\psi : f^*\mathcal{A} \to \bigoplus_{n \geq 0}\mathcal{L}^{\otimes n}$
を得る。この写像に付随する射 $r_{\mathcal{L}, \psi}$ がちょうど浸入
$i$ であることの確認は省略する。従って結果は補題
\ref{morphisms-lemma-characterize-relatively-ample} の (6) から従う。
\end{proof}

\noindent
この補題の正しい逆を得るため、相対的に豊富な可逆層
$\mathcal{L}$ が与えられたとき、整数 $n \geq 1$ で
$\mathcal{L}^{\otimes n}$ が相対的に非常に豊富となるものが
存在するかを問う。
一般にはこれは偽である。これが成り立つことを妨げるものがいくつかある。
\begin{enumerate}
\item $S$ がアフィンであっても、どの有限な整数
$n$ も使えないことがある。これは $X \to S$ が有限型でないからである。例
\ref{morphisms-example-not-finite-type-proj} 参照。
\item 基底が準コンパクトでなければ、$f$ が有限型であっても結果が
成り立たないことがある。実際、体 $k$ が与えられると、スキーム
$X_d$ で $k$ 上有限型のものと、その豊富な可逆層
$\mathcal{O}_{X_d}(1)$ で、
$\mathcal{O}_{X_d}(1)$ の非常に豊富なテンソル冪のうち最小のものが
$d$ 乗となるものが存在する。例 \ref{morphisms-example-not-bounded} 参照。
射 $f$ を、スキーム $X_d$ の非交和から $\Spec(k)$ のコピーの
非交和への射とすれば、例が得られる。
\end{enumerate}
逆の主張の我々の版については、下の補題
\ref{morphisms-lemma-finite-type-ample-very-ample} を見よ。
その証明の前に準備を行う。

\begin{example}
\label{morphisms-example-very-ample}
$S$ をスキームとする。$\mathcal{A}$ を準連接次数付き
$\mathcal{O}_S$-代数で、$\mathcal{A}_1$ により
$\mathcal{A}_0$ 上生成されるものとする。
$X = \underline{\text{Proj}}_S(\mathcal{A})$ とおく。この場合、
$\mathcal{O}_X(1)$ は可逆であり、$X/S$ 上非常に豊富である。実際、
次数付き $\mathcal{O}_S$-代数写像
$$
\text{Sym}_{\mathcal{O}_X}^*(\mathcal{A}_1)
\longrightarrow
\mathcal{A}
$$
に付随する射は閉埋め込み
$X \to \mathbf{P}(\mathcal{A}_1)$ であり、
$\mathcal{O}_{\mathbf{P}(\mathcal{A}_1)}(1)$ を
$\mathcal{O}_X(1)$ に引き戻す。構成、補題
\ref{constructions-lemma-surjective-generated-degree-1-map-relative-proj}
参照。
\end{example}

\begin{example}
\label{morphisms-example-not-finite-type-proj}
$k$ を体とする。次数付き $k$-代数
$$
A = k[U, V, Z_1, Z_2, Z_3, \ldots]/I
\quad
\text{with}
\quad
I = (U^2 - Z_1^2, U^4 - Z_2^2, U^6 - Z_3^2, \ldots)
$$
を考える。次数付けは
$\deg(U) = \deg(V) = \deg(Z_1) = 1$ および $\deg(Z_d) = d$
で与える。$X = \text{Proj}(A)$ は $D_{+}(U)$ と $D_{+}(V)$ で
被覆されることに注意せよ。従って層 $\mathcal{O}_X(n)$ はすべて
可逆であり、$\mathcal{O}_X(1)^{\otimes n}$ と同型である。
特に、$\mathcal{O}_X(1)$ は豊富であり、$f$-豊富でもある。
ここで射は $f : X \to \Spec(k)$ である。
$\mathcal{O}_X(1)$ のどの冪も $f$-相対的に非常に豊富ではないと
主張する。実際、$\Gamma(X, \mathcal{O}_X(n))$ は次数
$n$ の成分（代数 $A$ のもの）であることが容易にわかる。従って
$\mathcal{O}_X(n)$ が非常に豊富ならば、$X$ は $k$ 上の射影空間の
閉部分スキームとなり、従って $k$ 上有限型となる。一方、
$D_{+}(V)$ は
$k[t, t_1, t_2, \ldots]/(t^2 - t_1^2, t^4 - t_2^2, t^6 - t_3^2, \ldots)$
のスペクトルであり、これは $k$ 上有限型ではない。
\end{example}

\begin{example}
\label{morphisms-example-not-bounded}
$k$ を無限体とする。$\lambda_1, \lambda_2, \lambda_3, \ldots$ を
$k^*$ の相異なる元とする。（これは厳密には必要でなく、実際には
すべての $\lambda_i$ が $1$ に等しくても、この例は完全に成り立つ。）
次数付き $k$-代数
$$
A_d = k[U, V, Z]/I_d
\quad
\text{with}
\quad
I_d = (Z^2 - \prod\nolimits_{i = 1}^{2d} (U - \lambda_i V)).
$$
を考える。次数付けは $\deg(U) = \deg(V) = 1$ および
$\deg(Z) = d$ で与える。このとき
$X_d = \text{Proj}(A_d)$ は豊富な可逆層
$\mathcal{O}_{X_d}(1)$ をもつ。
$\mathcal{O}_{X_d}(n)$ が非常に豊富ならば $n \geq d$ であると
主張する。その理由は、$Z$ の次数が $d$ であり、従って
$\Gamma(X_d, \mathcal{O}_{X_d}(n)) =
k[U, V]_n$ が $n < d$ に対して成り立つからである。詳細は省略する。
\end{example}

\begin{lemma}
\label{morphisms-lemma-relatively-very-ample-separated}
$f : X \to S$ をスキームの射とし、$\mathcal{L}$ を $X$ 上の
可逆層とする。$\mathcal{L}$ が $X/S$ 上相対的に非常に豊富ならば、
$f$ は分離的である。
\end{lemma}

\begin{proof}
分離的であることは基底上局所的である（スキーム、節
\ref{schemes-section-separation-axioms} 参照）。浸入は分離的である
（スキーム、補題 \ref{schemes-lemma-immersions-monomorphisms} 参照）。
従って、局所的に $X$ は分離的な次数付き環の斉次スペクトルに
浸入するので、この補題は従う。構成、補題
\ref{constructions-lemma-proj-separated} 参照。
\end{proof}

\begin{lemma}
\label{morphisms-lemma-relatively-very-ample}
$f : X \to S$ をスキームの射とし、$\mathcal{L}$ を $X$ 上の
可逆層とする。$f$ は準コンパクトであると仮定する。次は同値である。
\begin{enumerate}
\item $\mathcal{L}$ は $X/S$ 上相対的に非常に豊富である。
\item 開被覆 $S = \bigcup V_j$ が存在して、
$\mathcal{L}|_{f^{-1}(V_j)}$ は $f^{-1}(V_j)/V_j$ 上相対的に
非常に豊富であり、これはすべての $j$ に対して成り立つ。
\item 準連接次数付き $\mathcal{O}_S$-代数の層 $\mathcal{A}$ で、
次数 $1$ で $\mathcal{O}_S$ 上生成されるものと、次数付き
$\mathcal{O}_X$-代数の写像
$\psi : f^*\mathcal{A} \to \bigoplus_{n \geq 0} \mathcal{L}^{\otimes n}$
が存在して、$f^*\mathcal{A}_1 \to \mathcal{L}$ は全射であり、
付随する射
$r_{\mathcal{L}, \psi} : X \to \underline{\text{Proj}}_S(\mathcal{A})$
は浸入である。
\item $f$ は準分離的であり、標準写像
$\psi : f^*f_*\mathcal{L} \to \mathcal{L}$ は全射であり、付随する写像
$r_{\mathcal{L}, \psi} : X \to \mathbf{P}(f_*\mathcal{L})$ は浸入である。
\end{enumerate}
\end{lemma}

\begin{proof}
(1) ならば (2) であることは明らかである。(4) ならば (1) である
ことも明らかである。(4) の準分離性の仮定は、スキーム、補題
\ref{schemes-lemma-push-forward-quasi-coherent} によって
$f_*\mathcal{L}$ が準連接であることを保証するために用いられる。

\medskip\noindent
(2) を仮定し、(4) を証明する。$S = \bigcup V_j$ を (2) の開被覆と
する。$X_j = f^{-1}(V_j)$ とおき、$f_j : X_j \to V_j$ を $f$ の
制限とする。補題 \ref{morphisms-lemma-relatively-very-ample-separated} により
$f$ は分離的である（分離的であることは基底上局所的だから）。
仮定により、準連接 $\mathcal{O}_{V_j}$-加群 $\mathcal{E}_j$ と浸入
$i_j : X_j \to \mathbf{P}(\mathcal{E}_j)$ が存在して、
$\mathcal{L}|_{X_j} \cong i_j^*\mathcal{O}_{\mathbf{P}(\mathcal{E}_j)}(1)$
となる。射 $i_j$ は全射
$f_j^*\mathcal{E}_j \to \mathcal{L}|_{X_j}$ に対応する。構成、節
\ref{constructions-section-projective-bundle} 参照。この写像は写像
$\mathcal{E}_j \to f_*\mathcal{L}|_{V_j}$ に随伴し、その合成
$$
f_j^*\mathcal{E}_j \to (f^*f_*\mathcal{L})|_{X_j} \to \mathcal{L}|_{X_j}
$$
は全射である。従って
$\psi : f^*f_*\mathcal{L} \to \mathcal{L}$ は全射である。付随する射を
$r_{\mathcal{L}, \psi} : X \to \mathbf{P}(f_*\mathcal{L})$ とする。
$r_{\mathcal{L}, \psi}$ が浸入であることを示す必要が残るが、
読者自身でこれを証明されたい。$\mathcal{O}_{V_j}$-加群の写像
$\mathcal{E}_j \to f_*\mathcal{L}|_{V_j}$ は対称代数上の準同型を定め、
これはさらに射
$$
\mathbf{P}(f_*\mathcal{L}|_{V_j}) \supset U_j \longrightarrow
\mathbf{P}(\mathcal{E}_j)
$$
を定める。ここで $U_j$ は構成、補題
\ref{constructions-lemma-morphism-relative-proj} の開部分スキームである。
$\psi$ と $\mathcal{E}_j \to f_*\mathcal{L}|_{V_j}$ の両立性から、
$r_{\mathcal{L}, \psi}(X_j) \subset U_j$ であり、分解
$$
\xymatrix{
X_j \ar[r]^-{r_{\mathcal{L}, \psi}} & U_j \ar[r] &
\mathbf{P}(\mathcal{E}_j)
}
$$
が存在することがわかる。確認は省略する。
これは $r_{\mathcal{L}, \psi}$ が浸入であることを示す。

\medskip\noindent
ここまでで (1)、(2)、(4) が同値であることがわかった。
(4) ならば (3) であることは明らかである。(3) を仮定し、(1) を証明する。
$\mathcal{A}$ を準連接次数付き $\mathcal{O}_S$-代数の層で、
次数 $1$ で $\mathcal{O}_S$ 上生成されるものとする。次数付き
$\mathcal{O}_S$-代数の写像
$\text{Sym}(\mathcal{A}_1) \to \mathcal{A}$ を考える。
これは仮定により全射であり、従って閉埋め込み
$$
\underline{\text{Proj}}_S(\mathcal{A})
\longrightarrow
\mathbf{P}(\mathcal{A}_1)
$$
を誘導し、この埋め込みは $\mathcal{O}(1)$ を $\mathcal{O}(1)$ に
引き戻す。構成、補題
\ref{constructions-lemma-surjective-generated-degree-1-map-relative-proj}
参照。従って (3) ならば (1) であることは明らかである。
\end{proof}

\begin{lemma}
\label{morphisms-lemma-very-ample-base-change}
$f : X \to S$ をスキームの射とし、$\mathcal{L}$ を可逆
$\mathcal{O}_X$-加群とする。$S' \to S$ をスキームの射とする。
$f' : X' \to S'$ を $f$ の基底変換とし、$\mathcal{L}'$ を
$\mathcal{L}$ の $X'$ への引き戻しと記す。
$\mathcal{L}$ が $f$-非常に豊富ならば、$\mathcal{L}'$ は
$f'$-非常に豊富である。
\end{lemma}

\begin{proof}
定義 \ref{morphisms-definition-very-ample} により、準連接
$\mathcal{O}_S$-加群 $\mathcal{E}$ と浸入
$i : X \to \mathbf{P}(\mathcal{E})$（$S$ 上のもの）が存在して、
$\mathcal{L} \cong i^*\mathcal{O}_{\mathbf{P}(\mathcal{E})}(1)$
となる。$\mathbf{P}(\mathcal{E})$ の $S'$ への基底変換は、
引き戻し $\mathcal{E}'$（$\mathcal{E}$ のもの）に付随する射影束であり、
$\mathcal{O}_{\mathbf{P}(\mathcal{E})}(1)$ の引き戻しは
$\mathcal{O}_{\mathbf{P}(\mathcal{E}')}(1)$ である。構成、補題
\ref{constructions-lemma-relative-proj-base-change} 参照。
最後に、浸入の基底変換は浸入である
（スキーム、補題 \ref{schemes-lemma-base-change-immersion}）。
\end{proof}






\section{有限型射に関する豊富な層と非常に豊富な層}
\label{morphisms-section-ample-finite-type}

\noindent
実際、本節の内容の大部分は、まだ定義していない（準）射影的射の
概念に関するものである。

\begin{lemma}
\label{morphisms-lemma-very-ample-finite-type-over-affine}
$f : X \to S$ をスキームの射とする。
$\mathcal{L}$ を $X$ 上の可逆層とする。
次を仮定する。
\begin{enumerate}
\item 可逆層 $\mathcal{L}$ は $X/S$ 上非常に豊富である。
\item 射 $X \to S$ は有限型である。
\item $S$ はアフィンである。
\end{enumerate}
このとき、ある $n \geq 0$ と浸入
$i : X \to \mathbf{P}^n_S$ が $S$ 上に存在して、
$\mathcal{L} \cong i^*\mathcal{O}_{\mathbf{P}^n_S}(1)$ となる。
\end{lemma}

\begin{proof}
(1)、(2)、(3) を仮定する。
条件 (3) は、$S = \Spec(R)$ がある環 $R$ について
成り立つことを意味する。
条件 (1) は、定義により、準連接 $\mathcal{O}_S$-加群
$\mathcal{E}$ と浸入 $\alpha : X \to \mathbf{P}(\mathcal{E})$
が存在して
$\mathcal{L} = \alpha^*\mathcal{O}_{\mathbf{P}(\mathcal{E})}(1)$
となることを意味する。
$\mathcal{E} = \widetilde{M}$ と書く。ここで、これはある
$R$-加群 $M$ によるものとする。
従って
$$
\mathbf{P}(\mathcal{E}) = \text{Proj}(\text{Sym}_R(M)).
$$
を得る。$\alpha$ は浸入であり、かつ
$\text{Proj}(\text{Sym}_R(M))$ の位相は標準開集合
$D_{+}(f)$、$f \in \text{Sym}_R^d(M)$、$d \geq 1$ により生成されるので、
各 $x \in X$ に対して、ある
$f \in \text{Sym}_R^d(M)$、$d \geq 1$ で
$\alpha(x) \in D_{+}(f)$ を満たし、
$$
\alpha|_{\alpha^{-1}(D_{+}(f))} : \alpha^{-1}(D_{+}(f)) \to D_{+}(f)
$$
が閉埋め込みとなるものを見つけられる。
条件 (2) から $X$ は準コンパクトである。従って有限個の元
$f_j \in \text{Sym}_R^{d_j}(M)$、$d_j \geq 1$ で、各 $f = f_j$ に対して
上の表示された写像が閉埋め込みであり、かつ
$\alpha(X) \subset \bigcup D_{+}(f_j)$ となるものを選べる。
$U_j = \alpha^{-1}(D_{+}(f_j))$ と書く。$U_j$ はアフィンスキーム
$D_{+}(f_j)$ の閉部分スキームとしてアフィンであることに注意する。
$U_j = \Spec(A_j)$ と書く。条件 (2) からさらに、補題
\ref{morphisms-lemma-locally-finite-type-characterize} により
$A_j$ は $R$ 上有限型である。
有限個の $x_{j, k} \in A_j$ で、$A_j$ を $R$-代数として
生成するものを選ぶ。
$\alpha|_{U_j}$ は閉埋め込みなので、$x_{j, k}$ はある元
$$
f_{j, k}/f_j^{e_{j, k}} \in \text{Sym}_R(M)_{(f_j)}
=
\Gamma(D_{+}(f_j), \mathcal{O}_{\text{Proj}(\text{Sym}_R(M))}).
$$
の像である。最後に $n \geq 1$ と元 $y_0, \ldots, y_n \in M$ を、
各多項式 $f_j, f_{j, k} \in \text{Sym}_R(M)$ が
$y_t$ の多項式で、その係数が $R$ に属するものとなるように選ぶ。
次数付き環写像
$$
\psi : R[Y_0, \ldots, Y_n] \longrightarrow \text{Sym}_R(M),
\quad Y_i \longmapsto y_i.
$$
を考える。$F_j$, $F_{j, k}$ を
$R[Y_0, \ldots, Y_n]$ の元で、
$\psi(F_j) = f_j$ および $\psi(F_{j, k}) = f_{j, k}$
を満たすものと記す。
構成、補題 \ref{constructions-lemma-morphism-proj} により、開部分スキーム
$$
U(\psi) \subset \text{Proj}(\text{Sym}_R(M))
$$
と射 $r_\psi : U(\psi) \to \mathbf{P}^n_R$ を得る。この射は
$r_\psi^{-1}(D_{+}(F_j)) = D_{+}(f_j)$ を満たすので、
$\alpha(X) \subset U(\psi)$ であることがわかる。さらに、
$$
i = r_\psi \circ \alpha : X \longrightarrow \mathbf{P}^n_R
$$
もなお浸入である。実際、構成により
$i^\sharp(F_{j, k}/F_j^{e_{j, k}}) = x_{j, k} \in
A_j = \Gamma(U_j, \mathcal{O}_X)$
だからである。また、射 $r_\psi$ には写像
$\theta : r_\psi^*\mathcal{O}_{\mathbf{P}^n_R}(1)
\to \mathcal{O}_{\text{Proj}(\text{Sym}_R(M))}(1)|_{U(\psi)}$
が備わっており、この場合それは同型である（$\theta$ の構成については
上で引用した補題を参照せよ。いくつかの詳細は省略する）。
もとの写像 $\alpha$ は
$\mathcal{L} = \alpha^*\mathcal{O}_{\text{Proj}(\text{Sym}_R(M))}(1)$
を満たすと仮定したので、望む結論を得る。
\end{proof}

\begin{lemma}
\label{morphisms-lemma-quasi-affine-finite-type-over-S}
$\pi : X \to S$ をスキームの射とする。
$X$ が準アフィンであり、$\pi$ が局所有限型であると仮定する。
このとき、ある $n \geq 0$ と浸入
$i : X \to \mathbf{A}^n_S$ が $S$ 上に存在する。
\end{lemma}

\begin{proof}
$A = \Gamma(X, \mathcal{O}_X)$ とおく。仮定により $X$ は
準コンパクトであり、性質、補題 \ref{properties-lemma-quasi-affine} により
$\Spec(A)$ の開部分スキームと同一視される。
さらに、$X_f$ がアフィンとなるような $f \in A$ に対する開集合
$X_f$ 全体は $X$ の位相の基をなす。性質、補題
\ref{properties-lemma-quasi-affine} の証明を参照せよ。
従って、有限個の $f_j \in A$、$j = 1, \ldots, m$ で
$X = \bigcup X_{f_j}$ となり、$\pi(X_{f_j}) \subset V_j$ が
あるアフィン開集合 $V_j \subset S$ に対して成り立つものを選べる。
補題 \ref{morphisms-lemma-locally-finite-type-characterize} により、環写像
$\mathcal{O}(V_j) \to \mathcal{O}(X_{f_j}) = A_{f_j}$ は有限型である。
従って $a_1, \ldots, a_N \in A$ を、元
$a_1, \ldots, a_N, 1/f_j$ が $A_{f_j}$ を
$\mathcal{O}(V_j)$ 上生成することが各 $j$ について成り立つように選べる。
$n = m + N$ とおき、大域切断
$f_1, \ldots, f_m, a_1, \ldots, a_N$（$X$ の構造層のもの）
によって与えられる射を
$$
i : X \longrightarrow \mathbf{A}^n_S
$$
とする。$D(x_j) \subset \mathbf{A}^n_S$ を、第 $j$ 座標関数が
零でない開部分スキームとする。
すると $1 \leq j \leq m$ に対して
$i^{-1}(D(x_j)) = X_{f_j}$ であり、誘導される射
$X_{f_j} \to D(x_j)$ は、アフィン開集合
$\Spec(\mathcal{O}(V_j)[x_1, \ldots, x_n, 1/x_j])$
（$D(x_j)$ のもの）を経由する。
環写像
$\mathcal{O}(V_j)[x_1, \ldots, x_n, 1/x_j] \to A_{f_j}$ は
構成により全射なので、$i^{-1}(D(x_j)) \to D(x_j)$ は望みどおり
浸入であると結論する。
\end{proof}

\begin{lemma}
\label{morphisms-lemma-quasi-projective-finite-type-over-S}
$f : X \to S$ をスキームの射とする。
$\mathcal{L}$ を $X$ 上の可逆層とする。
次を仮定する。
\begin{enumerate}
\item 可逆層 $\mathcal{L}$ は $X$ 上豊富である。
\item 射 $X \to S$ は局所有限型である。
\end{enumerate}
このとき、ある $d_0 \geq 1$ が存在して、すべての $d \geq d_0$ に対し
ある $n \geq 0$ と浸入
$i : X \to \mathbf{P}^n_S$ が $S$ 上に存在し、
$\mathcal{L}^{\otimes d} \cong i^*\mathcal{O}_{\mathbf{P}^n_S}(1)$
となる。
\end{lemma}

\begin{proof}
$A = \Gamma_*(X, \mathcal{L}) =
\bigoplus_{d \geq 0} \Gamma(X, \mathcal{L}^{\otimes d})$ とおく。
性質、命題 \ref{properties-proposition-characterize-ample} により、
$X_a$ がアフィン開集合となるような斉次元 $a \in A_{+}$ 全体は
$X$ の位相の基をなす。従って有限個のそのような元
$a_0, \ldots, a_n \in A_{+}$ で次を満たすものを選べる。
\begin{enumerate}
\item $X = \bigcup_{i = 0, \ldots, n} X_{a_i}$ である。
\item 各 $X_{a_i}$ はアフィンである。
\item 各 $X_{a_i}$ はあるアフィン開集合 $V_i \subset S$ に写る。
\end{enumerate}
補題 \ref{morphisms-lemma-locally-finite-type-characterize} により、環写像
$\mathcal{O}_S(V_i) \to \mathcal{O}_X(X_{a_i})$ は有限型である。
従って有限個の元
$f_{ij} \in \mathcal{O}_X(X_{a_i})$、$j = 1, \ldots, n_i$ で、
$\mathcal{O}_X(X_{a_i})$ を $\mathcal{O}_S(V_i)$-代数として
生成するものを選べる。性質、補題
\ref{properties-lemma-invert-s-sections} により、各
$f_{ij}$ を $a_{ij}/a_i^{e_{ij}}$ と書ける。ここで
$a_{ij} \in A_{+}$ はある斉次元である。$N$ を、元 $a_i$、$a_{ij}$ の
すべての次数の正の共通倍数とする。元
$$
a_i^{N/\deg(a_i)}, \ a_{ij}a_i^{(N/\deg(a_i)) - e_{ij}} \in A_N.
$$
を考える。構成により、これらは可逆層
$\mathcal{L}^{\otimes N}$ を $X$ 上生成する。従って、構成、補題
\ref{constructions-lemma-projective-space} および定義
\ref{constructions-definition-projective-space} により、これらから
$S$ 上の射
$$
j : X \longrightarrow
\mathbf{P}_S^{m}
\quad
\text{with } m = n + \sum n_i
$$
が生じる。さらに
$j^*\mathcal{O}_{\mathbf{P}_S}(1) = \mathcal{L}^{\otimes N}$ である。
斉次座標を $T_0, \ldots, T_n, T_{ij}$ と名づけ、
$T_0, \ldots, T_m$ とはしない。
$i = 0, \ldots, n$ に対して
$j^{-1}(D_{+}(T_i)) = X_{a_i}$ である。
さらに、元 $T_{ij}/T_i$ を $j^\sharp$ によって引き戻すと、
元 $f_{ij} \in \mathcal{O}_X(X_{a_i})$ を得る。
従って、$j$ の $X_{a_i}$ への制限は、$X_{a_i}$ からアフィン開集合
$D_{+}(T_i) \cap \mathbf{P}^m_{V_i}$（$\mathbf{P}^N_S$ のもの）
への閉埋め込みを与える。
従って射 $j$ は浸入であると結論する。
これは、ある $d$ と $n$ に対して補題が成り立つことを意味し、
実質上すべての応用にはこれで十分である。

\medskip\noindent
以上により、ある $d_2 \geq 1$（すなわち上の $d_2 = N$）と
ある $m \geq 0$ に対して、浸入
$j : X \to \mathbf{P}^m_S$ が存在し、これは大域切断
$s'_0, \ldots, s'_m \in \Gamma(X, \mathcal{L}^{\otimes d_2})$
により与えられることを示した。性質、命題
\ref{properties-proposition-characterize-ample} により、
ある整数 $d_1$ が存在して、$\mathcal{L}^{\otimes d}$ は
すべての $d \geq d_1$ に対して大域生成される。
$d_0 = d_1 + d_2$ とおく。この $d_0$ によって補題が成り立つと主張する。
実際、整数 $d \geq d_0$ が与えられたとき、
$s''_1, \ldots, s''_t
\in \Gamma(X, \mathcal{L}^{\otimes d - d_2})$ で、
$\mathcal{L}^{\otimes d - d_2}$ を $X$ 上生成するものを選べる。
$k = (m + 1)t$ とおき、$s_0, \ldots, s_k$ を切断
$s'_\alpha s''_\beta$、$\alpha = 0, \ldots, m$、
$\beta = 1, \ldots, t$ の族と記す。
これらは $\mathcal{L}^{\otimes d}$ を $X$ 上生成するので、射
$$
i : X \longrightarrow \mathbf{P}^{k - 1}_S
$$
で $i^*\mathcal{O}_{\mathbf{P}^n_S}(1) \cong \mathcal{L}^{\otimes d}$
を満たすものを定める。
$i$ が浸入であることを見るには、$i$ が合成
$$
X \longrightarrow \mathbf{P}^m_S \times_S \mathbf{P}^{t - 1}_S
\longrightarrow \mathbf{P}^{k - 1}_S
$$
であることに注意すればよい。ここで第一の射は $(j, j')$ であり、
$j'$ は $s''_1, \ldots, s''_t$ により与えられる。
第二の射は Segre 埋め込み
（構成、補題 \ref{constructions-lemma-segre-embedding}）である。
$j$ は浸入なので $(j, j')$ も浸入である
（$X \to \mathbf{P}^m_S \times_S \mathbf{P}^{t - 1}_S
\to \mathbf{P}^m_S$ に補題 \ref{morphisms-lemma-immersion-permanence} を適用せよ）。
従って $i$ は浸入の合成であり、それゆえ浸入である
（スキーム、補題 \ref{schemes-lemma-composition-immersion}）。
\end{proof}

\begin{lemma}
\label{morphisms-lemma-finite-type-over-affine-ample-very-ample}
$f : X \to S$ をスキームの射とする。
$\mathcal{L}$ を可逆 $\mathcal{O}_X$-加群とする。
$S$ はアフィンであり、$f$ は有限型であると仮定する。
次は同値である。
\begin{enumerate}
\item $\mathcal{L}$ は $X$ 上豊富である。
\item $\mathcal{L}$ は $f$-豊富である。
\item $\mathcal{L}^{\otimes d}$ が $f$-非常に豊富となる
$d \geq 1$ が存在する。
\item $\mathcal{L}^{\otimes d}$ は $f$-非常に豊富であり、
これはすべての $d \gg 1$ に対して成り立つ。
\item ある $d \geq 1$ に対して、ある $n \geq 1$ と浸入
$i : X \to \mathbf{P}^n_S$ が存在して
$\mathcal{L}^{\otimes d} \cong i^*\mathcal{O}_{\mathbf{P}^n_S}(1)$
となる。
\item すべての $d \gg 1$ に対して、ある $n \geq 1$ と浸入
$i : X \to \mathbf{P}^n_S$ が存在して
$\mathcal{L}^{\otimes d} \cong i^*\mathcal{O}_{\mathbf{P}^n_S}(1)$
となる。
\end{enumerate}
\end{lemma}

\begin{proof}
(1) と (2) の同値性は補題 \ref{morphisms-lemma-ample-over-affine} である。
(2) $\Rightarrow$ (6) は補題
\ref{morphisms-lemma-quasi-projective-finite-type-over-S} である。
(6) が (5) を含意することは自明である。
$\mathbf{P}^n_S$ は $S$ 上の射影束なので（構成、補題
\ref{constructions-lemma-projective-space-bundle} 参照）、
相対的に非常に豊富な層の定義から、(5) は (3) を含意し、
(6) は (4) を含意する。(4) が (3) を含意することは自明である。
最後に (3) が (2) を含意することを示せばよいが、これは補題
\ref{morphisms-lemma-ample-very-ample} と補題
\ref{morphisms-lemma-ample-power-ample} から従う。
\end{proof}

\begin{lemma}
\label{morphisms-lemma-finite-type-ample-very-ample}
$f : X \to S$ をスキームの射とする。
$\mathcal{L}$ を可逆 $\mathcal{O}_X$-加群とする。
$S$ は準コンパクトであり、$f$ は有限型であると仮定する。
次は同値である。
\begin{enumerate}
\item $\mathcal{L}$ は $f$-豊富である。
\item $\mathcal{L}^{\otimes d}$ が $f$-非常に豊富となる
$d \geq 1$ が存在する。
\item $\mathcal{L}^{\otimes d}$ は $f$-非常に豊富であり、
これはすべての $d \gg 1$ に対して成り立つ。
\end{enumerate}
\end{lemma}

\begin{proof}
(3) が (2) を含意することは自明である。補題
\ref{morphisms-lemma-ample-very-ample} により、(2) は (1) を含意する。
これは、有限型射が定義により準コンパクトだからである。
$\mathcal{L}$ が $f$-豊富であると仮定する。有限アフィン開被覆
$S = V_1 \cup \ldots \cup V_m$ を選ぶ。$X_i = f^{-1}(V_i)$ と書く。
上の補題 \ref{morphisms-lemma-finite-type-over-affine-ample-very-ample} により、
ある $d_0$ が存在して、$\mathcal{L}^{\otimes d}$ は
$X_i/V_i$ 上相対的に非常に豊富であり、これはすべての
$d \geq d_0$ に対して成り立つ。
従って補題 \ref{morphisms-lemma-relatively-very-ample} により、
(1) が (3) を含意すると結論する。
\end{proof}

\noindent
次の二つの補題は、射が有限型であるときの、相対的に非常に豊富な
可逆層と相対的に豊富な可逆層について、最もよく使われ、最も有用な
特徴づけを与える。

\begin{lemma}
\label{morphisms-lemma-characterize-very-ample-on-finite-type}
$f : X \to S$ をスキームの射とする。
$\mathcal{L}$ を $X$ 上の可逆層とする。
$f$ が有限型であると仮定する。
次は同値である。
\begin{enumerate}
\item $\mathcal{L}$ は $f$-相対的に非常に豊富である。
\item 開被覆 $S = \bigcup V_j$、各 $j$ に対する整数 $n_j$、
および $V_j$ 上の浸入
$$
i_j :
X_j = f^{-1}(V_j) = V_j \times_S X
\longrightarrow
\mathbf{P}^{n_j}_{V_j}
$$
が存在して、
$\mathcal{L}|_{X_j} \cong i_j^*\mathcal{O}_{\mathbf{P}^{n_j}_{V_j}}(1)$
となる。
\end{enumerate}
\end{lemma}

\begin{proof}
$S$ のアフィン開被覆を取り、$f$ と $\mathcal{L}$ の各制限に
補題 \ref{morphisms-lemma-very-ample-finite-type-over-affine} を適用すれば、
(1) が (2) を含意することがわかる。
補題 \ref{morphisms-lemma-relatively-very-ample} により、
(2) が (1) を含意することがわかる。
\end{proof}

\begin{lemma}
\label{morphisms-lemma-characterize-ample-on-finite-type}
$f : X \to S$ をスキームの射とする。
$\mathcal{L}$ を $X$ 上の可逆層とする。
$f$ が有限型であると仮定する。
次は同値である。
\begin{enumerate}
\item $\mathcal{L}$ は $f$-相対的に豊富である。
\item 開被覆 $S = \bigcup V_j$、各 $j$ に対する整数
$d_j \geq 1$、$n_j \geq 0$、および $V_j$ 上の浸入
$$
i_j :
X_j = f^{-1}(V_j) = V_j \times_S X
\longrightarrow
\mathbf{P}^{n_j}_{V_j}
$$
が存在して、
$\mathcal{L}^{\otimes d_j}|_{X_j} \cong
i_j^*\mathcal{O}_{\mathbf{P}^{n_j}_{V_j}}(1)$
となる。
\end{enumerate}
\end{lemma}

\begin{proof}
$S$ のアフィン開被覆を取り、$f$ と $\mathcal{L}$ の各制限に
補題 \ref{morphisms-lemma-finite-type-over-affine-ample-very-ample} を適用すれば、
(1) が (2) を含意することがわかる。
補題 \ref{morphisms-lemma-characterize-relatively-ample} により、
(2) が (1) を含意することがわかる。
\end{proof}

\begin{lemma}
\label{morphisms-lemma-invertible-add-enough-ample-very-ample}
$f : X \to S$ をスキームの射とする。
$\mathcal{N}$、$\mathcal{L}$ を可逆 $\mathcal{O}_X$-加群とする。
$S$ は準コンパクトで、$f$ は有限型であり、$\mathcal{L}$ は
$f$-豊富であると仮定する。このとき
$\mathcal{N} \otimes_{\mathcal{O}_X} \mathcal{L}^{\otimes d}$
は $f$-非常に豊富であり、これはすべての $d \gg 1$ に対して成り立つ。
\end{lemma}

\begin{proof}
補題 \ref{morphisms-lemma-characterize-very-ample-on-finite-type} により、
$S$ がアフィンの場合に帰着する。補題
\ref{morphisms-lemma-finite-type-over-affine-ample-very-ample} と性質、命題
\ref{properties-proposition-characterize-ample} を組み合わせると、
ある整数 $d_0$ で
$\mathcal{N} \otimes \mathcal{L}^{\otimes d_0}$ が
大域生成されるものを選べる。
$s_0, \ldots, s_n$ を、$\mathcal{N} \otimes \mathcal{L}^{\otimes d_0}$
を生成する大域切断として選ぶ。これにより射
$j : X \to \mathbf{P}^n_S$ が $S$ 上に定まる。
補題 \ref{morphisms-lemma-finite-type-over-affine-ample-very-ample} により、
ある整数 $d_1$ を、すべての $d \geq d_1$ に対して切断
$t_{d, 0}, \ldots, t_{d, n(d)}$ で
$\mathcal{L}^{\otimes d}$ を生成し、浸入
$$
j_d = \varphi_{\mathcal{L}^{\otimes d}, t_{d, 0}, \ldots, t_{d, n(d)}} :
X
\longrightarrow
\mathbf{P}^{n(d)}_S
$$
を $S$ 上に定めるものが存在するように選ぶこともできる。
そこで $d \geq d_0 + d_1$ に対して射
$$
\varphi_{\mathcal{N} \otimes \mathcal{L}^{\otimes d},
s_j \otimes t_{d - d_0, i}} :
X
\longrightarrow
\mathbf{P}^{(n + 1)(n(d - d_0) + 1) - 1}_S
$$
を考えることができる。この射は合成
$$
X \to \mathbf{P}^n_S \times_S \mathbf{P}^{n(d - d_0)}_S
\to \mathbf{P}^{(n + 1)(n(d - d_0) + 1) - 1}_S
$$
なので浸入である。ここで第一の射は $(j, j_{d - d_0})$ であり、
第二の射は Segre 埋め込み
（構成、補題 \ref{constructions-lemma-segre-embedding}）である。
$j$ は浸入なので $(j, j_{d - d_0})$ も浸入である
（補題 \ref{morphisms-lemma-immersion-permanence} を適用せよ）。
従って浸入の合成を得ており、それゆえこれは浸入である
（スキーム、補題 \ref{schemes-lemma-composition-immersion}）。
\end{proof}


\section{準射影的射}
\label{morphisms-section-quasi-projective}

\noindent
前節の議論は次の定義を示唆する。準射影的という我々の定義は
\cite{EGA} から取る。文字「H」を伴う版は \cite{H} の定義である。

\begin{definition}
\label{morphisms-definition-quasi-projective}
\begin{reference}
\cite[II, Definition 5.3.1]{EGA} and \cite[page 103]{H}
\end{reference}
$f : X \to S$ をスキームの射とする。
\begin{enumerate}
\item $f$ が {\it 準射影的}であるとは、$f$ が有限型であり、
$f$-相対的に豊富な可逆 $\mathcal{O}_X$-加群が存在することをいう。
\item $f$ が {\it H-準射影的}であるとは、準コンパクトな浸入
$X \to \mathbf{P}^n_S$ が $S$ 上に存在し、これがある $n$ に対して
成り立つことをいう。
\footnote{これは Hartshorne における定義と厳密には同じではない。
すなわち Hartshorne の定義（第8訂正版、1997年）では、$f$ は開埋め込みと
それに続く H-射影的射（定義 \ref{morphisms-definition-projective} 参照）の合成で
なければならないが、これは $f$ が準コンパクトであることを含意しない。
逆方向の含意については補題
\ref{morphisms-lemma-H-quasi-projective-open-H-projective} を参照せよ。}
\item $f$ が {\it 局所準射影的}であるとは、開被覆
$S = \bigcup V_j$ で、各 $f^{-1}(V_j) \to V_j$ が
準射影的となるものが存在することをいう。
\end{enumerate}
\end{definition}

\noindent
この定義が示唆するとおり、準射影的であるという性質は $S$ 上
局所的ではない。後に準射影的スキームの圏についてさらに述べられる。
射の続き、節 \ref{more-morphisms-section-quasi-projective} を参照せよ。

\begin{lemma}
\label{morphisms-lemma-base-change-quasi-projective}
準射影的射の基底変換は準射影的である。
\end{lemma}

\begin{proof}
これは補題 \ref{morphisms-lemma-base-change-finite-type} と
\ref{morphisms-lemma-ample-base-change} から従う。
\end{proof}

\begin{lemma}
\label{morphisms-lemma-composition-quasi-projective}
$f : X \to Y$ および $g : Y \to S$ をスキームの射とする。
$S$ が準コンパクトで、$f$ と $g$ が準射影的ならば、
$g \circ f$ は準射影的である。
\end{lemma}

\begin{proof}
これは補題 \ref{morphisms-lemma-composition-finite-type} と
\ref{morphisms-lemma-ample-composition} から従う。
\end{proof}

\begin{lemma}
\label{morphisms-lemma-quasi-projective-properties}
$f : X \to S$ をスキームの射とする。$f$ が準射影的、
H-準射影的、または局所準射影的ならば、$f$ は分離的かつ有限型である。
\end{lemma}

\begin{proof}
省略する。
\end{proof}

\begin{lemma}
\label{morphisms-lemma-H-quasi-projective-quasi-projective}
H-準射影的射は準射影的である。
\end{lemma}

\begin{proof}
省略する。
\end{proof}

\begin{lemma}
\label{morphisms-lemma-characterize-locally-quasi-projective}
$f : X \to S$ をスキームの射とする。
次は同値である。
\begin{enumerate}
\item 射 $f$ は局所準射影的である。
\item 開被覆 $S = \bigcup V_j$ で、各
$f^{-1}(V_j) \to V_j$ が H-準射影的となるものが存在する。
\end{enumerate}
\end{lemma}

\begin{proof}
補題 \ref{morphisms-lemma-H-quasi-projective-quasi-projective} により、
(2) が (1) を含意することがわかる。(1) を仮定する。
問題は $S$ 上局所的なので、$S$ はアフィンであり、$X$ は $S$ 上
有限型であり、$\mathcal{L}$ は $X/S$ 上相対的に豊富な可逆層であると
仮定してよい。補題
\ref{morphisms-lemma-finite-type-over-affine-ample-very-ample} により、
$\mathcal{L}$ は $X$ 上豊富であると仮定してよい。
補題 \ref{morphisms-lemma-quasi-projective-finite-type-over-S} により、
$X$ から $S$ 上の射影空間への浸入が存在すること、すなわち
$X$ は $S$ 上 H-準射影的であることがわかり、望む結論を得る。
\end{proof}

\begin{lemma}
\label{morphisms-lemma-quasi-affine-finite-type-quasi-projective}
\begin{reference}
\cite[II, Proposition 5.3.4 (i)]{EGA}
\end{reference}
有限型の準アフィン射は準射影的である。
\end{lemma}

\begin{proof}
これは補題 \ref{morphisms-lemma-quasi-affine-O-ample} から従う。
\end{proof}

\begin{lemma}
\label{morphisms-lemma-quasi-projective-permanence}
$g : Y \to S$ および $f : X \to Y$ をスキームの射とする。
$g \circ f$ が準射影的であり、$f$ が準コンパクト
\footnote{$g$ が準分離的ならば、これはスキーム、補題
\ref{schemes-lemma-quasi-compact-permanence} から従う。}
ならば、$f$ は準射影的である。
\end{lemma}

\begin{proof}
補題 \ref{morphisms-lemma-permanence-finite-type} により
$f$ は有限型であることに注意する。
従ってこの補題は補題 \ref{morphisms-lemma-ample-permanence} と定義から従う。
\end{proof}





\section{固有射}
\label{morphisms-section-proper}

\noindent
固有射の概念は代数幾何学において重要な役割を果たす。
固有射の重要な例は構造射 $\mathbf{P}^n_S \to S$
（射影 $n$-空間のもの）であり、実際、これが定義を導く動機となる例である。

\begin{definition}
\label{morphisms-definition-proper}
$f : X \to S$ をスキームの射とする。
$f$ が {\it 固有}であるとは、$f$ が分離的、有限型、かつ普遍閉であることをいう。
\end{definition}

\noindent
原点を二重化したアフィン直線からアフィン直線への射は有限型かつ
普遍閉であるから、上の定義では分離性の条件が必要である。
本節の残りでは、固有射および普遍閉射の基本的性質のいくつかを証明する。

\begin{lemma}
\label{morphisms-lemma-universally-closed-local-on-the-base}
$f : X \to S$ をスキームの射とする。
次は同値である。
\begin{enumerate}
\item 射 $f$ は普遍閉である。
\item 開被覆 $S = \bigcup V_j$ で、
$f^{-1}(V_j) \to V_j$ が普遍閉であることがすべての添字 $j$ に
対して成り立つものが存在する。
\end{enumerate}
\end{lemma}

\begin{proof}
定義から明らかである。
\end{proof}

\begin{lemma}
\label{morphisms-lemma-proper-local-on-the-base}
$f : X \to S$ をスキームの射とする。
次は同値である。
\begin{enumerate}
\item 射 $f$ は固有である。
\item 開被覆 $S = \bigcup V_j$ で、
$f^{-1}(V_j) \to V_j$ が固有であることがすべての添字 $j$ に
対して成り立つものが存在する。
\end{enumerate}
\end{lemma}

\begin{proof}
省略する。
\end{proof}

\begin{lemma}
\label{morphisms-lemma-composition-proper}
固有射の合成は固有である。
同じことが普遍閉射についても成り立つ。
\end{lemma}

\begin{proof}
閉射の合成は閉である。
$X \to Y \to Z$ が普遍閉射であり、
$Z' \to Z$ が任意の射ならば、
$Z' \times_Z X = (Z' \times_Z Y) \times_Y X  \to Z' \times_Z Y$
は閉であり、$Z' \times_Z Y \to Z'$ も閉である。
従って普遍閉射についての結果を得る。
「分離的」と「有限型」が合成で保たれることはすでに見た
（スキーム、補題 \ref{schemes-lemma-separated-permanence} および補題
\ref{morphisms-lemma-composition-finite-type}）。従って固有射についての結果も得る。
\end{proof}

\begin{lemma}
\label{morphisms-lemma-base-change-proper}
固有射の基底変換は固有である。
同じことが普遍閉射についても成り立つ。
\end{lemma}

\begin{proof}
普遍閉射については定義から正しい。
分離的射についても正しい
（スキーム、補題 \ref{schemes-lemma-separated-permanence}）。
有限型射についても正しい
（補題 \ref{morphisms-lemma-base-change-finite-type}）。
従って固有射についても正しい。
\end{proof}

\begin{lemma}
\label{morphisms-lemma-closed-immersion-proper}
閉埋め込みは固有であり、従って特に普遍閉である。
\end{lemma}

\begin{proof}
閉埋め込みの基底変換は閉埋め込みである
（スキーム、補題 \ref{schemes-lemma-base-change-immersion}）。
従ってそれは普遍閉である。
閉埋め込みは分離的である
（スキーム、補題 \ref{schemes-lemma-immersions-monomorphisms}）。
閉埋め込みは有限型である
（補題 \ref{morphisms-lemma-immersion-locally-finite-type}）。
従って閉埋め込みは固有である。
\end{proof}

\begin{lemma}
\label{morphisms-lemma-image-proper-scheme-closed}
スキームの可換図式
$$
\xymatrix{
X \ar[rr] \ar[rd] & &
Y \ar[ld] \\
& S &
}
$$
が与えられ、$Y$ は $S$ 上分離的であると仮定する。
\begin{enumerate}
\item $X \to S$ が普遍閉ならば、射 $X \to Y$ は普遍閉である。
\item $X$ が $S$ 上固有ならば、射 $X \to Y$ は固有である。
\end{enumerate}
特に、どちらの場合にも $X$ の $Y$ における像は閉である。
\end{lemma}

\begin{proof}
$X \to S$ が普遍閉（resp.\ 固有）であると仮定する。
射を $X \to X \times_S Y \to Y$ と分解する。
第一の射は閉埋め込みである。スキーム、補題
\ref{schemes-lemma-semi-diagonal} 参照。
従って第一の射は固有である（補題
\ref{morphisms-lemma-closed-immersion-proper}）。
射影 $X \times_S Y \to Y$ は普遍閉（resp.\ 固有）射の基底変換であり、
従って普遍閉（resp.\ 固有）である。補題
\ref{morphisms-lemma-base-change-proper} 参照。
従って $X \to Y$ は、普遍閉（resp.\ 固有）射の合成として
普遍閉（resp.\ 固有）である
（補題 \ref{morphisms-lemma-composition-proper}）。
\end{proof}

\noindent
次の補題の証明は Bjorn Poonen による。 
\href{https://mathoverflow.net/questions/23337/is-a-universally-closed-morphism-of-schemes-quasi-compact/23528#23528}{この箇所}
を参照せよ。

\begin{lemma}
\label{morphisms-lemma-universally-closed-quasi-compact}
\begin{reference}
Bjorn Poonen による。
\end{reference}
スキームの普遍閉射は準コンパクトである。
\end{lemma}

\begin{proof}
$f : X \to S$ を射とする。$f$ が準コンパクトでないと仮定する。
目標は $f$ が普遍閉でないことを示すことである。スキーム、補題
\ref{schemes-lemma-quasi-compact-affine} により、アフィン開集合
$V \subset S$ で $f^{-1}(V)$ が準コンパクトでないものが存在する。
目標を達成するには $f^{-1}(V) \to V$ が普遍閉でないことを
示せば十分なので、$S = \Spec(A)$ がある環 $A$ に対して成り立つと
仮定してよい。

\medskip\noindent
$X = \bigcup_{i \in I} X_i$ と書く。ここで $X_i$ は $X$ の
アフィン開部分スキームである。$T = \Spec(A[y_i ; i \in I])$ とおく。
$T_i = D(y_i) \subset T$ とおく。$Z$ を閉集合
$(X \times_S T) - \bigcup_{i \in I} (X_i \times_S T_i)$ とする。
像 $f_T(Z)$（$Z$ の $f_T : X \times_S T \to T$ によるもの）が
閉でないことを証明すれば十分である。

\medskip\noindent
ある点 $s \in S$ で、近傍 $U$（$s$ の $S$ におけるもの）で
$X_U$ が準コンパクトとなるものが存在しないものがある。
そうでなければ、$S$ を有限個のそのような $U$ で覆うことができ、
スキーム、補題 \ref{schemes-lemma-quasi-compact-affine} により
$f$ は準コンパクトとなる。そのような $s \in S$ を固定する。

\medskip\noindent
まず $f_T(Z_s) \ne T_s$ を確認する。$t \in T$ を、$s$ 上にあり、
$\kappa(t) = \kappa(s)$ を満たし、$y_i = 1$ が
$\kappa(t)$ においてすべての $i$ に対して成り立つ点とする。
すると $t \in T_i$ がすべての $i$ に対して成り立ち、
$Z_s \to T_s$ の $t$ 上のファイバーは
$(X - \bigcup_{i \in I} X_i)_s$ と同型であるが、これは空である。
従って $t \in T_s - f_T(Z_s)$ である。

\medskip\noindent
$f_T(Z)$ が $T$ で閉であると仮定する。このとき元
$g \in A[y_i; i \in I]$ で
$f_T(Z) \subset V(g)$ かつ $t \not \in V(g)$ を満たすものが存在する。
従って $g$ の像は $\kappa(t)$ において零でない。特に、$g$ のある係数は
$\kappa(s)$ における像が零でない。従ってこの係数は
ある近傍 $U$（$s$ のもの）上で可逆である。$J$ を有限集合
$j \in I$ で $y_j$ が $g$ に現れるもの全体とする。
$X_U$ は準コンパクトでないので、点
$x \in X - \bigcup_{j \in J} X_j$ で、ある $u \in U$ の上に
あるものを選べる。$g$ は $U$ 上可逆な係数をもつので、点
$t' \in T$ で $u$ 上にあり、$t' \not \in V(g)$ かつ
$t' \in V(y_i)$ となることがすべての $i \notin J$ に対して
成り立つものを見つけられる。これは
$V(y_i; i \in I, i \not\in J) = \Spec(A[y_j; j\in J])$
であり、このスキームの $u$ 上の点の集合が
$\Spec(\kappa(u)[y_j; j \in J])$ と全単射だからである。
言い換えれば、$t' \notin T_i$ が各 $i \notin J$ に対して成り立つ。
スキーム、補題 \ref{schemes-lemma-points-fibre-product} により、
点 $z$（$X \times_S T$ のもの）で、$x \in X$ と $t' \in T$ に写るものを
見つけられる。$x \not \in X_j$ が $j \in J$ に対して成り立ち、
$t' \not \in T_i$ が $i \in I \setminus J$ に対して成り立つので
$z \in Z$ である。一方、
$f_T(z) = t' \not \in V(g)$ であり、
これは $f_T(Z) \subset V(g)$ に矛盾する。
従って「$f_T(Z)$ は閉である」という仮定は誤りであり、
望みどおり $f_T$ は閉でないと結論する。
\end{proof}

\noindent
次の補題は、固有スキームの像（基底上有限型の分離的スキームにおけるもの）
が固有であることを述べる。

\begin{lemma}
\label{morphisms-lemma-image-proper-is-proper}
$S$ をスキームとする。$f : X \to Y$ を $S$ 上のスキームの射とする。
$X$ が $S$ 上普遍閉であり、$f$ が全射ならば、$Y$ は $S$ 上普遍閉である。
特に、さらに $Y$ が $S$ 上分離的かつ局所有限型ならば、
$Y$ は $S$ 上固有である。
\end{lemma}

\begin{proof}
$X$ が普遍閉であり、$f$ が全射であると仮定する。
構造射を $p : X \to S$、$q : Y \to S$ と記す。
$S' \to S$ をスキームの射とする。基底変換
$f' : X_{S'} \to Y_{S'}$ は全射であり
（補題 \ref{morphisms-lemma-base-change-surjective}）、基底変換
$p' : X_{S'} \to S'$ は閉である。
$T \subset Y_{S'}$ が閉ならば、
$(f')^{-1}(T) \subset X_{S'}$ は閉であり、従って
$p'((f')^{-1}(T)) = q'(T)$ は閉である。
ゆえに $q'$ は閉である。これで第一の主張が証明された。
従って補題 \ref{morphisms-lemma-universally-closed-quasi-compact} により
$Y \to S$ は準コンパクトであり、それゆえさらに
$Y \to S$ が局所有限型かつ分離的ならば、定義により
$Y \to S$ は固有である。
\end{proof}

\begin{lemma}
\label{morphisms-lemma-scheme-theoretic-image-is-proper}
スキームの可換図式
$$
\xymatrix{
X \ar[rr]_h \ar[rd]_f & & Y \ar[ld]^g \\
& S
}
$$
が与えられたとする。次を仮定する。
\begin{enumerate}
\item $X \to S$ は普遍閉（例えば固有）射である。
\item $Y \to S$ は分離的かつ局所有限型である。
\end{enumerate}
このとき、スキーム論的像 $Z \subset Y$（$h$ のもの）は $S$ 上固有であり、
$X \to Z$ は全射である。
\end{lemma}

\begin{proof}
$h$ のスキーム論的像は節
\ref{morphisms-section-scheme-theoretic-image} で構成した。
$f$ は準コンパクトなので（補題
\ref{morphisms-lemma-universally-closed-quasi-compact}）、
$h$ は準コンパクトである（スキーム、補題
\ref{schemes-lemma-quasi-compact-permanence}）。
従って $h(X) \subset Z$ は稠密である
（補題 \ref{morphisms-lemma-quasi-compact-scheme-theoretic-image}）。
一方、$h(X)$ は $Y$ で閉である
（補題 \ref{morphisms-lemma-image-proper-scheme-closed}）。
従って $X \to Z$ は全射である。
ゆえに $Z \to S$ は固有である
（補題 \ref{morphisms-lemma-image-proper-is-proper}）。
\end{proof}

\noindent
全射な普遍閉射による分離的スキームの像は分離的である。

\begin{lemma}
\label{morphisms-lemma-image-universally-closed-separated}
$S$ をスキームとする。$f : X \to Y$ を $S$ 上のスキームの
全射な普遍閉射とする。
\begin{enumerate}
\item $X$ が準分離的ならば、$Y$ は準分離的である。
\item $X$ が分離的ならば、$Y$ は分離的である。
\item $X$ が $S$ 上準分離的ならば、$Y$ は $S$ 上準分離的である。
\item $X$ が $S$ 上分離的ならば、$Y$ は $S$ 上分離的である。
\end{enumerate}
\end{lemma}

\begin{proof}
(1) と (2) は、$S = \Spec(\mathbf{Z})$ に対する (3) と (4) の帰結である
（スキーム、定義 \ref{schemes-definition-separated} 参照）。
可換図式
$$
\xymatrix{
X \ar[d] \ar[rr]_{\Delta_{X/S}} & & X \times_S X \ar[d] \\
Y \ar[rr]^{\Delta_{Y/S}} & & Y \times_S Y
}
$$
を考える。左の垂直矢印は全射（すなわち普遍全射）である。
右の垂直矢印は普遍閉射
$X \times_S X \to X \times_S Y \to Y \times_S Y$ の合成として
普遍閉である。従ってそれは準コンパクトでもある。補題
\ref{morphisms-lemma-universally-closed-quasi-compact} 参照。

\medskip\noindent
$X$ が $S$ 上準分離的、すなわち $\Delta_{X/S}$ が
準コンパクトであると仮定する。
$V \subset Y \times_S Y$ が準コンパクト開集合ならば、
$V \times_{Y \times_S Y} X \to \Delta_{Y/S}^{-1}(V)$ は全射であり、
$V \times_{Y \times_S Y} X$ は上の注意により準コンパクトである。
従って $\Delta_{Y/S}$ は準コンパクト、すなわち $Y$ は $S$ 上
準分離的である。

\medskip\noindent
$X$ が $S$ 上分離的、すなわち $\Delta_{X/S}$ が閉埋め込みであると
仮定する。このとき $X \to Y \times_S Y$ は閉射の合成として閉である。
$X \to Y$ は全射なので、$\Delta_{Y/S}(Y)$ は $Y \times_S Y$ で
閉であることが従う。従ってスキーム、定義
\ref{schemes-definition-separated} に続く議論により
$Y$ は $S$ 上分離的である。
\end{proof}







\section{付値判定法}
\label{morphisms-section-valuative-criteria}

\noindent
普遍閉性および分離性の付値判定法については、スキームの章、節
\ref{schemes-section-valuative-criterion-universal-closedness} および
\ref{schemes-section-valuative-separatedness} ですでに論じた。
本節では、そのいくつかの帰結と変形を論じる。
極限の章、節 \ref{limits-section-Noetherian-valuative-criterion} では、
局所 Noether スキームおよび有限型射を扱う場合には離散付値環だけを
考えれば十分であることを示す。

\begin{lemma}[固有性の付値判定法]
\label{morphisms-lemma-characterize-proper}
\begin{reference}
\cite[II Theorem 7.3.8]{EGA}
\end{reference}
$S$ をスキームとする。$f : X \to Y$ を $S$ 上のスキームの射とする。
$f$ は有限型かつ準分離的であると仮定する。このとき、次は同値である。
\begin{enumerate}
\item $f$ は固有である。
\item $f$ は付値判定法
（スキームの章、定義 \ref{schemes-definition-valuative-criterion}）を満たす。
\item 任意の実線部分からなる可換図式
$$
\xymatrix{
\Spec(K) \ar[r] \ar[d] & X \ar[d] \\
\Spec(A) \ar[r] \ar@{-->}[ru] & Y
}
$$
で、$A$ が分数体 $K$ をもつ付値環であるものが与えられると、
図式を可換にする点線矢印が一意に存在する。
\end{enumerate}
\end{lemma}

\begin{proof}
(3) は (2) の言い換えである。したがって、この補題は
スキームの章、命題 \ref{schemes-proposition-characterize-universally-closed}
および補題 \ref{schemes-lemma-separated-implies-valuative}、
補題 \ref{schemes-lemma-valuative-criterion-separatedness} と
諸定義から形式的に従う。
\end{proof}

\noindent
付値判定法を検証する際、通常は可能な図式をすべて考える必要はない。
次の補題に現れるような付値判定法を「精密化された付値判定法」と呼ぶことにする。

\begin{lemma}
\label{morphisms-lemma-refined-valuative-criterion-universally-closed}
$f : X \to S$ および $h : U \to X$ をスキームの射とする。
$f$ と $h$ は準コンパクトであり、$h(U)$ は $X$ で稠密であると仮定する。
任意の実線部分からなる可換図式
$$
\xymatrix{
\Spec(K) \ar[r] \ar[d] & U \ar[r]^h & X \ar[d]^f \\
\Spec(A) \ar[rr] \ar@{-->}[rru] & & S
}
$$
で、$A$ が分数体 $K$ をもつ付値環であるものが与えられたとき、
図式を可換にする点線矢印が一意に存在するならば、$f$ は普遍閉である。
さらに $f$ が準分離的ならば、$f$ は分離的である。
\end{lemma}

\begin{proof}
$f$ が普遍閉であることを示すため、$f$ に対する付値判定法の存在部分を
検証する。スキームの章、命題
\ref{schemes-proposition-characterize-universally-closed} により、
これは十分である。そのため、可換図式
$$
\xymatrix{
\Spec(K) \ar[r] \ar[d] & X \ar[d] \\
\Spec(A) \ar[r] & S
}
$$
を考える。ここで $A$ は付値環で、$K$ は $A$ の分数体である。
付値環と体は被約なので、スキームの章、補題
\ref{schemes-lemma-map-into-reduction} により、$U$、$X$、$S$ を
それぞれの被約化で置き換えてよい。この場合、$h(U)$ が稠密であるという
仮定は、$h : U \to X$ のスキーム論的像が $X$ であることを意味する。
補題 \ref{morphisms-lemma-scheme-theoretic-image-reduced} を参照せよ。
また、$S$ を、射 $\Spec(A) \to S$ が経由するアフィン開集合で
置き換えてよい。したがって $S = \Spec(R)$ と仮定してよい。

\medskip\noindent
$\Spec(B) \subset X$ を、射 $\Spec(K) \to X$ が経由する
アフィン開集合とする。多項式代数 $P$ を $B$ 上に選び、
$B$-代数の全射 $P \to K$ を選ぶ。このとき $\Spec(P) \to X$ は
平坦である。したがって補題
\ref{morphisms-lemma-flat-base-change-scheme-theoretic-image} により、射
$U \times_X \Spec(P) \to \Spec(P)$ のスキーム論的像は $\Spec(P)$ である。
補題 \ref{morphisms-lemma-reach-points-scheme-theoretic-image} により、可換図式
$$
\xymatrix{
\Spec(K') \ar[r] \ar[d] & U \times_X \Spec(P) \ar[d] \\
\Spec(A') \ar[r] & \Spec(P)
}
$$
を見いだせる。ここで $A'$ は付値環、$K'$ は $A'$ の分数体であり、
$\Spec(A')$ の閉点は $\Spec(K) \subset \Spec(P)$ に写る。
言い換えると、$B$-代数写像
$\varphi : K \to A'/\mathfrak m_{A'}$ が存在する。
付値環 $A'' \subset A'/\mathfrak m_{A'}$ で、
$\varphi(A)$ を支配し、分数体が
$K'' = A'/\mathfrak m_{A'}$ であるものを選ぶ
（代数の章、補題 \ref{algebra-lemma-dominate}）。次のようにおく。
$$
C = \{\lambda \in A' \mid \lambda \bmod \mathfrak m_{A'} \in A''\}.
$$
これは代数の章、補題
\ref{algebra-lemma-stack-valuation-rings} により付値環である。
$C$ は $R$-代数で、分数体が $K'$ なので、補題の主張に現れる形の可換図式
$$
\xymatrix{
\Spec(K') \ar[r] \ar[d] & U \ar[r] & X \ar[d] \\
\Spec(C) \ar[rr] \ar@{-->}[rru] & & S
}
$$
を得る。したがって、図示されたとおりに図式へ入る点線矢印が存在する。
補題の一意性の仮定により、合成
$\Spec(A') \to \Spec(C) \to X$ は、与えられた射
$\Spec(A') \to \Spec(P) \to \Spec(B) \subset X$ と一致する。
したがって、この射を $C/\mathfrak m_{A'} = A''$ のスペクトルに
制限すると、与えられた射
$\Spec(K'') = \Spec(A'/\mathfrak m_{A'}) \to \Spec(K) \to X$
が誘導される。$x \in X$ を、$\Spec(A'') \to X$ の閉点の像とする。
誘導された環写像 $\mathcal{O}_{X, x} \to A''$ の像は、
$K \subset K''$ に含まれる局所部分環である。
$A$ は $K$ の中で支配関係に関して極大であり、
$A \subset A''$ なので、$A = K \cap A''$ である。
したがって $\mathcal{O}_{X, x} \to A''$ は $A \subset A''$ を経由する。
これにより、求める矢印 $\Spec(A) \to X$ を得る。

\medskip\noindent
最後に、$f$ は準分離的であると仮定する。このとき
$\Delta : X \to X \times_S X$ は準コンパクトである。
実線部分からなる図式
$$
\xymatrix{
\Spec(K) \ar[r] \ar[d] & U \ar[r]^h & X \ar[d]^\Delta \\
\Spec(A) \ar[rr] \ar@{-->}[rru] & & X \times_S X
}
$$
で、$A$ が分数体 $K$ をもつ付値環であるものが与えられると、
図式を可換にする点線矢印が一意に存在する。実際、下側の水平矢印は、
補題の図式で点線矢印となりうる一対の射 $\Spec(A) \to X$ と同じデータである。
したがって、要求された一意性から、下側の水平矢印が $\Delta$ を
経由することがわかる。よって、直前に証明した結果を
$\Delta : X \to X \times_S X$ および $h : U \to X$ に適用して、
$\Delta$ が普遍閉であると結論できる。これは明らかに $f$ が分離的であることを
意味する。
\end{proof}

\begin{remark}
\label{morphisms-remark-check-val-on-open}
補題 \ref{morphisms-lemma-refined-valuative-criterion-universally-closed} における
点線矢印の一意性の仮定は必要である（詳細は省略する）。もちろん、
$f$ が分離的ならば一意性は保証される
（スキームの章、補題 \ref{schemes-lemma-separated-implies-valuative}）。
\end{remark}

\begin{lemma}
\label{morphisms-lemma-morphism-defined-local-ring}
$S$ をスキームとする。$X$、$Y$ を $S$ 上のスキームとする。
$s \in S$ とし、$x \in X$、$y \in Y$ を $s$ 上の点とする。
\begin{enumerate}
\item $f, g : X \to Y$ を $S$ 上の射で、
$f(x) = g(x) = y$ かつ
$f^\sharp_x = g^\sharp_x : \mathcal{O}_{Y, y} \to \mathcal{O}_{X, x}$
を満たすものとする。次の場合には、ある開近傍 $U \subset X$ 上で
$f|_U = g|_U$ となる。
\begin{enumerate}
\item $Y$ は $S$ 上局所有限型である。
\item $X$ は整である。
\item $X$ は局所 Noether である。または、
\item $X$ は被約で、既約成分を有限個しかもたない。
\end{enumerate}
\item $\varphi : \mathcal{O}_{Y, y} \to \mathcal{O}_{X, x}$ を
局所 $\mathcal{O}_{S, s}$-代数写像とする。次の場合には、ある開近傍
$U \subset X$（$x$ を含むもの）と射 $f : U \to Y$ が存在し、
後者は $x$ を $y$ に写して $f^\sharp_x = \varphi$ を満たす。
\begin{enumerate}
\item $Y$ は $S$ 上局所有限表示である。
\item $Y$ は局所有限型で、$X$ は整である。
\item $Y$ は局所有限型で、$X$ は局所 Noether である。または、
\item $Y$ は局所有限型で、$X$ は被約かつ既約成分を有限個しかもたない。
\end{enumerate}
\end{enumerate}
\end{lemma}

\begin{proof}
(1) の証明。$X$、$Y$、$S$ を $x$、$y$、$s$ の適切な
アフィン開近傍で置き換えることにより、次の代数の問題に帰着できる。
環 $R$ と二つの $R$-代数写像 $\varphi, \psi : B \to A$ が与えられ、
\begin{enumerate}
\item $R \to B$ が有限型であるか、$A$ が整域であるか、$A$ が
Noether であるか、または $A$ が被約で極小素イデアルを有限個しかもたず、
\item 二つの写像 $B \to A_\mathfrak p$ が、ある素イデアル
$\mathfrak p \subset A$ に対して一致する
\end{enumerate}
と仮定する。このとき、$\varphi, \psi$ が同じ写像 $B \to A_g$ を
定めることを示せばよい。ただし $g \in A$、$g \not \in \mathfrak p$ は
適切に選ぶ。
$R \to B$ が有限型ならば、$t_1, \ldots, t_m \in B$ を
$B$ の $R$-代数としての生成元とする。各 $j$ に対して
$g_j \in A$、$g_j \not \in \mathfrak p$ で、
$\varphi(t_j)$ と $\psi(t_j)$ の $A_{g_j}$ における像が一致するものを
見つけられる。このとき $g = \prod g_j$ とおく。
他の場合（$A$ が整域、Noether、または被約で極小素イデアルを有限個しか
もたない場合）には、$g \in A$、$g \not \in \mathfrak p$ で
$A_g \subset A_\mathfrak p$ を満たすものを見つけられる。
代数の章、補題 \ref{algebra-lemma-subring-of-local-ring} を参照せよ。
したがって望みどおり、写像 $B \to A_g$ は一致する。

\medskip\noindent
(2) の証明。このため、$X$、$Y$、$S$ を適切なアフィン開集合で
置き換えてよい。$X = \Spec(A)$、$Y = \Spec(B)$、
$S = \Spec(R)$ とする。$\mathfrak p \subset A$ を $x$ に対応する
素イデアルとする。$\mathfrak q \subset B$ を $y$ に対応する
素イデアルとする。このとき $\varphi$ は局所 $R$-代数写像
$\varphi : B_\mathfrak q \to A_\mathfrak p$ である。
$R \to B$ が有限表示の環写像ならば、ある
$g \in A \setminus \mathfrak p$ と $R$-代数写像 $B \to A_g$ が存在し、
$$
\xymatrix{
B_\mathfrak q \ar[r]_\varphi & A_\mathfrak p \\
B \ar[u] \ar[r] & A_g \ar[u]
}
$$
は可換となる。代数の章、補題
\ref{algebra-lemma-characterize-finite-presentation} および
\ref{algebra-lemma-localization-colimit} を参照せよ。
誘導される射 $\Spec(A_g) \to \Spec(B)$ が求めるものである。
$B$ が $R$ 上有限型ならば、$t_1, \ldots, t_m \in B$ を
$B$ の $R$-代数としての生成元とする。このとき
$g_j \in A$、$g_j \not \in \mathfrak p$ で、
$\varphi(t_j) \in \Im(A_{g_j} \to A_\mathfrak p)$ を満たすものを選べる。
したがって $A$ を $A[1/\prod g_j]$ で置き換えることにより、
$B$ が $A \to A_\mathfrak p$ の像に写ると仮定してよい。
$g \in A$、$g \not \in \mathfrak p$ で、$A_g \to A_\mathfrak p$ が
単射となるものを見つけられれば、求める $R$-代数写像 $B \to A_g$ を得る。
したがって、代数の章、補題
\ref{algebra-lemma-subring-of-local-ring} をもう一度適用すれば証明が完了する。
\end{proof}

\begin{lemma}
\label{morphisms-lemma-extend-across}
$S$ をスキームとする。$X$、$Y$ を $S$ 上のスキームとする。
$x \in X$ とする。$U \subset X$ を開集合とし、
$f : U \to Y$ を $S$ 上の射とする。次を仮定する。
\begin{enumerate}
\item $x$ は $U$ の閉包に属する。
\item $X$ は被約で既約成分を有限個しかもたないか、または
$X$ は Noether である。
\item $\mathcal{O}_{X, x}$ は付値環である。
\item $Y \to S$ は固有である。
\end{enumerate}
このとき、開集合 $U \subset U' \subset X$（$x$ を含むもの）と、
$S$-射 $f' : U' \to Y$ で $f$ を延長するものが存在する。
\end{lemma}

\begin{proof}
$X$ を $x$ の $X$ における開近傍で置き換えても差し支えない
（細部を一つ省略する）。性質の章、補題
\ref{properties-lemma-ring-affine-open-injective-local-ring} により、
$X$ はアフィンで、
$\Gamma(X, \mathcal{O}_X) \subset \mathcal{O}_{X, x}$ であると
仮定してよい。特に $X$ は整であり、一意な生成点 $\xi$ をもつ。
その剰余体は分数体 $K$、すなわち付値環
$\mathcal{O}_{X, x}$ の分数体である。
$x$ は $U$ の閉包に属するので $U$ は空でなく、したがって $U$ は
$\xi$ を含む。よって固有性の付値判定法
（補題 \ref{morphisms-lemma-characterize-proper}）により、可換図式
$$
\xymatrix{
\Spec(K) \ar[d]_\xi \ar[r] & \Spec(\mathcal{O}_{X, x}) \ar[d]_t \\
U \ar[r]^f & Y
}
$$
に入る射 $t : \Spec(\mathcal{O}_{X, x}) \to Y$ が存在する。
この図式の射はすべて $S$ 上のスキームの射である。
補題 \ref{morphisms-lemma-morphism-defined-local-ring} を
$y = t(x)$ および $\varphi = t^\sharp_x$ として適用すると、
開近傍 $V \subset X$（$x$ のもの）と射 $g : V \to Y$ を得る。
この射は $S$ 上で、$x$ を $y$ に写し、
$g^\sharp_x = t^\sharp_x$ を満たす。
$Y \to S$ は分離的なので、等化子 $E$（$f|_{U \cap V}$ と
$g|_{U \cap V}$ のもの）は $U \cap V$ の閉部分スキームである。
スキームの章、補題 \ref{schemes-lemma-where-are-they-equal} を参照せよ。
構成により $f$ と $g$ は同じ射 $\Spec(K) \to Y$ を定めるので、
$E$ は整スキーム $U \cap V$ の生成点を含む。
したがって $E = U \cap V$ であり、$f$ と $g$ は貼り合わさって
求める射 $U' = U \cup V \to Y$ を与える。
\end{proof}






\section{射影的射}
\label{morphisms-section-projective}

\noindent
射影的射については \cite{EGA} の定義を用いる。
「H」を伴う定義の版は \cite{H} による。
その結果得られる二つの定義は異なるが、どちらも有用である。

\begin{definition}
\label{morphisms-definition-projective}
$f : X \to S$ をスキームの射とする。
\begin{enumerate}
\item $f$ が {\it 射影的}であるとは、$X$ が $S$-スキームとして、
ある射影束 $\mathbf{P}(\mathcal{E})$ の閉部分スキームと同型であることをいう。
ここで $\mathcal{O}_S$-加群 $\mathcal{E}$ は準連接かつ有限型である。
\item $f$ が {\it H-射影的}であるとは、ある整数 $n$ と
閉埋め込み $X \to \mathbf{P}^n_S$ が $S$ 上に存在することをいう。
\item $f$ が {\it 局所射影的}であるとは、開被覆
$S = \bigcup U_i$ で、各 $f^{-1}(U_i) \to U_i$ が射影的となるものが
存在することをいう。
\end{enumerate}
\end{definition}

\noindent
予想されるとおり、射影的射は準射影的である。補題
\ref{morphisms-lemma-projective-quasi-projective} を参照せよ。
逆に、準射影的射はしばしば開埋め込みと射影的射の合成である。
補題 \ref{morphisms-lemma-quasi-projective-open-projective} を参照せよ。
準射影的な基底上の射影的射の性質の概観については、
射についてさらに、節 \ref{more-morphisms-section-projective} を参照せよ。

\begin{example}
\label{morphisms-example-projective}
$S$ をスキームとする。$\mathcal{A}$ を準連接な次数付き
$\mathcal{O}_S$-代数で、$\mathcal{A}_1$ によって $\mathcal{A}_0$ 上
生成されるものとする。さらに $\mathcal{A}_1$ は
$\mathcal{O}_S$ 上有限型であると仮定する。
$X = \underline{\text{Proj}}_S(\mathcal{A})$ とおく。
この場合、$X \to S$ は射影的である。実際、次数付き
$\mathcal{O}_S$-代数写像
$$
\text{Sym}_{\mathcal{O}_X}^*(\mathcal{A}_1)
\longrightarrow
\mathcal{A}
$$
に付随する射は閉埋め込みである。構成の章、補題
\ref{constructions-lemma-surjective-generated-degree-1-map-relative-proj}
を参照せよ。
\end{example}
 
\begin{lemma}
\label{morphisms-lemma-H-projective}
H-射影的射は H-準射影的である。
H-射影的射は射影的である。
\end{lemma}

\begin{proof}
第一の主張は定義から直ちに従う。
第二の主張は、$\mathbf{P}^n_S$ が $S$ 上の射影束であることから従う。
構成の章、補題 \ref{constructions-lemma-projective-space-bundle} を参照せよ。
\end{proof}

\begin{lemma}
\label{morphisms-lemma-characterize-locally-projective}
$f : X \to S$ をスキームの射とする。次は同値である。
\begin{enumerate}
\item 射 $f$ は局所射影的である。
\item 開被覆 $S = \bigcup U_i$ で、各
$f^{-1}(U_i) \to U_i$ が H-射影的となるものが存在する。
\end{enumerate}
\end{lemma}

\begin{proof}
補題 \ref{morphisms-lemma-H-projective} により (2) は (1) を含意する。(1) を仮定する。
各点 $s \in S$ に対して、アフィン開集合
$\Spec(R) = U \subset S$ で、$s$ の近傍となり、$X_U$ が
射影束 $\mathbf{P}(\mathcal{E})$ の閉部分スキームと同型となるものを
見つけられる。ここで $\mathcal{O}_U$-加群 $\mathcal{E}$ は
有限型かつ準連接である。
$\mathcal{E} = \widetilde{M}$ と書く。ここで $R$-加群 $M$ は有限型である
（性質の章、補題 \ref{properties-lemma-finite-type-module} 参照）。
$x_0, \ldots, x_n \in M$ を $M$ の $R$-加群としての生成元に選ぶ。
次数付き $R$-代数の全射
$$
R[X_0, \ldots, X_n] \longrightarrow \text{Sym}_R(M).
$$
を考える。構成の章、補題
\ref{constructions-lemma-surjective-graded-rings-map-proj} によれば、
対応する射
$$
\mathbf{P}(\mathcal{E}) \to \mathbf{P}^n_R
$$
は閉埋め込みである。したがって、$f^{-1}(U)$ は
$\mathbf{P}^n_U$ の閉部分スキームと同型である
（$U$ 上のスキームとして）。言い換えると、(2) が成り立つ。
\end{proof}

\begin{lemma}
\label{morphisms-lemma-locally-projective-proper}
局所射影的射は固有である。
\end{lemma}

\begin{proof}
$f : X \to S$ を局所射影的とする。$f$ が固有であることを示すには
基底上局所的に作業してよい。補題
\ref{morphisms-lemma-proper-local-on-the-base} を参照せよ。
したがって、上の補題 \ref{morphisms-lemma-characterize-locally-projective} により、
閉埋め込み $X \to \mathbf{P}^n_S$ が存在すると仮定してよい。
補題 \ref{morphisms-lemma-composition-proper} および
\ref{morphisms-lemma-closed-immersion-proper} により、
$\mathbf{P}^n_S \to S$ が固有であることを示せば十分である。
$\mathbf{P}^n_S \to S$ は
$\mathbf{P}^n_{\mathbf{Z}} \to \Spec(\mathbf{Z})$ の基底変換なので、
補題 \ref{morphisms-lemma-base-change-proper} により、
$\mathbf{P}^n_{\mathbf{Z}} \to \Spec(\mathbf{Z})$ が固有であることを
示せば十分である。構成の章、補題
\ref{constructions-lemma-proj-separated} により、スキーム
$\mathbf{P}^n_{\mathbf{Z}}$ は分離的である。
構成の章、補題 \ref{constructions-lemma-proj-quasi-compact} により、
スキーム $\mathbf{P}^n_{\mathbf{Z}}$ は準コンパクトである。
$\mathbf{P}^n_{\mathbf{Z}} \to \Spec(\mathbf{Z})$ が局所有限型であることは
明らかである。実際、$\mathbf{P}^n_{\mathbf{Z}}$ はアフィン開集合
$D_{+}(X_i)$ で覆われ、それぞれは有限型 $\mathbf{Z}$-代数
$$
\mathbf{Z}[X_0/X_i, \ldots, X_n/X_i].
$$
のスペクトルである。最後に、
$\mathbf{P}^n_{\mathbf{Z}} \to \Spec(\mathbf{Z})$ が普遍閉であることを
示さなければならない。これは構成の章、補題
\ref{constructions-lemma-proj-valuative-criterion} と付値判定法
（スキームの章、命題
\ref{schemes-proposition-characterize-universally-closed} 参照）から従う。
\end{proof}

\begin{lemma}
\label{morphisms-lemma-proper-ample-locally-projective}
$f : X \to S$ をスキームの固有射とする。$f$-豊富な可逆層が
$X$ 上に存在するならば、$f$ は局所射影的である。
\end{lemma}

\begin{proof}
$f$-豊富な可逆層が存在するならば、$S$ 上局所的に浸入
$i : X \to \mathbf{P}^n_S$ を見つけられる。補題
\ref{morphisms-lemma-finite-type-over-affine-ample-very-ample} を参照せよ。
$X \to S$ は固有なので、射 $i$ は閉埋め込みである。
補題 \ref{morphisms-lemma-image-proper-scheme-closed} を参照せよ。
\end{proof}

\begin{lemma}
\label{morphisms-lemma-H-projective-composition}
H-射影的射の合成は H-射影的である。
\end{lemma}

\begin{proof}
$X \to Y$ および $Y \to Z$ が H-射影的であると仮定する。
このとき、閉埋め込み $X \to \mathbf{P}^n_Y$ が $Y$ 上に存在し、
閉埋め込み $Y \to \mathbf{P}^m_Z$ が $Z$ 上に存在する。
次の図式を考える。
$$
\xymatrix{
X \ar[r] \ar[d] &
\mathbf{P}^n_Y \ar[r] \ar[dl] &
\mathbf{P}^n_{\mathbf{P}^m_Z} \ar[dl] \ar@{=}[r] &
\mathbf{P}^n_Z \times_Z \mathbf{P}^m_Z \ar[r] &
\mathbf{P}^{nm + n + m}_Z \ar[ddllll] \\
Y \ar[r] \ar[d] & \mathbf{P}^m_Z \ar[dl] & \\
Z & &
}
$$
ここで、最上段の最も右の水平矢印は Segre 埋め込みである。
構成の章、補題 \ref{constructions-lemma-segre-embedding} を参照せよ。
図式は、望みどおり $X$ を $\mathbf{P}^{nm + n + m}_Z$ の
閉部分スキームと同一視する。
\end{proof}

\begin{lemma}
\label{morphisms-lemma-H-projective-base-change}
H-射影的射の基底変換は H-射影的である。
\end{lemma}

\begin{proof}
これは、スキーム上の射影空間の基底変換が射影空間であることと、
閉埋め込みの基底変換が閉埋め込みであることから従う。
スキームの章、補題 \ref{schemes-lemma-base-change-immersion} を参照せよ。
\end{proof}

\begin{lemma}
\label{morphisms-lemma-base-change-projective}
（局所）射影的射の基底変換は（局所）射影的である。
\end{lemma}

\begin{proof}
これは、スキーム上の射影束の基底変換が射影束であること、
有限型 $\mathcal{O}$-加群の引き戻しが有限型であること
（加群の章、補題 \ref{modules-lemma-pullback-finite-type}）、
および閉埋め込みの基底変換が閉埋め込みであることから従う。
スキームの章、補題 \ref{schemes-lemma-base-change-immersion} を参照せよ。
いくつかの詳細は省略する。
\end{proof}

\begin{lemma}
\label{morphisms-lemma-projective-quasi-projective}
射影的射は準射影的である。
\end{lemma}

\begin{proof}
$f : X \to S$ を射影的射とする。閉埋め込み
$i : X \to \mathbf{P}(\mathcal{E})$ を選ぶ。ここで $\mathcal{E}$ は
準連接有限型 $\mathcal{O}_S$-加群である。このとき
$\mathcal{L} = i^*\mathcal{O}_{\mathbf{P}(\mathcal{E})}(1)$ は
$f$-非常に豊富である。$f$ は固有なので
（補題 \ref{morphisms-lemma-locally-projective-proper}）、準コンパクトである。
したがって補題 \ref{morphisms-lemma-ample-very-ample} により、
$\mathcal{L}$ は $f$-豊富である。$f$ は固有なので有限型である。
これで準射影的射の定義に含まれるすべての性質が成り立つことを確認した。
\end{proof}

\begin{lemma}
\label{morphisms-lemma-H-quasi-projective-open-H-projective}
$f : X \to S$ を H-準射影的射とする。このとき $f$ は
$X \to X' \to S$ と分解する。ここで $X \to X'$ は開埋め込み、
$X' \to S$ は H-射影的である。
\end{lemma}

\begin{proof}
定義により、$f$ を準コンパクトな浸入
$i : X \to \mathbf{P}^n_S$ と、それに続く射影
$\mathbf{P}^n_S \to S$ に分解できる。
補題 \ref{morphisms-lemma-quasi-compact-immersion} により、閉部分スキーム
$X' \subset \mathbf{P}^n_S$ で、$i$ が開埋め込み
$X \to X'$ を経由するものが存在する。これで補題が従う。
\end{proof}

\begin{lemma}
\label{morphisms-lemma-quasi-projective-open-projective}
$f : X \to S$ を準射影的射とし、$S$ は準コンパクトかつ
準分離的であるとする。このとき $f$ は $X \to X' \to S$ と分解する。
ここで $X \to X'$ は開埋め込み、$X' \to S$ は射影的である。
\end{lemma}

\begin{proof}
$\mathcal{L}$ を $f$-豊富とする。$f$ は有限型で $S$ は
準コンパクトなので、$\mathcal{L}^{\otimes n}$ が
$f$-非常に豊富となる $n > 0$ が存在する。補題
\ref{morphisms-lemma-finite-type-ample-very-ample} を参照せよ。
$\mathcal{L}$ を $\mathcal{L}^{\otimes n}$ で置き換える。
$\mathcal{F} = f_*\mathcal{L}$ と書く。これは
$\mathcal{O}_S$-準連接加群である。スキームの章、補題
\ref{schemes-lemma-push-forward-quasi-coherent} を参照せよ
（準射影的射は準コンパクトかつ分離的である。補題
\ref{morphisms-lemma-quasi-projective-properties} を参照せよ）。
性質の章、補題
\ref{properties-lemma-directed-colimit-finite-presentation} により、
有向集合 $I$ と、有限型準連接 $\mathcal{O}_S$-加群
$\mathcal{E}_i$ の $I$ 上の系で、$\mathcal{F} = \colim \mathcal{E}_i$ となるものを
見つけられる。合成
$\psi_i : f^*\mathcal{E}_i \to f^*\mathcal{F} \to \mathcal{L}$
を考える。有限アフィン開被覆
$S = \bigcup_{j = 1, \ldots, m} V_j$ を選ぶ。各 $j$ に対して、
切断
$$
s_{j, 0}, \ldots, s_{j, n_j} \in
\Gamma(f^{-1}(V_j), \mathcal{L}) = f_*\mathcal{L}(V_j) = \mathcal{F}(V_j)
$$
で、$\mathcal{L}$ を $f^{-1}V_j$ 上生成し、浸入
$$
f^{-1}V_j \longrightarrow \mathbf{P}^{n_j}_{V_j},
$$
を定めるものを選べる。補題
\ref{morphisms-lemma-very-ample-finite-type-over-affine} を参照せよ。
$i$ を、切断 $e_{j, t} \in \mathcal{E}_i(V_j)$ で、
$s_{j, t}$ に $\mathcal{F}$ において写るものが、すべての
$j = 1, \ldots, m$ および $t = 1, \ldots, n_j$ に対して
存在するように選ぶ。
このとき写像 $\psi_i$ は全射である。実際、切断
$f^*e_{j, t}$ は切断 $s_{j, t}$ と同じ像をもち、後者は
$\mathcal{L}|_{f^{-1}V_j}$ を生成する。したがって、$S$ 上の射
$$
r_{\mathcal{L}, \psi_i} : X \longrightarrow \mathbf{P}(\mathcal{E}_i)
$$
で、$V_j$ 上では上で与えた浸入が
$$
f^{-1}V_j \to \mathbf{P}(\mathcal{E}_i)|_{V_j} \to \mathbf{P}^{n_j}_{V_j}
$$
と分解するものを得る。このことから
$r_{\mathcal{L}, \psi_i}|_{V_j}$ は浸入である。補題
\ref{morphisms-lemma-immersion-permanence} を参照せよ。
$S = \bigcup V_j$ なので、$r_{\mathcal{L}, \psi_i}$ は浸入である。
$r_{\mathcal{L}, \psi_i}$ は、$X \to S$ が準コンパクトで
$\mathbf{P}(\mathcal{E}_i) \to S$ が分離的なので
（スキームの章、補題 \ref{schemes-lemma-quasi-compact-permanence} 参照）、
準コンパクトであることに注意する。
補題 \ref{morphisms-lemma-quasi-compact-immersion} により、閉部分スキーム
$X' \subset \mathbf{P}(\mathcal{E}_i)$ で、$i$ が開埋め込み
$X \to X'$ を経由するものが存在する。このとき
$X' \to S$ は定義により射影的であり、結論を得る。
\end{proof}

\begin{lemma}
\label{morphisms-lemma-projective-is-quasi-projective-proper}
$S$ を準コンパクトかつ準分離的なスキームとする。
$f : X \to S$ をスキームの射とする。このとき
\begin{enumerate}
\item $f$ が射影的であることと、$f$ が準射影的かつ固有であることは
同値である。また、
\item $f$ が H-射影的であることと、$f$ が H-準射影的かつ固有であることは
同値である。
\end{enumerate}
\end{lemma}

\begin{proof}
$f$ が射影的ならば、補題 \ref{morphisms-lemma-projective-quasi-projective} により
$f$ は準射影的であり、補題 \ref{morphisms-lemma-locally-projective-proper} により
固有である。逆に、$X \to S$ が準射影的かつ固有ならば、
補題 \ref{morphisms-lemma-quasi-projective-open-projective} により、
開埋め込み $X \to X'$ で、$X' \to S$ が射影的となるものを選べる。
$X \to S$ は固有なので、$X$ は $X'$ で閉である
（補題 \ref{morphisms-lemma-image-proper-scheme-closed}）。
すなわち、$X \to X'$ は（開かつ）閉埋め込みである。
$X'$ は $S$ 上の射影束の閉部分スキームと同型なので
（定義 \ref{morphisms-definition-projective}）、$X$ についても同じことが成り立つ。
すなわち、$X \to S$ は射影的射である。これで (1) が証明された。
(2) の証明も同じであるが、補題 \ref{morphisms-lemma-H-projective} および
\ref{morphisms-lemma-H-quasi-projective-open-H-projective} を用いる。
\end{proof}

\begin{lemma}
\label{morphisms-lemma-composition-projective}
$f : X \to Y$ および $g : Y \to S$ をスキームの射とする。
$S$ が準コンパクトかつ準分離的で、$f$ と $g$ が射影的ならば、
$g \circ f$ は射影的である。
\end{lemma}

\begin{proof}
補題 \ref{morphisms-lemma-projective-quasi-projective} および
\ref{morphisms-lemma-locally-projective-proper} により、$f$ と $g$ は
準射影的かつ固有である。補題 \ref{morphisms-lemma-composition-proper} および
\ref{morphisms-lemma-composition-quasi-projective} により、$g \circ f$ は
固有かつ準射影的である。したがって補題
\ref{morphisms-lemma-projective-is-quasi-projective-proper} により、
$g \circ f$ は射影的である。
\end{proof}

\begin{lemma}
\label{morphisms-lemma-projective-permanence}
$g : Y \to S$ および $f : X \to Y$ をスキームの射とする。
$g \circ f$ が射影的で $g$ が分離的ならば、$f$ は射影的である。
\end{lemma}

\begin{proof}
閉埋め込み $X \to \mathbf{P}(\mathcal{E})$ を選ぶ。ここで
$\mathcal{E}$ は準連接有限型 $\mathcal{O}_S$-加群である。
このとき射 $X \to \mathbf{P}(\mathcal{E}) \times_S Y$ を得る。
この射は、合成
$$
X \to X \times_S Y \to \mathbf{P}(\mathcal{E}) \times_S Y
$$
として閉埋め込みである。第一の射はスキームの章、補題
\ref{schemes-lemma-semi-diagonal} により閉埋め込みであり
（ここで $g$ が分離的であることを使う）、第二の射は閉埋め込みの
基底変換として閉埋め込みである。最後に、ファイバー積
$\mathbf{P}(\mathcal{E}) \times_S Y$ は
$\mathbf{P}(g^*\mathcal{E})$ と同型であり、引き戻しは準連接有限型加群を保つ。
\end{proof}

\begin{lemma}
\label{morphisms-lemma-projective-over-quasi-projective-is-H-projective}
$S$ を豊富な可逆層をもつスキームとする。このとき
\begin{enumerate}
\item 任意の射影的射 $X \to S$ は H-射影的であり、
\item 任意の準射影的射 $X \to S$ は H-準射影的である。
\end{enumerate}
\end{lemma}

\begin{proof}
$S$ に関する仮定は、$S$ が準コンパクトかつ分離的であることを含意する。
性質の章、定義 \ref{properties-definition-ample} と補題
\ref{properties-lemma-ample-immersion-into-proj}、および構成の章、補題
\ref{constructions-lemma-proj-separated} を参照せよ。
したがって補題 \ref{morphisms-lemma-quasi-projective-open-projective} が適用でき、
(1) が (2) を含意することがわかる。
$\mathcal{E}$ を有限型準連接 $\mathcal{O}_S$-加群とする。
射影的射の定義により、
$\mathbf{P}(\mathcal{E}) \to S$ が H-射影的であることを示せば十分である。
$\mathcal{E}$ が有限個の大域切断で生成されるならば、対応する全射
$\mathcal{O}_S^{\oplus n} \to \mathcal{E}$ は閉埋め込み
$$
\mathbf{P}(\mathcal{E}) \longrightarrow
\mathbf{P}(\mathcal{O}_S^{\oplus n}) = \mathbf{P}^{n - 1}_S
$$
を誘導し、これは望むものである。一般の場合、$\mathcal{L}$ を
$S$ 上の豊富な可逆層とする。性質の章、命題
\ref{properties-proposition-characterize-ample} により、ある整数 $n$ に対して
$\mathcal{E} \otimes_{\mathcal{O}_S} \mathcal{L}^{\otimes n}$
は有限個の大域切断で生成される。構成の章、補題
\ref{constructions-lemma-twisting-and-proj} により
$\mathbf{P}(\mathcal{E}) =
\mathbf{P}(\mathcal{E} \otimes_{\mathcal{O}_S} \mathcal{L}^{\otimes n})$
なので、証明が完了する。
\end{proof}

\begin{lemma}
\label{morphisms-lemma-proper-ample-is-proj}
$f : X \to S$ を普遍閉射とする。
$\mathcal{L}$ を $f$-豊富な可逆 $\mathcal{O}_X$-加群とする。
このとき、補題 \ref{morphisms-lemma-characterize-relatively-ample} の標準射
$$
r : X
\longrightarrow
\underline{\text{Proj}}_S
\left(
\bigoplus\nolimits_{d \geq 0} f_*\mathcal{L}^{\otimes d}
\right)
$$
は同型である。
\end{lemma}

\begin{proof}
$f$-豊富な可逆加群の存在は $f$ が準コンパクトであることを
強制するので、$f$ は準コンパクトである。引用した補題により射 $r$ は
開埋め込みである。一方、補題
\ref{morphisms-lemma-image-proper-scheme-closed} により $r$ の像は閉である
（構成の章、補題 \ref{constructions-lemma-relative-proj-separated} により、
$r$ の標的は $S$ 上分離的である）。最後に、性質の章、補題
\ref{properties-lemma-ample-immersion-into-proj} により $r$ の像は稠密である
（ここでは、補題 \ref{morphisms-lemma-characterize-relatively-ample} の証明で、
射 $r$ が $S$ のアフィン開集合上で、性質の章、補題
\ref{properties-lemma-map-into-proj} の標準射によって与えられることが
示されたことも用いる）。したがって $r$ は全射な開埋め込み、
すなわち同型である。
\end{proof}

\begin{lemma}
\label{morphisms-lemma-proper-ample-delete-affine}
$f : X \to S$ を普遍閉射とする。$\mathcal{L}$ を
$f$-豊富な可逆 $\mathcal{O}_X$-加群とする。
$s \in \Gamma(X, \mathcal{L})$ とする。このとき $X_s \to S$ は
アフィン射である。
\end{lemma}

\begin{proof}
問題は $S$ 上局所的である（補題 \ref{morphisms-lemma-characterize-affine}）。
したがって $S$ はアフィンと仮定してよい。
補題 \ref{morphisms-lemma-proper-ample-is-proj} により
$X = \text{Proj}(A)$ と書ける。ここで $A$ は次数付き環で、
$s$ は $f \in A_1$ に対応し、
$X_s = D_+(f)$ である
（性質の章、補題 \ref{properties-lemma-map-into-proj}）。
これは $\text{Proj}(A)$ の構成により補題を証明する。
構成の章、節 \ref{constructions-section-proj} を参照せよ。
\end{proof}











\section{整射と有限射}
\label{morphisms-section-integral}

\noindent
環準同型 $R \to A$ は、$A$ の各元が $R$ に係数をもつモニックな
方程式を満たすとき {\it 整}であるということを思い出そう。
また、環準同型 $R \to A$ は、$A$ が $R$-加群として有限生成である
とき {\it 有限}であるという。代数、定義
\ref{algebra-definition-integral-ring-map} を参照せよ。

\begin{definition}
\label{morphisms-definition-integral}
$f : X \to S$ をスキームの射とする。
\begin{enumerate}
\item $f$ が {\it 整}であるとは、$f$ がアフィンであり、任意のアフィン開集合
$\Spec(R) = V \subset S$ に対し、その逆像が
$\Spec(A) = f^{-1}(V) \subset X$ であるとき、対応する環準同型
$R \to A$ が整であることをいう。
\item $f$ が {\it 有限}であるとは、$f$ がアフィンであり、任意のアフィン開集合
$\Spec(R) = V \subset S$ に対し、その逆像が
$\Spec(A) = f^{-1}(V) \subset X$ であるとき、対応する環準同型
$R \to A$ が有限であることをいう。
\end{enumerate}
\end{definition}

\noindent
整射および有限射が分離的かつ準コンパクトであることは明らかである。
また、有限射が有限型の射であることも明らかである。本節の補題の
ほとんどはまったく標準的である。ただし、本節末尾の興味深い補題
\ref{morphisms-lemma-integral-universally-closed} に注意されたい。

\begin{lemma}
\label{morphisms-lemma-integral-local}
$f : X \to S$ をスキームの射とする。次は同値である。
\begin{enumerate}
\item 射 $f$ は整である。
\item アフィン開被覆 $S = \bigcup U_i$ であって、各
$f^{-1}(U_i)$ がアフィンであり、
$\mathcal{O}_S(U_i) \to \mathcal{O}_X(f^{-1}(U_i))$ が整となるものが
存在する。
\item 開被覆 $S = \bigcup U_i$ であって、各
$f^{-1}(U_i) \to U_i$ が整となるものが存在する。
\end{enumerate}
さらに、$f$ が整ならば、任意の開部分スキーム
$U \subset S$ に対し、射 $f : f^{-1}(U) \to U$ は整である。
\end{lemma}

\begin{proof}
代数、補題 \ref{algebra-lemma-integral-local} を参照せよ。
いくつかの詳細は省略する。
\end{proof}

\begin{lemma}
\label{morphisms-lemma-finite-local}
$f : X \to S$ をスキームの射とする。次は同値である。
\begin{enumerate}
\item 射 $f$ は有限である。
\item アフィン開被覆 $S = \bigcup U_i$ であって、各
$f^{-1}(U_i)$ がアフィンであり、
$\mathcal{O}_S(U_i) \to \mathcal{O}_X(f^{-1}(U_i))$ が有限となるものが
存在する。
\item 開被覆 $S = \bigcup U_i$ であって、各
$f^{-1}(U_i) \to U_i$ が有限となるものが存在する。
\end{enumerate}
さらに、$f$ が有限ならば、任意の開部分スキーム
$U \subset S$ に対し、射 $f : f^{-1}(U) \to U$ は有限である。
\end{lemma}

\begin{proof}
代数、補題 \ref{algebra-lemma-integral-local} を参照せよ。
いくつかの詳細は省略する。
\end{proof}

\begin{lemma}
\label{morphisms-lemma-finite-integral}
有限射は整である。局所有限型の整射は有限である。
\end{lemma}

\begin{proof}
代数、補題 \ref{algebra-lemma-finite-is-integral} および補題
\ref{algebra-lemma-characterize-finite-in-terms-of-integral} を参照せよ。
\end{proof}

\begin{lemma}
\label{morphisms-lemma-composition-finite}
有限射の合成は有限である。整射についても同様である。
\end{lemma}

\begin{proof}
代数、補題 \ref{algebra-lemma-finite-transitive} および
\ref{algebra-lemma-integral-transitive} を参照せよ。
\end{proof}

\begin{lemma}
\label{morphisms-lemma-base-change-finite}
有限射の基底変換は有限である。整射についても同様である。
\end{lemma}

\begin{proof}
代数、補題 \ref{algebra-lemma-base-change-integral} を参照せよ。
\end{proof}

\begin{lemma}
\label{morphisms-lemma-integral-universally-closed}
\begin{slogan}
整 $=$ アフィン $+$ 普遍閉
\end{slogan}
$f : X \to S$ をスキームの射とする。次は同値である。
\begin{enumerate}
\item $f$ は整であり、
\item $f$ はアフィンかつ普遍閉である。
\end{enumerate}
\end{lemma}

\begin{proof}
(1) を仮定する。整射は定義によりアフィンである。整射の基底変換は
整であるから、(2) を示すには整射が閉であることを示せば十分である。
これは代数、補題 \ref{algebra-lemma-integral-going-up} および
\ref{algebra-lemma-going-up-closed} から従う。

\medskip\noindent
(2) を仮定する。$f$ は射 $f : \Spec(A) \to \Spec(R)$ であり、
環準同型 $R \to A$ から得られるとしてよい。$a$ を $A$ の元とする。
$a$ が $R$ 上整であること、すなわち核 $I$ にモニック多項式が
存在することを示さなければならない。ここで準同型 $R[x] \to A$ は
$x$ を $a$ に送る。環 $B = A[x]/(ax -1)$ を考え、$J$ を合成
$R[x]\to A[x] \to B$ の核とする。$f\in J$ ならば、
$q\in A[x]$ が存在して、$f = (ax-1)q$ が $A[x]$ において成り立つ。
したがって、$f = \sum_i f_ix^i$ かつ $q = \sum_iq_ix^i$ ならば、
すべての $i \geq 0$ に対して $f_i = aq_{i-1} - q_i$ である。
$n \geq \deg q + 1$ に対し、多項式
$$
\sum\nolimits_{i \geq 0} f_i x^{n - i} =
\sum\nolimits_{i \geq 0} (a q_{i - 1} - q_i) x^{n - i} =
(a - x) \sum\nolimits_{i \geq 0} q_i x^{n - i - 1}
$$
は明らかに $I$ に属する。$f_0 = 1$ ならばこの多項式はモニックでも
あるから、$J$ が定数項 $1$ の多項式を含むことを示す問題に帰着する。
これを、$\Spec(R[x]/(J + (x))$ が空であることを示して証明する。


\medskip\noindent
$f$ は普遍閉であるから、基底変換 $\Spec(A[x]) \to \Spec(R[x])$ は
閉である。したがって、閉部分集合
$\Spec(B) \subset \Spec(A[x])$ の像は閉部分集合
$\Spec(R[x]/J) \subset \Spec(R[x])$ である。
例 \ref{morphisms-example-scheme-theoretic-image} および補題
\ref{morphisms-lemma-quasi-compact-scheme-theoretic-image} を参照せよ。
特に $\Spec(B) \to \Spec(R[x]/J)$ は全射である。すべての正方形が
引き戻しである次の図式を考える。
$$
\xymatrix{
\Spec(B) \ar@{->>}[r]^g &
\Spec(R[x]/J)  \ar[r] &
\Spec(R[x])\\
\emptyset \ar[u] \ar[r] &
\Spec(R[x]/(J + (x)))\ar[u] \ar[r] &
\Spec(R) \ar[u]^0
}
$$
左下隅は空である。実際、これは $R\otimes_{R[x]} B$ のスペクトルであり、
ここで準同型 $R[x]\to B$ は $x$ を可逆元に送り、準同型
$R[x]\to R$ は $x$ を $0$ に送る。$g$ は全射なので、望んだとおり
$\Spec(R[x]/(J + (x)))$ は空である。
\end{proof}

\begin{lemma}
\label{morphisms-lemma-integral-fibres}
$f : X \to S$ を整射とする。このとき $X$ の各点はそのファイバー内で
閉である。
\end{lemma}

\begin{proof}
代数、補題 \ref{algebra-lemma-integral-no-inclusion} を参照せよ。
\end{proof}

\begin{lemma}
\label{morphisms-lemma-integral-dimension}
$f : X \to Y$ を整射とする。このとき $\dim(X) \leq \dim(Y)$ である。
$f$ が全射ならば $\dim(X) = \dim(Y)$ である。
\end{lemma}

\begin{proof}
$X$ と $Y$ の次元はアフィン開被覆の各要素の次元の上限であるから、
$Y$ と $X$ はアフィンであるとしてよい。不等式は代数、補題
\ref{algebra-lemma-integral-dim-up} から従う。等式はさらに代数、補題
\ref{algebra-lemma-dimension-going-up} および
\ref{algebra-lemma-integral-going-up} から従う。
\end{proof}

\begin{lemma}
\label{morphisms-lemma-finite-quasi-finite}
有限射は準有限である。
\end{lemma}

\begin{proof}
これは代数、補題 \ref{algebra-lemma-quasi-finite} および補題
\ref{morphisms-lemma-quasi-finite-locally-quasi-compact} から従う。
別の証明として、補題 \ref{morphisms-lemma-integral-fibres}（および有限射は
整であるという事実）によりファイバー内のすべての点は閉点である。
そこで補題 \ref{morphisms-lemma-quasi-finite-at-point-characterize} (3) を用いれば、
$f$ は $x$ において準有限であり、これはすべての $x \in X$ に対して
成り立つとわかる。
\end{proof}

\begin{lemma}
\label{morphisms-lemma-finite-proper}
$f : X \to S$ をスキームの射とする。次は同値である。
\begin{enumerate}
\item $f$ は有限であり、
\item $f$ はアフィンかつ固有である。
\end{enumerate}
\end{lemma}

\begin{proof}
これは補題 \ref{morphisms-lemma-integral-universally-closed}、有限射が整かつ分離的で
あるという事実、固有射が有限型かつ分離的かつ普遍閉な射にほかならない
という事実、および有限型の整射が有限であるという事実（補題
\ref{morphisms-lemma-finite-integral}）から形式的に従う。
\end{proof}

\begin{lemma}
\label{morphisms-lemma-closed-immersion-finite}
閉埋め込みは有限である（したがって、とりわけ整である）。
\end{lemma}

\begin{proof}
閉埋め込みはアフィンであり（補題
\ref{morphisms-lemma-closed-immersion-affine}）、全射環準同型は有限かつ整だからである。
\end{proof}

\begin{lemma}
\label{morphisms-lemma-finite-union-finite}
$X_i \to Y$, $i = 1, \ldots, n$ をスキームの有限射とする。このとき
$X_1 \amalg \ldots \amalg X_n \to Y$ も有限である。
\end{lemma}

\begin{proof}
$R \to A_i$, $i = 1, \ldots, n$ が有限環準同型ならば、
$R \to A_1 \times \ldots \times A_n$ も有限であるという代数的事実から従う。
\end{proof}

\begin{lemma}
\label{morphisms-lemma-finite-permanence}
$f : X \to Y$ および $g : Y \to Z$ を射とする。
\begin{enumerate}
\item $g \circ f$ が有限であり $g$ が分離的ならば、$f$ は有限である。
\item $g \circ f$ が整であり $g$ が分離的ならば、$f$ は整である。
\end{enumerate}
\end{lemma}

\begin{proof}
$g \circ f$ が有限（それぞれ整）であり、$g$ が分離的であると仮定する。
補題 \ref{morphisms-lemma-base-change-finite} により、基底変換
$X \times_Z Y \to Y$ は有限（それぞれ整）である。
射 $X \to X \times_Z Y$ は、$Y \to Z$ が分離的なので閉埋め込みである。
スキーム、補題 \ref{schemes-lemma-section-immersion} を参照せよ。
閉埋め込みは有限（それぞれ整）である。補題
\ref{morphisms-lemma-closed-immersion-finite} を参照せよ。
有限射（それぞれ整射）の合成は有限（それぞれ整）である。
補題 \ref{morphisms-lemma-composition-finite} を参照せよ。これでよい。
\end{proof}

\begin{lemma}
\label{morphisms-lemma-finite-monomorphism-closed}
$f : X \to Y$ をスキームの射とする。$f$ が有限かつモノ射ならば、
$f$ は閉埋め込みである。
\end{lemma}

\begin{proof}
これは代数、補題
\ref{algebra-lemma-finite-epimorphism-surjective} に帰着する。
\end{proof}

\begin{lemma}
\label{morphisms-lemma-finite-projective}
有限射は射影的である。
\end{lemma}

\begin{proof}
$f : X \to S$ を有限射とする。このとき $f_*\mathcal{O}_X$ は
準連接 $\mathcal{O}_S$-加群であり（補題
\ref{morphisms-lemma-affine-equivalence-algebras}）、有限型である
（有限射の定義および性質、補題
\ref{properties-lemma-finite-type-module} による）。証明を終えるため、
$S$ 上の閉埋め込み
$$
\sigma :
X
\longrightarrow
\mathbf{P}(f_*\mathcal{O}_X) =
\underline{\text{Proj}}_S(\text{Sym}^*_{\mathcal{O}_S}(f_*\mathcal{O}_X))
$$
が存在すると主張する。すなわち、$\sigma$ を（構成、補題
\ref{constructions-lemma-apply-relative} を通じて）全射
$$
f^*f_*\mathcal{O}_X \longrightarrow \mathcal{O}_X
$$
に対応する射とする。この全射は随伴射 $f^*f_* \to \text{id}$ から
得られる。このとき $\sigma$ は閉埋め込みである。
実際、$S$ 上アフィン局所的に $X = \Spec(A)$ および
$S = \Spec(R)$ と書ける。$X$ は $S$ 上有限なので、
$a_1, \ldots, a_n \in A$ を、$A$ を $R$-加群として生成するように
選べる。このとき $R$-代数準同型
$R[T_0, T_1, \ldots, T_n] \to \text{Sym}_R^*(A)$ で、$T_0$ を $1$ に、
$T_i$ を $a_i$ に送るもの（$1 \leq i \leq n$）は全射である。
したがって、構成、補題
\ref{constructions-lemma-surjective-graded-rings-map-proj} により、
$\mathbf{P}(f_*\mathcal{O}_X)$ は $\mathbf{P}^n_S$ の閉部分スキームである。
ゆえに、誘導される射 $X \to \mathbf{P}^n_S$ が閉埋め込みであることを
示せば十分である（例えば、補題
\ref{morphisms-lemma-closed-immersion-ideals} を参照せよ）。
読者は、$X \to \mathbf{P}^n_S$ の像が開集合
$D_+(T_0) \cong \mathbf{A}^n_S$ に含まれ、$X \to \mathbf{A}^n_S$ が、
全射 $R$-代数準同型
$R[x_1, \ldots, x_n] \to \text{Sym}_R^*(A)$ で $x_i$ を $a_i$ に送る
ものに対応することを確認できる。
したがって $X \to \mathbf{P}^n_S$ は（閉埋め込みと開埋め込みの合成として）
埋め込みである。$X$ は $S$ 上有限なので、この埋め込みの像は閉である
（ここでは補題 \ref{morphisms-lemma-finite-proper}、補題
\ref{morphisms-lemma-image-proper-scheme-closed}、および構成、補題
\ref{constructions-lemma-projective-space-separated} を用いる）。
スキーム、補題 \ref{schemes-lemma-immersion-when-closed} により結論を得る。
\end{proof}






\section{普遍同相射}
\label{morphisms-section-universal-homeomorphisms}

\noindent
次の定義は実のところ不要である。実際、普遍同相射とはまさに整で、
普遍的単射かつ全射な射にほかならない。補題
\ref{morphisms-lemma-universal-homeomorphism} を参照せよ。

\begin{definition}
\label{morphisms-definition-universal-homeomorphism}
スキームの射 $f : X \to Y$ は、基底変換
$f' : Y' \times_Y X \to Y'$ が任意の射 $Y' \to Y$ に対して同相写像となるとき
{\it 普遍同相射}と呼ばれる。
\end{definition}

\noindent
まず、当然必要となる補題を述べる。

\begin{lemma}
\label{morphisms-lemma-base-change-universal-homeomorphism}
スキームの普遍同相射をスキームの任意の射によって基底変換したものは
普遍同相射である。
\end{lemma}

\begin{proof}
定義から直ちに従う。
\end{proof}

\begin{lemma}
\label{morphisms-lemma-composition-universal-homeomorphism}
スキームの二つの普遍同相射の合成は普遍同相射である。
\end{lemma}

\begin{proof}
省略する。
\end{proof}

\noindent
次の簡単な補題が、普遍同相射を特徴づける鍵である。

\begin{lemma}
\label{morphisms-lemma-homeomorphism-affine}
$f : X \to Y$ をスキームの射とする。$f$ が $Y$ の閉部分集合への
同相写像ならば、$f$ はアフィンである。
\end{lemma}

\begin{proof}
$y \in Y$ を点とする。$y \not \in f(X)$ ならば、$y$ のアフィン近傍で
$f(X)$ と交わらないものが存在する。$y \in f(X)$ ならば、$x \in X$ を
$X$ の点のうち $y$ に写る一意なものとする。$y \in V$ をアフィン開近傍とする。
$U \subset X$ を $x$ のアフィン開近傍で、$V$ に写るものとする。
仮定により $f(U) \subset V \cap f(X)$ は誘導位相に関して開である。
$f$ に関するこの仮定から、
$h \in \Gamma(V, \mathcal{O}_Y)$ であって $y \in D(h)$ かつ
$D(h) \cap f(X) \subset f(U)$ となるものを選べる。
$h' \in \Gamma(U, \mathcal{O}_X)$ で $f^\sharp(h)$ の $U$ への制限を表す。
すると $D(h') \subset U$ は $f^{-1}(D(h))$ に等しい。換言すれば、
$Y$ の各点には逆像がアフィンとなる開近傍がある。したがって $f$ は
アフィンである。補題 \ref{morphisms-lemma-characterize-affine} を参照せよ。
\end{proof}

\begin{lemma}
\label{morphisms-lemma-universal-homeomorphism}
$f : X \to Y$ をスキームの射とする。次は同値である。
\begin{enumerate}
\item $f$ は普遍同相射であり、
\item $f$ は整、普遍的単射、かつ全射である。
\end{enumerate}
\end{lemma}

\begin{proof}
$f$ が普遍同相射であると仮定する。補題
\ref{morphisms-lemma-homeomorphism-affine} により $f$ はアフィンである。
$f$ は明らかに普遍閉なので、補題
\ref{morphisms-lemma-integral-universally-closed} により $f$ は整である。
さらに、$f$ が普遍的単射かつ全射であることも明らかである。

\medskip\noindent
$f$ が整、普遍的単射、かつ全射であると仮定する。補題
\ref{morphisms-lemma-integral-universally-closed} により $f$ は普遍閉である。
さらに普遍的全単射でもあるから（補題
\ref{morphisms-lemma-base-change-surjective} を参照）、普遍同相射である。
\end{proof}

\begin{lemma}
\label{morphisms-lemma-reduction-universal-homeomorphism}
$X$ をスキームとする。標準的な閉埋め込み $X_{red} \to X$
（スキーム、定義
\ref{schemes-definition-reduced-induced-scheme} を参照）は普遍同相射である。
\end{lemma}

\begin{proof}
省略する。
\end{proof}

\begin{lemma}
\label{morphisms-lemma-check-closed-infinitesimally}
$f : X \to S$ および $S' \to S$ をスキームの射とする。次を仮定する。
\begin{enumerate}
\item $S' \to S$ は閉埋め込みである。
\item $S' \to S$ は点の上で全単射である。
\item $X \times_S S' \to S'$ は閉埋め込みである。
\item $X \to S$ は有限型であるか、または $S' \to S$ は有限表示である。
\end{enumerate}
このとき $f : X \to S$ は閉埋め込みである。
\end{lemma}

\begin{proof}
仮定 (1) と (2) から $S' \to S$ は普遍同相射である
（例えば、$S_{red} = S'_{red}$ であることと補題
\ref{morphisms-lemma-reduction-universal-homeomorphism} を用いればよい）。
従って (3) から $X \to S$ は $S$ の閉部分集合への同相写像である。
すると補題 \ref{morphisms-lemma-homeomorphism-affine} により $X \to S$ は
アフィンである。$U \subset S$ をアフィン開集合とし、
$U = \Spec(A)$ と書く。このとき (1) により $S' = \Spec(A/I)$ であり、
(2) により $I$ は局所冪零イデアルである。$f$ はアフィンなので
$f^{-1}(U) = \Spec(B)$ と書ける。仮定 (4) から、$B$ は有限型
$A$-代数である（補題 \ref{morphisms-lemma-locally-finite-type-characterize}）か、
または $I$ は有限生成である（補題
\ref{morphisms-lemma-closed-immersion-finite-presentation}）。
仮定 (3) は $A/I \to B/IB$ が全射であるということである。
$A \to B$ が有限型ならば代数、補題
\ref{algebra-lemma-surjective-mod-locally-nilpotent} により、また
$I$ が有限生成で従って冪零ならば代数、補題
\ref{algebra-lemma-NAK} により、$A \to B$ が全射であると結論できる。
これは $f$ が閉埋め込みであることを意味する。補題
\ref{morphisms-lemma-closed-immersion} を参照せよ。
\end{proof}

\begin{lemma}
\label{morphisms-lemma-permanence-universal-homeomorphism}
$f : X \to Z$ を二つの射 $g : X \to Y$ および
$h : Y \to Z$ の合成とする。射 $\{f, g, h\}$ のうち二つが
普遍同相射ならば、残りの一つも普遍同相射である。
\end{lemma}

\begin{proof}
$g$ と $h$ がともに普遍同相射ならば、補題
\ref{morphisms-lemma-composition-universal-homeomorphism} により $f$ もそうである。

\medskip\noindent
$f$ と $g$ がともに普遍同相射であると仮定する。$h$ もそうであることを
示したい。スキームの任意の射 $\alpha : Z' \to Z$ に沿って図式を基底変換
すると、すべての正方形がカルテジアンである次の図式を得る。
$$
\xymatrix{
X' \ar[r]^{g'} \ar[d] & Y' \ar[r]^{h'} \ar[d] & Z' \ar[d] \\
X \ar[r]^{g} & Y \ar[r]^{h} & Z.
}
$$
仮定により、合成 $f'= h' \circ g' : X' \to Z'$ と
$g' : X' \to Y'$ は同相写像であるから、$h'$ も同相写像である。
これで $h$ が普遍同相射であることが証明された。

\medskip\noindent
最後に、$f$ と $h$ が普遍同相射であると仮定する。$g$ が普遍同相射で
あることを示したい。$\beta : Y' \to Y$ をスキームの任意の射とする。
すべての正方形がカルテジアンである次の図式を得る。
$$
\xymatrix{
X' \ar[r]^{g'} \ar[d] & Y' \ar[d]^{\gamma} & \\
X'' \ar[r]^{g''} \ar[d] & Y'' \ar[r]^{h''} \ar[d] & Y' \ar[d]^{h \circ \beta} \\
X \ar[r]^{g} & Y \ar[r]^{h} & Z.
}
$$
ここで射 $\gamma : Y' \to Y''$ は、ファイバー積の普遍性と二つの射
$id_{Y'} : Y' \to Y'$ および $\beta : Y' \to Y$ によって定義される。
$g'$ が同相写像であることを証明しよう。同相写像であるという性質は
三つのうち二つから残りの一つが従う性質をもつので、$g''$ は同相写像である。
上の正方形を見ると、$\gamma$ が普遍同相射であることを示せば十分である。
$h''$ は同相写像なので、補題 \ref{morphisms-lemma-homeomorphism-affine} により
アフィン射であり、従って特に分離的である（補題
\ref{morphisms-lemma-affine-separated}）。$h'' \circ \gamma$ は恒等射であるから、
スキーム、補題 \ref{schemes-lemma-section-immersion} により $\gamma$ は
閉埋め込みである。$h''$ は全単射なので、$\gamma$ は全単射な閉埋め込み、
従って普遍同相射である（例えば、補題
\ref{morphisms-lemma-universal-homeomorphism} の特徴づけによる）。これでよい。
\end{proof}



\section{アフィンスキームの普遍同相射}
\label{morphisms-section-universal-homeomorphisms-affine}

\noindent
本節では、アフィンスキームの普遍同相射を特徴づける。

\begin{lemma}
\label{morphisms-lemma-subalgebra-inherits}
$A \to B$ を環準同型とし、誘導されるスキームの射
$f : \Spec(B) \to \Spec(A)$ が普遍同相射、resp.\ 剰余体上で同型を
誘導する普遍同相射、resp.\ 普遍閉、resp.\ 普遍閉かつ普遍的単射であるとする。
このとき、任意の $A$-部分代数 $B' \subset B$ に対し、
$f' : \Spec(B') \to \Spec(A)$ について同じことが成り立つ。
\end{lemma}

\begin{proof}
$f$ が普遍閉ならば、補題 \ref{morphisms-lemma-integral-universally-closed} により
$B$ は $A$ 上整である。従って $B'$ は $A$ 上整であり、$f'$ は
（同じ補題により）普遍閉である。これで $f$ が普遍閉の場合が証明された。

\medskip\noindent
続いて、$B$ は $B'$ 上整であることがわかる
（代数、補題 \ref{algebra-lemma-integral-permanence}）。
従って $\Spec(B) \to \Spec(B')$ は全射である
（代数、補題 \ref{algebra-lemma-integral-overring-surjective}）。
ゆえに $A \to B$ が剰余体の純非分離拡大を誘導するならば、
$A \to B'$ についても同じことが成り立つ。これで、$f$ が普遍閉かつ
普遍的単射である場合が証明された。補題
\ref{morphisms-lemma-universally-injective} を参照せよ。

\medskip\noindent
$f$ が普遍同相射である場合は、以上の考察、補題
\ref{morphisms-lemma-universal-homeomorphism}、および $f$ が全射ならば $f'$ も
全射であるという明らかな事実から従う。

\medskip\noindent
$A \to B$ が剰余体上で同型を誘導するならば、$A \to B'$ についても
同じである（第二段落の議論を参照）。以上により残る場合にも補題が成り立つ。
\end{proof}

\begin{lemma}
\label{morphisms-lemma-colimit-inherits}
$A$ を環とする。$B = \colim B_\lambda$ を $A$-代数のフィルター余極限とする。
各 $f_\lambda : \Spec(B_\lambda) \to \Spec(A)$ が普遍同相射、
resp.\ 剰余体上で同型を誘導する普遍同相射、resp.\ 普遍閉、
resp.\ 普遍閉かつ普遍的単射ならば、
$f : \Spec(B) \to \Spec(A)$ についても同じことが成り立つ。
\end{lemma}

\begin{proof}
$f_\lambda$ が普遍閉ならば、補題
\ref{morphisms-lemma-integral-universally-closed} により $B_\lambda$ は $A$ 上整である。
従って $B$ は $A$ 上整であり、$f$ は（同じ補題により）普遍閉である。
これで各 $f_\lambda$ が普遍閉である場合が証明された。

\medskip\noindent
$\mathfrak q \subset B$ を $\mathfrak p \subset A$ の上にある素イデアルとし、
その逆像を $\mathfrak q_\lambda \subset B_\lambda$ と表す。このとき
$\kappa(\mathfrak q) = \colim \kappa(\mathfrak q_\lambda)$ である。
従って $A \to B_\lambda$ が剰余体の純非分離拡大を誘導するならば、
$A \to B$ についても同じである。これで、$f_\lambda$ が普遍閉かつ
普遍的単射である場合が証明された。補題
\ref{morphisms-lemma-universally-injective} を参照せよ。

\medskip\noindent
$f$ が普遍同相射である場合は、以上の考察と補題
\ref{morphisms-lemma-universal-homeomorphism}、および $B$ の素イデアルが
すべての $B_\lambda$ の素イデアルの整合的な族にほかならないという
事実を組み合わせれば従う。

\medskip\noindent
$A \to B_\lambda$ が剰余体上で同型を誘導するならば、
$A \to B$ についても同じである（第二段落の議論を参照）。
以上により残る場合にも補題が成り立つ。
\end{proof}

\begin{lemma}
\label{morphisms-lemma-special-elements-and-localization}
$A \subset B$ を環拡大とする。$S \subset A$ を乗法的部分集合とする。
$n \geq 1$ とし、$b_i \in B$ を $1 \leq i \leq n$ に対して選ぶ。
条件
$$
x \not \in S^{-1}A\text{ and } b_i x^i \in S^{-1}A\text{ for }i = 1, \ldots, n
$$
を満たす任意の $x \in S^{-1}B$ は $s^{-1}y$ に等しい。ここで
$s \in S$ および $y \in B$ は
$$
y \not \in A\text{ and } b_i y^i \in A\text{ for }i = 1, \ldots, n
$$
を満たす。
\end{lemma}

\begin{proof}
省略する。ヒント：分母を払え。
\end{proof}

\begin{lemma}
\label{morphisms-lemma-nth-and-nplusone-implies-square-and-cube}
$A \subset B$ を環拡大とする。$b \in B$、$b \not \in A$、および整数
$n \geq 2$ で $b^n \in A$ かつ $b^{n + 1} \in A$ となるものが存在するならば、
$b' \in B$、$b' \not \in A$ で $(b')^2 \in A$ かつ
$(b')^3 \in A$ となるものが存在する。
\end{lemma}

\begin{proof}
$b$ と $n$ を補題のとおりとする。このとき $b$ の十分高い冪はすべて
$A$ に属する。実際、
$(b^n)^k(b^{n + 1})^i = b^{(k + i)n + i}$ であり、これから
任意の冪 $b^m$ で $m \geq n^2$ を満たすものは $A$ に属する。従って
$i \geq 1$ を $b^i \not \in A$ となる最大の整数とすれば、
$(b^i)^2 \in A$ かつ $(b^i)^3 \in A$ である。
\end{proof}

\begin{lemma}
\label{morphisms-lemma-square-and-cube}
$A \subset B$ を環拡大とし、$\Spec(B) \to \Spec(A)$ は剰余体上で
同型を誘導する普遍同相射であるとする。$A \not = B$ ならば、
$b \in B$、$b \not \in A$ で $b^2 \in A$ かつ $b^3 \in A$ となるものが存在する。
\end{lemma}

\begin{proof}
$A \subset B$ は整であることを思い出そう
（補題 \ref{morphisms-lemma-integral-universally-closed}）。
補題 \ref{morphisms-lemma-subalgebra-inherits} により、$B$ は $A$ 上一元で生成されると
仮定してよい。従って $B$ は $A$ 上有限である
（代数、補題 \ref{algebra-lemma-characterize-finite-in-terms-of-integral}）。
ゆえに、$B/A$ の $A$-加群としての台は閉で空ではない
（代数、補題 \ref{algebra-lemma-support-closed} および
\ref{algebra-lemma-support-zero}）。$\mathfrak p \subset A$ をこの台の
極小素イデアルとする。$A \subset B$ を
$A_\mathfrak p \subset B_\mathfrak p$ で置き換えれば
（補題 \ref{morphisms-lemma-special-elements-and-localization} により許される）、
$(A, \mathfrak m)$ は局所環、$B$ は $A$ 上有限、そして $B/A$ の
台は $\{\mathfrak m\}$ であると仮定してよい（これは $A$-加群としての台である）。
$B/A$ は有限加群なので、
$I = \text{Ann}_A(B/A)$ は $\mathfrak m = \sqrt{I}$ を満たす
（代数、補題 \ref{algebra-lemma-support-closed}）。
$\mathfrak m' \subset B$ を $\mathfrak m$ の上にある一意な素イデアルとする。
$\Spec(B) \to \Spec(A)$ は同相写像なので
$\mathfrak m' = \sqrt{IB}$ である。$f \in \mathfrak m'$ に対し、
$n \geq 1$ を $f^n \in IB$ となるよう選ぶ。このとき
$f^{n + 1} \in IB$ でもある。$IB \subset A$ が $I$ の選び方により成り立つので、
$f^n, f^{n + 1} \in A$ と結論できる。補題
\ref{morphisms-lemma-nth-and-nplusone-implies-square-and-cube} により、
$\mathfrak m' \not \subset A$ ならば補題の結論が成り立つ。一方、
$\mathfrak m' \subset A$ ならば $\mathfrak m' = \mathfrak m$ であり、
仮定により剰余体も同型なので $A = B$ と結論される。この矛盾により
証明が完了する。
\end{proof}

\begin{lemma}
\label{morphisms-lemma-pth-power-and-multiple}
$A \subset B$ を環拡大とし、$\Spec(B) \to \Spec(A)$ は普遍同相射であるとする。
$A \not = B$ ならば、$b \in B$、$b \not \in A$ で
$b^2 \in A$ かつ $b^3 \in A$ となるものが存在するか、または素数 $p$ と
$b \in B$、$b \not \in A$ で $pb \in A$ かつ $b^p \in A$ となるものが存在する。
\end{lemma}

\begin{proof}
議論は補題 \ref{morphisms-lemma-square-and-cube} の証明とほぼまったく同じだが、
確実に機能することがわかるよう、すべてを書き下す。

\medskip\noindent
$A \subset B$ は整であることを思い出そう
（補題 \ref{morphisms-lemma-integral-universally-closed}）。
補題 \ref{morphisms-lemma-subalgebra-inherits} により、$B$ は $A$ 上一元で生成されると
仮定してよい。従って $B$ は $A$ 上有限である
（代数、補題 \ref{algebra-lemma-characterize-finite-in-terms-of-integral}）。
ゆえに、$B/A$ の $A$-加群としての台は閉で空ではない
（代数、補題 \ref{algebra-lemma-support-closed} および
\ref{algebra-lemma-support-zero}）。$\mathfrak p \subset A$ をこの台の
極小素イデアルとする。$A \subset B$ を
$A_\mathfrak p \subset B_\mathfrak p$ で置き換えれば
（補題 \ref{morphisms-lemma-special-elements-and-localization} により許される）、
$(A, \mathfrak m)$ は局所環、$B$ は $A$ 上有限、そして $B/A$ の
台は $\{\mathfrak m\}$ であると仮定してよい（これは $A$-加群としての台である）。
$B/A$ は有限加群なので、
$I = \text{Ann}_A(B/A)$ は $\mathfrak m = \sqrt{I}$ を満たす
（代数、補題 \ref{algebra-lemma-support-closed}）。
$\mathfrak m' \subset B$ を $\mathfrak m$ の上にある一意な素イデアルとする。
$\Spec(B) \to \Spec(A)$ は同相写像なので
$\mathfrak m' = \sqrt{IB}$ である。$f \in \mathfrak m'$ に対し、
$n \geq 1$ を $f^n \in IB$ となるよう選ぶ。このとき
$f^{n + 1} \in IB$ でもある。$IB \subset A$ が $I$ の選び方により成り立つので、
$f^n, f^{n + 1} \in A$ と結論できる。補題
\ref{morphisms-lemma-nth-and-nplusone-implies-square-and-cube} により、
$\mathfrak m' \not \subset A$ ならば補題の結論が成り立つ。
$\mathfrak m' \subset A$ ならば $\mathfrak m' = \mathfrak m$ である。
$A \not = B$ なので、剰余体の写像
$\kappa = A/\mathfrak m \to B/\mathfrak m' = \kappa'$ は同型ではありえない。
補題 \ref{morphisms-lemma-universally-injective} により、$\kappa$ の標数は素数 $p$ であり、
拡大 $\kappa'/\kappa$ は純非分離である。$b \in B$ を、その $\kappa'$ に
おける像が $\kappa$ に属さない一方、その $p$ 乗が $\kappa$ に属するように選ぶ。
すると $b \not \in A$、$b^p \in A$、かつ $pb \in A$ である
（なぜなら $pb \in \mathfrak m' = \mathfrak m \subset A$）。これが求める結論である。
\end{proof}

\begin{proposition}
\label{morphisms-proposition-universal-homeomorphism-equal-residue-fields}
$A \subset B$ を環拡大とする。次は同値である。
\begin{enumerate}
\item $\Spec(B) \to \Spec(A)$ は剰余体上で同型を誘導する普遍同相射であり、
\item 任意の有限部分集合 $E \subset B$ は拡大
$$
A[b_1, \ldots, b_n] \subset B
$$
に含まれ、$b_i^2, b_i^3 \in A[b_1, \ldots, b_{i - 1}]$ が
$i = 1, \ldots, n$ に対して成り立つ。
\end{enumerate}
\end{proposition}

\begin{proof}
(1) を仮定する。超限再帰を用い、各順序数 $\alpha$ に対して
$A$-部分代数 $B_\alpha \subset B$ を次のように構成する。
$B_0 = A$ と置く。$\alpha$ が極限順序数ならば
$B_\alpha = \colim_{\beta < \alpha} B_\beta$ と置く。
$\alpha = \beta + 1$ ならば、$B_\beta = B$ の場合には
$B_\alpha = B_\beta$ と置き、$B_\beta \not = B$ の場合には補題
\ref{morphisms-lemma-square-and-cube} を適用し、$b_\alpha \in B$、
$b_\alpha \not \in B_\beta$ かつ $b_\alpha^2, b_\alpha^3 \in B_\beta$
となるものを選び、$B_\alpha = B_\beta[b_\alpha] \subset B$ と置く。
明らかに $B = \colim B_\alpha$ である（実際、濃度を考えれば
$B = B_\alpha$ となる順序数 $\alpha$ が存在する）。
$A \to B_\alpha$ が各順序数 $\alpha$ に対し (2) を満たすことを
超限帰納法で証明する。$\alpha = 0$ のときは明らかである。
すべての $\beta < \alpha$ に対して主張が成り立つと仮定し、
$E \subset B_\alpha$ を有限部分集合とする。$\alpha$ が極限順序数ならば
$B_\alpha = \bigcup_{\beta < \alpha} B_\beta$ であり、
$E \subset B_\beta$ となる $\beta < \alpha$ が存在するから、この場合の
結果を得る。$\alpha = \beta + 1$ ならば
$B_\alpha = B_\beta[b_\alpha]$ である。従って任意の $e \in E$ は、
多項式 $e = \sum d_{e, i}b_\alpha^i$ と書け、ここで
$d_{e, i} \in B_\beta$ である。
$D \subset B_\beta$ を
$D = \{d_{e, i}\} \cup \{b_\alpha^2, b_\alpha^3\}$ と置く。
帰納法の仮定により、補題の主張にある形の $A$-部分代数
$A[b_1, \ldots, b_n] \subset B_\beta$ で $D$ を含むものが存在する。
すると $A[b_1, \ldots, b_n, b_\alpha] \subset B_\alpha$ は、
補題の主張にある形の $A$-部分代数で、$B_\alpha$ に含まれ、かつ $E$ を含む。

\medskip\noindent
(2) を仮定する。$B = \colim B_\lambda$ を、その有限
$A$-部分代数の余極限として書く。補題 \ref{morphisms-lemma-colimit-inherits} により、
$\Spec(B_\lambda) \to \Spec(A)$ が剰余体上で同型を誘導する普遍同相射である
ことを示せば十分である。普遍閉射の合成は普遍閉であり、剰余体上で
同型を誘導する射についても同じことが成り立つ。従って、
$A \subset B$ で $B$ が一つの元 $b$ によって生成され、
$b^2, b^3 \in A$ であるときに (1) が成り立つことを示せば十分である。
このような拡大は整であるから、$\Spec(B) \to \Spec(A)$ は普遍閉かつ全射である
（補題 \ref{morphisms-lemma-integral-universally-closed} および代数、補題
\ref{algebra-lemma-integral-overring-surjective}）。
$(b^2)^3 = (b^3)^2$ が $A$ において成り立つことに注意する。
任意の環準同型 $\varphi : A \to K$ であって $K$ が体であるものに対し、
$\lambda \in K$ であって $\varphi(b^2) = \lambda^2$ かつ
$\varphi(b^3) = \lambda^3$ となるものが存在する。実際、$\lambda = 0$ と
置くのは $\varphi(b^2) = 0$ の場合であり、そうでなければ
$\lambda = \varphi(b^3)/\varphi(b^2)$ とすればよい。従って
$B \otimes_A K$ は $K[x]/(x^2 - \lambda^2, x^3 - \lambda^3)$ の商である。
この環は剰余体 $K$ をもつ素イデアルをちょうど一つもつ。
これは $\Spec(B) \to \Spec(A)$ が全単射で、剰余体上で同型を誘導することを
意味する。普遍閉性と組み合わせれば (1) が従う。補題
\ref{morphisms-lemma-universal-homeomorphism} および
\ref{morphisms-lemma-universally-injective} を参照せよ。
\end{proof}

\begin{proposition}
\label{morphisms-proposition-universal-homeomorphism}
$A \subset B$ を環拡大とする。次は同値である。
\begin{enumerate}
\item $\Spec(B) \to \Spec(A)$ は普遍同相射であり、
\item 任意の有限部分集合 $E \subset B$ は拡大
$$
A[b_1, \ldots, b_n] \subset B
$$
に含まれ、$i = 1, \ldots, n$ に対して次のいずれかが成り立つ。
\begin{enumerate}
\item $b_i^2, b_i^3 \in A[b_1, \ldots, b_{i - 1}]$、または
\item 素数 $p$ が存在して
$pb_i, b_i^p \in A[b_1, \ldots, b_{i - 1}]$。
\end{enumerate}
\end{enumerate}
\end{proposition}

\begin{proof}
証明は命題 \ref{morphisms-proposition-universal-homeomorphism-equal-residue-fields}
の証明とまったく同じであるが、次の変更を行う。
\begin{enumerate}
\item 補題 \ref{morphisms-lemma-pth-power-and-multiple} を補題
\ref{morphisms-lemma-square-and-cube} の代わりに用いる。従って各後続順序数
$\alpha = \beta + 1$ に対し、$b_\alpha^2, b_\alpha^3 \in B_\beta$
であるか、または素数 $p$ が存在して
$pb_\alpha, b_\alpha^p \in B_\beta$ である。
\item $\alpha$ が後続順序数ならば、
$D = \{d_{e, i}\} \cup \{b_\alpha^2, b_\alpha^3\}$ と取るか、または
$D = \{d_{e, i}\} \cup \{pb_\alpha, b_\alpha^p\}$ と取る。この選択は
$\alpha$ がどちらの場合に属するかによる。
\item (2) $\Rightarrow$ (1) の証明では、$B$ が $A$ 上一つの元 $b$ で生成され、
$pb, b^p \in B$ となる素数 $p$ が存在する場合も考える必要がある。
このとき $A \subset B$ は、例えば代数、補題
\ref{algebra-lemma-p-ring-map} により、スペクトル上の普遍同相射を誘導する。
\end{enumerate}
これで証明が完了する。
\end{proof}

\begin{lemma}
\label{morphisms-lemma-universal-homeo-iso-if-invert-p}
$p$ を素数とする。$A \to B$ を、同型 $A[1/p] \to B[1/p]$ を誘導する
環準同型とする（例えば $p$ が $A$ において冪零である場合）。次は同値である。
\begin{enumerate}
\item $\Spec(B) \to \Spec(A)$ は普遍同相射であり、
\item $A \to B$ の核は局所冪零イデアルであり、任意の $b \in B$ に対して
$p$ の冪 $q$ が存在し、$qb$ と $b^q$ は $A \to B$ の像に属する。
\end{enumerate}
\end{lemma}

\begin{proof}
(2) が成り立つならば、代数、補題 \ref{algebra-lemma-p-ring-map} により (1) が成り立つ。
(1) を仮定する。このとき、代数、補題
\ref{algebra-lemma-image-dense-generic-points} により、$A \to B$ の核は
冪零元からなる。従って $A$ を $A \to B$ の像で置き換え、$A \subset B$ と
仮定してよい。代数、補題 \ref{algebra-lemma-help-with-powers} により、集合
$$
B' = \{b \in B \mid p^nb, b^{p^n} \in A\text{ for some }n \geq 0\}
$$
は $A$-部分代数として $B$ に含まれる（積について閉じていることは明らかである）。
$B' = B$ を示さなければならない。そうでなければ、補題
\ref{morphisms-lemma-pth-power-and-multiple} によれば、$b \in B$、$b \not \in B'$ で
$b^2, b^3 \in B'$ となるものが存在するか、または素数 $\ell$ が存在して
$\ell b, b^\ell \in B'$ となる。この二つの場合がとも矛盾を導くことを示せば、
補題が証明される。

\medskip\noindent
$A[1/p] = B[1/p]$ なので、$p$ の冪 $q$ で $qb \in A$ となるものを選べる。

\medskip\noindent
$b^2, b^3 \in B'$ ならば $b^q \in B'$ でもある。$B'$ の定義により、
$(b^q)^{q'} \in A$ となることが、ある $p$ の冪 $q'$ について成り立つ。すると
$qq'b, b^{qq'} \in A$ であり、従って $b \in B'$ となって矛盾する。

\medskip\noindent
今度は素数 $\ell$ が存在して $\ell b, b^\ell \in B'$ となると仮定する。
$\ell \not = p$ ならば、$\ell b \in B'$ と $qb \in A \subset B'$ から
$b \in B'$ となり矛盾する。従って $\ell = p$ であり、$b^p \in B'$ であるから、
前とまったく同様に矛盾を得る。
\end{proof}

\begin{lemma}
\label{morphisms-lemma-make-universal-homeo}
$A$ を環とし、$x, y \in A$ とする。
\begin{enumerate}
\item $x^3 = y^2$ が $A$ において成り立つならば、
$A \to B = A[t]/(t^2 - x, t^3 - y)$ は剰余体上で全単射を誘導し、
スペクトル上で普遍同相射を誘導する。
\item 素数 $p$ が存在して $p^px = y^p$ が $A$ において成り立つならば、
$A \to B = A[t]/(t^p - x, pt - y)$ はスペクトル上で普遍同相射を誘導する。
\end{enumerate}
\end{lemma}

\begin{proof}
補題 \ref{morphisms-lemma-universal-homeomorphism} の判定法を用いてこれを確認する。
いずれの場合にも環準同型は整である。従って、体 $k$ と環準同型
$\varphi : A \to k$ が与えられたとき、$k$-代数 $B \otimes_A k$ が一意な
素イデアルをもち、その剰余体が (1) の場合には $k$ に等しく、(2) の場合には
$k$ 上純非分離であることを示せば十分である。補題
\ref{morphisms-lemma-universally-injective} を参照せよ。

\medskip\noindent
(1) の場合、$\lambda = 0$ と置くのは $\varphi(x) = 0$ のときであり、そうでなければ
$\lambda = \varphi(y)/\varphi(x)$ と置く。このとき
$B = k[t]/(t^2 - \lambda^2, t^3 - \lambda^2)$ である。従って結果は明らかである。

\medskip\noindent
(2) の場合、$k$ の標数が $p$ ならば、$\varphi(y) = 0$ および
$B = k[t]/(t^p - \varphi(x))$ を得る。これは局所アルティン $k$-代数であり、
その剰余体は $k$ であるか、または次数 $p$ の $k$ 上純非分離拡大である。
$k$ の標数が $p$ でなければ、$\lambda = \varphi(y)/p$ と置くことにより
$B = k[t]/(t - \lambda) = k$ がわかり、この場合も結論を得る。
\end{proof}

\begin{lemma}
\label{morphisms-lemma-universal-homeo-limit}
$A \to B$ を環準同型とする。
\begin{enumerate}
\item $A \to B$ がスペクトル上で普遍同相射を誘導するならば、
$B = \colim B_i$ は有限表示 $A$-代数 $B_i$ のフィルター余極限として書け、
$A \to B_i$ はスペクトル上で普遍同相射を誘導する。
\item $A \to B$ が剰余体上で同型かつスペクトル上で普遍同相射を誘導するならば、
$B = \colim B_i$ は有限表示 $A$-代数 $B_i$ のフィルター余極限として書け、
$A \to B_i$ は剰余体上で同型かつスペクトル上で普遍同相射を誘導する。
\end{enumerate}
\end{lemma}

\begin{proof}
(1) の証明。代数、補題 \ref{algebra-lemma-when-colimit} の判定法を用いる。
$A \to C$ を有限表示とし、$\varphi : C \to B$ を $A$-代数準同型とする。
$B' = \varphi(C) \subset B$ をその像とする。このとき補題
\ref{morphisms-lemma-subalgebra-inherits} により、$A \to B'$ はスペクトル上で
普遍同相射を誘導する。代数、補題 \ref{algebra-lemma-ring-colimit-fp} により
$B' = \colim_{i \in I} B_i$ と書け、$A \to B_i$ は有限表示で、遷移写像は全射である。
代数、補題 \ref{algebra-lemma-characterize-finite-presentation} により、添字
$0 \in I$ と、分解 $C \to B_0 \to B'$ を写像 $C \to B'$ に対して選べる。
$\Spec(B_i) \to \Spec(A)$ は $i$ が十分大きいとき普遍同相射であると主張する。
この主張から (1) の証明が完了する。

\medskip\noindent
主張の証明。補題 \ref{morphisms-lemma-reduction-universal-homeomorphism} により、環準同型
$A_{red} \to B'_{red}$ はスペクトル上で普遍同相射を誘導する。従って、代数、補題
\ref{algebra-lemma-image-dense-generic-points} により $A_{red} \subset B'_{red}$ である。
$A' = \Im(A \to B')$ と置くと、全射 $A \to A' \to A_{red}$ があり、
$\Spec(A_{red}) = \Spec(A') = \Spec(A)$ という全単射を誘導する。
従って $A' \subset B'$ はスペクトル上で普遍同相射を誘導する。命題
\ref{morphisms-proposition-universal-homeomorphism} と、$B'$ が $A'$ 上有限型であることから、
$n$ および $b'_1, \ldots, b'_n \in B'$ で、
$B' = A'[b'_1, \ldots, b'_n]$ となり、かつ $j = 1, \ldots, n$ に対して
次のいずれかが成り立つものを見つけられる。
\begin{enumerate}
\item $(b'_j)^2, (b'_j)^3 \in A'[b'_1, \ldots, b'_{j - 1}]$、または
\item 素数 $p$ が存在して
$pb'_j, (b'_j)^p \in A'[b'_1, \ldots, b'_{j - 1}]$。
\end{enumerate}
$b_1, \ldots, b_n \in B_0$ を $b'_1, \ldots, b'_n$ の持ち上げとして選ぶ。
$i \geq 0$ に対し、$b_{j, i}$ で $b_j$ の $B_i$ における像を表す。
$i$ が十分大きければ、$j = 1, \ldots, n$ に対して次のいずれかが成り立つ。
\begin{enumerate}
\item $b_{j, i}^2, b_{j, i}^3 \in A_i[b_{1, i}, \ldots, b_{j - 1, i}]$、または
\item 素数 $p$ が存在して
$pb_{j, i}, b_{j, i}^p \in A_i[b_{1, i}, \ldots, b_{j - 1, i}]$。
\end{enumerate}
ここで $A_i \subset B_i$ は $A \to B_i$ の像である。$A \to A_i$ は、核が
局所冪零イデアルである全射環準同型であることに注意する。必要ならばさらに
$i$ を大きくすることにより、$B_i$ は $b_1, \ldots, b_n$ により $A_i$ 上生成される、
すなわち $B_i = A_i[b_1, \ldots, b_n]$ と仮定してよい。代数、補題
\ref{algebra-lemma-p-ring-map} および \ref{algebra-lemma-2-3-ring-map} により、
$A \to A_i \to A_i[b_1] \to \ldots \to A_i[b_1, \ldots, b_n] = B_i$
はスペクトル上で普遍同相射を誘導すると結論できる。これで主張の証明が完了する。

\medskip\noindent
(2) の証明もまったく同じである。
\end{proof}






\section{絶対弱正規化と半正規化}
\label{morphisms-section-seminormalization}

\noindent
前節で証明した結果に動機づけられ、次の定義を与える。

\begin{definition}
\label{morphisms-definition-seminormal-ring}
$A$ を環とする。
\begin{enumerate}
\item $A$ が {\it 半正規} であるとは、すべての $x, y \in A$ で
$x^3 = y^2$ を満たすものに対し、一意な $a \in A$ が存在して
$x = a^2$ かつ $y = a^3$ となることをいう。
\item $A$ が {\it 絶対弱正規} であるとは、(a) $A$ が半正規であり、かつ
(b) 任意の素数 $p$ と $x, y \in A$ で $p^px = y^p$ を満たすものに対し、
一意な $a \in A$ が存在して $x = a^p$ かつ $y = pa$ となることをいう。
\end{enumerate}
\end{definition}

\noindent
半正規環の定義では存在を仮定すれば十分であるという、興味深い観察がある
（\cite{Costa}）。
\footnote{$A$ を、すべての $x, y \in A$ で $x^3 = y^2$ を満たすものに対し、
$a \in A$ が存在して $x = a^2$ かつ $y = a^3$ となるような環とする。
このとき $A$ は被約である。実際、$x^2 = 0$ ならば $x^2 = x^3$ であり、従って
$a$ が存在して $x = a^3$ かつ $x = a^2$ となる。すると
$x = a^3 = ax = a^4 = x^2 = 0$ である。最後に、
$a_1^2 = a_2^2$ かつ $a_1^3 = a_2^3$ が被約環の $a_1, a_2$ に対して成り立つならば、
$(a_1 - a_2)^3 = a_1^3 - 3a_1^2a_2 +3a_1a_2^2 - a_2^3 =
(1 - 3 + 3 - 1)a_1^3 = 0$ であり、従って $a_1 = a_2$ である。}
従って $a$ の存在だけを仮定すればよい。絶対弱正規スキームは
\cite[Appendix B]{rydh_descent} で定義された。

\begin{lemma}
\label{morphisms-lemma-seminormal-local-property}
半正規であること、または絶対弱正規であることは環の局所的性質である。
性質、定義 \ref{properties-definition-property-local} を参照せよ。
\end{lemma}

\begin{proof}
$A$ が半正規で $f \in A$ であると仮定する。$x', y' \in A_f$ を
$(x')^3 = (y')^2$ を満たすものとする。$x' = x/f^{2n}$、
$y' = y/f^{3n}$ と書く。ここである $n \geq 0$ および $x, y \in A$ を取る。$x, y$ を
$f^{2m}x, f^{3m}y$ で、$n$ を $n + m$ で置き換えれば、
$x^3 = y^2$ が $A$ において成り立つ。そこで一意な $a \in A$ が見つかり、
$x = a^2$ かつ $y = a^3$ となる。$a' = a/f^n$ と置けば、求めるとおり
$x' = (a')^2$ かつ $y' = (a')^3$ となる。$a'$ の一意性は $a$ の一意性から従う。
まったく同様に、$A$ が絶対弱正規ならば $A_f$ も絶対弱正規であることを
読者は示せる。

\medskip\noindent
$A$ を環とし、$f_1, \ldots, f_n \in A$ は単位イデアルを生成すると仮定する。
各 $A_{f_i}$ が $i$ ごとに半正規であると仮定する。$x, y \in A$ が
$x^3 = y^2$ を満たすとする。各 $i$ に対し、一意な $a_i \in A_{f_i}$ が
見つかり、$x = a_i^2$ かつ $y = a_i^3$ が
$A_{f_i}$ において成り立つ。
一意性と第一段落の結果（これは $A_{f_if_j}$ が半正規であることを示す）により、
$a_i$ と $a_j$ は $A_{f_if_j}$ の同じ元へ写る。代数、補題
\ref{algebra-lemma-standard-covering} により、一意な $a \in A$ が見つかり、
$a_i$ へ $A_{f_i}$ において各 $i$ に対して写る。同じ理由により
$x = a^2$ かつ $y = a^3$ である。明らかにこの $a$ は一意である。
従って $A$ は半正規である。$A_{f_i}$ が絶対弱正規であると仮定すれば、
まったく同じ議論により $A$ は絶対弱正規である。
\end{proof}

\noindent
次に半正規スキームと絶対弱正規スキームを定義する。

\begin{definition}
\label{morphisms-definition-seminormal}
$X$ をスキームとする。
\begin{enumerate}
\item $X$ が {\it 半正規} であるとは、各 $x \in X$ がアフィン開近傍
$\Spec(R) = U \subset X$ をもち、その環 $R$ が半正規であることをいう。
\item $X$ が {\it 絶対弱正規} であるとは、各 $x \in X$ がアフィン開近傍
$\Spec(R) = U \subset X$ をもち、その環 $R$ が絶対弱正規であることをいう。
\end{enumerate}
\end{definition}

\noindent
ここで定番の補題を述べる。

\begin{lemma}
\label{morphisms-lemma-locally-seminormal}
$X$ をスキームとする。次は同値である。
\begin{enumerate}
\item スキーム $X$ は半正規である。
\item 任意のアフィン開集合 $U \subset X$ に対し、環 $\mathcal{O}_X(U)$ は半正規である。
\item アフィン開被覆 $X = \bigcup U_i$ が存在し、各 $\mathcal{O}_X(U_i)$ は半正規である。
\item 開被覆 $X = \bigcup X_j$ が存在し、各開部分スキーム $X_j$ は半正規である。
\end{enumerate}
さらに、$X$ が半正規ならば、すべての開部分スキームは半正規である。
「半正規」を「絶対弱正規」で置き換えても同じ主張が成り立つ。
\end{lemma}

\begin{proof}
性質、補題 \ref{properties-lemma-locally-P} と補題
\ref{morphisms-lemma-seminormal-local-property} を組み合わせればよい。
\end{proof}

\begin{lemma}
\label{morphisms-lemma-seminormal-reduced}
半正規なスキームまたは環は被約である。従って絶対弱正規なスキームまたは環も被約である。
\end{lemma}

\begin{proof}
$A$ を環とする。$a \in A$ が非零であるが $a^2 = 0$ ならば、
$a^2 = 0^2$ かつ $a^3 = 0^3$ であるから、$A$ は半正規ではない。
\end{proof}

\begin{lemma}
\label{morphisms-lemma-seminormalization-ring}
$A$ を環とする。
\begin{enumerate}
\item スペクトル上で普遍同相射を誘導する環準同型 $A \to B$ の圏は、
終対象 $A \to A^{awn}$ をもつ。
\item (1) の圏における $A \to B$ が与えられたとき、得られる写像
$B \to A^{awn}$ が同型であるための必要十分条件は $B$ が絶対弱正規であることである。
\item 剰余体上で同型を誘導し、スペクトル上で普遍同相射を誘導する環準同型
$A \to B$ の圏は、終対象 $A \to A^{sn}$ をもつ。
\item (3) の圏における $A \to B$ が与えられたとき、得られる写像
$B \to A^{sn}$ が同型であるための必要十分条件は $B$ が半正規であることである。
\end{enumerate}
任意の環準同型 $\varphi : A \to A'$ に対し、一意な写像
$\varphi^{awn} : A^{awn} \to (A')^{awn}$ および
$\varphi^{sn} : A^{sn} \to (A')^{sn}$ が存在し、これらは $\varphi$ と可換する。
\end{lemma}

\begin{proof}
(1) と (2) を証明し、(3)、(4)、および最後の主張の証明は省略する。
次の形の $A$-代数の圏を考える。
$$
B = A[x_1, \ldots, x_n]/J
$$
ここで $J$ は有限生成イデアルであり、$A \to B$ はスペクトル上で
普遍同相射を定める。この圏は有向であると主張する
（圏、定義 \ref{categories-definition-directed}）。実際、
$$
B = A[x_1, \ldots, x_n]/J
\quad\text{and}\quad
B' = A[x_1, \ldots, x_{n'}]/J'
$$
が与えられたとき、
$$
B'' = A[x_1, \ldots, x_{n + n'}]/J''
$$
を考えられる。ここで $J''$ は $J$ の元と、
$f(x_{n + 1}, \ldots, x_{n + n'})$（$f \in J'$）という元によって生成される。
すると $A$-代数準同型 $B \to B''$ および $B' \to B''$ があり、これらは同型
$B \otimes_A B' \to B''$ を誘導する。補題
\ref{morphisms-lemma-base-change-universal-homeomorphism} および
\ref{morphisms-lemma-composition-universal-homeomorphism} から、
$\Spec(B'') \to \Spec(A)$ は普遍同相射であり、従って $A \to B''$ はこの圏に属する。
最後に、この圏における射 $\varphi, \varphi' : B \to B'$ が与えられ、
$B$ は上に表示したものとする。このとき $B''$ を、$B'$ を
$\varphi(x_i) - \varphi'(x_i)$（$i = 1, \ldots, n$）で生成されるイデアルで割った商とする。
$\Spec(B') = \Spec(B)$ なので、
$\Spec(B'') \to \Spec(B')$ は全単射な閉埋め込み、従って普遍同相射である。
ゆえに $B''$ はこの圏に属し、$\varphi, \varphi'$ は $B' \to B''$ によって等化される。
これで主張の証明が完了する。次のように置く。
$$
A^{awn} = \colim B
$$
ここで余極限は今述べた圏上で取る。補題 \ref{morphisms-lemma-colimit-inherits} により、
$A \to A^{awn}$ はスペクトル上で普遍同相射を誘導することに注意する
（ここで圏が有向であることを用いる）。

\medskip\noindent
スペクトル上で普遍同相射を誘導する有限表示の環準同型 $A \to B$ が与えられると、
標準的な写像 $B \to A^{awn}$ を $A^{awn}$ の構成そのものから得る。
(1) にある任意の $A \to B$ は、(1) にある有限表示の $A \to B$ の
フィルター余極限であるから（補題 \ref{morphisms-lemma-universal-homeo-limit}）、
$A \to A^{awn}$ は圏 (1) の終対象である。

\medskip\noindent
$x, y \in A^{awn}$ を $x^3 = y^2$ を満たす元とする。このとき補題
\ref{morphisms-lemma-make-universal-homeo} により、
$A^{awn} \to A^{awn}[t]/(t^2 - x, t^3 - y)$ はスペクトル上で
普遍同相射を誘導する。従って
$A \to A^{awn}[t]/(t^2 - x, t^3 - y)$ は圏 (1) に属し、一意な $A$-代数準同型
$A^{awn}[t]/(t^2 - x, t^3 - y) \to A^{awn}$ を得る。ゆえにその像
$a \in A^{awn}$（$t$ の像）は、$a^2 = x$ かつ $a^3 = y$ を
$A^{awn}$ において満たす一意な元である。
まったく同様に、素数 $p$ と $x, y \in A^{awn}$ が与えられ、
$p^px = y^p$ ならば、一意な $a \in A^{awn}$ で $a^p = x$ かつ
$pq = y$ を満たすものが見つかる。従って定義により $A^{awn}$ は絶対弱正規である。

\medskip\noindent
最後に、$A \to B$ は圏 (1) に属し、$B$ は絶対弱正規であるとする。
$A^{awn} \to B^{awn}$ はスペクトル上で普遍同相射を誘導し、かつ
$A^{awn}$ は被約であるから（補題 \ref{morphisms-lemma-seminormal-reduced}）、
$A^{awn} \subset B^{awn}$ を得る
（代数、補題 \ref{algebra-lemma-image-dense-generic-points} を参照）。
この包含が等号でなければ、補題 \ref{morphisms-lemma-pth-power-and-multiple} により、
$b \in B^{awn}$、$b \not \in A^{awn}$ で、$b^2, b^3 \in A^{awn}$ または
$pb, b^p \in A^{awn}$ を満たすものが、ある素数 $p$ に対して存在する。
しかし定義 \ref{morphisms-definition-seminormal-ring} の存在と一意性により
$b \in A^{awn}$ と強制され、証明を終える矛盾を得る。
\end{proof}

\begin{lemma}
\label{morphisms-lemma-seminormalization}
$X$ をスキームとする。
\begin{enumerate}
\item 普遍同相射 $Y \to X$ の圏は、始対象 $X^{awn} \to X$ をもつ。
\item (1) の圏における $Y \to X$ が与えられたとき、得られる射
$X^{awn} \to Y$ が同型であるための必要十分条件は、$Y$ が絶対弱正規であることである。
\item 剰余体上で同型を誘導する普遍同相射 $Y \to X$ の圏は、
始対象 $X^{sn} \to X$ をもつ。
\item (3) の圏における $Y \to X$ が与えられたとき、得られる射
$X^{sn} \to Y$ が同型であるための必要十分条件は、$Y$ が半正規であることである。
\end{enumerate}
任意のスキームの射 $h : X' \to X$ に対し、一意な射
$h^{awn} : (X')^{awn} \to X^{awn}$ および $h^{sn} : (X')^{sn} \to X^{sn}$ が存在し、これらは $h$ と可換する。
\end{lemma}

\begin{proof}
(1) と (2) を証明し、(3) と (4) の証明は省略する。$h : X' \to X$ を
スキームの射とする。(1) が $X$ と $X'$ に対して成り立つならば、
$X' \times_X X^{awn} \to X'$ は普遍同相射であり、従って
$(X')^{awn} \to X' \times_X X^{awn}$ という $X'$ 上の一意な射を、
$(X')^{awn} \to X'$ の普遍性により得る。これと射影
$X' \times_X X^{awn} \to X^{awn}$ を合成して $h^{awn}$ を得る。
さらに (2) が $X$ と $X'$ に対して成り立ち、$h$ が開埋め込みならば、
$X' \times_X X^{awn}$ は絶対弱正規であり（補題
\ref{morphisms-lemma-locally-seminormal}）、$(X')^{awn} \to X' \times_X X^{awn}$ は同型である。

\medskip\noindent
任意の普遍同相射がアフィンであることを思い出そう。補題
\ref{morphisms-lemma-homeomorphism-affine} を参照せよ。従って $X$ がアフィンならば、
(1) と (2) は補題 \ref{morphisms-lemma-seminormalization-ring} から直ちに従う。
$X$ をスキームとし、$\mathcal{B}$ を $X$ のアフィン開集合全体とする。
各 $U \in \mathcal{B}$ に対して $U^{awn} \to U$ を得る。また
$V \subset U$ かつ $V, U \in \mathcal{B}$ ならば、前段落の議論により標準的な同型
$\rho_{V, U} : V^{awn} \to V \times_U U^{awn}$ を得る。従って相対的貼り合わせ
（構成、補題 \ref{constructions-lemma-relative-glueing}）により、射
$X^{awn} \to X$ を得る。この射の $U^{awn}$ への制限は $U$ 上で
$\rho_{V, U}$ と可換する。次に $Y \to X$ を普遍同相射とする。
このとき $U \times_X Y \to U$ は $U \in \mathcal{B}$ に対して普遍同相射であり、
一意な射 $g_U : U^{awn} \to U \times_X Y$ を $U$ 上で得る。
これらの $g_U$ は射 $\rho_{V, U}$ と可換する。詳細は省略する。
従って一意な射 $g : X^{awn} \to Y$ が存在し、これは $X$ 上のもので、$g_U$ と $U$ 上で一致する。
構成、注意 \ref{constructions-remark-relative-glueing-functorial} を参照せよ。
これで $X$ に対する (1) が証明された。(2) はアフィン局所的に成り立つことから従う。
\end{proof}

\begin{definition}
\label{morphisms-definition-seminormalization}
$X$ をスキームとする。
\begin{enumerate}
\item 補題 \ref{morphisms-lemma-seminormalization} で構成した射 $X^{sn} \to X$ を
$X$ の {\it 半正規化} という。
\item 補題 \ref{morphisms-lemma-seminormalization} で構成した射 $X^{awn} \to X$ を
$X$ の {\it 絶対弱正規化} という。
\end{enumerate}
\end{definition}

\noindent
確かに、半正規化 $X^{sn}$、すなわち $X$ の半正規化は半正規スキームであり、
絶対弱正規化 $X^{awn}$ は絶対弱正規スキームである。さらに、任意のスキームの射
$h : Y \to X$ に対して、標準的な可換図
$$
\xymatrix{
Y^{awn} \ar[d]^{h^{awn}} \ar[r] &
Y^{sn} \ar[d]^{h^{sn}} \ar[r] &
Y \ar[d]^h \\
X^{awn} \ar[r] &
X^{sn} \ar[r] &
X
}
$$
を得る。射 $h^{sn}$ および $h^{awn}$ は $h$ と可換する一意な射である。

\begin{lemma}
\label{morphisms-lemma-characterize-seminormal}
$X$ をスキームとする。次は同値である。
\begin{enumerate}
\item $X$ は半正規である。
\item $X$ はそれ自身の半正規化に等しい、すなわち射 $X^{sn} \to X$ は同型である。
\item $\pi : Y \to X$ が剰余体上で同型を誘導する普遍同相射で、$Y$ が被約ならば、
$\pi$ は同型である。
\end{enumerate}
\end{lemma}

\begin{proof}
(1) と (2) の同値性は補題 \ref{morphisms-lemma-seminormalization} から明らかである。
(3) が成り立つならば、$X^{sn} \to X$ は同型であり、(2) が成り立つとわかる。

\medskip\noindent
(2) が成り立つと仮定し、$\pi : Y \to X$ を剰余体上で同型を誘導する普遍同相射で、
$Y$ が被約であるものとする。補題 \ref{morphisms-lemma-seminormalization} により、
分解 $X \to Y \to X$ が存在し、その合成は $\text{id}_X$ である。このとき $X \to Y$ は閉埋め込みである
（スキーム、補題 \ref{schemes-lemma-section-immersion}、および $\pi$ が例えば補題
\ref{morphisms-lemma-universally-injective-separated} により分離的であることによる）。
$X \to Y$ は点上でも全単射なので、$Y$ の被約性から同型でなければならない。
これで証明が完了する。
\end{proof}

\begin{lemma}
\label{morphisms-lemma-characterize-absolutely-weakly-normal}
$X$ をスキームとする。次は同値である。
\begin{enumerate}
\item $X$ は絶対弱正規である。
\item $X$ はそれ自身の絶対弱正規化に等しい、すなわち射 $X^{awn} \to X$ は同型である。
\item $\pi : Y \to X$ が普遍同相射で、$Y$ が被約ならば、$\pi$ は同型である。
\end{enumerate}
\end{lemma}

\begin{proof}
これは補題 \ref{morphisms-lemma-characterize-seminormal} とまったく同様に証明される。
\end{proof}





\section{有限局所自由射}
\label{morphisms-section-finite-locally-free}

\noindent
多くの論文では、実際には有限局所自由射を意味するところで、著者は有限平坦射を用いている。
その理由は、基底が局所 Noetherian ならば両者が同じだからである。しかし一般には同じではない。
演習、演習 \ref{exercises-exercise-flat-not-projective} を参照せよ。

\begin{definition}
\label{morphisms-definition-finite-locally-free}
$f : X \to S$ をスキームの射とする。$f$ が {\it 有限局所自由} であるとは、$f$ が
アフィンであり、$f_*\mathcal{O}_X$ が有限局所自由な
$\mathcal{O}_S$-加群であることをいう。この場合、$f$ が
{\it 階数} または {\it 次数} $d$ をもつとは、層 $f_*\mathcal{O}_X$ が
次数 $d$ の有限局所自由層であることをいう。
\end{definition}

\noindent
$f : X \to S$ が有限局所自由ならば、$S$ は開かつ閉な部分スキーム $S_d$ の
非交和であり、$f^{-1}(S_d) \to S_d$ は次数 $d$ の有限局所自由射となることに注意する。

\begin{lemma}
\label{morphisms-lemma-finite-flat}
$f : X \to S$ をスキームの射とする。次は同値である。
\begin{enumerate}
\item $f$ は有限局所自由である。
\item $f$ は有限、平坦、かつ有限表示局所的である。
\end{enumerate}
$S$ が局所 Noetherian ならば、これらはさらに次とも同値である。
\begin{enumerate}
\item[(3)] $f$ は有限かつ平坦である。
\end{enumerate}
\end{lemma}

\begin{proof}
$V \subset S$ をアフィン開集合とする。三つのいずれの場合にも射はアフィンであるから、
$f^{-1}(V)$ はアフィンである。従って
$V = \Spec(R)$ および $f^{-1}(V) = \Spec(A)$ と書け、ここである $R$-代数 $A$ を取る。
(1) を仮定する。これは、$S$ をアフィン開集合 $V = \Spec(R)$ で被覆し、
$A$ が有限自由 $R$-加群となるようにできることを意味する。このとき、代数、補題
\ref{algebra-lemma-finite-finite-type} により $R \to A$ は有限表示である。
従って (2) が成り立つ。逆に (2) を仮定する。任意のアフィン開集合
$V = \Spec(R)$ を $S$ の中に取ると、環準同型 $R \to A$ は有限かつ有限表示であり、$A$ は
$R$-加群として平坦である。代数、補題
\ref{algebra-lemma-finite-finitely-presented-extension} により、
$A$ は有限表示 $R$-加群である。従って代数、補題
\ref{algebra-lemma-finite-projective} から $A$ は有限局所自由である。(1) が成り立つ。
Noetherian の場合は、Noetherian 環上の有限加群が有限表示加群であることから従う。
代数、補題 \ref{algebra-lemma-Noetherian-finite-type-is-finite-presentation} を参照せよ。
\end{proof}

\begin{lemma}
\label{morphisms-lemma-composition-finite-locally-free}
有限局所自由射の合成は有限局所自由である。
\end{lemma}

\begin{proof}
省略する。
\end{proof}

\begin{lemma}
\label{morphisms-lemma-base-change-finite-locally-free}
有限局所自由射の基底変換は有限局所自由である。
\end{lemma}

\begin{proof}
省略する。
\end{proof}

\begin{lemma}
\label{morphisms-lemma-finite-locally-free}
$f : X \to S$ をスキームの有限局所自由射とする。開かつ閉な部分スキームによる
非交和分解 $S = \coprod_{d \geq 0} S_d$ が存在し、$X_d = f^{-1}(S_d)$ と置けば、
制限 $f|_{X_d}$ は有限局所自由射 $X_d \to S_d$ であり、その次数は $d$ である。
\end{lemma}

\begin{proof}
有限局所自由層が局所的によく定まった階数をもつことから従う。詳細は省略する。
\end{proof}

\begin{lemma}
\label{morphisms-lemma-massage-finite}
$f : Y \to X$ を有限射とし、$X$ はアフィンであるとする。次の図式が存在する。
$$
\xymatrix{
Z' \ar[rd] &
Y' \ar[l]^i \ar[d] \ar[r] &
Y \ar[d] \\
 & X' \ar[r] & X
}
$$
ここで
\begin{enumerate}
\item $Y' \to Y$ および $X' \to X$ は全射な有限局所自由射である。
\item $Y' = X' \times_X Y$ である。
\item $i : Y' \to Z'$ は閉埋め込みである。
\item $Z' \to X'$ は有限局所自由である。
\item $Z' = \bigcup_{j = 1, \ldots, m} Z'_j$ は閉部分スキームの
（集合論的）有限合併であり、各部分スキームは $X'$ へ同型に写る。
\end{enumerate}
\end{lemma}

\begin{proof}
$X = \Spec(A)$ および $Y = \Spec(B)$ と書く。代数詳論、節
\ref{more-algebra-section-descent-flatness-integral} も参照せよ。
$x_1, \ldots, x_n \in B$ を $B$ の $A$-代数としての生成元とする。各 $i$ に対して、
ある首一多項式 $P_i(T) \in A[T]$ を $P(x_i) = 0$ が $B$ で成り立つように選べる。
代数、補題 \ref{algebra-lemma-adjoin-roots} を（$n$ 回）適用すると、
有限局所自由な環拡大 $A \subset A'$ が存在し、各 $P_i$ は完全に分解する。
$$
P_i(T) = \prod\nolimits_{k = 1, \ldots, d_i} (T - \alpha_{ik})
$$
ここで $\alpha_{ik} \in A'$ である。次のように置く。
$$
C = A'[T_1, \ldots, T_n]/(P_1(T_1), \ldots, P_n(T_n))
$$
また $B' = A' \otimes_A B$ と置く。写像 $C \to B'$、$T_i \mapsto 1 \otimes x_i$ は
$A'$-代数の全射である。$X' = \Spec(A')$、$Y' = \Spec(B')$、
$Z' = \Spec(C)$ と置けば、(1) -- (4) が成り立つ。(5) は、集合論的に
$\Spec(C)$ が次のイデアルで切り出される閉部分スキームの和だから成り立つ。
$$
(T_1 - \alpha_{1k_1}, T_2 - \alpha_{2k_2}, \ldots, T_n - \alpha_{nk_n})
$$
ここで任意の $1 \leq k_i \leq d_i$ を取る。
\end{proof}

\noindent
次の補題は、後の補題 \ref{morphisms-lemma-image-nowhere-dense-quasi-finite} で
適切な一般性のもとに述べられる。

\begin{lemma}
\label{morphisms-lemma-image-nowhere-dense-finite}
$f : Y \to X$ をスキームの有限射とする。$T \subset Y$ を $Y$ の閉じた疎部分集合とする。
このとき $f(T) \subset X$ は $X$ の閉じた疎部分集合である。
\end{lemma}

\begin{proof}
補題 \ref{morphisms-lemma-finite-proper} により $f(T) \subset X$ は閉じている。
$X = \bigcup X_i$ をアフィン被覆とする。$T$ は $Y$ で疎であるから、
$T \cap f^{-1}(X_i)$ も $f^{-1}(X_i)$ で疎である。従って定理をアフィンの場合に
証明できれば、$f(T) \cap X_i$ は疎であるとわかる。すると位相、補題
\ref{topology-lemma-nowhere-dense-local} により、$T$ は $X$ で疎であることが従う。

\medskip\noindent
$X$ がアフィンであると仮定する。補題 \ref{morphisms-lemma-massage-finite} にあるような図式
$$
\xymatrix{
Z' \ar[rd] &
Y' \ar[l]^i \ar[d]^{f'} \ar[r]_a &
Y \ar[d]^f \\
 & X' \ar[r]^b & X
}
$$
を選ぶ。射 $a$、$b$ は有限局所自由であるから開である
（補題 \ref{morphisms-lemma-finite-flat} および \ref{morphisms-lemma-fppf-open}）。
従って $T' = a^{-1}(T)$ は疎である。位相、補題
\ref{topology-lemma-open-inverse-image-closed-nowhere-dense} を参照せよ。
射 $b$ は全射かつ開である。従って、$f'(T') = b^{-1}(f(T))$ が疎であることを
証明できれば、$f(T)$ は疎である。位相、補題
\ref{topology-lemma-open-inverse-image-closed-nowhere-dense} を参照せよ。
$i$ は閉埋め込みなので、位相、補題
\ref{topology-lemma-closed-image-nowhere-dense} により
$i(T') \subset Z'$ は閉じており疎である。従って問題は次段落で論じる場合に帰着された。

\medskip\noindent
$Y = \bigcup_{i = 1, \ldots, n} Y_i$ が閉部分集合の有限合併であり、各部分集合が
$X$ へ同型に写ると仮定する。$T_i = Y_i \cap T$ と置く。各 $T_i$ が $Y_i$ で疎ならば、
各 $f(T_i)$ は $X$ で疎である。これは $Y_i \to X$ が同型だからである。従って
$f(T) = f(T_i)$ は $X$ の閉じた疎部分集合の有限合併であり、結論を得る。
位相、補題 \ref{topology-lemma-nowhere-dense} を参照せよ。
そうでないとし、例えば $T_1$ が空でない開集合 $V \subset Y_1$ を含むとする。
これが矛盾を導くことを示す。$Y_2 \cap V \subset V$ を考える。これは真の閉部分集合であるか、
$V$ に等しい。前者ならば $V$ を $V \setminus V \cap Y_2$ で置き換えると、
$V \subset T_1$ は $Y_1$ で開であり、$Y_2$ と交わらない。後者ならば
$V \subset Y_1 \cap Y_2$ は $Y_1$ と $Y_2$ の双方で開である。
$i = 3, \ldots, n$ に対して順に同じ操作を繰り返す。その結果、非交和分解
$$
\{1, \ldots, n\} = I_1 \amalg I_2, \quad 1 \in I_1
$$
と、$V$ という $Y_1$ の開集合で $T_1$ に含まれ、$V \subset Y_i$ が $i \in I_1$ に対して、
$V \cap Y_i = \emptyset$ が $i \in I_2$ に対して成り立つものを得る。
$U = f(V)$ と置く。これは $X$ の開集合である。実際、$f|_{Y_1} : Y_1 \to X$ は同型である。
このとき
$$
f^{-1}(U) = V\ \amalg\ \bigcup\nolimits_{i \in I_2} (Y_i \cap f^{-1}(U))
$$
である。$\bigcup_{i \in I_2} Y_i$ は閉じているので、$V \subset f^{-1}(U)$ は開であり、
従って $V \subset Y$ は開である。これは $T$ が $Y$ で疎であるという仮定に反する。
\end{proof}





\section{有理写像}
\label{morphisms-section-rational-maps}

\noindent
$X$ をスキームとする。$U$、$V$ が $X$ の稠密開集合ならば、
$U \cap V$ も稠密開集合であることに注意する。

\begin{definition}
\label{morphisms-definition-rational-map}
$X$、$Y$ をスキームとする。
\begin{enumerate}
\item $f : U \to Y$、$g : V \to Y$ を、稠密開部分集合 $U$、$V$ という
$X$ の部分集合上で定義されたスキームの射とする。$f$ が $g$ と {\it 同値} であるとは、
$f|_W = g|_W$ を満たすある $W \subset U \cap V$ が存在し、$X$ で稠密開であることをいう。
\item {\it $X$ から $Y$ への有理写像} とは、(1) で定義した同値関係に関する同値類である。
\item $X$、$Y$ が基底スキーム $S$ 上のスキームであるとする。$X$ から $Y$ への有理写像を
{\it $S$-有理写像（$X$ から $Y$ への）} とよぶのは、その同値類が
代表 $f : U \to Y$ をもち、それが $S$-射であるときである。
\end{enumerate}
\end{definition}

\noindent
定義の (1) にある二つの射 $f$、$g$ について、同値であるという代わりに、
同じ有理写像を定めるということもある。場合によっては、有理写像は
生成点における局所環上の写像によって決定される。

\begin{lemma}
\label{morphisms-lemma-rational-map-finite-presentation}
$S$ をスキームとする。$X$ と $Y$ を $S$ 上のスキームとする。$X$ の既約成分は有限個で、
その生成点を $x_1, \ldots, x_n$ と仮定する。$s_i \in S$ を $x_i$ の像とする。写像
$$
\left\{
\begin{matrix}
S\text{-rational maps} \\
\text{from }X\text{ to }Y
\end{matrix}
\right\}
\longrightarrow
\left\{
\begin{matrix}
(y_1, \varphi_1, \ldots, y_n, \varphi_n)\text{ where }
y_i \in Y\text{ lies over }s_i\text{ and}\\
\varphi_i : \mathcal{O}_{Y, y_i} \to \mathcal{O}_{X, x_i}
\text{ is a local }\mathcal{O}_{S, s_i}\text{-algebra map}
\end{matrix}
\right\}
$$
で、$f : U \to Y$ を $2n$-組へ送り、その成分を
$y_i = f(x_i)$ および $\varphi_i = f^\sharp_{x_i}$ とするものを考える。このとき
\begin{enumerate}
\item $Y \to S$ が有限型局所的ならば、この写像は単射である。
\item $Y \to S$ が有限表示局所的ならば、この写像は全単射である。
\item $Y \to S$ が有限型局所的で $X$ が被約ならば、この写像は全単射である。
\end{enumerate}
\end{lemma}

\begin{proof}
$X$ の任意の稠密開集合は点 $x_i$ を含むので、この構成は意味をもつ。
(1) または (2) を証明するため、$X$ を任意の稠密開集合で置き換えてよい。従って
$Z_1, \ldots, Z_n$ が $X$ の既約成分ならば、$X$ を
$X \setminus \bigcup_{i \not = j} Z_i \cap Z_j$ で置き換えてよい。
こうした後、$X$ はその既約成分（開かつ閉な部分スキームと見る）の非交和となる。
すると矢印の右辺と左辺はともに既約成分にわたる積であり、$X$ が既約な場合に帰着する。

\medskip\noindent
$X$ は既約で、生成点 $x$ は $s \in S$ の上にあると仮定する。(1) は補題
\ref{morphisms-lemma-morphism-defined-local-ring} の (1) から従う。(2) と (3) は
同じ補題の (2) から従う。
\end{proof}

\begin{definition}
\label{morphisms-definition-rational-function}
$X$ をスキームとする。$X$ 上の {\it 有理関数} とは、
$X$ から $\mathbf{A}^1_{\mathbf{Z}}$ への有理写像である。
\end{definition}

\noindent
アフィン直線 $\mathbf{A}^1$ の定義については、構成、定義
\ref{constructions-definition-affine-n-space} を参照せよ。$X$ を $S$ 上のスキームとする。
任意の開集合 $U \subset X$ に対し、射 $U \to \mathbf{A}^1_{\mathbf{Z}}$ は
射 $U \to \mathbf{A}^1_S$ で $S$ 上のものと同じである。従って有理関数は、
$S$-有理写像（$X$ から $\mathbf{A}^1_S$ への）とも同じである。

\medskip\noindent
標準的な同一視
$\Mor(T, \mathbf{A}^1_{\mathbf{Z}}) = \Gamma(T, \mathcal{O}_T)$
が任意のスキーム $T$ に対してあることを思い出そう。スキーム、例
\ref{schemes-example-global-sections} を参照せよ。従って
$\mathbf{A}^1_{\mathbf{Z}}$ はスキームの圏における環対象である。より正確には、射
\begin{eqnarray*}
+ : \mathbf{A}^1_{\mathbf{Z}} \times \mathbf{A}^1_{\mathbf{Z}}
& \longrightarrow &
\mathbf{A}^1_{\mathbf{Z}} \\
(f, g) & \longmapsto & f + g \\
* : \mathbf{A}^1_{\mathbf{Z}} \times \mathbf{A}^1_{\mathbf{Z}}
& \longrightarrow &
\mathbf{A}^1_{\mathbf{Z}} \\
(f, g) & \longmapsto & fg
\end{eqnarray*}
は環の加法と乗法のすべての公理を満たす（いつものように $1$ をもつ可換環である）。
従って $\mathbf{A}^1_{\mathbf{Z}}$ への有理写像の集合も自然な環構造をもつ。

\begin{definition}
\label{morphisms-definition-ring-of-rational-functions}
$X$ をスキームとする。$X$ 上の {\it 有理関数環} とは、元が有理関数であり、
上で述べた加法と乗法をもつ環 $R(X)$ のことである。
\end{definition}

\noindent
既約成分が有限個のスキームについては、この環を計算できる。

\begin{lemma}
\label{morphisms-lemma-integral-scheme-rational-functions}
$X$ を有限個の既約成分 $X_1, \ldots, X_n$ をもつスキームとする。
$\eta_i \in X_i$ が生成点ならば、
$$
R(X) = \mathcal{O}_{X, \eta_1} \times \ldots \times \mathcal{O}_{X, \eta_n}
$$
$X$ が被約ならば、これは $\prod \kappa(\eta_i)$ に等しい。
$X$ が整ならば、$R(X) = \mathcal{O}_{X, \eta} = \kappa(\eta)$ は体である。
\end{lemma}

\begin{proof}
$U \subset X$ を稠密開部分集合とする。このとき
$U_i = (U \cap X_i) \setminus (\bigcup_{j \not = i} X_j)$
は $\eta_i$ を含むので空でない開集合であり、$X_i$ に含まれ、
$\bigcup U_i \subset U \subset X$ は稠密である。従って補題の同一視は次の等式列から得られる。
\begin{align*}
R(X)
& =
\colim_{U \subset X\text{ open dense}} \Mor(U, \mathbf{A}^1_\mathbf{Z}) \\
& =
\colim_{U \subset X\text{ open dense}} \mathcal{O}_X(U) \\
& =
\colim_{\eta_i \in U_i \subset X\text{ open}} \prod \mathcal{O}_X(U_i) \\
& =
\prod \colim_{\eta_i \in U_i \subset X\text{ open}} \mathcal{O}_X(U_i) \\
& =
\prod \mathcal{O}_{X, \eta_i}
\end{align*}
ここで第二の等号は、スキーム、例
\ref{schemes-example-global-sections} による。最後の主張は、代数、補題
\ref{algebra-lemma-minimal-prime-reduced-ring} から従う。
\end{proof}

\begin{definition}
\label{morphisms-definition-function-field}
$X$ を整スキームとする。$X$ の {\it 関数体}、または {\it 有理関数体} とは体 $R(X)$ のことである。
\end{definition}

\noindent
この体を $k(X)$ で表し、$R(X)$ を用いないこともある。関数体の概念を用いると、
整スキーム上の分離条件を明らかにできる。補題
\ref{morphisms-lemma-integral-scheme-rational-functions} により、整スキーム上では各局所環
$\mathcal{O}_{X, x}$ を $R(X)$ の局所部分環とみなせることに注意する。

\begin{lemma}
\label{morphisms-lemma-distinct-local-rings}
$X$ を整分離スキームとする。$Z_1$、$Z_2$ を $X$ の異なる既約閉部分集合とし、
$\eta_i$ を $Z_i$ の生成点とする。$Z_1 \not\subset Z_2$ ならば、
$\mathcal{O}_{X, \eta_1} \not \subset \mathcal{O}_{X, \eta_2}$
が $R(X)$ の部分環として成り立つ。特に $Z_1 = \{x\}$ が一つの閉点 $x$ からなるならば、
$x$ の近傍で正則だが $\mathcal{O}_{X, \eta_{2}}$ に属さない関数が存在する。
\end{lemma}

\begin{proof}
まず、$X$ が分離的であるという仮定のもとで、スキームの一意な写像
$\Spec(\mathcal{O}_{X, \eta_2}) \to X$ が $X$ 上に存在し、その合成
$$
\Spec(R(X)) \longrightarrow
\Spec(\mathcal{O}_{X, \eta_2}) \longrightarrow X
$$
が標準写像 $\Spec(R(X)) \to X$ となることを観察する。実際、標準写像
$can : \Spec(\mathcal{O}_{X, \eta_2}) \to X$ がある。スキーム、式
(\ref{schemes-equation-canonical-morphism}) を参照せよ。第二の射 $a$ で $X$ へのものが
与えられると、仮定により $a$ は $can$ と
$\Spec(\mathcal{O}_{X, \eta_2})$ の生成点上で一致する。ここで $X$ が分離的であることから、
$\Spec(\mathcal{O}_{X, \eta_2})$ のうち二つの写像が一致する部分集合は閉じている。
スキーム、補題 \ref{schemes-lemma-where-are-they-equal} を参照せよ。
従って $a = can$ が $\Spec(\mathcal{O}_{X, \eta_2})$ 全体で成り立つ。

\medskip\noindent
$Z_1 \not \subset Z_2$ を仮定し、反対に
$\mathcal{O}_{X, \eta_{1}} \subset \mathcal{O}_{X, \eta_{2}}$ が
$R(X)$ の部分環として成り立つと仮定する。このとき第二の射
$$
\Spec(\mathcal{O}_{X, \eta_{2}}) \longrightarrow
\Spec(\mathcal{O}_{X, \eta_{1}}) \longrightarrow
X.
$$
を得る。上の議論により、この合成は $can$ に等しくなければならない。
これは $\eta_2$ が $\eta_1$ へ特殊化することを意味する
（スキーム、補題 \ref{schemes-lemma-specialize-points} を参照）。
しかしこれは仮定 $Z_1 \not \subset Z_2$ に反する。
\end{proof}

\begin{definition}
\label{morphisms-definition-domain-of-definition}
$\varphi$ を二つのスキーム $X$ と $Y$ の間の有理写像とする。
$\varphi$ が {\it 点 $x \in X$ で定義される} とは、$(U, f)$ という $\varphi$ の代表で
$x \in U$ となるものが存在することをいう。$\varphi$ の {\it 定義域} とは、
$\varphi$ が定義される点すべての集合である。
\end{definition}

\noindent
この定義では、一般に $\varphi$ が定義域全体で定義された代表をもつとは限らない。

\begin{lemma}
\label{morphisms-lemma-rational-map-from-reduced-to-separated}
$X$ と $Y$ をスキームとする。$X$ は被約で $Y$ は分離的であると仮定する。
$\varphi$ を $X$ から $Y$ への有理写像とし、その定義域を $U \subset X$ とする。
このとき一意な射 $f : U \to Y$ が存在し、これは $\varphi$ を表す。
$X$ と $Y$ が分離スキーム $S$ 上のスキームであり、$\varphi$ が
$S$-有理写像ならば、$f$ は $S$ 上の射である。
\end{lemma}

\begin{proof}
$(V, g)$ と $(V', g')$ を $\varphi$ の代表とする。このとき $g, g'$ は
稠密開部分スキーム $W \subset V \cap V'$ 上で一致する。一方、等化子 $E$、すなわち
$g|_{V \cap V'}$ と $g'|_{V \cap V'}$ の等化子は $V \cap V'$ の閉部分スキームである
（スキーム、補題 \ref{schemes-lemma-where-are-they-equal}）。
ここで $W \subset E$ から $E = V \cap V'$ が集合論的に従う。
$V \cap V'$ は被約なので、$E = V \cap V'$ がスキーム論的にも成り立つと結論できる。
すなわち $g|_{V \cap V'} = g'|_{V \cap V'}$ である。従って
$g : V \to Y$ という $\varphi$ の代表たちを貼り合わせ、射 $f : U \to Y$ を得ることができる。
スキーム、補題 \ref{schemes-lemma-glue} を参照せよ。最後の主張の証明は省略する。
\end{proof}

\noindent
一般には有理写像を合成することに意味はない。その理由は、第一の有理写像の代表の像が、
第二の有理写像の定義域と空の交わりしかもたないことがあるからである。しかし、
スキームが既約であると仮定して支配的有理写像を考えれば、有理写像を合成できる。

\begin{definition}
\label{morphisms-definition-dominant-rational}
$X$ と $Y$ を既約スキームとする。$X$ から $Y$ への有理写像を
{\it 支配的} とよぶのは、任意の代表 $f : U \to Y$ がスキームの支配的射であるときである。
\end{definition}

\noindent
補題 \ref{morphisms-lemma-dominant-finite-number-irreducible-components} により、これは
生成点 $\eta \in X$ が生成点 $\xi$（$Y$ の）へ写ること、すなわち
$f(\eta) = \xi$ が任意の代表 $f : U \to Y$ に対して成り立つことと同値である。
支配的有理写像 $\varphi$（既約スキーム $X$ と $Y$ の間の）と、
任意の有理写像 $\psi$（$Y$ から $Z$ への）を合成できる。実際、代表
$f : U \to Y$ で $U \subset X$ が稠密開であるものと、代表
$g : V \to Z$ で $V \subset Y$ が稠密開であるものを選ぶ。このとき
$W = f^{-1}(V) \subset X$ は空でない開集合である（$X$ の生成点を含むからである）。
$\psi \circ \varphi$ を $g \circ f|_W : W \to Z$ の同値類と定める。
これがよく定義されることの確認は省略する。

\medskip\noindent
このようにして、対象が既約スキームで、射が支配的有理写像である圏を得る。
基底スキーム $S$ が与えられれば、同様に $S$ 上の既約スキームを対象とし、
支配的 $S$-有理写像を射とする圏を定義できる。

\begin{definition}
\label{morphisms-definition-birational-schemes}
$X$ と $Y$ を既約スキームとする。
\begin{enumerate}
\item $X$ と $Y$ が {\it 双有理} であるとは、$X$ と $Y$ が既約スキームと
支配的有理写像の圏で同型であることをいう。
\item $X$ と $Y$ が基底スキーム $S$ 上のスキームであると仮定する。
$X$ と $Y$ が {\it $S$-双有理} であるとは、$X$ と $Y$ が $S$ 上の既約スキームと
支配的 $S$-有理写像の圏で同型であることをいう。
\end{enumerate}
\end{definition}

\noindent
$X$ と $Y$ が双有理な既約スキームならば、$X$ から $Z$ への有理写像の集合と
$Y$ から $Z$ への有理写像の集合は、任意のスキーム $Z$ に対して全単射である
（この $Z$ に関して関手的に）。「一般の」既約スキームに対して、これは可能な定義の一つにすぎない。
別の定義として、$X$ と $Y$ の有理関数環が同型であることを要求できる。
多様体については、これらの条件は同値である。補題
\ref{morphisms-lemma-criterion-birational-finite-presentation} を参照せよ。

\begin{lemma}
\label{morphisms-lemma-birational-integral}
$X$ と $Y$ を既約スキームとする。
\begin{enumerate}
\item スキーム $X$ と $Y$ が双有理であるための必要十分条件は、同型な空でない開集合をもつことである。
\item $X$ と $Y$ が基底スキーム $S$ 上のスキームであると仮定する。このとき
$X$ と $Y$ が $S$-双有理であるための必要十分条件は、空でない開集合
$U \subset X$ と $V \subset Y$ が存在し、これらが $S$-同型であることである。
\end{enumerate}
\end{lemma}

\begin{proof}
$X$ と $Y$ が双有理であると仮定する。$f : U \to Y$ と $g : V \to X$ が、
$X$ から $Y$ へ、および $Y$ から $X$ への互いに逆な支配的有理写像を定めるとする。
$V$ はアフィンであると仮定してよい。$U$ を $f^{-1}(V)$ のアフィン開集合で置き換えてよい。
$g \circ f$ は支配的有理写像として恒等写像なので、合成 $U \to V \to X$ は
$U$ の稠密開集合上で恒等写像である。従って $U$ をより小さいアフィン開集合で
置き換えることにより、$U \to V \to X$ が $U$ の $X$ への包含写像であると仮定できる。
従って $U \to V$ は局所閉埋め込みである
（スキーム、補題 \ref{schemes-lemma-section-immersion} を
$U \to g^{-1}(U) \to U$ に適用する）。
しかし $U$ と $V$ の役割を交換して上の議論を繰り返すと、
空でないアフィン開集合 $V' \subset V$ で、その包含が $V' \to U \to V$ と
分解するものが存在するとわかる。このとき $V' \to U$ は必ず開埋め込みである。実際、
$V' \to f^{-1}(V') \to V'$ はモノ射であり
（スキーム、補題 \ref{schemes-lemma-immersions-monomorphisms}）、合成が恒等射なので同型である。
従って $V'$ は $X$ と $Y$ の双方の開集合と同型である。
$S$-有理写像の場合にも、まったく同じ議論が成り立つ。
\end{proof}

\begin{remark}
\label{morphisms-remark-generalize-category}
既約スキームと支配的有理写像の圏には次の一般化がある。スキーム $X$ に対し、
$X^0$ で点 $x \in X$ のうち $\dim(\mathcal{O}_{X, x}) = 0$ となるものの集合を表す。
言い換えれば、$X^0$ は $X$ の既約成分の生成点全体の集合である。
このとき次の圏を考えられる。
\begin{enumerate}
\item 対象は、各準コンパクト開集合が有限個の既約成分をもつスキーム $X$ である。
\item $X$ から $Y$ への射は、有理写像 $f : U \to Y$（$X$ から $Y$ への）で
$f(U^0) = Y^0$ を満たすものである。
\end{enumerate}
$U \subset X$ がスキームの稠密開集合ならば、$U^0 \subset X^0$ は
等号であるとは限らないが、$X$ がこの圏の対象ならば等号となる。
従ってこの圏の二つの射が与えられると、その合成はよく定義され、この圏の射となる。
\end{remark}

\begin{remark}
\label{morphisms-remark-pseudo-morphisms}
定義 \ref{morphisms-definition-rational-map} には、射 $U \to Y$ で、スキーム論的に稠密な開部分スキーム
$U \subset X$ 上で定義されたものだけを考える変種がある。
補題 \ref{morphisms-lemma-intersection-scheme-theoretically-dense} により、これは同値関係を与える。
これらの同値類を {\it $X$ から $Y$ への擬射} という。
$X$ が被約ならば二つの概念は一致する。
\end{remark}





\section{双有理射}
\label{morphisms-section-birational}

\noindent
多様体の双有理写像が、標的のある開部分集合上では同型であるという概念に
馴染みがあるかもしれない。しかし一般には、双有理射はどの空でない開集合上でも
同型でないことがある。例 \ref{morphisms-example-birational-not-iso-over-open} を参照せよ。
以下が正式な定義である。

\begin{definition}
\label{morphisms-definition-birational}
\begin{reference}
\cite[(2.2.9)]{EGA1}
\end{reference}
$X$, $Y$ をスキームとする。$X$ と $Y$ は有限個の既約成分をもつと仮定する。
射 $f : X \to Y$ が次を満たすとき、この射を {\it 双有理} であるという。
\begin{enumerate}
\item $f$ は、$X$ の既約成分の生成点の集合と $Y$ の既約成分の
生成点の集合との間に全単射を誘導する。
\item 生成点 $\eta \in X$ で $X$ の既約成分のものすべてに対して、局所環の写像
$\mathcal{O}_{Y, f(\eta)} \to \mathcal{O}_{X, \eta}$
は同型である。
\end{enumerate}
\end{definition}

\noindent
以下で、双有理射の生成点上のファイバーが一点集合であることを見る。
さらに、実際に現れるほとんどの場合には、既約スキーム $X$ と $Y$ の間に
双有理射が存在すれば、$X$ と $Y$ は双有理スキームとなることを見る。

\begin{lemma}
\label{morphisms-lemma-birational-dominant}
$f : X \to Y$ を、有限個の既約成分をもつスキームの射とする。
$f$ が双有理ならば、$f$ は支配的である。
\end{lemma}

\begin{proof}
補題 \ref{morphisms-lemma-generic-points-in-image-dominant} と定義から従う。
\end{proof}

\begin{lemma}
\label{morphisms-lemma-birational-generic-fibres}
$f : X \to Y$ を、有限個の既約成分をもつスキームの双有理射とする。
$y \in Y$ がある既約成分の生成点ならば、基底変換
$X \times_Y \Spec(\mathcal{O}_{Y, y}) \to \Spec(\mathcal{O}_{Y, y})$
は同型である。
\end{lemma}

\begin{proof}
$Y = \Spec(B)$ はアフィンかつ既約であるとしてよい。
このとき $X$ も既約である。任意の空でないアフィン開集合 $U \subset X$ に対して
結果を証明すれば、$X$ に対する結果が従う（短い議論は省略する）。
従って $X$ もアフィンであるとしてよく、$X = \Spec(A)$ と書く。
$y \in Y$ が極小素イデアル $\mathfrak q \subset B$ に対応するとする。
仮定により、$A$ は唯一の極小素イデアル $\mathfrak p$ をもち、これは
$\mathfrak q$ の上にあり、$B_\mathfrak q \to A_\mathfrak p$ は同型である。
従って $A_\mathfrak q \to \kappa(\mathfrak p)$ は全射であり、
$\mathfrak p A_\mathfrak q$ は極大イデアルである。一方、
$\mathfrak p A_\mathfrak q$ は $A_\mathfrak q$ の唯一の極小素イデアルである。
よって $\mathfrak p A_\mathfrak q$ は $A_\mathfrak q$ の唯一の素イデアルであり、
$A_\mathfrak q = A_\mathfrak p$ である。$A_\mathfrak q = A \otimes_B B_\mathfrak q$
であるから、補題が従う。
\end{proof}

\begin{example}
\label{morphisms-example-birational-not-iso-over-open}
標的のどの開集合上でも同型でない双有理射の例を二つ挙げる。

\medskip\noindent
第一の例。$k$ を無限体とする。$A = k[x]$ とおく。さらに
$B = k[x, \{y_{\alpha}\}_{\alpha \in k}]/
((x-\alpha)y_\alpha, y_\alpha y_\beta)$
とおく。包含 $A \subset B$ とレトラクション $B \to A$ が存在し、
後者はすべての $y_\alpha$ を零にする。
射 $\Spec(A) \to \Spec(B)$ と射 $\Spec(B) \to \Spec(A)$ は
どちらも双有理だが、どの開集合上でも同型ではない。

\medskip\noindent
第二の例。$A$ を整域とする。$S \subset A$ を $0$ を含まない
乗法的部分集合とする。$B = S^{-1}A$ とおけば、射
$f : \Spec(B) \to \Spec(A)$ は双有理である。
開集合 $U$ が $\Spec(A)$ に存在して $f^{-1}(U) \to U$ が同型ならば、
ある $a \in A$ が存在して、$S$ のすべての元が主局所化 $A_a$ で可逆になる。
$A = \mathbf{Z}$ とし、$S$ を奇整数の集合とすれば反例を得る。
\end{example}

\begin{lemma}
\label{morphisms-lemma-birational-birational}
$f : X \to Y$ を、基底スキーム $S$ 上で有限個の既約成分をもつスキームの
双有理射とする。次の条件のいずれかが満たされると仮定する。
\begin{enumerate}
\item $f$ は局所有限型であり、$Y$ は被約である。
\item $f$ は局所有限表示である。
\end{enumerate}
このとき稠密開集合 $U \subset X$ と $V \subset Y$ で、
$f(U) \subset V$ かつ $f|_U : U \to V$ が同型となるものが存在する。
特に $X$ と $Y$ が既約ならば、$X$ と $Y$ は $S$-双有理である。
\end{lemma}

\begin{proof}
$X$ と $Y$ が既約である場合へ直ちに帰着できるが、その議論は省略する。
さらに縮小することにより、$X$ と $Y$ はアフィンであるとしてよい。
$X = \Spec(A)$, $Y = \Spec(B)$ と書く。
仮定により $A$, resp.\ $B$ は唯一の極小素イデアル $\mathfrak p$,
resp.\ $\mathfrak q$ をもち、素イデアル $\mathfrak p$ は
$\mathfrak q$ の上にあり、$B_\mathfrak q = A_\mathfrak p$ である。
補題 \ref{morphisms-lemma-birational-generic-fibres} により、
$B_\mathfrak q = A_\mathfrak q = A_\mathfrak p$ である。

\medskip\noindent
$B \to A$ が有限型であるとし、$A = B[x_1, \ldots, x_n]$ と書く。
$b_i \in B$ と $g_i \in B \setminus \mathfrak q$ で、$b_i/g_i$ が
$x_i$ の $A_\mathfrak q$ における像へ写るものが存在する。
従って $b_i  - g_ix_i$ は $A_{g_i'}$ で零へ写る。ただし、ある
$g_i' \in B \setminus \mathfrak q$ を選んだ。
$g = \prod g_i g'_i$ とおくと、$B_g \to A_g$ は全射である。
さらに $Y$ が被約ならば、写像 $B_g \to B_\mathfrak q$ は単射なので、
$B_g \to A_g$ も単射である。これで (1) が証明された。

\medskip\noindent
(2) の証明。前段落の議論により、$B \to A$ は全射であるとしてよい。
$f$ は局所有限表示だから、核 $J \subset B$ は有限生成イデアルである。
$J = (b_1, \ldots, b_r)$ と書く。$B_\mathfrak q = A_\mathfrak q$ なので、
$g_i \in B \setminus \mathfrak q$ で $g_i b_i = 0$ を満たすものが存在する。
$g = \prod g_i$ とおけば、$B_g \to A_g$ は同型である。
\end{proof}

\begin{lemma}
\label{morphisms-lemma-criterion-birational-finite-presentation}
$S$ をスキームとする。$X$ と $Y$ を、$S$ 上局所有限表示な既約スキームとする。
$x \in X$ と $y \in Y$ を生成点とする。次は同値である。
\begin{enumerate}
\item $X$ と $Y$ は $S$-双有理である。
\item $X$ と $Y$ には、$S$-同型な空でない開集合が存在する。
\item $x$ と $y$ は同じ点 $s$ へ写り、その点は $S$ に属し、
$\mathcal{O}_{X, x}$ と $\mathcal{O}_{Y, y}$ は
$\mathcal{O}_{S, s}$-代数として同型である。
\end{enumerate}
\end{lemma}

\begin{proof}
(1) と (2) の同値性は補題 \ref{morphisms-lemma-birational-integral} で見た。
(2) が (3) を含意することは直ちに分かる。
最後に (3) が成り立つと仮定し、(1) を証明する。
補題 \ref{morphisms-lemma-rational-map-finite-presentation} により、有理写像
$f : U \to Y$ で、$x \in U$ を $y$ へ送り、与えられた同型
$\mathcal{O}_{Y, y} \cong \mathcal{O}_{X, x}$ を誘導するものが存在する。
従って $f$ は双有理射であり、
補題 \ref{morphisms-lemma-birational-birational} により空でない開集合上の
同型を誘導する。これで証明が完了する。
\end{proof}

\begin{lemma}
\label{morphisms-lemma-common-open}
$S$ をスキームとする。$X$ と $Y$ を、$S$ 上局所有限型な整スキームとする。
$x \in X$ と $y \in Y$ を生成点とする。次は同値である。
\begin{enumerate}
\item $X$ と $Y$ は $S$-双有理である。
\item $X$ と $Y$ には、$S$-同型な空でない開集合が存在する。
\item $x$ と $y$ は同じ点 $s \in S$ へ写り、
$\kappa(x) \cong \kappa(y)$ は $\kappa(s)$-拡大としての同型である。
\end{enumerate}
\end{lemma}

\begin{proof}
(1) と (2) の同値性は補題 \ref{morphisms-lemma-birational-integral} で見た。
(2) が (3) を含意することは直ちに分かる。
最後に (3) が成り立つと仮定し、(1) を証明する。
代数、補題 \ref{algebra-lemma-minimal-prime-reduced-ring} により、
$\mathcal{O}_{X, x} = \kappa(x)$ かつ
$\mathcal{O}_{Y, y} = \kappa(y)$ である。
補題 \ref{morphisms-lemma-rational-map-finite-presentation} により、有理写像
$f : U \to Y$ で、$x \in U$ を $y$ へ送り、与えられた同型
$\mathcal{O}_{Y, y} \cong \mathcal{O}_{X, x}$ を誘導するものが存在する。
従って $f$ は双有理射であり、
補題 \ref{morphisms-lemma-birational-birational} により空でない開集合上の
同型を誘導する。これで証明が完了する。
\end{proof}





\section{生成的有限射}
\label{morphisms-section-generically-finite}

\noindent
この節では、局所有限型であって、ある意味で「生成的有限」である
スキーム間の写像を特徴づける。

\begin{lemma}
\label{morphisms-lemma-generically-finite}
$X$, $Y$ をスキームとする。
$f : X \to Y$ を局所有限型とする。
$\eta \in Y$ を $Y$ の既約成分の生成点とする。次は同値である。
\begin{enumerate}
\item 集合 $f^{-1}(\{\eta\})$ は有限である。
\item アフィン開集合 $U_i \subset X$, $i = 1, \ldots, n$ と
$V \subset Y$ で、$f(U_i) \subset V$,
$\eta \in V$, $f^{-1}(\{\eta\}) \subset \bigcup U_i$ を満たし、
各 $f|_{U_i} : U_i \to V$ が有限となるものが存在する。
\end{enumerate}
$f$ が準分離的ならば、これらは次とも同値である。
\begin{enumerate}
\item[(3)] アフィン開集合 $V \subset Y$ と
$U \subset X$ で、$f(U) \subset V$,
$\eta \in V$, $f^{-1}(\{\eta\}) \subset U$ を満たし、
$f|_U : U \to V$ が有限となるものが存在する。
\end{enumerate}
$f$ が準コンパクトかつ準分離的ならば、これらは次とも同値である。
\begin{enumerate}
\item[(4)] アフィン開集合 $V \subset Y$, $\eta \in V$ で、
$f^{-1}(V) \to V$ が有限となるものが存在する。
\end{enumerate}
\end{lemma}

\begin{proof}
問題は基底上局所的である。従って $Y$ を $\eta$ のアフィン近傍で
置き換えることができ、以下の証明を通じて $Y$ はアフィンであると仮定する。
$Y = \Spec(R)$ と書く。

\medskip\noindent
(2) が (1) を含意することは明らかである。
$f^{-1}(\{\eta\}) = \{\xi_1, \ldots, \xi_n\}$ が有限であると仮定する。
アフィン開集合 $U_i \subset X$ で $\xi_i \in U_i$ を満たすものを選ぶ。
代数、補題 \ref{algebra-lemma-generically-finite} により、
$Y$ を $Y$ の標準開集合で置き換えた後、各射 $U_i \to Y$ は有限となる。
言い換えれば (2) が成り立つ。

\medskip\noindent
(3) が (1) を含意することは明らかである。$f$ は準分離的であり、(1) が
成り立つと仮定する。$f^{-1}(\{\eta\}) = \{\xi_1, \ldots, \xi_n\}$ と書く。
補題 \ref{morphisms-lemma-finite-fibre} により、$\xi_i$ の間に特殊化は存在しない。
各 $\xi_i$ は生成点 $\eta$ へ写り、これは $Y$ の既約成分のものなので、
非自明な特殊化 $\xi \leadsto \xi_i$ は $X$ に存在しえない
（その場合 $\xi$ も $\eta$ へ写るからである）。
従って各 $\xi_i$ は $X$ の既約成分の生成点である。
$Y$ はアフィンで $f$ は準分離的だから、$X$ は準分離的である。
性質、補題 \ref{properties-lemma-maximal-points-affine} により、
アフィン開集合 $U \subset X$ で各 $\xi_i$ を含むものを見つけられる。
代数、補題 \ref{algebra-lemma-generically-finite} により、
$Y$ を $Y$ の標準開集合で置き換えた後、射 $U \to Y$ は有限となる。
言い換えれば (3) が成り立つ。

\medskip\noindent
(4) は、$f$ に追加の仮定を置かずに (1) -- (3) のすべてを含意する。
$f$ は準コンパクトかつ準分離的であると仮定する。
同値な条件 (1) -- (3) が (4) を含意することを示せばよい。
$U$, $V$ を (3) のようにとる。$Y$ を $V$ で置き換える。
$f$ は準コンパクトで、$Y$ は（アフィンなので）準コンパクトだから、
$X$ は準コンパクトである。従って $Z = X \setminus U$ は準コンパクトであり、
射 $f|_Z : Z \to Y$ も準コンパクトである。$Z$ の構成により
$\eta \not \in f(Z)$ である。従って補題
\ref{morphisms-lemma-quasi-compact-generic-point-not-in-image} により、
$V'$ で、$\eta$ の $Y$ におけるアフィン開近傍となり、
$f^{-1}(V') \cap Z = \emptyset$ を満たすものが存在する。
このとき $f^{-1}(V') \subset U$ であり、これは
$f^{-1}(V') \to V'$ が有限であることを意味する。
\end{proof}

\begin{example}
\label{morphisms-example-counter-generically-finite}
$A = \prod_{n \in \mathbf{N}} \mathbf{F}_2$ とおく。
$A$ のすべての元は冪等元である。従ってすべての素イデアルは極大であり、
剰余体は $\mathbf{F}_2$ である。
よって $X = \Spec(A)$ の位相は全不連結かつ準コンパクトである。
射影写像 $A \to \mathbf{F}_2$ は $\Spec(A)$ の開点を定める。
$X$ は準コンパクトだから、$X$ のすべての点が開であることはありえない。
$x \in X$ を、開でない閉点とする。このときスキーム $Y$ を、
$X$ の二つのコピーを $X \setminus \{x\}$ に沿って貼り合わせて作れる。
言い換えれば、これは $X$ で $x$ を二重化したものである。スキーム、例
\ref{schemes-example-affine-space-zero-doubled} と比較せよ。
射 $f : Y \to X$ は準コンパクトかつ有限型で、有限ファイバーをもつが、
準分離的ではない。
点 $x \in X$ は $X$ の既約成分の生成点である
（$X$ は全不連結だからである）。しかし補題
\ref{morphisms-lemma-generically-finite} の性質 (3) と (4) は成り立たない。
その理由は、任意の開近傍 $x \in U \subset X$ に対して、逆像
$f^{-1}(U)$ はアフィンでないからである。実際、$f^{-1}(U)$ 上の関数は
$x$ 上にある二つの点を分離できない（証明は省略する。これはよい演習である）。
従って補題の (3) と (4) では、$f$ が準分離的であるという条件が必要である。
\end{example}

\begin{remark}
\label{morphisms-remark-quasi-finite-finite-over-dense-open}
補題 \ref{morphisms-lemma-generically-finite} に代わるものとして、
準有限射は標的の稠密開集合上で有限であるという命題がある。
これは射についてさらに、補題
\ref{more-morphisms-lemma-quasi-finite-finite-over-dense-open}
で示される。
\end{remark}

\begin{lemma}
\label{morphisms-lemma-quasi-finiteness-over-generic-point}
$X$, $Y$ をスキームとする。$f : X \to Y$ を局所有限型とする。
$X^0$, resp.\ $Y^0$ で、それぞれ $X$, resp.\ $Y$ の既約成分の
生成点の集合を表す。$\eta \in Y^0$ とする。次は同値である。
\begin{enumerate}
\item $f^{-1}(\{\eta\}) \subset X^0$ である。
\item $f$ は $\eta$ 上にあるすべての点で準有限である。
\item $f$ はすべての $\xi \in X^0$ で準有限であり、これらは $\eta$ 上にある。
\end{enumerate}
\end{lemma}

\begin{proof}
条件 (1) により、ファイバー $X_\eta$ の点の間に特殊化は存在しない。
従って補題 \ref{morphisms-lemma-quasi-finite-at-point-characterize} により (2) が成り立つ。
(2) $\Rightarrow$ (3) は直ちに分かる。
$\eta$ は $Y$ の生成点だから、$X_\eta$ の生成点は $X$ の生成点である。
従って (3) と補題 \ref{morphisms-lemma-quasi-finite-at-point-characterize} により、
$X_\eta$ の生成点は閉でもある。よって $X_\eta$ のすべての点は生成点であり、
(1) が成り立つ。
\end{proof}

\begin{lemma}
\label{morphisms-lemma-finite-over-dense-open}
$X$, $Y$ をスキームとする。$f : X \to Y$ を局所有限型とする。
$X^0$, resp.\ $Y^0$ で、それぞれ $X$, resp.\ $Y$ の既約成分の
生成点の集合を表す。次を仮定する。
\begin{enumerate}
\item $X^0$ と $Y^0$ は有限であり、$f^{-1}(Y^0) = X^0$ である。
\item $f$ は準コンパクトであるか、または $f$ は分離的である。
\end{enumerate}
このとき稠密開集合 $V \subset Y$ で、$f^{-1}(V) \to V$ が有限となるものが存在する。
\end{lemma}

\begin{proof}
$Y$ は有限個の既約成分をもつので、その既約成分の非交和である稠密開集合を
見つけられる。従って $Y$ は生成点 $\eta$ をもつ既約アフィンスキームであるとしてよい。
$\eta$ 上のファイバーは有限である。実際、$X^0$ は有限である。

\medskip\noindent
$f$ は分離的で、$Y$ は既約アフィンであると仮定する。
補題 \ref{morphisms-lemma-generically-finite} の (3) のように $V \subset Y$ と
$U \subset X$ を選ぶ。$f|_U : U \to V$ は有限だから、
$U \subset f^{-1}(V)$ は開かつ閉である（補題
\ref{morphisms-lemma-image-proper-scheme-closed} と \ref{morphisms-lemma-finite-proper}）。
従って $f^{-1}(V) = U \amalg W$ である。ただし $W$ は $X$ の
ある開部分スキームである。しかし $U$ は $X$ のすべての生成点を
含むので、望む通り $W = \emptyset$ である。

\medskip\noindent
$f$ は準コンパクトで、$Y$ は既約アフィンであると仮定する。
このとき $X$ は準コンパクトだから、分離的な稠密開部分スキーム
$U \subset X$ が存在する
（性質、補題 \ref{properties-lemma-quasi-compact-dense-open-separated}）。
生成点の集合 $X^0$ は有限だから、$X^0 \subset U$ である。
従って $\eta \not \in f(X \setminus U)$ である。
$X \setminus U \to Y$ は準コンパクトだから、空でない開集合
$V \subset Y$ で $f^{-1}(V) \subset U$ を満たすものが存在する。
補題 \ref{morphisms-lemma-quasi-compact-dominant} を参照せよ。
$X$ を $f^{-1}(V)$ で、$Y$ を $V$ で置き換えると、
前段落で扱った分離的な場合へ帰着する。
\end{proof}

\begin{lemma}
\label{morphisms-lemma-birational-isomorphism-over-dense-open}
$X$, $Y$ をスキームとする。$f : X \to Y$ を、
有限個の既約成分をもつスキーム間の双有理射とする。次を仮定する。
\begin{enumerate}
\item $f$ は準コンパクトであるか、または $f$ は分離的である。
\item $f$ は局所有限型かつ $Y$ は被約であるか、または
$f$ は局所有限表示である。
\end{enumerate}
このとき稠密開集合 $V \subset Y$ で、$f^{-1}(V) \to V$ が同型となるものが存在する。
\end{lemma}

\begin{proof}
補題 \ref{morphisms-lemma-finite-over-dense-open} により、$f$ は有限であるとしてよい。
$Y$ は有限個の既約成分をもつので、その既約成分の非交和である稠密開集合を
見つけられる。従って $Y$ は既約であるとしてよい。
補題 \ref{morphisms-lemma-birational-birational} により、空でない開集合 $U \subset X$ で
$f|_U : U \to Y$ が開埋め込みとなるものが見つかる。
$f$ は有限なので閉である部分集合 $f(X \setminus U)$ を
$Y$ から除くと、$f$ は同型となる。
\end{proof}

\begin{lemma}
\label{morphisms-lemma-finite-degree}
$X$, $Y$ を整スキームとする。
$f : X \to Y$ を局所有限型とする。
$f$ は支配的であると仮定する。次は同値である。
\begin{enumerate}
\item 拡大 $R(Y) \subset R(X)$ の超越次数は $0$ である。
\item 拡大 $R(Y) \subset R(X)$ は有限である。
\item 空でないアフィン開集合 $U \subset X$ と
$V \subset Y$ で、$f(U) \subset V$ かつ
$f|_U : U \to V$ が有限となるものが存在する。
\item $X$ の生成点は、$X$ の点で $Y$ の生成点へ写る唯一のものである。
\end{enumerate}
$f$ が分離的であるか、または $f$ が準コンパクトならば、これらは次とも同値である。
\begin{enumerate}
\item[(5)] 空でないアフィン開集合 $V \subset Y$ で、
$f^{-1}(V) \to V$ が有限となるものが存在する。
\end{enumerate}
\end{lemma}

\begin{proof}
任意のアフィン開集合 $\Spec(A) = U \subset X$ と
$\Spec(R) = V \subset Y$ で $f(U) \subset V$ を満たすものを選ぶ。
定義により $R$ と $A$ は整域である。環準同型 $R \to A$ は有限型である
（補題 \ref{morphisms-lemma-locally-finite-type-characterize}）。
補題 \ref{morphisms-lemma-dominant-finite-number-irreducible-components} により、
$X$ の生成点は $Y$ の生成点へ写るので、$R \to A$ は単射である。
$K = R(Y)$ を $R$ の分数体、$L = R(X)$ を $A$ の分数体とする。
このとき $L/K$ は有限生成体拡大である。従って (1) と (2) は同値である。

\medskip\noindent
(2) が成り立つと仮定する。$x_1, \ldots, x_n \in A$ を
$A$ の $R$ 上の生成元とする。仮定により、零でない多項式
$P_i(X) \in R[X]$ で $P_i(x_i) = 0$ を満たすものが存在する。$f_i \in R$ を
$P_i$ の最高次係数とする。このとき
$R_{f_1 \ldots f_n} \to A_{f_1 \ldots f_n}$ は有限、すなわち (3) が成り立つ。
(3) が (2) を含意することに注意せよ。従って (1), (2), (3) はすべて同値である。

\medskip\noindent
$\eta$ を $X$ の生成点、$\eta' \in Y$ を $Y$ の生成点とする。
(4) を仮定する。このとき $\dim_\eta(X_{\eta'}) = 0$ であり、
補題 \ref{morphisms-lemma-dimension-fibre-at-a-point} により
$R(X) = \kappa(\eta)$ は超越次数 $0$ を $R(Y) = \kappa(\eta')$ 上にもつ。
言い換えれば (1) が成り立つ。同値な条件 (1), (2), (3) を仮定する。
$x \in X$ が $\eta'$ へ写る点であるとする。
$x$ は $\eta$ の特殊化だから、包含
$R(Y) \subset \mathcal{O}_{X, x} \subset R(X)$ を得る。
これは $\mathcal{O}_{X, x}$ が体であることを含意する。代数、補題
\ref{algebra-lemma-integral-over-field} を参照せよ。
従って $x = \eta$ である。よって (1) -- (4) はすべて同値である。

\medskip\noindent
(5) が、$f$ に追加の仮定を置かずに (3) を含意することは明らかである。
残るのは、$f$ が分離的または準コンパクトならば、同値な条件
(1) -- (4) が (5) を含意することを証明することである。
これは補題 \ref{morphisms-lemma-finite-over-dense-open} から従う。
\end{proof}

\begin{definition}
\label{morphisms-definition-degree}
$X$ と $Y$ を整スキームとする。
$f : X \to Y$ を局所有限型かつ支配的とする。
$[R(X) : R(Y)] < \infty$、または補題 \ref{morphisms-lemma-finite-degree} の
同値な条件 (1) -- (4) のいずれかを仮定する。このとき正の整数
$$
\deg(X/Y) = [R(X) : R(Y)]
$$
を {\it $X$ の $Y$ 上の次数} という。
\end{definition}

\noindent
この概念は射 $f : X \to Y$ に拡張できる。ただし、(a) $Y$ は
生成点 $\eta$ をもつ整スキームであり、(b) $f$ は局所有限型であり、
(c) $f^{-1}(\{\eta\})$ は有限であるとする。この場合、次で定義できる。
$$
\deg(X/Y)
=
\sum\nolimits_{\xi \in X, \ f(\xi) = \eta}
\dim_{R(Y)} (\mathcal{O}_{X, \xi}).
$$
すなわち、$R(Y) = \kappa(\eta) = \mathcal{O}_{Y, \eta}$ であること
（補題 \ref{morphisms-lemma-integral-scheme-rational-functions}）を用いると、
上の次元は補題 \ref{morphisms-lemma-generically-finite} により有限である。
しかし多くの応用では、先に与えた定義が適切である。

\begin{lemma}
\label{morphisms-lemma-degree-composition}
$X$, $Y$, $Z$ を整スキームとする。
$f : X \to Y$ と $g : Y \to Z$ を局所有限型な支配的射とする。
$[R(X) : R(Y)] < \infty$ かつ $[R(Y) : R(Z)] < \infty$ と仮定する。このとき
$$
\deg(X/Z) = \deg(X/Y) \deg(Y/Z).
$$
\end{lemma}

\begin{proof}
これは体の有限拡大の塔における次数の乗法性から従う。
体、補題 \ref{fields-lemma-multiplicativity-degrees} を参照せよ。
\end{proof}

\begin{remark}
\label{morphisms-remark-definition-generically-finite}
$f : X \to Y$ を局所有限型なスキームの射とする。
{\it 生成的有限} 射を定義するために使える性質が（少なくとも）二つある。
これらは、その性質を源上局所的にしたいか標的上局所的にしたいかに対応する。
\begin{enumerate}
\item （標的上局所的。Ravi Vakil の提案による。）
$Y$ のすべての準コンパクト開集合は有限個の既約成分をもつと仮定する
（例えば $Y$ が局所 Noether 的ならば成り立つ）。
各生成点の逆像が有限であることを要請する。補題
\ref{morphisms-lemma-generically-finite} を参照せよ。
\item （源上局所的。）稠密開集合 $U \subset X$ で、
$U \to Y$ が局所準有限となるものが存在することを要請する。
\end{enumerate}
(1) の場合、$Y$ に対する補助条件を置かずに要請を定式化できるが、
一般のスキームに対する正しい概念にはならないと思われる。
上の形の性質 (2) は、生成点上のファイバーが有限であることを含意しない。
しかし $f$ が準コンパクトで、$Y$ が (1) のようなものならば含意する。
\end{remark}

\begin{definition}
\label{morphisms-definition-modification}
$X$ を整スキームとする。{\it $X$ の修正} とは、
双有理固有射 $f : X' \to X$ で $X'$ が整であるもののことである。
\end{definition}

\noindent
$f : X' \to X$ を定義のような修正とする。補題
\ref{morphisms-lemma-finite-degree} により、空でない $U \subset X$ で
$f^{-1}(U) \to U$ が有限となるものが存在する。生成的平坦性
（命題 \ref{morphisms-proposition-generic-flatness}）により、
$f^{-1}(U) \to U$ は平坦かつ有限表示であるとしてよい。
従って $f^{-1}(U) \to U$ は有限局所自由である
（補題 \ref{morphisms-lemma-finite-flat}）。$f$ は双有理だから、
$X'$ の $X$ 上の次数は $1$ である。
従って $f^{-1}(U) \to U$ は次数 $1$ の有限局所自由射、すなわち同型である。
よって修正を、整スキームの固有射 $f : X' \to X$ で、
$f^{-1}(U) \to U$ が同型となるような空でない開集合
$U \subset X$ が存在するものと {\it 再定義} できる。

\begin{definition}
\label{morphisms-definition-alteration}
\begin{reference}
\cite[Definition 2.20]{alterations}
\end{reference}
$X$ を整スキームとする。{\it $X$ のオルタレーション} とは、
固有支配的射 $f : Y \to X$ で、$Y$ が整であり、
$f^{-1}(U) \to U$ が有限となるような空でない開集合
$U \subset X$ が存在するもののことである。
\end{definition}

\noindent
これは \cite{alterations} で与えられた定義であるが、ここでは
$X$ と $Y$ が Noether 的であることを要請しない。
上と同様に論じると、オルタレーションは、関数体の有限拡大を誘導する
整スキームの固有支配的射 $f : Y \to X$、すなわち補題
\ref{morphisms-lemma-finite-degree} の同値な条件を満たす射であることが分かる。



\section{次元公式}
\label{morphisms-section-dimension-formula}

\noindent
Noether スキーム間の射については、局所環の次元に関してもう少し詳しく
述べることができる。以下は重要な（しかも証明はそれほど難しくない）結果である。
$R(X)$ は整スキーム $X$ の関数体を表すことを思い出そう。

\begin{lemma}
\label{morphisms-lemma-dimension-formula}
$S$ をスキームとする。
$f : X \to S$ をスキームの射とする。
$x \in X$ とし、$s = f(x)$ とおく。次を仮定する。
\begin{enumerate}
\item $S$ は局所 Noether 的である。
\item $f$ は局所有限型である。
\item $X$ と $S$ は整である。
\item $f$ は支配的である。
\end{enumerate}
このとき
\begin{equation}
\label{morphisms-equation-dimension-formula}
\dim(\mathcal{O}_{X, x})
\leq
\dim(\mathcal{O}_{S, s}) + \text{trdeg}_{R(S)}R(X)
- \text{trdeg}_{\kappa(s)} \kappa(x).
\end{equation}
さらに、$S$ が普遍鎖状ならば等号が成り立つ。
\end{lemma}

\begin{proof}
対応する代数的な主張は、代数、補題
\ref{algebra-lemma-dimension-formula} である。
\end{proof}

\begin{lemma}
\label{morphisms-lemma-dimension-formula-general}
$S$ をスキームとする。$f : X \to S$ をスキームの射とする。
$x \in X$ とし、$s = f(x)$ とおく。$S$ は局所 Noether 的であり、
$f$ は局所有限型であると仮定する。このとき
\begin{equation}
\label{morphisms-equation-dimension-formula-general}
\dim(\mathcal{O}_{X, x})
\leq
\dim(\mathcal{O}_{S, s}) + E - \text{trdeg}_{\kappa(s)} \kappa(x).
\end{equation}
ここで $E$ は $\text{trdeg}_{\kappa(f(\xi))}(\kappa(\xi))$ の最大値である。
ただし $\xi$ は $X$ の既約成分の生成点を動き、その成分は $x$ を含む。
\end{lemma}

\begin{proof}
$X_1, \ldots, X_n$ を、$X$ の既約成分で $x$ を含むものに被約誘導スキーム構造を
入れたものとする。これらは極小素イデアル $\mathfrak q_i$（
$\mathcal{O}_{X, x}$ のもの）に対応するので、有限個しかない
（スキーム、補題 \ref{schemes-lemma-specialize-points} および
代数、補題 \ref{algebra-lemma-Noetherian-irreducible-components}）。
従って
$\dim(\mathcal{O}_{X, x}) = \max \dim(\mathcal{O}_{X, x}/\mathfrak q_i)
= \max \dim(\mathcal{O}_{X_i, x})$ である。
ここで現れる $\xi$、すなわち $E$ の定義に現れるものは、
ちょうど生成点 $\xi_i \in X_i$ である。
$Z_i = \overline{\{f(\xi_i)\}} \subset S$ に被約誘導スキーム構造を入れる。
合成 $X_i \to X \to S$ は $Z_i$ を経由する
（スキーム、補題 \ref{schemes-lemma-map-into-reduction}）。従って次元公式
（補題 \ref{morphisms-lemma-dimension-formula}）を適用すると
$\dim(\mathcal{O}_{X_i, x}) \leq \dim(\mathcal{O}_{Z_i, x}) +
\text{trdeg}_{\kappa(f(\xi))}(\kappa(\xi)) -
\text{trdeg}_{\kappa(s)} \kappa(x)$ を得る。以上を合わせれば補題が従う。
\end{proof}

\noindent
一つの応用は、次元関数を備えた普遍鎖状スキーム上の任意の有限型スキームに
次元関数を構成することである。次元関数の定義については、位相、定義
\ref{topology-definition-dimension-function} を参照せよ。

\begin{lemma}
\label{morphisms-lemma-dimension-function-propagates}
$S$ を局所 Noether 的かつ普遍鎖状なスキームとする。
$\delta : S \to \mathbf{Z}$ を次元関数とする。
$f : X \to S$ をスキームの射とする。
$f$ は局所有限型であると仮定する。このとき写像
\begin{align*}
\delta = \delta_{X/S} : X & \longrightarrow \mathbf{Z} \\
x & \longmapsto \delta(f(x)) + \text{trdeg}_{\kappa(f(x))} \kappa(x)
\end{align*}
は $X$ 上の次元関数である。
\end{lemma}

\begin{proof}
$f : X \to S$ を局所有限型とする。
$x \leadsto y$, $x \not = y$ を $X$ における特殊化とする。
$\delta_{X/S}(x) > \delta_{X/S}(y)$ であること、および
$\delta_{X/S}(x) = \delta_{X/S}(y) + 1$ であることを、$y$ が $x$ の
直後の特殊化である場合に示さなければならない。

\medskip\noindent
アフィン開集合 $V \subset S$ で $y$ の像を含むものを選び、アフィン開集合
$U \subset X$ で $V$ へ写り $y$ を含むものを選ぶ。$X$ を $U$ で、$S$ を $V$ で
置き換えてよいことは明らかである。従って $X = \Spec(A)$、
$S = \Spec(R)$ であり、$f$ は環準同型 $R \to A$ により与えられるとしてよい。
環 $R$ は普遍鎖状であり（補題 \ref{morphisms-lemma-universally-catenary-local}）、
準同型 $R \to A$ は有限型である
（補題 \ref{morphisms-lemma-locally-finite-type-characterize}）。

\medskip\noindent
素イデアル $\mathfrak q \subset A$ を点 $x$ に対応するものとし、素イデアル
$\mathfrak p \subset R$ を $f(x)$ に対応するものとする。制限 $\delta'$、
すなわち $\delta$ の $S' = \Spec(R/\mathfrak p) \subset S$ への制限は
次元関数である。
環 $R/\mathfrak p$ は普遍鎖状である。$\delta_{X/S}$ の
$X' = \Spec(A/\mathfrak q)$ への制限は関数 $\delta_{X'/S'}$ に明らかに等しい。
この関数は次元関数 $\delta'$ を用いて構成したものである。従って上の仮定に加えて、
$R \subset A$ は整域である、言い換えれば $X$ と $S$ は整スキームであり、
$x$ は $X$ の生成点、$f(x)$ は $S$ の生成点であるとしてよい。

\medskip\noindent
$\mathcal{O}_{X, x} = R(X)$ であることに注意する。また $x \leadsto y$、
$x \not = y$ なので、$\mathcal{O}_{X, y}$ のスペクトルは少なくとも二点をもつ
（スキーム、補題 \ref{schemes-lemma-specialize-points}）。従って
$\dim(\mathcal{O}_{X, y}) > 0$ である。
$y$ が $x$ の直後の特殊化ならば、
$\Spec(\mathcal{O}_{X, y}) = \{x, y\}$ かつ
$\dim(\mathcal{O}_{X, y}) = 1$ である。

\medskip\noindent
$s = f(x)$、$t = f(y)$ と書く。計算すると
\begin{align*}
\delta_{X/S}(x) - \delta_{X/S}(y)
& =
\delta(s) + \text{trdeg}_{\kappa(s)} \kappa(x)
- \delta(t) - \text{trdeg}_{\kappa(t)} \kappa(y) \\
& =
\delta(s) - \delta(t) +
\text{trdeg}_{R(S)} R(X) - \text{trdeg}_{\kappa(t)} \kappa(y) \\
& =
\delta(s) - \delta(t) + \dim(\mathcal{O}_{X, y})
- \dim(\mathcal{O}_{S, t})
\end{align*}
を得る。最後の段階では (\ref{morphisms-equation-dimension-formula}) の等号を用いた。
$\delta$ はスキーム $S$ 上の次元関数であり、$s \in S$ は生成点だから、差
$\delta(s) - \delta(t)$ は、位相、補題
\ref{topology-lemma-dimension-function-catenary} により
$\text{codim}(\overline{\{t\}}, S)$ に等しい。これは性質、補題
\ref{properties-lemma-codimension-local-ring} により
$\dim(\mathcal{O}_{S, t})$ に等しい。従って
$$
\delta_{X/S}(x) - \delta_{X/S}(y) = \dim(\mathcal{O}_{X, y})
$$
となり、$\dim(\mathcal{O}_{X, y})$ について上で述べたことから補題が従う。
\end{proof}

\noindent
次元公式のもう一つの応用は、「オルタレーション」（後で定義する）によって
次元が変わらないことである。

\begin{lemma}
\label{morphisms-lemma-alteration-dimension}
$f : X \to Y$ をスキームの射とする。次を仮定する。
\begin{enumerate}
\item $Y$ は局所 Noether 的である。
\item $X$ と $Y$ は整スキームである。
\item $f$ は支配的である。
\item $f$ は局所有限型である。
\end{enumerate}
このとき
$$
\dim(X) \leq \dim(Y) + \text{trdeg}_{R(Y)} R(X).
$$
$f$ が閉である\footnote{例えば $f$ が固有である場合。定義
\ref{morphisms-definition-proper} を参照せよ。}ならば等号が成り立つ。
\end{lemma}

\begin{proof}
$f : X \to Y$ を補題のような射とする。
$\xi_0 \leadsto \xi_1 \leadsto \ldots \leadsto \xi_e$ を
$X$ における特殊化の列とする。$x = \xi_e$、$y = f(x)$ とおく。
$e \leq \dim(\mathcal{O}_{X, x})$ であることに注意せよ。実際、与えられた特殊化は
$\mathcal{O}_{X, x}$ のスペクトルに現れる。スキーム、補題
\ref{schemes-lemma-specialize-points} を参照せよ。
次元公式、補題 \ref{morphisms-lemma-dimension-formula} により
\begin{align*}
e & \leq \dim(\mathcal{O}_{X, x}) \\
& \leq \dim(\mathcal{O}_{Y, y}) +
\text{trdeg}_{R(Y)} R(X) - \text{trdeg}_{\kappa(y)} \kappa(x) \\
& \leq \dim(\mathcal{O}_{Y, y}) + \text{trdeg}_{R(Y)} R(X)
\end{align*}
を得る。従って望む通り $e \leq \dim(Y) + \text{trdeg}_{R(Y)} R(X)$ である。

\medskip\noindent
次に、$f$ は閉でもあると仮定する。
$\overline{\xi}_0 \leadsto \overline{\xi}_1 \leadsto \ldots
\leadsto \overline{\xi}_d$ を $Y$ における特殊化の列とする。
$\dim(X) \geq d + r$ を示したい。
$\overline{\xi}_0 = \eta$ は $Y$ の生成点であるとしてよい。
一般ファイバー $X_\eta$ は $\kappa(\eta) = R(Y)$ 上局所有限型のスキームである。
$f$ は支配的なので、これは空でない。従って補題
\ref{morphisms-lemma-ubiquity-Jacobson-schemes} により Jacobson スキームである。
よって補題 \ref{morphisms-lemma-jacobson-finite-type-points} により、閉点
$\xi_0 \in X_\eta$ で、拡大 $\kappa(\eta) \subset \kappa(\xi_0)$ が
有限となるものを見つけられる。
$\mathcal{O}_{X, \xi_0} = \mathcal{O}_{X_\eta, \xi_0}$ であることに注意する。
これは $\eta$ が $Y$ の生成点だからである。従って
$\dim(\mathcal{O}_{X, \xi_0}) = r$ である。実際、補題
\ref{morphisms-lemma-dimension-formula} を、普遍鎖状スキーム上のスキーム $X_\eta$、
すなわち $\Spec(\kappa(\eta))$ 上のもの（補題
\ref{morphisms-lemma-ubiquity-uc} を参照）と点 $\xi_0$ に適用すればよい。
これは $\xi_{-r} \leadsto \ldots \leadsto \xi_{-1} \leadsto \xi_0$ を
$X$ において見つけられることを意味する。一方、$f$ は閉なので、特殊化は $f$ に沿って
持ち上がる。位相、補題
\ref{topology-lemma-closed-open-map-specialization} を参照せよ。
従って $\xi_0$ は $\eta = \overline{\xi}_0$ 上にあるから、特殊化
$\xi_0 \leadsto \xi_1 \leadsto \ldots \leadsto \xi_d$ で、
$\overline{\xi}_0 \leadsto \overline{\xi}_1 \leadsto \ldots
\leadsto \overline{\xi}_d$ 上にあるものを見つけられる。
言い換えれば
$$
\xi_{-r} \leadsto \ldots \leadsto \xi_{-1} \leadsto \xi_0
\leadsto \xi_1 \leadsto \ldots \leadsto \xi_d
$$
であり、望む通り $\dim(X) \geq d + r$ となる。
\end{proof}

\begin{lemma}
\label{morphisms-lemma-alteration-dimension-general}
$f : X \to Y$ をスキームの射とする。$Y$ は局所 Noether 的であり、
$f$ は局所有限型であると仮定する。このとき
$$
\dim(X) \leq \dim(Y) + E
$$
である。ここで $E$ は $\text{trdeg}_{\kappa(f(\xi))}(\kappa(\xi))$ の
上限であり、$\xi$ は $X$ の既約成分の生成点を動く。
\end{lemma}

\begin{proof}
補題 \ref{morphisms-lemma-dimension-formula-general} と性質、補題
\ref{properties-lemma-dimension} から直ちに従う。
\end{proof}


\section{相対正規化}
\label{morphisms-section-normalization-X-in-Y}

\noindent
この節では、{\it 一つのスキームの別のスキームにおける正規化}を構成する。

\begin{lemma}
\label{morphisms-lemma-integral-closure}
$X$ をスキームとする。$\mathcal{A}$ を準連接
$\mathcal{O}_X$-代数の層とする。規則
$$
U \longmapsto \{f \in \mathcal{A}(U) \mid
f_x \in \mathcal{A}_x \text{ integral over } \mathcal{O}_{X, x}
\text{ for all }x \in U\}
$$
で定められる部分層 $\mathcal{A}' \subset \mathcal{A}$ は準連接
$\mathcal{O}_X$-代数である。茎 $\mathcal{A}'_x$ は
$\mathcal{O}_{X, x}$ の整閉包であり、これは $\mathcal{A}_x$ の中でとる。
任意のアフィン開集合 $U \subset X$ に対し、環
$\mathcal{A}'(U) \subset \mathcal{A}(U)$ は
$\mathcal{O}_X(U)$ の整閉包であり、これは $\mathcal{A}(U)$ の中でとる。
\end{lemma}

\begin{proof}
条件は局所的なので、これは部分層である。
代数、補題 \ref{algebra-lemma-integral-closure-is-ring} により
$\mathcal{O}_X$-代数である。
$U \subset X$ をアフィン開集合とする。$U = \Spec(R)$ と書き、
$\mathcal{A}$ は $R$-代数 $A$ に付随する準連接層であるとする。
代数、補題 \ref{algebra-lemma-integral-closure-stalks} によれば、
$\mathcal{A}'$ の $U$ 上の値は $A'$ により与えられる。これは
$R$ の $A$ における整閉包である。これで補題の最後の主張が証明された。
$\mathcal{A}'$ が準連接であることを示すには、
$\mathcal{A}'(D(f)) = A'_f$ を示せば十分である。
これは整閉包と局所化が可換であることから従う。代数、補題
\ref{algebra-lemma-integral-closure-localize} を参照せよ。
同じ事実により、茎についての主張も従う。
\end{proof}

\begin{definition}
\label{morphisms-definition-integral-closure}
$X$ をスキームとする。$\mathcal{A}$ を準連接
$\mathcal{O}_X$-代数の層とする。{\it $\mathcal{O}_X$ の整閉包
（$\mathcal{A}$ におけるもの）}とは、上の補題
\ref{morphisms-lemma-integral-closure} で構成した準連接
$\mathcal{O}_X$-部分代数 $\mathcal{A}' \subset \mathcal{A}$ のことである。
\end{definition}

\noindent
上の定義の状況では、相対スペクトルの射
$$
\xymatrix{
Y = \underline{\Spec}_X(\mathcal{A}) \ar[rr] \ar[rd] & &
X' = \underline{\Spec}_X(\mathcal{A}') \ar[ld] \\
& X &
}
$$
を考えることができる。補題 \ref{morphisms-lemma-affine-equivalence-algebras} を
参照せよ。スキーム $X' \to X$ が $X$ の正規化となり、これはスキーム
$Y$ においてとる。以下では少し一般的な状況を考える。
スキームの準コンパクトかつ準分離的な射 $f : Y \to X$ が
与えられたとする。この場合、$\mathcal{O}_X$-代数の層
$f_*\mathcal{O}_Y$ は準連接である。スキーム、補題
\ref{schemes-lemma-push-forward-quasi-coherent} を参照せよ。
整閉包 $\mathcal{O}' \subset f_*\mathcal{O}_Y$ をとると、
準連接 $\mathcal{O}_X$-代数の層が得られ、その相対スペクトルが
$X$ の正規化となる。これは $Y$ においてとる。正式な定義は次の通りである。

\begin{definition}
\label{morphisms-definition-normalization-X-in-Y}
$f : Y \to X$ をスキームの準コンパクトかつ準分離的な射とする。
$\mathcal{O}'$ を $\mathcal{O}_X$ の整閉包とし、これは
$f_*\mathcal{O}_Y$ の中でとる。{\it $X$ の正規化
（$Y$ におけるもの）}とは、次のスキーム
\footnote{スキーム $X'$ は正規とは限らない。例えば $Y = X$ かつ
$f = \text{id}_X$ ならば $X' = X$ である。}
$$
\nu : X' = \underline{\Spec}_X(\mathcal{O}') \to X
$$
のことであり、これは $X$ 上にある。さらに、もとの射 $f$ の自然な分解
$$
Y \xrightarrow{f'} X' \xrightarrow{\nu} X
$$
を備える。
\end{definition}

\noindent
この分解は、標準射
$Y \to \underline{\Spec}_X(f_*\mathcal{O}_Y)$（構成、補題
\ref{constructions-lemma-canonical-morphism} を参照）と、包含写像
$\mathcal{O}' \to f_*\mathcal{O}_Y$ から生じる相対スペクトルの射との
合成である。正規化は次のように特徴づけられる。

\begin{lemma}
\label{morphisms-lemma-characterize-normalization}
$f : Y \to X$ をスキームの準コンパクトかつ準分離的な射とする。
$f = \nu \circ f'$ をその分解とし、$\nu : X' \to X$ を
$X$ の正規化（$Y$ におけるもの）とする。この分解は次の二つの性質で
特徴づけられる。
\begin{enumerate}
\item 射 $\nu$ は整である。
\item 任意の分解 $f = \pi \circ g$ で $\pi : Z \to X$ が整であるものに対し、
可換図式
$$
\xymatrix{
Y \ar[d]_{f'} \ar[r]_g & Z \ar[d]^\pi \\
X' \ar[ru]^h \ar[r]^\nu & X
}
$$
を満たす一意な射 $h : X' \to Z$ が存在する。
\end{enumerate}
さらに、射 $f' : Y \to X'$ は支配的であり、(2) における射
$h : X' \to Z$ は $Z$ の正規化（$Y$ におけるもの）である。
\end{lemma}

\begin{proof}
定義 \ref{morphisms-definition-normalization-X-in-Y} のように、
$\mathcal{O}' \subset f_*\mathcal{O}_Y$ を $\mathcal{O}_X$ の
整閉包とする。構成により射 $\nu$ は整であり、(1) が従う。
(2) のような分解 $f = \pi \circ g$ で、$\pi : Z \to X$ が整であるものが
与えられたとする。定義 \ref{morphisms-definition-integral} により $\pi$ は
アフィンである。従って $Z$ は、ある準連接 $\mathcal{O}_X$-代数の層
$\mathcal{B}$ の相対スペクトルである。
射 $g : Y \to Z$ は $\mathcal{O}_X$-代数の写像
$\chi : \mathcal{B} \to f_*\mathcal{O}_Y$ に対応する。
任意の場合に $\mathcal{B}(U)$ は $\mathcal{O}_X(U)$ 上整である。
ここで $U \subset X$ はアフィン開集合であるから
（定義 \ref{morphisms-definition-integral}）、補題
\ref{morphisms-lemma-integral-closure} により
$\chi(\mathcal{B}) \subset \mathcal{O}'$ である。
相対スペクトルの関手性、補題
\ref{morphisms-lemma-affine-equivalence-algebras} により、これは一意な射
$h : X' \to Z$ を与える。図式が可換であることの確認は省略する。

\medskip\noindent
(1) と (2) が分解 $f = \nu \circ f'$ を特徴づけることは明らかである。
実際、ある圏の始対象として特徴づけている。

\medskip\noindent
(2) の普遍性から、$f'$ は $X'$ の真の閉部分スキームを経由しない。
従って $f'$ のスキーム論的像は $X'$ である。$f'$ は準コンパクトである
（スキーム、補題 \ref{schemes-lemma-quasi-compact-permanence} と、
$\nu$ がアフィン射なので分離的であることによる）から、
$f'(Y)$ は $X'$ で稠密である。従って $f'$ は支配的である。

\medskip\noindent
$g$ は、スキーム、補題 \ref{schemes-lemma-compose-after-separated} と
\ref{schemes-lemma-quasi-compact-permanence} により準コンパクトかつ
準分離的であることに注意せよ。従って補題の最後の主張は意味をもつ。
(2) の射 $h$ は補題 \ref{morphisms-lemma-finite-permanence} により整である。
分解 $g = \pi' \circ g'$ で $\pi' : Z' \to Z$ が整であるものが
与えられると、分解 $f = (\pi \circ \pi') \circ g'$ が得られ、射
$h' : X' \to Z'$ が得られる。一意性により $\pi' \circ h' = h$ である。
従って特徴づけ (1), (2) を射 $h : X' \to Z$ に適用でき、
補題の最後の主張が従う。
\end{proof}

\begin{lemma}
\label{morphisms-lemma-functoriality-normalization}
可換図式
$$
\xymatrix{
Y_2 \ar[d]_{f_2} \ar[r] & Y_1 \ar[d]^{f_1} \\
X_2 \ar[r] & X_1
}
$$
をスキームの射のものとする。
$f_1$, $f_2$ は準コンパクトかつ準分離的であると仮定する。
$f_i = \nu_i \circ f_i'$, $i = 1, 2$ を標準分解とし、ここで
$\nu_i : X_i' \to X_i$ は $X_i$ の正規化（$Y_i$ におけるもの）とする。
このとき一意な矢印 $X'_2 \to X'_1$ が存在し、可換図式
$$
\xymatrix{
Y_2 \ar[d]_{f_2'} \ar[r] & Y_1 \ar[d]^{f_1'} \\
X_2' \ar[d]_{\nu_2} \ar[r] & X_1' \ar[d]^{\nu_1} \\
X_2 \ar[r] & X_1
}
$$
に組み込まれる。
\end{lemma}

\begin{proof}
補題 \ref{morphisms-lemma-characterize-normalization} (1) と
\ref{morphisms-lemma-base-change-finite} により、基底変換
$X_2 \times_{X_1} X'_1 \to X_2$ は整である。
$f_2$ はこの射を経由することに注意せよ。従って補題
\ref{morphisms-lemma-characterize-normalization} (2) により、一意な射
$X'_2 \to X_2 \times_{X_1} X'_1$ を得る。
これは可換図式に組み込まれる矢印 $X'_2 \to X'_1$ を与え、
一意性も同時に従う。
\end{proof}


\begin{lemma}
\label{morphisms-lemma-normalization-localization}
$f : Y \to X$ をスキームの準コンパクトかつ準分離的な射とする。
$U \subset X$ を開部分スキームとし、$V = f^{-1}(U)$ とおく。
このとき、$U$ の $V$ における正規化は $U$ の逆像である。ここで逆像は、
$X$ の $Y$ における正規化においてとる。
\end{lemma}

\begin{proof}
構成から明らかである。
\end{proof}

\begin{lemma}
\label{morphisms-lemma-normalization-is-normalization}
$f : Y \to X$ をスキームの準コンパクトかつ準分離的な射とする。
$X'$ を $X$ の $Y$ における正規化とする。このとき、$X'$ の $Y$ における
正規化は $X'$ である。
\end{lemma}

\begin{proof}
$Y \to X'' \to X'$ が $X'$ の $Y$ における正規化であるならば、
補題 \ref{morphisms-lemma-characterize-normalization} を合成 $X'' \to X$ に適用して、
標準射 $h : X' \to X''$ を $X$ 上で得る。射 $h$ と $X'' \to X'$ が
互いに逆であることの確認は省略する（補題における分解の一意性を用いる）。
\end{proof}

\begin{lemma}
\label{morphisms-lemma-normalization-in-reduced}
$f : Y \to X$ をスキームの準コンパクトかつ準分離的な射とする。
$X' \to X$ を $X$ の $Y$ における正規化とする。$Y$ が被約ならば、
$X'$ も被約である。
\end{lemma}

\begin{proof}
これは被約環の部分環が被約であるという事実から従う。
いくつかの詳細は省略する。
\end{proof}

\begin{lemma}
\label{morphisms-lemma-normalization-generic}
$f : Y \to X$ をスキームの準コンパクトかつ準分離的な射とする。
$X' \to X$ を $X$ の $Y$ における正規化とする。$X'$ の既約成分の
任意の生成点は、$Y$ の既約成分のある生成点の像である。
\end{lemma}

\begin{proof}
補題 \ref{morphisms-lemma-normalization-localization} により、$X = \Spec(A)$ は
アフィンであると仮定してよい。有限アフィン開被覆
$Y = \bigcup \Spec(B_i)$ を選ぶ。このとき $X' = \Spec(A')$ であり、射
$\Spec(B_i) \to Y \to X'$ は合わせて単射な $A$-代数写像
$A' \to \prod B_i$ を定める。従って補題は、代数、補題
\ref{algebra-lemma-injective-minimal-primes-in-image} から従う。
\end{proof}

\begin{lemma}
\label{morphisms-lemma-normalization-in-disjoint-union}
$f : Y \to X$ をスキームの準コンパクトかつ準分離的な射とする。
$Y = Y_1 \amalg Y_2$ が二つのスキームの非交和であると仮定する。
$f_i = f|_{Y_i}$ と書く。$X_i'$ を $X$ の $Y_i$ における正規化とする。
このとき $X_1' \amalg X_2'$ は $X$ の $Y$ における正規化である。
\end{lemma}

\begin{proof}
整閉包の言葉では、これは次の事実に対応する。
$A \to B$ を環準同型とする。$B = B_1 \times B_2$ と仮定する。
$A_i'$ を $A$ の $B_i$ における整閉包とする。このとき
$A_1' \times A_2'$ は $A$ の $B$ における整閉包である。
これが成り立つ理由は、元 $(1, 0)$ と $(0, 1)$ が $B$ の冪等元であり、
従って $A$ 上整だからである。従って $A'$（$A$ の $B$ における整閉包）は
直積であり、その因子が上で述べた整閉包 $A'_i$ であることは難しくない
（いくつかの詳細は省略する）。
\end{proof}

\begin{lemma}
\label{morphisms-lemma-normalization-in-universally-closed}
$f : X \to S$ をスキームの準コンパクト、準分離的かつ普遍閉な射とする。
このとき $f_*\mathcal{O}_X$ は $\mathcal{O}_S$ 上整である。言い換えれば、
$S$ の正規化（$X$ におけるもの）は、構成、補題
\ref{constructions-lemma-canonical-morphism} の分解
$$
X \longrightarrow \underline{\Spec}_S(f_*\mathcal{O}_X)
\longrightarrow S
$$
に等しい。
\end{lemma}

\begin{proof}
問題は $S$ 上局所的なので、$S = \Spec(R)$ はアフィンであると仮定してよい。
$h \in \Gamma(X, \mathcal{O}_X)$ とする。$h$ が $R$ 上のモニック方程式を
満たすことを示さなければならない。次の可換図式のように、$h$ を射と考える。
$$
\xymatrix{
X \ar[rr]_h \ar[rd]_f & & \mathbf{A}^1_S \ar[ld] \\
& S &
}
$$
$Z \subset \mathbf{A}^1_S$ を $h$ のスキーム論的像とする。定義
\ref{morphisms-definition-scheme-theoretic-image} を参照せよ。
射 $h$ は、$f$ が準コンパクトであり、$\mathbf{A}^1_S \to S$ が分離的なので、
準コンパクトである。スキーム、補題
\ref{schemes-lemma-quasi-compact-permanence} を参照せよ。補題
\ref{morphisms-lemma-quasi-compact-scheme-theoretic-image} により射 $X \to Z$ は
支配的である。補題 \ref{morphisms-lemma-image-proper-scheme-closed} により射
$X \to Z$ は閉である。従って $h(X) = Z$ である（集合論的に）。
従って補題 \ref{morphisms-lemma-image-proper-is-proper} を用いて、$Z \to S$ は
普遍閉（さらには固有）であると結論できる。
$Z \subset \mathbf{A}^1_S$ なので、$Z \to S$ はアフィンかつ固有であり、
従って補題 \ref{morphisms-lemma-integral-universally-closed} により整である。
$\mathbf{A}^1_S = \Spec(R[T])$ と書くと、イデアル $I \subset R[T]$、
すなわち $Z$ のイデアルはモニック多項式 $P(T) \in R[T]$ を含むと結論できる。
従って $P(h) = 0$ となり、主張を得る。
\end{proof}

\begin{lemma}
\label{morphisms-lemma-normalization-in-integral}
$f : Y \to X$ を整射とする。
このとき $X$ の正規化（$Y$ におけるもの）は $Y$ に等しい。
\end{lemma}

\begin{proof}
補題 \ref{morphisms-lemma-integral-universally-closed} により、これは補題
\ref{morphisms-lemma-normalization-in-universally-closed} の特別な場合である。
\end{proof}

\begin{lemma}
\label{morphisms-lemma-normal-normalization}
$f : Y \to X$ をスキームの準コンパクトかつ準分離的な射とする。
$X'$ を $X$ の $Y$ における正規化とする。次を仮定する。
\begin{enumerate}
\item $Y$ は正規スキームである。
\item $Y$ の準コンパクトな開集合は有限個の既約成分をもつ。
\end{enumerate}
このとき $X'$ は正規整スキームの非交和である。さらに、射
$Y \to X'$ は支配的であり、既約成分の間の全単射を誘導する。
\end{lemma}

\begin{proof}
$U \subset X$ をアフィン開集合とする。$U'$ を $U$ の $X'$ における逆像として
考える。$V = f^{-1}(U)$ とおく。補題
\ref{morphisms-lemma-normalization-localization} により、$V \to U' \to U$ は
$U$ の正規化（$V$ におけるもの）である。$U = \Spec(A)$ と書く。
このとき $V$ は準コンパクトであり、従って仮定により有限個の既約成分をもつ。
従って
$V = \coprod_{i = 1, \ldots n} V_i$ は正規整スキームの有限非交和である。
性質、補題
\ref{properties-lemma-normal-locally-finite-nr-irreducibles} を参照せよ。
補題 \ref{morphisms-lemma-normalization-in-disjoint-union} により
$U' = \coprod_{i = 1, \ldots, n} U_i'$ である。ここで $U'_i$ は
$U$ の正規化（$V_i$ におけるもの）である。性質、補題
\ref{properties-lemma-normal-integral-sections} により
$B_i = \Gamma(V_i, \mathcal{O}_{V_i})$ は正規整域である。
$U_i' = \Spec(A_i')$ であることに注意せよ。ここで
$A_i' \subset B_i$ は $A$ の $B_i$ における整閉包である。補題
\ref{morphisms-lemma-integral-closure} を参照せよ。代数、補題
\ref{algebra-lemma-integral-closure-in-normal} により、包含
$A_i' \subset B_i$ は正規整域であることが分かる。
従って $U' = \coprod U_i'$ は正規整スキームの有限合併であり、従って正規である。

\medskip\noindent
$X'$ はスキーム $U'$ による開被覆をもつので、性質、補題
\ref{properties-lemma-locally-normal} から $X'$ は正規であると結論する。
一方、各 $U'$ は有限個の既約スキームの非交和なので、$X'$ の各準コンパクト
開集合は有限個の既約成分をもつ（省略する位相的議論による）。従って $X'$ は、
性質、補題
\ref{properties-lemma-normal-locally-finite-nr-irreducibles} により
正規整スキームの非交和である。上の $X'$ の記述から、$Y \to X'$ が
支配的であり、$V \to U'$ が既約成分上の全単射を誘導することは、
任意のアフィン開集合 $U \subset X$ に対して明らかである。
射 $Y \to X'$ に対する既約成分の全単射は、これから位相的議論により従う
（省略する）。
\end{proof}

\begin{lemma}
\label{morphisms-lemma-nagata-normalization-finite-general}
$f : X \to S$ を射とする。次を仮定する。
\begin{enumerate}
\item $S$ は永田スキームである。
\item $f$ は準コンパクトかつ準分離的である。
\item $X$ の準コンパクトな開集合は有限個の既約成分をもつ。
\item $x \in X$ が既約成分の生成点ならば、体拡大
$\kappa(x)/\kappa(f(x))$ は有限生成である。
\item $X$ は被約である。
\end{enumerate}
このとき、正規化 $\nu : S' \to S$、すなわち $S$ の $X$ における正規化は
有限である。
\end{lemma}

\begin{proof}
仮定 (1) により、$S = \Spec(R)$ かつ $R$ が永田環である場合へ直ちに
帰着できる。整閉包 $A$、すなわち $R$ の $\Gamma(X, \mathcal{O}_X)$ におけるものが
$R$ 上有限であることを示さなければならない。仮定 (2) により $f$ は
準コンパクトなので、$X = \bigcup_{i = 1, \ldots, n} U_i$ と書け、各 $U_i$ は
アフィンである。$U_i = \Spec(B_i)$ とする。
各 $B_i$ は仮定 (5) により被約であり、仮定 (3) と代数、補題
\ref{algebra-lemma-irreducible} により有限個の極小素イデアル
$\mathfrak q_{i1}, \ldots, \mathfrak q_{im_i}$ をもつ。次の包含がある。
$$
\Gamma(X, \mathcal{O}_X) \subset B_1 \times \ldots \times B_n
\subset
\prod\nolimits_{i = 1, \ldots, n}
\prod\nolimits_{j = 1, \ldots, m_i} (B_i)_{\mathfrak q_{ij}}
$$
第二の包含は、代数、補題
\ref{algebra-lemma-reduced-ring-sub-product-fields} による。
代数、補題 \ref{algebra-lemma-minimal-prime-reduced-ring} により
$\kappa(\mathfrak q_{ij}) = (B_i)_{\mathfrak q_{ij}}$ である。
従って整閉包 $A$、すなわち $R$ の $\Gamma(X, \mathcal{O}_X)$ におけるものは、
整閉包 $A_{ij}$、すなわち $R$ の $\kappa(\mathfrak q_{ij})$ におけるものたちの
直積に含まれる。$R$ は Noether 環なので、各 $A_{ij}$ が有限 $R$-加群で
あることを各 $i, j$ について示せば十分である。$\mathfrak p_{ij} \subset R$ を
$\mathfrak q_{ij}$ の像とする。仮定 (4) により
$\kappa(\mathfrak q_{ij})/\kappa(\mathfrak p_{ij})$ は有限生成体拡大なので、
$R \to \kappa(\mathfrak q_{ij})$ は本質的有限型であることが分かる。
従って、代数、補題
\ref{algebra-lemma-nagata-in-reduced-finite-type-finite-integral-closure}
により $R \to A_{ij}$ は有限である。
\end{proof}

\begin{lemma}
\label{morphisms-lemma-nagata-normalization-finite}
$f : X \to S$ を射とする。次を仮定する。
\begin{enumerate}
\item $S$ は永田スキームである。
\item $f$ は有限型である。
\item $X$ は被約である。
\end{enumerate}
このとき、正規化 $\nu : S' \to S$、すなわち $S$ の $X$ における正規化は
有限である。
\end{lemma}

\begin{proof}
これは補題 \ref{morphisms-lemma-nagata-normalization-finite-general} の特別な場合である。
実際、有限型射 $f$ は定義により準コンパクトであり、補題
\ref{morphisms-lemma-finite-type-Noetherian-quasi-separated} により準分離的なので、
(2) が成り立つ。$X$ は補題 \ref{morphisms-lemma-finite-type-noetherian} により
局所 Noether 的なので、条件 (3) が成り立つ。最後に、有限型射は有限生成な
剰余体拡大を誘導するので、条件 (4) が成り立つ。
\end{proof}

\begin{lemma}
\label{morphisms-lemma-relative-normalization-normal-codim-1}
$f : Y \to X$ を有限型なスキームの射とし、$Y$ は被約、$X$ は永田スキームで
あるとする。$X'$ を $X$ の $Y$ における正規化とする。$x' \in X'$ を
次を満たす点とする。
\begin{enumerate}
\item $\dim(\mathcal{O}_{X', x'}) = 1$ である。
\item $Y \to X'$ の $x'$ 上のファイバーは空である。
\end{enumerate}
このとき $\mathcal{O}_{X', x'}$ は離散付値環である。
\end{lemma}

\begin{proof}
$X$ を $x'$ の像のアフィン近傍で置き換えることができる。
従って、$X = \Spec(A)$ とし、$A$ は永田環であると仮定してよい。
補題 \ref{morphisms-lemma-nagata-normalization-finite} により射 $X' \to X$ は有限である。
従って $X' = \Spec(A')$ と書け、ある有限 $A$-代数を $A'$ と表す。
補題 \ref{morphisms-lemma-normalization-is-normalization} により、$X$ を $X'$ で
置き換えた後、次の段落で述べる場合に帰着する。

\medskip\noindent
$X = X' = \Spec(A)$ であり、$A$ が Noether 環である場合。
$\mathfrak p \subset A$ を点 $x'$ に対応する素イデアルとする。
$g \in \mathfrak p$ を選び、これは $A$ のどの極小素イデアルにも含まれないとする
（素イデアル回避と、$A$ が有限個の極小素イデアルをもつという事実を使う。
代数、補題 \ref{algebra-lemma-silly} および
\ref{algebra-lemma-Noetherian-irreducible-components} を参照せよ）。
$Z = f^{-1}V(g) \subset Y$ とおく。これは $Y$ の閉部分スキームである。
$f(Z)$ は $g$ の選び方によりどの生成点も含まず、$x'$ を含まない。
これは $x'$ が $f$ の像にないからである。補題
\ref{morphisms-lemma-reach-points-scheme-theoretic-image} により、
$f(Z)$ の閉包は $f(Z)$ の点の特殊化全体である。従って $f(Z)$ の閉包は
$x'$ を含まない。なぜなら、条件 $\dim(\mathcal{O}_{X', x'}) = 1$ は、
$X = X'$ の生成点のみが $x'$ に特殊化することを含意するからである。
言い換えれば、$X$ を $x'$ のアフィン開近傍で置き換えた後、
$f^{-1}V(g) = \emptyset$ と仮定してよい。従って $g$ は $Y$ 上の可逆な
大域関数へ写り、分解
$$
A \to A_g \to \Gamma(Y, \mathcal{O}_Y)
$$
を得る。$X = X'$ なので、これは $A$ が $A$ の $A_g$ における整閉包に
等しいことを含意する。代数、補題
\ref{algebra-lemma-integral-closure-localize} により、$A_\mathfrak p$ は
$A_\mathfrak p$ の $A_\mathfrak p[1/g]$ における整閉包であると結論する。
$g$ の選び方により、$\dim(A_\mathfrak p) = 1$ であり、$A$ は被約なので、
$A_\mathfrak p[1/g]$ は体の有限直積（$\mathfrak p$ に含まれる極小素イデアルの
剰余体の直積）であることが分かる。従って $A_\mathfrak p$ は正規であり
（代数、補題 \ref{algebra-lemma-characterize-reduced-ring-normal}）、証明が完了する。
いくつかの詳細は省略する。
\end{proof}


\section{正規化}
\label{morphisms-section-normalization}

\noindent
次に、スキーム $X$ の正規化を扱う。
$X$ が局所有限個の既約成分をもつ場合にのみ、これを定義し構成する。
$X$ を、各準コンパクト開集合が有限個の既約成分をもつスキームとする。
$X^{(0)} \subset X$ を $X$ の既約成分の生成点全体とする。写像
\begin{equation}
\label{morphisms-equation-generic-points}
f :
Y = \coprod\nolimits_{\eta \in X^{(0)}} \Spec(\kappa(\eta))
\longrightarrow
X
\end{equation}
を、スキーム、第 \ref{schemes-section-points} 節の標準写像を用いた、生成点の
$X$ への包含とする。この射は仮定により準コンパクトであり、$Y$ が分離的なので
準分離的であることに注意せよ（スキーム、第
\ref{schemes-section-separation-axioms} 節を参照）。

\begin{definition}
\label{morphisms-definition-normalization}
$X$ を、各準コンパクト開集合が有限個の既約成分をもつスキームとする。
$X$ の {\it 正規化}を、射
$$
\nu : X^\nu \longrightarrow X
$$
として定義する。これは $X$ の正規化であり、上で構成した射
$f : Y \to X$（\ref{morphisms-equation-generic-points}）においてとる。
\end{definition}

\noindent
任意の局所 Noether スキームでは既約成分全体は局所有限であるから、
この定義を適用できる。通常、正規化は被約スキームに対してのみ定義される。
上の定義では、$X$ の正規化は、被約化 $X_{red}$、すなわち $X$ の被約化の
正規化と同じである。

\begin{lemma}
\label{morphisms-lemma-normalization-reduced}
$X$ を、各準コンパクト開集合が有限個の既約成分をもつスキームとする。
正規化射 $\nu$ は被約化 $X_{red}$ を経由し、$X^\nu \to X_{red}$ は
$X_{red}$ の正規化である。
\end{lemma}

\begin{proof}
$f : Y \to X$ を射（\ref{morphisms-equation-generic-points}）とする。スキーム、補題
\ref{schemes-lemma-map-into-reduction} から、分解
$Y \to X_{red} \to X$ を得る。これは $f$ の分解である。補題
\ref{morphisms-lemma-characterize-normalization} により標準射 $X^\nu \to X_{red}$ を得て、
$X^\nu$ は $X_{red}$ の正規化（$Y$ におけるもの）である。射
$Y \to X_{red}$ は $X_{red}$ に対して構成した射
（\ref{morphisms-equation-generic-points}）と同一なので、補題が従う。
\end{proof}

\noindent
$X$ が被約ならば、$X$ の正規化は、$\mathcal{O}_X$ の有理型関数の層
$\mathcal{K}_X$ における整閉包の相対スペクトルと同じである
（因子、第 \ref{divisors-section-meromorphic-functions} 節を参照）。
実際、この場合には $\mathcal{K}_X = f_*\mathcal{O}_Y$ である。因子、補題
\ref{divisors-lemma-reduced-finite-irreducible} とその証明を参照せよ。
ここでこれを明示的に記述する。

\begin{lemma}
\label{morphisms-lemma-description-normalization}
$X$ を、各準コンパクト開集合が有限個の既約成分をもつ被約スキームとする。
$\Spec(A) = U \subset X$ をアフィン開集合とする。このとき
\begin{enumerate}
\item $A$ は有限個の極小素イデアル
$\mathfrak q_1, \ldots, \mathfrak q_t$ をもつ。
\item $Q(A)$、すなわち $A$ の全商環は
$Q(A/\mathfrak q_1) \times \ldots \times Q(A/\mathfrak q_t)$ である。
\item 整閉包 $A'$、すなわち $A$ の $Q(A)$ における整閉包は、整域
$A/\mathfrak q_i$ の体 $Q(A/\mathfrak q_i)$ における整閉包の直積である。
\item $\nu^{-1}(U)$ は $A'$ のスペクトルと同一視される。ここで
$\nu : X^\nu \to X$ は正規化射である。
\end{enumerate}
\end{lemma}

\begin{proof}
極小素イデアルは既約成分に対応するので（代数、補題
\ref{algebra-lemma-irreducible}）、仮定から (1) を得る。さらに
$(0) = \mathfrak q_1 \cap \ldots \cap \mathfrak q_t$ である。これは $A$ が被約だからである
（代数、補題 \ref{algebra-lemma-Zariski-topology}）。従って
$Q(A) = \prod A_{\mathfrak q_i} = \prod \kappa(\mathfrak q_i)$ である。代数、補題
\ref{algebra-lemma-total-ring-fractions-no-embedded-points} および
\ref{algebra-lemma-minimal-prime-reduced-ring} を参照せよ。これで (2) が証明された。
(3) は代数、補題 \ref{algebra-lemma-characterize-reduced-ring-normal} または補題
\ref{morphisms-lemma-normalization-in-disjoint-union} から従う。(4) が成り立つのは、
$f^{-1}(U) \to U$ が明らかに射
$$
\Spec\left(\prod \kappa(\mathfrak q_i)\right)
\longrightarrow
\Spec(A)
$$
だからである。ここで $f : Y \to X$ は射
（\ref{morphisms-equation-generic-points}）である。
\end{proof}

\begin{lemma}
\label{morphisms-lemma-stalk-normalization}
$X$ を、各準コンパクト開集合が有限個の既約成分をもつスキームとする。
$\nu : X^\nu \to X$ を $X$ の正規化とする。$x \in X$ とする。このとき、
次のものは $\mathcal{O}_{X, x}$-代数として標準的に同型である。
\begin{enumerate}
\item 茎 $(\nu_*\mathcal{O}_{X^\nu})_x$。
\item $\mathcal{O}_{X, x}$ の整閉包で、
$(\mathcal{O}_{X, x})_{red}$ の全商環においてとるもの。
\item $\mathcal{O}_{X, x}$ の極小素イデアルの剰余体の直積における
$\mathcal{O}_{X, x}$ の整閉包（このような極小素イデアルは有限個である）。
\end{enumerate}
\end{lemma}

\begin{proof}
$X$ を $x$ のアフィン開近傍で置き換えた後、$X$ は有限個の既約成分をもち、
$x$ はそのすべてに含まれると仮定してよい。このとき茎
$(\nu_*\mathcal{O}_{X^\nu})_x$ は、$A = \mathcal{O}_{X, x}$ の極小素イデアルの
剰余体の直積 $L$ における $A$ の整閉包である。これは正規化の構成と補題
\ref{morphisms-lemma-integral-closure} から従う。あるいは補題
\ref{morphisms-lemma-description-normalization} と、正規化が局所化と可換であるという事実
（代数、補題 \ref{algebra-lemma-integral-closure-localize}）を用いてもよい。
$A_{red}$ は有限個の極小素イデアルをもつので（これらは $X$ の既約成分で
$x$ を通るものの生成点とちょうど対応する）、$L$ は $A_{red}$ の全商環である
（代数、補題 \ref{algebra-lemma-total-ring-fractions-no-embedded-points}）。
従って、いまの環は $A$ の $A_{red}$ の全商環における整閉包でもある。
\end{proof}

\begin{lemma}
\label{morphisms-lemma-normalization-normal}
$X$ を、各準コンパクト開集合が有限個の既約成分をもつスキームとする。
\begin{enumerate}
\item 正規化 $X^\nu$ は正規整スキームの非交和である。
\item 射 $\nu : X^\nu \to X$ は整かつ全射であり、既約成分上の全単射を誘導する。
\item 任意の整射 $\alpha : X' \to X$ で、準コンパクト開集合
$U \subset X$ に対して逆像 $\alpha^{-1}(U)$ が有限個の既約成分をもち、
$\alpha|_{\alpha^{-1}(U)} : \alpha^{-1}(U) \to U$ が双有理であるものに対して
\footnote{この不格好な定式化が必要なのは、始域と終域が有限個の既約成分をもつ
場合についてしか、射が双有理であることの意味を定義していないからである。
$X'_{red} \to X_{red}$ が条件を満たせば十分である。}
分解 $X^\nu \to X' \to X$ が存在し、$X^\nu \to X'$ は $X'$ の正規化である。
\item 任意の射 $Z \to X$ で、$Z$ が正規スキームであり、$Z$ の各既約成分が
$X$ のある既約成分を支配するものに対して、一意な分解
$Z \to X^\nu \to X$ が存在する。
\end{enumerate}
\end{lemma}

\begin{proof}
$f : Y \to X$ を（\ref{morphisms-equation-generic-points}）のものとする。スキーム $X^\nu$ は、
$Y$ が正規であり、$Y$ の任意のアフィン開集合が有限個の既約成分をもつから、補題
\ref{morphisms-lemma-normal-normalization} により正規整スキームの
非交和である。これで (1) が証明された。別法として、補題
\ref{morphisms-lemma-normalization-reduced} と \ref{morphisms-lemma-description-normalization}
から (1) を導いてもよい。

\medskip\noindent
射 $\nu$ は補題 \ref{morphisms-lemma-characterize-normalization} により整である。
補題 \ref{morphisms-lemma-normal-normalization} により射 $Y \to X^\nu$ は既約成分上の
全単射を誘導し、$Y$ の構成により、これは $X^\nu \to X$ が既約成分上の
全単射を誘導することを含意する。構成により $f : Y \to X$ は支配的なので、
$\nu$ も支配的である。整射は閉であるから（補題
\ref{morphisms-lemma-integral-universally-closed}）、これは $\nu$ が全射であることを
含意する。これで (2) が証明された。

\medskip\noindent
$\alpha : X' \to X$ が (3) のような射であると仮定する。$X'$ が正規化を
定義するための仮定を満たすことは明らかである。$f' : Y' \to X'$ を、$X'$ から
出発して構成した射（\ref{morphisms-equation-generic-points}）とする。$\alpha$ は局所的に
双有理なので、$Y' = Y$ かつ $f = \alpha \circ f'$ であることは明らかである。
従って分解 $X^\nu \to X' \to X$ が存在し、補題
\ref{morphisms-lemma-characterize-normalization} により $X^\nu \to X'$ は $X'$ の
正規化である。これで (3) が証明された。

\medskip\noindent
$g : Z \to X$ を、始域が正規スキームであり、各既約成分が $X$ のある既約成分を
支配する射とする。補題 \ref{morphisms-lemma-normalization-reduced} により
$X^\nu = X_{red}^\nu$ であり、スキーム、補題
\ref{schemes-lemma-map-into-reduction} により $Z \to X$ は $X_{red}$ を経由する。
従って $X$ を $X_{red}$ で置き換え、$X$ は被約であると仮定してよい。
さらに、分解は一意なので、$Z$ 上局所的に構成すれば十分である。
$W \subset Z$ と $U \subset X$ を、$g(W) \subset U$ を満たすアフィン開集合とする。
$U = \Spec(A)$ および $W = \Spec(B)$ と書き、$g|_W$ は
$\varphi : A \to B$ で与えられるとする。補題
\ref{morphisms-lemma-description-normalization} の結果を自由に用いる。
$\mathfrak p_1, \ldots, \mathfrak p_t$ を $A$ の極小素イデアルとする。
$Z$ は正規なので、$B$ は正規環、特に被約環である。さらに仮定により、任意の
極小素イデアル $\mathfrak q \subset B$ に対して
$\varphi^{-1}(\mathfrak q)$ は $A$ の極小素イデアルである。従って
$x \in A$ が非零因子、すなわち $x \not \in \bigcup \mathfrak p_i$ ならば、
$\varphi(x)$ は $B$ の非零因子である。こうして標準環準同型
$Q(A) \to Q(B)$ を得る。$B$ は正規なので、$Q(B)$ における自身の整閉包に
等しい（代数、補題 \ref{algebra-lemma-normal-ring-integrally-closed} を参照）。
従って整閉包 $A' \subset Q(A)$、すなわち $A$ の整閉包は、$B$ へ標準写像
$Q(A) \to Q(B)$ により写る。$\nu^{-1}(U) = \Spec(A')$ なので、
これは標準分解 $W \to \nu^{-1}(U) \to U$ を与える。これは $\nu|_W$ の分解である。
一意性の確認は省略する。
\end{proof}


\begin{lemma}
\label{morphisms-lemma-normalization-in-terms-of-components}
$X$ を、各準コンパクト開集合が有限個の既約成分をもつスキームとする。
$Z_i \subset X$, $i \in I$ を、被約誘導構造を備えた $X$ の既約成分とする。
$Z_i^\nu \to Z_i$ を正規化とする。このとき
$\coprod_{i \in I} Z_i^\nu \to X$ は $X$ の正規化である。
\end{lemma}

\begin{proof}
補題 \ref{morphisms-lemma-normalization-reduced} により、$X$ は被約であると仮定してよい。
このとき補題は、補題 \ref{morphisms-lemma-description-normalization} の局所的記述から、
または補題 \ref{morphisms-lemma-normalization-normal} の (3) から従う。後者を用いる場合、
$\coprod Z_i \to X$ は整かつ局所双有理である（$X$ は被約であり、局所有限個の
既約成分をもつからである）。
\end{proof}

\begin{lemma}
\label{morphisms-lemma-normalization-birational}
$X$ を有限個の既約成分をもつ被約スキームとする。
このとき正規化射 $X^\nu \to X$ は双有理である。
\end{lemma}

\begin{proof}
補題 \ref{morphisms-lemma-normalization-normal} により、正規化は既約成分の全単射を誘導する。
$\eta \in X$ を $X$ のある既約成分の生成点とし、$\eta^\nu \in X^\nu$ を
$X^\nu$ の対応する既約成分の生成点とする。このとき $\eta^\nu \mapsto \eta$ であり、
証明を終えるには
$\mathcal{O}_{X, \eta} \to \mathcal{O}_{X^\nu, \eta^\nu}$ が同型であることを
示せばよい。定義 \ref{morphisms-definition-birational} を参照せよ。$X$ と $X^\nu$ は
被約なので、両方の局所環はそれぞれの剰余体に等しい
（代数、補題 \ref{algebra-lemma-minimal-prime-reduced-ring}）。一方、正規化を
$X$ の $Y = \coprod \Spec(\kappa(\eta))$ における正規化として構成したことから、
$\kappa(\eta) \subset \kappa(\eta^\nu) \subset \kappa(\eta)$ であり、証明が完了する。
\end{proof}

\begin{lemma}
\label{morphisms-lemma-finite-birational-over-normal}
整スキーム間の有限（あるいは整でさえある）双有理射 $f : X \to Y$ で、
$Y$ が正規であるものは同型である。
\end{lemma}

\begin{proof}
$V \subset Y$ をアフィン開集合とし、その逆像 $U \subset X$ もアフィン開集合とする。
$f$ は整スキーム間の双有理射なので、準同型
$\mathcal{O}_Y(V) \to \mathcal{O}_X(U)$ は整域間の単射であり、商体の同型を誘導する。
$Y$ は正規なので、環 $\mathcal{O}_Y(V)$ は商体において整閉である。
$f$ は有限（あるいは整）なので、$\mathcal{O}_X(U)$ の各元は
$\mathcal{O}_Y(V)$ 上整である。従って
$\mathcal{O}_Y(V) = \mathcal{O}_X(U)$ と結論する。これは望む通り $f$ が同型で
あることを証明する。
\end{proof}

\begin{lemma}
\label{morphisms-lemma-normalization-trivial}
$X$ を局所有限個の既約成分をもつスキームとする。正規化射
$\nu : X^\nu \to X$ が同型であることと $X$ が正規であることは同値である。
\end{lemma}

\begin{proof}
$\nu$ が同型ならば、補題 \ref{morphisms-lemma-normalization-normal} により $X$ は正規である。
逆に $X$ が正規であると仮定する。補題 \ref{morphisms-lemma-normalization-localization} と
性質、補題 \ref{properties-lemma-normal-locally-finite-nr-irreducibles} により、
$X$ は整であると仮定してよい。補題 \ref{morphisms-lemma-normalization-normal} により射
$\nu$ は整であり、$X^\nu$ はただ一つの既約成分をもつ（従って整スキームである）。
補題 \ref{morphisms-lemma-normalization-birational} と
\ref{morphisms-lemma-finite-birational-over-normal} から結論が従う。
\end{proof}

\begin{lemma}
\label{morphisms-lemma-Japanese-normalization}
$X$ を整な日本スキームとする。
正規化 $\nu : X^\nu \to X$ は有限射である。
\end{lemma}

\begin{proof}
定義（性質、定義 \ref{properties-definition-nagata}）と補題
\ref{morphisms-lemma-description-normalization} から従う。実際、この場合、補題の主張は
$\nu^{-1}(\Spec(A))$ が $A$ の商体における整閉包のスペクトルであるということである。
\end{proof}

\begin{lemma}
\label{morphisms-lemma-nagata-normalization}
$X$ を永田スキームとする。
正規化 $\nu : X^\nu \to X$ は有限射である。
\end{lemma}

\begin{proof}
永田スキームは局所 Noether 的なので、定義 \ref{morphisms-definition-normalization} を適用できる
ことに注意せよ。補題は補題
\ref{morphisms-lemma-nagata-normalization-finite-general} の特別な場合であるが、次のように
直接証明することもできる。$X^\nu \to X$ を合成
$X^\nu \to X_{red} \to X$ として書く。$X_{red} \to X$ は閉埋め込みなので有限である。
従って、補題 \ref{morphisms-lemma-composition-finite} により、被約な永田スキームの場合を
証明すれば十分である。$\Spec(A) = U \subset X$ をアフィン開集合とする。
補題 \ref{morphisms-lemma-description-normalization} により
$\nu^{-1}(U) = \Spec(\prod A_i')$ である。ここで $A_i'$ は
$A/\mathfrak q_i$ の商体における整閉包である。$A$ は永田環なので
（性質、補題 \ref{properties-lemma-locally-nagata} を参照）、各環拡大
$A/\mathfrak q_i \subset A'_i$ は有限である。従って $A \to \prod A'_i$ は
有限環準同型であり、主張を得る。
\end{proof}

\begin{lemma}
\label{morphisms-lemma-normalization-geom-unibranch-univ-homeo}
$X$ を既約かつ幾何的単枝なスキームとする。
正規化射 $\nu : X^\nu \to X$ は普遍同相射である。
\end{lemma}

\begin{proof}
補題 \ref{morphisms-lemma-universal-homeomorphism} により、$\nu$ が整、普遍的単射、かつ全射で
あることを示さなければならない。補題 \ref{morphisms-lemma-normalization-normal} により射
$\nu$ は整である。$x \in X$ とし、$A = \mathcal{O}_{X, x}$ とおく。
$X$ は既約なので、$A$ はただ一つの極小素イデアル $\mathfrak p$ をもち、
$A_{red} = A/\mathfrak p$ である。補題 \ref{morphisms-lemma-stalk-normalization} により、茎
$A' = (\nu_*\mathcal{O}_{X^\nu})_x$ は $A$ の $A_{red}$ の商体における整閉包である。
発展代数、定義 \ref{more-algebra-definition-unibranch} により、$A'$ は
ただ一つの素イデアル $\mathfrak m'$ をもち、これは $\mathfrak m_x \subset A$ 上にあり、
$\kappa(\mathfrak m')/\kappa(x)$ は純非分離である。従って $\nu$ は全単射
（特に全射）であり、補題 \ref{morphisms-lemma-universally-injective} により普遍的単射である。
\end{proof}


\section{弱正規化}
\label{morphisms-section-weak-normalization}

\noindent
スキームが局所有限個の既約成分をもつ場合にのみ、その弱正規化を定義する。
これは正規化の場合と同様である。

\begin{lemma}
\label{morphisms-lemma-relative-weak-normalization-algebra}
$A \to B$ を、スペクトルの支配的射 $\Spec(B) \to \Spec(A)$ を誘導する
環準同型とする。次を満たす $A$-部分代数 $B' \subset B$ が存在する。
\begin{enumerate}
\item $\Spec(B') \to \Spec(A)$ は普遍同相射である。
\item 分解 $A \to C \to B$ で、$\Spec(C) \to \Spec(A)$ が普遍同相射であるものが
与えられると、$C \to B$ の像は $B'$ に含まれる。
\end{enumerate}
\end{lemma}

\begin{proof}
補題 \ref{morphisms-lemma-reduction-universal-homeomorphism} を今後断りなく用いる。
可換図式
$$
\xymatrix{
B \ar[r] & B_{red} \\
A \ar[u] \ar[r] & A_{red} \ar[u]
}
$$
を考える。(2) のような任意の分解 $A \to C \to B$、すなわち $A \to B$ の分解に対し、
$A_{red} \to C_{red} \to B_{red}$ は (2) のような
$A_{red} \to B_{red}$ の分解である。従って補題が $A_{red} \to B_{red}$ に対して
成り立ち、$A_{red}$-部分代数 $B'_{red} \subset B_{red}$ を与えるなら、
$B' \subset B$ を $B'_{red}$ の逆像とおくことで $A \to B$ に対する補題が解ける。
これにより次の段落で扱う場合に帰着する。

\medskip\noindent
$A$ と $B$ は被約であると仮定する。この場合、代数、補題
\ref{algebra-lemma-image-dense-generic-points} により $A \subset B$ である。
$A \to C \to B$ を (2) のような分解とする。このとき命題
\ref{morphisms-proposition-universal-homeomorphism} を $A \subset C$ に適用すると、$C$ の
各元は、拡大 $A[c_1, \ldots, c_n] \subset C$ に含まれ、$i = 1, \ldots, n$ に
対して次のいずれかを満たすことが分かる。
\begin{enumerate}
\item $c_i^2, c_i^3 \in A[c_1, \ldots, c_{i - 1}]$ である。または
\item 素数 $p$ が存在して
$pc_i, c_i^p \in A[c_1, \ldots, c_{i - 1}]$ である。
\end{enumerate}
従って $B' \subset B$ を、元 $b \in B$ で、次の拡大
$
A[b_1, \ldots, b_n] \subset B
$
に含まれるもの全体として定義すれば、
次の条件 (*) のもとで
性質 (2) が成り立つ。すなわち、$i = 1, \ldots, n$ に対して次のいずれかが成り立つ。
\begin{enumerate}
\item $b_i^2, b_i^3 \in A[b_1, \ldots, b_{i - 1}]$ である。または
\item 素数 $p$ が存在して
$pb_i, b_i^p \in A[b_1, \ldots, b_{i - 1}]$ である。
\end{enumerate}
確認すべきことは二つだけである。(a) $B'$ が $A$-部分代数であること、そして
(b) $\Spec(B') \to \Spec(A)$ が普遍同相射であることである。(a) は次から従う。
$n \geq 0$ と $b_1, \ldots, b_n \in B$ が (*) を満たし、$m \geq 0$ と
$b'_1, \ldots, b'_m \in B$ が (*) を満たすとき、整数 $n + m$ と
$b_1, \ldots, b_n, b'_1, \ldots, b'_m \in B$ も (*) を満たす。
最後に、(b) は命題 \ref{morphisms-proposition-universal-homeomorphism} と $B'$ の構成から従う。
\end{proof}

\begin{lemma}
\label{morphisms-lemma-relative-weak-normalization-localization}
$A \to B$ を、スペクトルの支配的射 $\Spec(B) \to \Spec(A)$ を誘導する
環準同型とする。補題 \ref{morphisms-lemma-relative-weak-normalization-algebra} の
$A$-部分代数 $B' \subset B$ の形成は局所化と可換である
（説明については証明を参照）。
\end{lemma}

\begin{proof}
$S \subset A$ を積閉部分集合とする。このとき $S^{-1}A \to S^{-1}B$ も、
支配的射 $\Spec(S^{-1}B) \to \Spec(S^{-1}A)$ を誘導する環準同型である
（補題 \ref{morphisms-lemma-base-change-dominant} および
\ref{morphisms-lemma-generalizations-lift-flat} を参照）。従って補題
\ref{morphisms-lemma-relative-weak-normalization-algebra} は $S^{-1}A$-部分代数
$(S^{-1}B)' \subset S^{-1}B$ を与える。主張は、等式
$S^{-1}B' = (S^{-1}B)'$ が $S^{-1}A$-部分代数として $S^{-1}B$ の中で
成り立つという意味である。

\medskip\noindent
これを示すため、補題 \ref{morphisms-lemma-relative-weak-normalization-algebra} の証明における
$B'$ と $(S^{-1}B)'$ の構成を用いる。第一段階で、$B'$ は $A_{red}$-部分代数
$B'_{red} \subset B_{red}$ の逆像である。これは環準同型
$A_{red} \to B_{red}$ に対して構成したものであり、$(S^{-1}B)'$ についても同様である。
$S^{-1}B_{red} = (S^{-1}B)_{red}$ であることに注意すれば、次の段落で扱う場合に
帰着する。

\medskip\noindent
$A$ と $B$ が被約ならば、$B'$ を部分代数 $A[b_1, \ldots, b_n]$ の合併として
構成した。ここで $i = 1, \ldots, n$ に対して次のいずれかを満たす。
\begin{enumerate}
\item $b_i^2, b_i^3 \in A[b_1, \ldots, b_{i - 1}]$ である。または
\item 素数 $p$ が存在して
$pb_i, b_i^p \in A[b_1, \ldots, b_{i - 1}]$ である。
\end{enumerate}
$(S^{-1}B)' \subset S^{-1}B$ についても同様である。従って
$B' \to B \to S^{-1}B$ の像が $(S^{-1}B)'$ に含まれることは明らかである。
対応する写像 $S^{-1}B' \to (S^{-1}B)'$ が全射であることを示すには、補題
\ref{morphisms-lemma-special-elements-and-localization} を用いて分母を順次払えばよい。
詳細は省略する。
\end{proof}

\begin{lemma}
\label{morphisms-lemma-relative-seminormalization-algebra}
$A \to B$ を、スペクトルの支配的射 $\Spec(B) \to \Spec(A)$ を誘導する
環準同型とする。次を満たす $A$-部分代数 $B' \subset B$ が存在する。
\begin{enumerate}
\item $\Spec(B') \to \Spec(A)$ は、剰余体上で同型を誘導する普遍同相射である。
\item 分解 $A \to C \to B$ で、$\Spec(C) \to \Spec(A)$ が剰余体上で同型を
誘導する普遍同相射であるものが与えられると、$C \to B$ の像は $B'$ に含まれる。
\end{enumerate}
\end{lemma}

\begin{proof}
証明は補題 \ref{morphisms-lemma-relative-weak-normalization-algebra} の証明とまったく同じである。
ただし命題 \ref{morphisms-proposition-universal-homeomorphism-equal-residue-fields} を命題
\ref{morphisms-proposition-universal-homeomorphism} の代わりに用いる。
\end{proof}

\begin{lemma}
\label{morphisms-lemma-relative-seminormalization-localization}
$A \to B$ を、スペクトルの支配的射 $\Spec(B) \to \Spec(A)$ を誘導する
環準同型とする。補題 \ref{morphisms-lemma-relative-seminormalization-algebra} の
$A$-部分代数 $B' \subset B$ の形成は局所化と可換である
（説明については証明を参照）。
\end{lemma}

\begin{proof}
証明は補題 \ref{morphisms-lemma-relative-weak-normalization-localization} の証明と同じである。
\end{proof}

\begin{lemma}
\label{morphisms-lemma-relative-seminormalization}
$f : Y \to X$ をスキームの準コンパクト、準分離的、かつ支配的な射とする。
\begin{enumerate}
\item 分解 $Y \to X' \to X$ で $X' \to X$ が普遍同相射であるものの圏は、
始対象 $Y \to X^{Y/wn} \to X$ をもつ。
\item 分解 $Y \to X' \to X$ で $X' \to X$ が剰余体上で同型を誘導する
普遍同相射であるものの圏は、始対象 $Y \to X^{Y/sn} \to X$ をもつ。
\end{enumerate}
さらに、分解 $Y \to X^{Y/wn} \to X$ および $Y \to X^{Y/sn} \to X$ の形成は、
$X$ の開部分スキームへの基底変換と可換である。
\end{lemma}

\begin{proof}
(1) を証明し、(2) の証明は省略する。また、最後の主張も分解の構成から従う。
補題 \ref{morphisms-lemma-universal-homeomorphism} を以後断りなく用いる。まず、
$Y \to X^{Y/n} \to X$ を $X$ の $Y$ における正規化とする。定義
\ref{morphisms-definition-normalization-X-in-Y} を参照せよ。(1) のような
$Y \to X' \to X$ に対し、与えられた射と整合する一意な射
$X^{Y/n} \to X'$ を得る。補題 \ref{morphisms-lemma-characterize-normalization} を参照せよ。
従って $f$ を $X^{Y/n} \to X$ で置き換えて補題を証明すれば十分である。
これにより、次の段落で考察する場合に帰着する。

\medskip\noindent
$f$ は整であると仮定する（証明の残りは、より一般に $f$ がアフィンの場合にも成り立つ）。
$U = \Spec(A)$ を $X$ のアフィン開集合とし、$V = f^{-1}(U) = \Spec(B)$ を
$Y$ における逆像とする。このとき $A \to B$ はスペクトル上の支配的射を誘導する
環準同型である。補題 \ref{morphisms-lemma-relative-weak-normalization-algebra} により、
$A$-部分代数 $B' \subset B$ で、$U^{V/wn} = \Spec(B')$ とおくと分解
$V \to U^{V/wn} \to U$ が、分解 $V \to U' \to U$ で $U' \to U$ が
普遍同相射であるものの圏において始対象となるものを得る。

\medskip\noindent
$U_1 \subset U_2 \subset X$ をアフィン開集合とし、$V_i = f^{-1}(U_i)$ とおく。
すると、次の標準射を得る。
$$
\rho_{U_1}^{U_2} : U_1^{V_1/wn} \to U_1 \times_{U_2} U_2^{V_2/wn}
$$
これは $U_1$ 上の射であり、$U_1^{V_1/wn}$ の普遍性によって得られる。
これらの射は自然な関手性を満たすが、その定式化と証明は読者に委ねる。
さらに射 $\rho_{U_1}^{U_2}$ は同型である。$U_1 \subset U_2$ が標準開集合なら、
これは補題 \ref{morphisms-lemma-relative-weak-normalization-localization} から従い、一般の場合は、
これらの写像の関手的性質とスキーム、補題
\ref{schemes-lemma-standard-open-two-affines} によりこの場合に帰着できる
（詳細は省略する）。従って相対的貼り合わせ
（構成、補題 \ref{constructions-lemma-relative-glueing}）により、射
$X^{Y/wn} \to X$ を得る。これは $U^{V/wn} \to U$ に制限される $U$ 上の射で、
$\rho_{U_1}^{U_2}$ と整合する。もちろん、射 $V \to U^{V/wn}$ は貼り合わさって射
$Y \to X^{Y/wn}$ を与える（構成、注意
\ref{constructions-remark-relative-glueing-functorial} を参照）。こうして分解
$Y \to X^{Y/wn} \to X$ を得、その第二の射は普遍同相射である。

\medskip\noindent
最後に、$Y \to X' \to X$ を (1) のような分解とする。上のような
$V \to U^{V/wn} \to U \subset X$ に対し、第二の矢印が普遍同相射である分解
$V \to U \times_X X' \to U$ を得る。次いで一意な射
$g_U : U^{V/wn} \to U \times_X X'$ を $U$ 上で得る。これは
$U^{V/wn}$ の普遍性による。これらの $g_U$ は射
$\rho_{U_1}^{U_2}$ と整合する。詳細は省略する。従って一意な射
$g : X^{Y/wn} \to X'$ が存在し、これは $X$ 上の射で、$g_U$ と $U$ 上で
一致する。構成、注意
\ref{constructions-remark-relative-glueing-functorial} を参照せよ。
これは $Y \to X^{Y/wn} \to X$ がこの圏の始対象であることを証明し、証明は完了する。
\end{proof}

\begin{definition}
\label{morphisms-definition-relative-seminormalization}
$f : Y \to X$ をスキームの準コンパクト、準分離的、かつ支配的な射とする。
\begin{enumerate}
\item 補題 \ref{morphisms-lemma-relative-seminormalization} の (2) で構成した分解
$Y \to X^{Y/sn} \to X$ を、{\it $X$ の $Y$ における半正規化}という。
\item 補題 \ref{morphisms-lemma-relative-seminormalization} の (1) で構成した分解
$Y \to X^{Y/wn} \to X$ を、{\it $X$ の $Y$ における弱正規化}という。
\end{enumerate}
\end{definition}

\noindent
局所有限個の既約成分をもつスキームの半正規化を解釈し直す方法を述べる。

\begin{lemma}
\label{morphisms-lemma-seminormal-and-normalization}
$X$ を、各準コンパクト開集合が有限個の既約成分をもつスキームとする。
$\nu : X^\nu \to X$ を $X$ の正規化とする。このとき $X$ の $X^\nu$ における
半正規化は $X$ の半正規化である。式で書けば、
$X^{sn} = X^{X^\nu/sn}$ である。
\end{lemma}

\begin{proof}
(\ref{morphisms-equation-generic-points}) のような $f : Y \to X$ をとり、$X^\nu$ が
$X$ の $Y$ における正規化となるようにする。半正規化
$X^{sn} \to X$ は、$X$ 上で剰余体上の同型を誘導する普遍同相射 $X' \to X$ の圏の
始対象である。$Y$ は $X$ の既約成分の生成点における剰余体のスペクトルの
非交和なので、この圏の任意の $X' \to X$ に対し、標準的な持ち上げ
$f' : Y \to X'$、すなわち $f$ の持ち上げを得る。次いで補題
\ref{morphisms-lemma-characterize-normalization} により、
標準射 $X^\nu \to X'$ を得る。従って今度は標準射
$X^{X^\nu/sn} \to X'$ を得る。これは $X^{X^\nu/sn}$ の普遍性による。
このようにして $X^{X^\nu/sn}$ が
$X^{sn}$ と同じ普遍性を満たすことが分かり、結論を得る。
\end{proof}

\noindent
補題 \ref{morphisms-lemma-seminormal-and-normalization} は次の定義を動機づける。
$X$ が局所有限個の既約成分をもつ場合にしか正規化を構成していないので、
弱正規化についてもこの場合に限定する。

\begin{definition}
\label{morphisms-definition-weak-normalization}
$X$ を、各準コンパクト開集合が有限個の既約成分をもつスキームとする。
$X$ の{\it 弱正規化}を、弱正規化
$$
X^\nu \longrightarrow X^{wn} \longrightarrow X
$$
、すなわち $X$ の、その $X^\nu$（$X$ の正規化）における弱正規化として定義する
（定義 \ref{morphisms-definition-normalization}）。式で書けば、
$X^{wn} = X^{X^\nu/wn}$ である。
\end{definition}

\noindent
補題 \ref{morphisms-lemma-seminormal-and-normalization} と合わせると、局所有限個の
既約成分をもつスキーム $X$ に対して標準射
$$
X^\nu \to X^{wn} \to X^{sn} \to X
$$
が存在することが分かる。この定義をしたので、（局所有限個の既約成分をもつという
条件のもとで）スキームが弱正規であることの意味を述べられる。

\begin{definition}
\label{morphisms-definition-weakly-normal}
$X$ を、各準コンパクト開集合が有限個の既約成分をもつスキームとする。
$X$ が{\it 弱正規}であるとは、弱正規化 $X^{wn} \to X$ が同型であることをいう
（定義 \ref{morphisms-definition-weak-normalization}）。
\end{definition}

\noindent
定義から直ちに、各準コンパクト開集合が有限個の既約成分をもつスキーム $X$ に対して
$$
X\text{ normal} \Rightarrow
X\text{ weakly normal} \Rightarrow
X\text{ seminormal}
$$
を得る。アフィンの場合の弱正規性の意味は次のように具体化できる。

\begin{lemma}
\label{morphisms-lemma-affine-weakly-normal}
$X = \Spec(A)$ を有限個の既約成分をもつアフィンスキームとする。このとき $X$ が
弱正規であるための必要十分条件は次の二条件である。
\begin{enumerate}
\item $A$ は半正規である（定義 \ref{morphisms-definition-seminormal-ring}）。
\item 素数 $p$ と $z, w \in A$ が、(a) $z$ は零因子でなく、
(b) $w^p$ は $z^p$ で割り切れ、(c) $pw$ は $z$ で割り切れる、を満たすならば、
$w$ は $z$ で割り切れる。
\end{enumerate}
\end{lemma}

\begin{proof}
$X$ は弱正規であると仮定する。弱正規なスキームは半正規なので、(1) が成り立つ
（これは弱正規スキームの定義による）。特に $A$ は被約である。
$p, z, w$ を (2) のようにとる。$x, y \in A$ を、$z^p x = w^p$ および
$zy = pw$ を満たすように選ぶ。このとき $p^p x = y^p$ である。環準同型
$A \to C = A[t]/(t^p - x, pt - y)$ はスペクトル上の普遍同相射を誘導する。
$X^\nu$ は $X$ の正規化であり、$A'$、すなわち $A$ の、$A$ の全商環における整閉包の
スペクトルである。
補題 \ref{morphisms-lemma-description-normalization} を参照せよ。
$a = w/z \in A'$ は $a^p = x$ により成り立つことに注意する。
従って、$A$-代数準同型 $A \to C \to A'$ で、$t$ を $a$ に送るものがある。
この時点で、弱正規であることの定義的性質 $X = X^{wn} = X^{X^\nu/wn}$ から、
$C \to A'$ の像が $A$ に含まれることが分かる。従って望む通り $a \in A$ を得る。

\medskip\noindent
逆に (1) と (2) を仮定する。$A'$ を前段落と同じものとする。
$X^{X^\nu/wn} = X$ を示さなければならない。補題
\ref{morphisms-lemma-relative-weak-normalization-algebra} の証明における構成により、
$X^{X^\nu/wn}$ は、$A'$ の部分環 $A[a_1, \ldots, a_n] \subset A'$ で、
$i = 1, \ldots, n$ に対して次のいずれかを満たすものの合併である部分環の
スペクトルである。
\begin{enumerate}
\item[(a)] $a_i^2, a_i^3 \in A[a_1, \ldots, a_{i - 1}]$ である。または
\item[(b)] 素数 $p$ が存在して
$pa_i, a_i^p \in A[a_1, \ldots, a_{i - 1}]$ である。
\end{enumerate}
このとき (1) と (2) を用いると、帰納的に $a_1, \ldots, a_n \in A$ と分かる。
詳細は省略する。従って $X = X^{X^\nu/wn}$ であり、ゆえに $X$ は弱正規である。
\end{proof}

\noindent
必要となる局所性の補題を述べる。

\begin{lemma}
\label{morphisms-lemma-locally-weakly-normal}
$X$ を、各準コンパクト開集合が有限個の既約成分をもつスキームとする。
次の条件は同値である。
\begin{enumerate}
\item スキーム $X$ は弱正規である。
\item 任意のアフィン開集合 $U \subset X$ に対して、環 $\mathcal{O}_X(U)$ は
補題 \ref{morphisms-lemma-affine-weakly-normal} の条件 (1) と (2) を満たす。
\item アフィン開被覆 $X = \bigcup U_i$ で、各環 $\mathcal{O}_X(U_i)$ が
補題 \ref{morphisms-lemma-affine-weakly-normal} の条件 (1) と (2) を満たすものが存在する。
\item 開被覆 $X = \bigcup X_j$ で、各開部分スキーム $X_j$ が弱正規であるものが存在する。
\end{enumerate}
さらに、$X$ が弱正規ならば、すべての開部分スキームも弱正規である。
\end{lemma}

\begin{proof}
$X$ が弱正規であるための条件は、射 $X^{wn} = X^{X^\nu/wn} \to X$ が
同型であることである。$X^\nu \to X$ の構成は開部分スキームへの基底変換と可換であり、
$X^{X^\nu/wn}$ の構成も $X$ の開部分スキームへの基底変換と可換である
（補題 \ref{morphisms-lemma-relative-seminormalization}）から、補題は明らかである。
\end{proof}


\section{ザリスキーの主定理（代数的版）}
\label{morphisms-section-Zariski}

\noindent
これは純粋に代数的な方法で証明できる版である。
（非アフィン射に対する）より強力な版を証明する前に、さらに道具を整備する必要がある。
スキームのコホモロジー、節
\ref{coherent-section-applications-formal-functions}
および
射についてさらに、節
\ref{more-morphisms-section-application-etale-neighbourhoods}
を参照せよ。

\begin{theorem}[ザリスキーの主定理の代数的版]
\label{morphisms-theorem-main-theorem}
$f : Y \to X$ をスキームのアフィン射とし、$f$ は有限型であると仮定する。
$X'$ を $X$ の $Y$ における正規化とする。図式は次の通りである。
$$
\xymatrix{
Y \ar[rd]_f \ar[rr]_{f'} & & X' \ar[ld]^\nu \\
& X &
}
$$
このとき、次を満たす開部分スキーム $U' \subset X'$ が存在する。
\begin{enumerate}
\item $(f')^{-1}(U') \to U'$ は同型である。
\item $(f')^{-1}(U') \subset Y$ は、$f$ が準有限である点全体の集合である。
\end{enumerate}
\end{theorem}

\begin{proof}
$X$、従って $Y$ がアフィンである場合に直ちに帰着する。
$X = \Spec(R)$ および $Y = \Spec(A)$ と書く。このとき $X' = \Spec(A')$ であり、
$A'$ は $R$ の $A$ における整閉包である。定義
\ref{morphisms-definition-integral-closure} および \ref{morphisms-definition-normalization-X-in-Y} を参照せよ。
代数、定理 \ref{algebra-theorem-main-theorem} により、各点 $y \in Y$ で $f$ が
準有限であるものに対し、開集合 $U'_y \subset X'$ が存在して
$(f')^{-1}(U'_y) \to U'_y$ は同型となる。$U' = \bigcup U'_y$ とおく。ここで
$y \in Y$ は $f$ の準有限な点全体を動く。残るのは、$f$ が $(f')^{-1}(U')$ のすべての点で
準有限であることを示すことである。$y \in (f')^{-1}(U')$ の像を $x \in X$ とする。
すると $Y_x \to X'_x$ は $y$ の近傍で同型であることが分かる。従って、$Y_x$ には
$y$ に特殊化する点が存在しない。これは $f'(y)$ について $X'_x$ で
成り立つからである。補題 \ref{morphisms-lemma-integral-fibres} を参照せよ。
補題 \ref{morphisms-lemma-quasi-finite-at-point-characterize} の (3) により、これは $f$ が
$y$ において準有限であることを意味する。
\end{proof}

\noindent
ザリスキーの主定理の代数的版を用いて、射が準有限である点の集合が開であることを示せる。

\begin{lemma}
\label{morphisms-lemma-quasi-finite-points-open}
\begin{slogan}
射の局所準有限軌跡は開である
\end{slogan}
$f : X \to S$ をスキームの射とする。$X$ の点で $f$ が準有限であるもの全体は、
開集合 $U \subset X$ である。誘導射 $U \to S$ は局所準有限である。
\end{lemma}

\begin{proof}
$f$ は $x$ において準有限であると仮定する。定義 \ref{morphisms-definition-quasi-finite} のような
アフィン開集合 $x \in U = \Spec(A) \subset X$、$V = \Spec(R) \subset S$ をとる。
上の定理 \ref{morphisms-theorem-main-theorem} または代数、補題
\ref{algebra-lemma-quasi-finite-open} のいずれによっても、素イデアル $\mathfrak q$ で
$R \to A$ が準有限となるもの全体は $\Spec(A)$ の開集合である。これらはすべて
$X$ の点で $f$ が準有限であるものに対応するので、第一の主張を得る。第二の主張は明らかである。
\end{proof}

\noindent
次の補題は、後に一般の準有限分離射へと強める。射についてさらに、補題
\ref{more-morphisms-lemma-quasi-finite-separated-pass-through-finite}
を参照せよ。

\begin{lemma}
\label{morphisms-lemma-quasi-finite-affine}
$f : Y \to X$ をスキームの射とする。次を仮定する。
\begin{enumerate}
\item $X$ と $Y$ はアフィンである。
\item $f$ は準有限である。
\end{enumerate}
このとき図式
$$
\xymatrix{
Y \ar[rd]_f \ar[rr]_j & & Z \ar[ld]^\pi \\
& X &
}
$$
で、$Z$ はアフィン、$\pi$ は有限、$j$ は開埋め込みであるものが存在する。
\end{lemma}

\begin{proof}
これは代数、補題 \ref{algebra-lemma-quasi-finite-open-integral-closure} を
スキームの言葉で言い換えたものである。
\end{proof}

\begin{lemma}
\label{morphisms-lemma-image-nowhere-dense-quasi-finite}
$f : Y \to X$ をスキームの準有限射とする。$T \subset Y$ を $Y$ の閉な疎部分集合とする。
このとき $f(T) \subset X$ は $X$ の疎部分集合である。
\end{lemma}

\begin{proof}
補題 \ref{morphisms-lemma-image-nowhere-dense-finite} の証明と同様に、基底 $X$ がアフィンである場合に
直ちに帰着する。この場合、$Y = \bigcup_{i = 1, \ldots, n} Y_i$ はアフィン開集合の
有限合併である（$f$ が準コンパクトだからである）。各 $T \cap Y_i$ は疎であり、
疎集合の有限合併も疎なので（位相、補題
\ref{topology-lemma-nowhere-dense} を参照）、像 $f(T \cap Y_i)$ が $X$ で疎であることを
証明すれば十分である。これにより $X$ と $Y$ がともにアフィンである場合に帰着する。
ここで上の補題 \ref{morphisms-lemma-quasi-finite-affine} を適用して図式
$$
\xymatrix{
Y \ar[rd]_f \ar[rr]_j & & Z \ar[ld]^\pi \\
& X &
}
$$
を得る。ここで $Z$ はアフィン、$\pi$ は有限、$j$ は開埋め込みである。
$\overline{T} = \overline{j(T)} \subset Z$ とおく。位相、補題
\ref{topology-lemma-image-nowhere-dense-open} により、$\overline{T}$ は $Z$ で疎である。
$f(T) \subset \pi(\overline{T})$ なので、有限の場合の対応する結果、すなわち補題
\ref{morphisms-lemma-image-nowhere-dense-finite} から補題が従う。
\end{proof}


\section{普遍有界なファイバー}
\label{morphisms-section-universally-bounded}

\noindent
$X$ を体 $k$ 上のスキームとする。$X$ が $k$ 上有限ならば、$X = \Spec(A)$ と書け、
$A$ は有限 $k$-代数である。これは、$X$ が $\Spec(k)$ 上有限局所自由であるともいえる。
定義 \ref{morphisms-definition-finite-locally-free} を参照せよ。従って $X \to \Spec(k)$ は
整数 $d \geq 0$ である{\it 次数}をもち、具体的には $d = \dim_k(A)$ である。
これを（有限）スキーム $X$ の $k$ 上の{\it 次数}と呼ぶこともある。

\begin{definition}
\label{morphisms-definition-universally-bounded}
$f : X \to Y$ をスキームの射とする。
\begin{enumerate}
\item 整数 $n$ が{\it $f$ のファイバーの次数を上から抑える}とは、任意の
$y \in Y$ に対し、ファイバー $X_y$ が $\kappa(y)$ 上の有限スキームで、
$\kappa(y)$ 上の次数が $\leq n$ であることをいう。
\item {\it $f$ のファイバーは普遍有界である}\footnote{これはおそらく標準的な用語ではない。}
とは、整数 $n$ で、$f$ のファイバーの次数を上から抑えるものが存在することをいう。
\end{enumerate}
\end{definition}

\noindent
特に、ファイバーの点の個数も $n$ で上から抑えられることに注意せよ。
（すべてのファイバーが有限被約スキームであっても、逆は成り立たない。）

\begin{lemma}
\label{morphisms-lemma-characterize-universally-bounded}
$f : X \to Y$ をスキームの射とし、$n \geq 0$ とする。次は同値である。
\begin{enumerate}
\item 整数 $n$ は $f$ のファイバーの次数を上から抑える。
\item 任意の射 $\Spec(k) \to Y$（ここで $k$ は体）に対し、ファイバー積
$X_k = \Spec(k) \times_Y X$ は $k$ 上有限で、その次数は $\leq n$ である。
\end{enumerate}
この場合、$f$ のファイバーは普遍有界であり、スキーム $X_k$ の点は高々 $n$ 個である。
より正確には、$X_k = \{x_1, \ldots, x_t\}$ ならば、
$$
n \geq \sum\nolimits_{i = 1, \ldots, t} [\kappa(x_i) : k]
$$
である。
\end{lemma}

\begin{proof}
(2) $\Rightarrow$ (1) は自明である。逆向きは、射 $\Spec(k) \to Y$ の像が $y$ ならば
$X_k = \Spec(k) \times_{\Spec(\kappa(y))} X_y$ であることから従う。定義により、
$f$ のファイバーが普遍有界であるとは、このような $n$ が存在することである。
最後に $X_k = \Spec(A)$ と仮定する。このとき $\dim_k A = n$ である。
従って $A$ は Artin 環で、すべての素イデアルは極大イデアル $\mathfrak m_i$ であり、
$A$ はこれらの極大イデアルにおける局所化の積である。代数、補題
\ref{algebra-lemma-finite-dimensional-algebra} および
\ref{algebra-lemma-artinian-finite-length} を参照せよ。すると $\mathfrak m_i$ は $x_i$ に対応し、
$A_{\mathfrak m_i} = \mathcal{O}_{X_k, x_i}$ であり、従って全射
$A \to  \bigoplus \kappa(\mathfrak m_i) = \bigoplus \kappa(x_i)$
が存在する。線形代数により、これは補題の主張の不等式を与える。
\end{proof}

\begin{lemma}
\label{morphisms-lemma-finite-locally-free-universally-bounded}
$f$ が次数 $d$ の有限局所自由射ならば、$d$ は $f$ のファイバーの次数を上から抑える。
\end{lemma}

\begin{proof}
$f$ の任意の基底変換は次数 $d$ の有限局所自由射であり
（補題 \ref{morphisms-lemma-base-change-finite-locally-free}）、従って $f$ のファイバーは
すべて次数 $d$ をもつので、主張は成り立つ。
\end{proof}

\begin{lemma}
\label{morphisms-lemma-composition-universally-bounded}
普遍有界なファイバーをもつ射の合成も、普遍有界なファイバーをもつ。
より正確には、$n$ が $f : X \to Y$ のファイバーの次数を上から抑え、$m$ が
$g : Y \to Z$ の次数を上から抑えると仮定する。このとき $nm$ は
$g \circ f : X \to Z$ のファイバーの次数を上から抑える。
\end{lemma}

\begin{proof}
$f : X \to Y$ と $g : Y \to Z$ が普遍有界なファイバーをもつとする。
$\deg(X_y/\kappa(y)) \leq n$ がすべての $y \in Y$ に対して成り立ち、
$\deg(Y_z/\kappa(z)) \leq m$ がすべての $z \in Z$ に対して成り立つとする。点 $z \in Z$ をとる。
仮定により、スキーム $Y_z$ は $\Spec(\kappa(z))$ 上有限である。特に $Y_z$ の基礎位相空間は
有限離散集合である。射 $f_z : X_z \to Y_z$ のファイバーは、$f$ の、対応する $Y$ の点における
ファイバーであり、上の議論により有限離散集合である。従って $X_z$ の基礎位相空間も
有限離散集合である。ゆえに $X_z$ はアフィンスキームである（これはよい演習であり、例えば
$X_z$ の全点からなる集合に性質、補題
\ref{properties-lemma-maximal-points-affine} を適用しても従う）。
$X_z = \Spec(A)$、$Y_z = \Spec(B)$、$k = \kappa(z)$ と書く。このとき
$k \to B \to A$ であり、(a) $\dim_k(B) \leq m$、(b) 各極大イデアル
$\mathfrak m \subset B$ に対して $\dim_{\kappa(\mathfrak m)}(A/\mathfrak mA) \leq n$
であることが分かっている。これから $\dim_k(A) \leq nm$ が従うと主張する。
$B$ はその局所化 $B_{\mathfrak m}$ の積であることに注意せよ。これは例えば $Y_z$ が
$1$ 点スキームの非交和であることから、または代数、補題
\ref{algebra-lemma-finite-dimensional-algebra} および
\ref{algebra-lemma-artinian-finite-length} から従う。従って
$\dim_k(B) = \sum_{\mathfrak m} \dim_k(B_{\mathfrak m})$ および
$\dim_k(A) = \sum_{\mathfrak m} \dim_k(A_{\mathfrak m})$ である。いずれの場合も
$\mathfrak m$ は $B$ の（$A$ のではない）極大イデアル全体を動く。上のことと中山の補題
（代数、補題 \ref{algebra-lemma-NAK}）により、各 $A_{\mathfrak m}$ は
$B_{\mathfrak m}^{\oplus n}$ の、$B_{\mathfrak m}$-加群としての商である。従って
$\dim_k(A_{\mathfrak m}) \leq n \dim_k(B_{\mathfrak m})$ である。すべてを合わせると
$$
\dim_k(A) = \sum\nolimits_{\mathfrak m} \dim_ta	(A_{\mathfrak m})
\leq \sum\nolimits_{\mathfrak m} n \dim_k(B_{\mathfrak m})
= n \dim_k(B) \leq nm
$$
となり、望む結果を得る。
\end{proof}

\begin{lemma}
\label{morphisms-lemma-base-change-universally-bounded}
普遍有界なファイバーをもつ射の基底変換は、普遍有界なファイバーをもつ射である。
より正確には、$n$ が $f : X \to Y$ のファイバーの次数を上から抑え、$Y' \to Y$ が
任意の射ならば、基底変換 $f' : Y' \times_Y X \to Y'$ のファイバーの次数も $n$ で
上から抑えられる。
\end{lemma}

\begin{proof}
補題 \ref{morphisms-lemma-characterize-universally-bounded} の結果から明らかである。
\end{proof}

\begin{lemma}
\label{morphisms-lemma-descent-universally-bounded}
$f : X \to Y$ をスキームの射とする。$Y' \to Y$ をスキームの射とし、
$f' : X' = X_{Y'} \to Y'$ を $f$ の基底変換とする。$Y' \to Y$ が全射で、$f'$ が
普遍有界なファイバーをもつならば、$f$ も普遍有界なファイバーをもつ。
より正確には、$n$ が $f'$ のファイバーの次数を上から抑えるならば、$n$ は $f$ の
ファイバーの次数も上から抑える。
\end{lemma}

\begin{proof}
$n \geq 0$ を $f'$ のファイバーの次数を上から抑える整数とする。$n$ は $f$ に対しても
使えると主張する。実際、点 $y \in Y$ に対して、点 $y' \in Y'$ で $y$ の上にあるものを選ぶと、
$$
X'_{y'} = \Spec(\kappa(y')) \times_{\Spec(\kappa(y))} X_y.
$$
$X'_{y'}$ は次数 $\leq n$ の、$\kappa(y')$ 上の有限スキームであると仮定しているので、
$X_y$ も次数 $\leq n$ の、$\kappa(y)$ 上の有限スキームである。（詳細の一部は省略する。）
\end{proof}

\begin{lemma}
\label{morphisms-lemma-immersion-universally-bounded}
埋め込みは普遍有界なファイバーをもつ。
\end{lemma}

\begin{proof}
定義において整数 $n = 1$ が使える。
\end{proof}

\begin{lemma}
\label{morphisms-lemma-etale-universally-bounded}
$f : X \to Y$ をスキームのエタール射とし、$n \geq 0$ とする。次は同値である。
\begin{enumerate}
\item 整数 $n$ はファイバーの次数を上から抑える。
\item 任意の体 $k$ と射 $\Spec(k) \to Y$ に対し、基底変換
$X_k = \Spec(k) \times_Y X$ は高々 $n$ 個の点をもつ。
\item 任意の $y \in Y$ と任意の分離代数閉包 $\kappa(y) \subset \kappa(y)^{sep}$ に対し、
スキーム $X_{\kappa(y)^{sep}}$ は高々 $n$ 個の点をもつ。
\end{enumerate}
\end{lemma}

\begin{proof}
これは補題 \ref{morphisms-lemma-characterize-universally-bounded} と、ファイバー $X_y$ が
$\kappa(y)$ の有限分離拡大のスペクトルの非交和であるという事実から従う。
補題 \ref{morphisms-lemma-etale-over-field} を参照せよ。
\end{proof}

\noindent
普遍有界なファイバーをもつことは絶対的な概念であり、相対的な概念ではない。
このため次の補題の条件は、$X$ が準コンパクトであることであって、$f$ が
準コンパクトであることではない。

\begin{lemma}
\label{morphisms-lemma-locally-quasi-finite-qc-source-universally-bounded}
$f : X \to Y$ をスキームの射とする。次を仮定する。
\begin{enumerate}
\item $f$ は局所準有限である。
\item $X$ は準コンパクトである。
\end{enumerate}
このとき $f$ は普遍有界なファイバーをもつ。
\end{lemma}

\begin{proof}
$X$ は準コンパクトなので、有限アフィン開被覆 $X = \bigcup_{i = 1, \ldots, n} U_i$ と、
アフィン開集合 $V_i \subset Y$（$i = 1, \ldots, n$）で $f(U_i) \subset V_i$ を
満たすものが存在する。「局所準有限性」の局所的性質（補題
\ref{morphisms-lemma-quasi-finite-at-point-characterize} の (4) を参照）により、射
$f|_{U_i} : U_i \to V_i$ はアフィン間の局所準有限射、従って準有限射であることが分かる。
補題 \ref{morphisms-lemma-quasi-finite-locally-quasi-compact} を参照せよ。
$y \in Y$ に対し、$X_y = \bigcup_{y \in V_i} (U_i)_y$ は明らかに開被覆である。
従ってアフィン間の準有限射について補題を証明すれば十分である。実際、$n_i$ が各射
$f|_{U_i} : U_i \to V_i$ に使えるならば、$\sum n_i$ が $f$ に使える。

\medskip\noindent
$f : X \to Y$ をアフィン間の準有限射と仮定する。補題
\ref{morphisms-lemma-quasi-finite-affine} により、図式
$$
\xymatrix{
X \ar[rd]_f \ar[rr]_j & & Z \ar[ld]^\pi \\
& Y &
}
$$
で、$Z$ はアフィン、$\pi$ は有限、$j$ は開埋め込みであるものを得る。
$j$ は普遍有界なファイバーをもつので（補題
\ref{morphisms-lemma-immersion-universally-bounded}）、$\pi$ が普遍有界なファイバーをもつことを
示す場合に帰着する（補題 \ref{morphisms-lemma-composition-universally-bounded}）。

\medskip\noindent
これにより、$\Spec(B) \to \Spec(A)$ という形の射で $A \to B$ が有限であるものに帰着する。
$B$ は $x_1, \ldots, x_n$ によって $A$-加群として生成されるとする。すると
$$
A^{\oplus n} \longrightarrow B, \quad
(a_1, \ldots, a_n) \longmapsto \sum a_i x_i
$$
は全射な $A$-加群準同型である。任意の素イデアル $\mathfrak p \subset A$ に対し、これは
$\kappa(\mathfrak p)$-ベクトル空間の全射
$$
\kappa(\mathfrak p)^{\oplus n} \longrightarrow B \otimes_A \kappa(\mathfrak p)
$$
を誘導する。言い換えれば、普遍有界なファイバーをもつ射の定義において整数 $n$ が使える。
\end{proof}

\begin{lemma}
\label{morphisms-lemma-universally-bounded-permanence}
スキームの射の可換図式
$$
\xymatrix{
X \ar[rd]_g \ar[rr]_f & & Y \ar[ld]^h \\
& Z &
}
$$
を考える。$g$ が普遍有界なファイバーをもち、$f$ が全射かつ平坦ならば、$h$ も
普遍有界なファイバーをもつ。より正確には、$n$ が $g$ のファイバーの次数を上から
抑えるならば、$n$ は $h$ のファイバーの次数も上から抑える。
\end{lemma}

\begin{proof}
$g$ が普遍有界なファイバーをもち、$f$ が全射かつ平坦であると仮定する。
$g$ のファイバーの次数は $n \in \mathbf{N}$ で上から抑えられるとする。
$n$ は $h$ にも使えると主張する。$z \in Z$ とし、スキームの射 $X_z \to Y_z$ を考える。
これは平坦かつ全射である。仮定により $X_z$ は $\kappa(z)$ 上の有限スキームであり、
特に Artin 環のスペクトルである（代数、補題
\ref{algebra-lemma-finite-dimensional-algebra} による）。補題 \ref{morphisms-lemma-Artinian-affine} により、
射 $X_z \to Y_z$ はアフィン、特に準コンパクトである。補題
\ref{morphisms-lemma-fpqc-quotient-topology} から、$Y_z$ は、$X_z$ についてそうであるため有限離散であると従う。
従って $Y_z$ はアフィンスキームである（これはよい演習であり、例えば $Y_z$ の
全点からなる集合に性質、補題 \ref{properties-lemma-maximal-points-affine} を適用しても従う）。
$Y_z = \Spec(B)$、$X_z = \Spec(A)$ と書く。このとき $A$ は $B$ 上忠実平坦なので、
$B \subset A$ である。従って望む通り $\dim_k(B) \leq \dim_k(A) \leq n$ である。
\end{proof}


\section{雑録}
\label{morphisms-section-misc}

\noindent
他の箇所に収まらない結果を集める。

\begin{lemma}
\label{morphisms-lemma-factor-through-point}
$f : Y \to X$ をスキームの射とし、$x \in X$ を点とする。$Y$ は被約で、
$f(Y)$ は集合論的に $\{x\}$ に含まれると仮定する。このとき $f$ は標準射
$x = \Spec(\kappa(x)) \to X$ を経由する。
\end{lemma}

\begin{proof}
省略する。ヒント：アフィン局所的に考えると可換環論の補題に帰着する。
環準同型 $A \to B$ で $B$ が被約であり、素イデアル $\mathfrak p \subset A$ が
$\Spec(B) \to \Spec(A)$ の像の中で一意に存在するならば、$A \to B$ は
$\kappa(\mathfrak p)$ を経由する。これはよい演習である。
\end{proof}

\begin{lemma}
\label{morphisms-lemma-factor-through-set-of-points}
$f : Y \to X$ をスキームの射とし、$E \subset X$ とする。$X$ は局所 Noether 的で、
$E$ の元の間に非自明な特殊化がなく、$Y$ は被約で、$f(Y) \subset E$ であると仮定する。
このとき $f$ は $\coprod_{x \in E} x \to X$ を経由する。
\end{lemma}

\begin{proof}
$E$ が一元集合なら、補題 \ref{morphisms-lemma-factor-through-point} から従う。$E$ が有限ならば、
$E$ は（$X$ から誘導される位相に関して）特殊化についての仮定により有限離散空間である。
従ってこの場合は一元集合の場合に帰着する。一般の場合、$X$ と $Y$ がアフィンスキームである
場合に帰着する。$f : Y \to X$ は環準同型 $\varphi : A \to B$ に対応するとする。
$A' \subset B$ を $\varphi$ の像とする。$E' \subset \Spec(A') \subset \Spec(A)$ を
$A'$ の極小素イデアル全体とする。代数、補題
\ref{algebra-lemma-injective-minimal-primes-in-image} により、集合 $E'$ は
$\Spec(B) \to \Spec(A') \subset \Spec(A)$ の像に含まれる。従って $E' \subset E$ である。
$A'$ は Noether 的なので、代数、補題
\ref{algebra-lemma-Noetherian-irreducible-components} により $E'$ は有限である。
$\Spec(B) \to \Spec(A)$ の像の他の各点は $E'$ の元の特殊化であり、かつ $E$ に属するので、
その像は $E'$ に含まれると結論する（$E$ の点の間の特殊化についての仮定による）。
従って、上で扱った $E$ が有限の場合に帰着する。
\end{proof}
